\input{preamble}

% Bien, commençons ici.
%
\begin{document}

\title{Espaces simpliciaux}


\maketitle

\phantomsection
\label{section-phantom}

\tableofcontents

\section{Introduction}
\label{section-introduction}

\noindent
Ce chapitre développe une théorie des espaces topologiques simpliciaux, des
espaces annelés simpliciaux, des schémas simpliciaux et des espaces
algébriques simpliciaux. Cette théorie permet parfois de démontrer des principes
du local au global qui paraissent difficiles à établir autrement.
On trouvera quelques exemples d'applications dans les articles
\cite{faltings_finiteness}, \cite{Kiehl} et \cite{HodgeIII}.

\medskip\noindent
Nous supposons dans tout ce chapitre que le lecteur connaît les notions et
résultats fondamentaux du chapitre Méthodes simpliciales, voir
Méthodes simpliciales, section \ref{simplicial-section-introduction}.
En particulier, nous continuons à écrire $X$, et non $X_\bullet$,
pour un objet simplicial.








\section{Espaces topologiques simpliciaux}
\label{section-simplicial-top}

\noindent
Un {\it espace simplicial} est un objet simplicial de la catégorie des
espaces topologiques, dont les morphismes sont les applications continues.
(Nous emploierons ``espace algébrique simplicial'' pour désigner les objets
simpliciaux de la catégorie des espaces algébriques.)
On peut représenter un espace simplicial $X$ comme suit :
$$
\xymatrix{
X_2
\ar@<2ex>[r]
\ar@<0ex>[r]
\ar@<-2ex>[r]
&
X_1
\ar@<1ex>[r]
\ar@<-1ex>[r]
\ar@<1ex>[l]
\ar@<-1ex>[l]
&
X_0
\ar@<0ex>[l]
}
$$
Il y a ici deux morphismes $d^1_0, d^1_1 : X_1 \to X_0$
et un unique morphisme $s^0_0 : X_0 \to X_1$, etc.
Il importe de garder à l'esprit que $d^n_i : X_n \to X_{n - 1}$
doit être considéré comme une ``projection qui oublie la
$i$-ième coordonnée'', et $s^n_j : X_n \to X_{n + 1}$ comme l'application
diagonale qui répète la $j$-ième coordonnée.

\medskip\noindent
Soit $X$ un espace simplicial. Nous lui associons un site
$X_{Zar}$\footnote{Cette notation est analogue à celle de
Sites, exemple \ref{sites-example-site-topological}
et de
Topologies, définition \ref{topologies-definition-big-small-Zariski}.}
à $X$ de la manière suivante.
\begin{enumerate}
\item Un objet de $X_{Zar}$ est un ouvert $U$ de $X_n$ pour un certain $n$,
\item un morphisme $U \to V$ de $X_{Zar}$ est donné par un
$\varphi : [m] \to [n]$, où $n, m$ sont tels que
$U \subset X_n$, $V \subset X_m$, et où $\varphi$ vérifie
$X(\varphi)(U) \subset V$, et
\item un recouvrement $\{U_i \to U\}$ dans $X_{Zar}$ signifie
que $U, U_i \subset X_n$ sont ouverts, que les applications $U_i \to U$ sont
données par $\text{id} : [n] \to [n]$, et que $U = \bigcup U_i$.
\end{enumerate}
Remarquons en particulier que, si $U \to V$ est un morphisme de $X_{Zar}$
donné par $\varphi$, alors $X(\varphi) : X_n \to X_m$ induit bien
une application continue $U \to V$ d'espaces topologiques.

\noindent
Il est clair que ce qui précède est un cas particulier d'une construction
qui associe un site à tout diagramme d'espaces topologiques. Énonçons
le lemme obligé.

\begin{lemma}
\label{lemma-simplicial-site}
Soit $X$ un espace simplicial. Alors $X_{Zar}$, tel qu'il est
défini ci-dessus, est un site.
\end{lemma}

\begin{proof}
Omis.
\end{proof}

\noindent
Soit $X$ un espace simplicial. Soit $\mathcal{F}$ un faisceau sur $X_{Zar}$.
Il résulte clairement de la définition des recouvrements que la restriction
de $\mathcal{F}$ aux ouverts de $X_n$ définit un faisceau $\mathcal{F}_n$
sur l'espace topologique $X_n$. Pour tout $\varphi : [m] \to [n]$, les
applications de restriction de $\mathcal{F}$ pour les couples $U \subset X_n$, $V \subset X_m$
tels que $X(\varphi)(U) \subset V$ définissent un $X(\varphi)$-morphisme
$\mathcal{F}(\varphi) : \mathcal{F}_m \to \mathcal{F}_n$, voir
Faisceaux, définition \ref{sheaves-definition-f-map}.
De plus, étant donnés $\varphi : [m] \to [n]$ et $\psi : [l] \to [m]$,
on a
$$
\mathcal{F}(\varphi) \circ \mathcal{F}(\psi) =
\mathcal{F}(\varphi \circ \psi)
$$
(dans le membre de gauche, on utilise la composition des $f$-morphismes, voir
Faisceaux, définition \ref{sheaves-definition-composition-f-maps}).
La réciproque est manifestement vraie elle aussi : si l'on dispose d'un système
$(\{\mathcal{F}_n\}_{n \geq 0},
\{\mathcal{F}(\varphi)\}_{\varphi \in \text{Flèches}(\Delta)})$
comme ci-dessus, satisfaisant les égalités affichées,
on obtient alors un faisceau sur $X_{Zar}$.

\begin{lemma}
\label{lemma-describe-sheaves-simplicial-site}
Soit $X$ un espace simplicial. Il existe une équivalence de
catégories entre
\begin{enumerate}
\item $\Sh(X_{Zar})$, et
\item la catégorie des systèmes $(\mathcal{F}_n, \mathcal{F}(\varphi))$
décrits ci-dessus.
\end{enumerate}
\end{lemma}

\begin{proof}
Voir la discussion ci-dessus.
\end{proof}

\begin{lemma}
\label{lemma-simplicial-space-site-functorial}
Soit $f : Y \to X$ un morphisme d'espaces simpliciaux.
Alors le foncteur $u : X_{Zar} \to Y_{Zar}$
qui associe à l'ouvert $U \subset X_n$ l'ouvert
$f_n^{-1}(U) \subset Y_n$ définit un morphisme de sites
$f_{Zar} : Y_{Zar} \to X_{Zar}$.
\end{lemma}

\begin{proof}
Il est clair que $u$ est un foncteur continu. On obtient donc les foncteurs
$f_{Zar, *} = u^s$ et $f_{Zar}^{-1} = u_s$, voir
Sites, section \ref{sites-section-morphism-sites}.
Pour voir que l'on obtient un morphisme de sites, il faut montrer
que $u_s$ est exact. Nous utiliserons
Sites, lemme \ref{sites-lemma-directed-morphism} à cette fin.
Soit $V \subset Y_n$ un ouvert. La catégorie
$\mathcal{I}_V^u$ (voir Sites, section \ref{sites-section-functoriality-PSh})
est formée des couples $(U, \varphi)$ où
$\varphi : [m] \to [n]$ et où $U \subset X_m$ est un ouvert tel que
$Y(\varphi)(V) \subset f_m^{-1}(U)$. De plus, un morphisme
$(U, \varphi) \to (U', \varphi')$ est donné par un
$\psi : [m'] \to [m]$ tel que $X(\psi)(U) \subset U'$
et $\varphi \circ \psi = \varphi'$.
Il s'agit de montrer que $\mathcal{I}_V^u$ est cofiltrante.

\medskip\noindent
Vérifions les conditions de
Catégories, définition \ref{categories-definition-codirected}.
La condition (1) est satisfaite parce que $(X_n, \text{id}_{[n]})$ est un objet.
Soit $(U, \varphi)$ un objet. La condition
$Y(\varphi)(V) \subset f_m^{-1}(U)$ équivaut à
$V \subset f_n^{-1}(X(\varphi)^{-1}(U))$. On obtient donc un morphisme
$(X(\varphi)^{-1}(U), \text{id}_{[n]}) \to (U, \varphi)$ en posant
$\psi = \varphi$. De plus, étant donné un couple d'objets
de la forme $(U, \text{id}_{[n]})$ et $(U', \text{id}_{[n]})$,
on voit qu'il existe un objet, à savoir $(U \cap U', \text{id}_{[n]})$,
qui se projette sur chacun d'eux. La condition (2) est donc satisfaite.
Pour vérifier la condition (3), supposons donnés deux morphismes
$a, a': (U, \varphi) \to (U', \varphi')$ définis par $\psi, \psi' : [m'] \to [m]$.
Leur précomposition par le morphisme
$(X(\varphi)^{-1}(U), \text{id}_{[n]}) \to (U, \varphi)$ défini
par $\varphi$ égalise alors $a, a'$, puisque
$\varphi \circ \psi = \varphi' = \varphi \circ \psi'$.
Ceci achève la démonstration.
\end{proof}

\begin{lemma}
\label{lemma-describe-functoriality}
Soit $f : Y \to X$ un morphisme d'espaces simpliciaux. En termes de la
description des faisceaux donnée dans le
lemme \ref{lemma-describe-sheaves-simplicial-site}, le
morphisme $f_{Zar}$ du lemme \ref{lemma-simplicial-space-site-functorial}
se décrit comme suit.
\begin{enumerate}
\item Si $\mathcal{G}$ est un faisceau sur $Y$, alors
$(f_{Zar, *}\mathcal{G})_n = f_{n, *}\mathcal{G}_n$.
\item Si $\mathcal{F}$ est un faisceau sur $X$, alors
$(f_{Zar}^{-1}\mathcal{F})_n = f_n^{-1}\mathcal{F}_n$.
\end{enumerate}
\end{lemma}

\begin{proof}
La première assertion résulte immédiatement des définitions. Pour la seconde,
remarquons que, dans la démonstration du
lemme \ref{lemma-simplicial-space-site-functorial},
nous avons montré que, pour un ouvert $V \subset Y_n$, la catégorie
$(\mathcal{I}_V^u)^{opp}$ contient comme sous-catégorie cofinale
la catégorie des ouverts $U \subset X_n$ tels que $f_n^{-1}(U) \supset V$,
les morphismes étant les inclusions. On voit donc que la restriction
de $u_p\mathcal{F}$ aux ouverts de $Y_n$ est le préfaisceau
$f_{n, p}\mathcal{F}_n$ défini dans
Faisceaux, lemme \ref{sheaves-lemma-pullback-presheaves}.
Puisque $f_{Zar}^{-1}\mathcal{F} = u_s\mathcal{F}$ est le faisceau associé à
$u_p\mathcal{F}$, que sa formation ne fait intervenir que les recouvrements et
que les recouvrements dans $Y_{Zar}$ ne font intervenir que des inclusions
entre ouverts d'un même $Y_n$; le résultat découle du fait que $f_n^{-1}\mathcal{F}_n$
est le faisceau associé à $f_{n, p}\mathcal{F}_n$, voir
Faisceaux, section \ref{sheaves-section-presheaves-functorial}.
\end{proof}

\noindent
Soit $X$ un espace topologique. Dans
Sites, exemple \ref{sites-example-site-topological},
nous avons noté $X_{Zar}$ le site dont les objets sont les ouverts de $X$,
les morphismes les inclusions et les recouvrements les recouvrements ouverts.
Nous identifions le topos $\Sh(X_{Zar})$ à la catégorie
des faisceaux sur $X$.

\begin{lemma}
\label{lemma-restriction-to-components}
Soit $X$ un espace simplicial. Le foncteur
$X_{n, Zar} \to X_{Zar}$, $U \mapsto U$, est continu
et cocontinu. Le morphisme de topos associé
$g_n : \Sh(X_n) \to \Sh(X_{Zar})$ vérifie les propriétés suivantes :
\begin{enumerate}
\item $g_n^{-1}$ associe au faisceau $\mathcal{F}$ sur $X$
le faisceau $\mathcal{F}_n$ sur $X_n$,
\item $g_n^{-1} : \Sh(X_{Zar}) \to \Sh(X_n)$ admet un adjoint à gauche $g^{Sh}_{n!}$,
\item $g^{Sh}_{n!}$ commute aux limites finies connexes,
\item $g_n^{-1} : \textit{Ab}(X_{Zar}) \to \textit{Ab}(X_n)$
admet un adjoint à gauche $g_{n!}$, et
\item $g_{n!}$ est exact.
\end{enumerate}
\end{lemma}

\begin{proof}
Outre les propriétés de notre foncteur mentionnées dans l'énoncé,
la catégorie $X_{n, Zar}$ possède des produits fibrés et des égalisateurs,
et le foncteur commute à ceux-ci (attention, $X_{Zar}$ ne possède pas
tous les produits fibrés). Le lemme résulte donc de la discussion de
Sites, sections \ref{sites-section-cocontinuous-functors} et
\ref{sites-section-cocontinuous-morphism-topoi}
et
Modules sur les sites, section \ref{sites-modules-section-exactness-lower-shriek}.
Plus précisément,
Sites, lemmes \ref{sites-lemma-cocontinuous-morphism-topoi},
\ref{sites-lemma-when-shriek} et
\ref{sites-lemma-preserve-equalizers},
ainsi que de
Modules sur les sites, lemmes
\ref{sites-modules-lemma-g-shriek-adjoint} et
\ref{sites-modules-lemma-exactness-lower-shriek}.
\end{proof}

\begin{lemma}
\label{lemma-restriction-injective-to-component}
Soit $X$ un espace simplicial. Si $\mathcal{I}$ est un faisceau abélien
injectif sur $X_{Zar}$, alors $\mathcal{I}_n$ est un faisceau abélien injectif
sur $X_n$.
\end{lemma}

\begin{proof}
Ceci résulte de
Homologie, lemme \ref{homology-lemma-adjoint-preserve-injectives}
et du
lemme \ref{lemma-restriction-to-components}.
\end{proof}

\begin{lemma}
\label{lemma-restriction-to-components-functorial}
Soit $f : Y \to X$ un morphisme d'espaces simpliciaux. Alors
$$
\xymatrix{
\Sh(Y_n) \ar[d] \ar[r]_{f_n} & \Sh(X_n) \ar[d] \\
\Sh(Y_{Zar}) \ar[r]^{f_{Zar}} & \Sh(X_{Zar})
}
$$
est un diagramme commutatif de topos.
\end{lemma}

\begin{proof}
Cela résulte directement de la description des foncteurs image inverse dans les
lemmes \ref{lemma-describe-functoriality} et
\ref{lemma-restriction-to-components}.
\end{proof}

\begin{lemma}
\label{lemma-augmentation}
Soit $Y$ un espace simplicial et soit $a : Y \to X$ une augmentation
(Méthodes simpliciales, définition \ref{simplicial-definition-augmentation}).
Soient $a_n : Y_n \to X$ les morphismes d'espaces topologiques correspondants.
Il existe un morphisme canonique de topos
$$
a : \Sh(Y_{Zar}) \to \Sh(X)
$$
qui possède les propriétés suivantes :
\begin{enumerate}
\item $a^{-1}\mathcal{F}$ est le faisceau dont la restriction est $a_n^{-1}\mathcal{F}$
sur $Y_n$,
\item $a_m \circ Y(\varphi) = a_n$ pour tout $\varphi : [m] \to [n]$,
\item $a \circ g_n = a_n$ comme morphismes de topos, où
$g_n$ est défini au lemme \ref{lemma-restriction-to-components},
\item $a_*\mathcal{G}$, pour $\mathcal{G} \in \Sh(Y_{Zar})$,
est l'égalisateur des deux applications
$a_{0, *}\mathcal{G}_0 \to a_{1, *}\mathcal{G}_1$.
\end{enumerate}
\end{lemma}

\begin{proof}
La partie (2) vaut pour les augmentations d'objets simpliciaux de toute catégorie.
Ainsi, $Y(\varphi)^{-1} a_m^{-1} \mathcal{F} = a_n^{-1}\mathcal{F}$,
ce qui définit un $Y(\varphi)$-morphisme de $a_m^{-1}\mathcal{F}$
vers $a_n^{-1}\mathcal{F}$.
On peut donc prendre (1) pour définition de $a^{-1}\mathcal{F}$ (en utilisant le
lemme \ref{lemma-describe-sheaves-simplicial-site}) et
(4) pour définition de $a_*$. Si ceci définit un morphisme de topos,
alors la partie (3) en résulte, car on aura $g_n^{-1} \circ a^{-1} = a_n^{-1}$
par construction. Pour vérifier que $a$ est un morphisme de topos, il faut montrer
que $a^{-1}$ est adjoint à gauche de $a_*$ et
que $a^{-1}$ est exact. Ce dernier fait résulte immédiatement de l'exactitude des
foncteurs $a_n^{-1}$.

\medskip\noindent
Soit $\mathcal{F}$ un objet de $\Sh(X)$ et soit $\mathcal{G}$
un objet de $\Sh(Y_{Zar})$. Étant donné
$\beta : a^{-1}\mathcal{F} \to \mathcal{G}$, considérons ses
composantes $\beta_n : a_n^{-1}\mathcal{F} \to \mathcal{G}_n$.
Ces applications sont adjointes aux applications
$\beta_n : \mathcal{F} \to a_{n, *}\mathcal{G}_n$.
La compatibilité à la structure simpliciale montre que
$\beta_0$ se factorise par $a_*\mathcal{G}$.
Réciproquement, supposons donnée une application $\alpha : \mathcal{F} \to a_*\mathcal{G}$.
Pour tout $n$, choisissons un $\varphi : [0] \to [n]$. On peut alors considérer
la composée
$$
\mathcal{F} \xrightarrow{\alpha} a_*\mathcal{G}
\to a_{0, *}\mathcal{G}_0 \xrightarrow{\mathcal{G}(\varphi)}
a_{n, *}\mathcal{G}_n
$$
Ses adjointes sont des applications $a_n^{-1}\mathcal{F} \to \mathcal{G}_n$
qui définissent un morphisme de faisceaux $a^{-1}\mathcal{F} \to \mathcal{G}$.
Nous omettons la vérification que les constructions ci-dessus définissent
des bijections inverses l'une de l'autre
$$
\Mor_{\Sh(Y_{Zar})}(a^{-1}\mathcal{F}, \mathcal{G}) =
\Mor_{\Sh(X)}(\mathcal{F}, a_*\mathcal{G})
$$
Ceci achève la démonstration. Une observation intéressante est que
ce morphisme de topos ne correspond à aucun foncteur géométrique évident
entre les sites qui définissent les topos.
\end{proof}

\begin{lemma}
\label{lemma-simplicial-resolution-Z}
Soit $X$ un espace topologique simplicial. Le complexe de
préfaisceaux abéliens sur $X_{Zar}$
$$
\ldots \to \mathbf{Z}_{X_2} \to \mathbf{Z}_{X_1} \to \mathbf{Z}_{X_0}
$$
dont le différentiel est $\sum (-1)^i d^n_i$ est une résolution
du préfaisceau constant $\mathbf{Z}$.
\end{lemma}

\begin{proof}
Soit $U \subset X_m$ un objet de $X_{Zar}$. La valeur
du complexe ci-dessus sur $U$ est le complexe de groupes abéliens
$$
\ldots \to
\mathbf{Z}[\Mor_\Delta([2], [m])] \to
\mathbf{Z}[\Mor_\Delta([1], [m])] \to
\mathbf{Z}[\Mor_\Delta([0], [m])]
$$
Autrement dit, c'est le complexe associé au
groupe abélien libre engendré par l'ensemble simplicial $\Delta[m]$, voir
Méthodes simpliciales, exemple \ref{simplicial-example-simplex-simplicial-set}.
Puisque $\Delta[m]$ est homotopiquement équivalent à $\Delta[0]$, voir
Méthodes simpliciales, exemple \ref{simplicial-example-simplex-contractible},
et puisque ``prendre les groupes abéliens libres'' est un foncteur,
on voit que le complexe ci-dessus est homotopiquement équivalent au
groupe abélien libre engendré par $\Delta[0]$
(Méthodes simpliciales, remarque \ref{simplicial-remark-homotopy-better} et
lemme \ref{simplicial-lemma-homotopy-equivalence-s-N}).
Ce complexe est acyclique en degrés strictement positifs
et égal à $\mathbf{Z}$ en degré $0$.
\end{proof}

\begin{lemma}
\label{lemma-simplicial-sheaf-cohomology}
Soit $X$ un espace topologique simplicial. Soit $\mathcal{F}$ un faisceau abélien
sur $X$. Il existe une suite spectrale $(E_r, d_r)_{r \geq 0}$ telle que
$$
E_1^{p, q} = H^q(X_p, \mathcal{F}_p)
$$
et qui aboutit à $H^{p + q}(X_{Zar}, \mathcal{F})$.
Cette suite spectrale est fonctorielle en $\mathcal{F}$.
\end{lemma}

\begin{proof}
Soit $\mathcal{F} \to \mathcal{I}^\bullet$ une résolution injective.
Considérons le complexe double de termes
$$
A^{p, q} = \mathcal{I}^q(X_p)
$$
et dont le premier différentiel est donné par la somme alternée des morphismes
$d^{p + 1}_i$ induits $\mathcal{I}_p^q \to \mathcal{I}_{p + 1}^q$, voir
le lemme \ref{lemma-describe-sheaves-simplicial-site}. Remarquons que
$$
A^{p, q} = \Gamma(X_p, \mathcal{I}_p^q) =
\Mor_{\textit{PSh}}(h_{X_p}, \mathcal{I}^q) =
\Mor_{\textit{PAb}}(\mathbf{Z}_{X_p}, \mathcal{I}^q)
$$
Il résulte donc du lemme \ref{lemma-simplicial-resolution-Z} et de
Cohomologie sur les sites, lemme
\ref{sites-cohomology-lemma-injective-abelian-sheaf-injective-presheaf}
que les lignes du complexe double sont exactes en degrés strictement positifs et
ont pour valeur $\Gamma(X_{Zar}, \mathcal{I}^q)$ en degré $0$.
D'autre part, puisque la restriction est exacte
(lemme \ref{lemma-restriction-to-components}),
l'application
$$
\mathcal{F}_p \to \mathcal{I}_p^\bullet
$$
est une résolution. Les faisceaux $\mathcal{I}_p^q$ sont des faisceaux
abéliens injectifs sur $X_p$
(lemme \ref{lemma-restriction-injective-to-component}).
La cohomologie des colonnes calcule donc les groupes
$H^q(X_p, \mathcal{F}_p)$. On conclut en appliquant
Homologie, lemmes \ref{homology-lemma-first-quadrant-ss} et
\ref{homology-lemma-double-complex-gives-resolution}.
\end{proof}

\begin{lemma}
\label{lemma-augmentation-pushforward-higher-direct-image}
Soit $X$ un espace simplicial et soit $a : X \to Y$
une augmentation. Soit $\mathcal{F}$ un faisceau abélien
sur $X_{Zar}$. Alors $R^na_*\mathcal{F}$ est le faisceau associé
au préfaisceau
$$
V \longmapsto H^n((X \times_Y V)_{Zar}, \mathcal{F}|_{(X \times_Y V)_{Zar}})
$$
\end{lemma}

\begin{proof}
C'est l'analogue de
Cohomologie, lemme \ref{cohomology-lemma-describe-higher-direct-images} ou de
Cohomologie sur les sites, lemme \ref{sites-cohomology-lemma-higher-direct-images},
et nous conseillons vivement au lecteur d'omettre la démonstration.
En choisissant une résolution injective de $\mathcal{F}$ sur
$X_{Zar}$ et en appliquant les définitions, on voit qu'il suffit de montrer :
(1) que la restriction d'un faisceau abélien injectif
sur $X_{Zar}$ à $(X \times_Y V)_{Zar}$ est un faisceau abélien injectif, et
(2) que $a_*\mathcal{F}$ est égal à la règle
$$
V \longmapsto H^0((X \times_Y V)_{Zar}, \mathcal{F}|_{(X \times_Y V)_{Zar}})
$$
La partie (2) résulte des faits suivants :
\begin{enumerate}
\item[(2a)] $a_*\mathcal{F}$ est l'égalisateur des deux applications
$a_{0, *}\mathcal{F}_0 \to a_{1, *}\mathcal{F}_1$
d'après le lemme \ref{lemma-augmentation},
\item[(2b)] $a_{0, *}\mathcal{F}_0(V) =
H^0(a_0^{-1}(V), \mathcal{F}_0)$ et
$a_{1, *}\mathcal{F}_1(V) = H^0(a_1^{-1}(V), \mathcal{F}_1)$,
\item[(2c)] $X_0 \times_Y V = a_0^{-1}(V)$ et $X_1 \times_Y V = a_1^{-1}(V)$,
\item[(2d)] $H^0((X \times_Y V)_{Zar}, \mathcal{F}|_{(X \times_Y V)_{Zar}})$
est l'égalisateur des deux applications
$H^0(X_0 \times_Y V, \mathcal{F}_0) \to H^0(X_1 \times_Y V, \mathcal{F}_1)$,
par exemple d'après le lemme \ref{lemma-simplicial-sheaf-cohomology}.
\end{enumerate}
La partie (1) résulte de ce que l'on définit un adjoint à gauche exact
$j_! : \textit{Ab}((X \times_Y V)_{Zar}) \to \textit{Ab}(X_{Zar})$
(extension par zéro) de la restriction
$\textit{Ab}(X_{Zar}) \to \textit{Ab}((X \times_Y V)_{Zar})$,
et que l'on applique Homologie, lemme \ref{homology-lemma-adjoint-preserve-injectives}.
\end{proof}

\noindent
Soit $X$ un espace topologique. Notons $X_\bullet$ l'espace topologique
simplicial constant de valeur $X$. D'après le
lemme \ref{lemma-describe-sheaves-simplicial-site},
se donner un faisceau sur $X_{\bullet, Zar}$ revient à se
donner un objet cosimplicial de la catégorie des faisceaux sur $X$.

\begin{lemma}
\label{lemma-constant-simplicial-space}
Soit $X$ un espace topologique. Soit $X_\bullet$ l'espace topologique
simplicial constant de valeur $X$. Le foncteur
$$
X_{\bullet, Zar} \longrightarrow X_{Zar},\quad
U \longmapsto U
$$
est continu et cocontinu et définit un morphisme de
topos $g : \Sh(X_{\bullet, Zar}) \to \Sh(X)$, ainsi qu'un adjoint à gauche
$g_!$ de $g^{-1}$. On a
\begin{enumerate}
\item $g^{-1}$ associe à un faisceau sur $X$ le faisceau cosimplicial
constant sur $X$,
\item $g_!$ associe à un faisceau $\mathcal{F}$ sur $X_{\bullet, Zar}$ le
faisceau $\mathcal{F}_0$, et
\item $g_*$ associe à un faisceau $\mathcal{F}$ sur $X_{\bullet, Zar}$
l'égalisateur des deux applications $\mathcal{F}_0 \to \mathcal{F}_1$.
\end{enumerate}
\end{lemma}

\begin{proof}
Les assertions concernant le foncteur se vérifient immédiatement.
L'existence de $g$ et de $g_!$ résulte de
Sites, lemmes \ref{sites-lemma-cocontinuous-morphism-topoi} et
\ref{sites-lemma-when-shriek}. La description de
$g^{-1}$ résulte immédiatement de Sites, lemme \ref{sites-lemma-when-shriek}.
La description de $g_*$ et de $g_!$ résulte de ce que les foncteurs indiqués sont
respectivement adjoint à droite et adjoint à gauche de $g^{-1}$.
\end{proof}








\section{Sites simpliciaux et topos}
\label{section-simplicial-sites}

\noindent
Il paraît naturel de définir un {\it site simplicial} comme un objet
simplicial de la (grande) catégorie dont les objets sont les sites
et dont les morphismes sont les morphismes de sites.
Voir Sites, définitions \ref{sites-definition-site} et
\ref{sites-definition-morphism-sites},
la composition des morphismes étant celle de
Sites, lemme \ref{sites-lemma-composition-morphisms-sites}.
Mais voici quelques variantes que l'on pourrait vouloir considérer :
(a) on pourrait plutôt travailler avec des foncteurs cocontinus
(voir Sites, sections \ref{sites-section-cocontinuous-functors} et
\ref{sites-section-cocontinuous-morphism-topoi}) entre sites,
(b) on pourrait travailler dans une $2$-catégorie convenable de sites où l'on introduit
la notion de $2$-morphisme entre morphismes de sites,
(c) on pourrait travailler dans une $2$-catégorie construite à partir de foncteurs
cocontinus. Plutôt que de choisir l'une de ces variantes comme définition,
nous développerons simplement la théorie au fur et à mesure des besoins.

\medskip\noindent
Un {\it topos simplicial} devrait sans doute être défini comme un
pseudo-foncteur de $\Delta^{opp}$ vers la $2$-catégorie des topos.
Voir Catégories, définition \ref{categories-definition-functor-into-2-category}
et Sites, section \ref{sites-section-topoi} et
\ref{sites-section-2-category}. Nous essaierons d'éviter, si possible,
de travailler avec un tel objet.

\medskip\noindent
{\bf Cas A.}
Soit $\mathcal{C}$ un objet simplicial de la catégorie dont les objets
sont les sites et dont les morphismes sont les morphismes de sites. Cela signifie que,
pour tout morphisme $\varphi : [m] \to [n]$ de $\Delta$, on dispose d'un morphisme
de sites $f_\varphi : \mathcal{C}_n \to \mathcal{C}_m$. Ce morphisme est
donné par un foncteur continu en sens opposé, que l'on notera
$u_\varphi : \mathcal{C}_m \to \mathcal{C}_n$.

\begin{lemma}
\label{lemma-simplicial-site-site}
Soit $\mathcal{C}$ un objet simplicial de la catégorie des sites.
Avec les notations ci-dessus, construisons un site $\mathcal{C}_{total}$ comme suit.
\begin{enumerate}
\item Un objet de $\mathcal{C}_{total}$ est un objet $U$ de
$\mathcal{C}_n$ pour un certain $n$,
\item un morphisme $(\varphi, f) : U \to V$ de $\mathcal{C}_{total}$
est donné par une application $\varphi : [m] \to [n]$ avec
$U \in \Ob(\mathcal{C}_n)$, $V \in \Ob(\mathcal{C}_m)$,
et un morphisme $f : U \to u_\varphi(V)$ de $\mathcal{C}_n$, et
\item un recouvrement $\{(\text{id}, f_i) :  U_i \to U\}$ dans $\mathcal{C}_{total}$
est donné par un entier $n$ et un recouvrement $\{f_i : U_i \to U\}$
de $\mathcal{C}_n$.
\end{enumerate}
\end{lemma}

\begin{proof}
La composée de $(\varphi, f) : U \to V$ et de $(\psi, g) : V \to W$
est donnée par $(\varphi \circ \psi, u_\varphi(g) \circ f)$.
On utilise ici l'égalité $u_\varphi \circ u_\psi = u_{\varphi \circ \psi}$.

\medskip\noindent
Soit $\{(\text{id}, f_i) :  U_i \to U\}$ un recouvrement comme en (3),
et soit $(\varphi, g) : W \to U$ un morphisme avec
$W \in \Ob(\mathcal{C}_m)$. Nous affirmons que
$$
W \times_{(\varphi, g), U, (\text{id}, f_i)} U_i =
W \times_{g, u_\varphi(U), u_\varphi(f_i)} u_\varphi(U_i)
$$
dans la catégorie $\mathcal{C}_{total}$. Cela a un sens car, d'après notre
définition des morphismes de sites, les produits fibrés requis
existent dans $\mathcal{C}_m$, puisque $u_\varphi$ transforme les recouvrements en
recouvrements. Le même raisonnement démontre l'affirmation (détails omis).
On voit donc que la famille des recouvrements est stable par changement
de base. Les autres axiomes d'un site sont immédiats.
\end{proof}

\noindent
{\bf Cas B.}
Soit $\mathcal{C}$ un objet simplicial de la catégorie dont les objets sont les
sites et dont les morphismes sont les foncteurs cocontinus. Cela signifie que, pour
tout morphisme $\varphi : [m] \to [n]$ de $\Delta$, on dispose d'un foncteur cocontinu
noté $u_\varphi : \mathcal{C}_n \to \mathcal{C}_m$. Le morphisme de topos
associé est noté
$f_\varphi : \Sh(\mathcal{C}_n) \to \Sh(\mathcal{C}_m)$.

\begin{lemma}
\label{lemma-simplicial-cocontinuous-site}
Soit $\mathcal{C}$ un objet simplicial de la catégorie dont les objets sont les
sites et dont les morphismes sont les foncteurs cocontinus. Avec les notations ci-dessus,
supposons que les foncteurs $u_\varphi : \mathcal{C}_n \to \mathcal{C}_m$
possèdent la propriété $P$ de Sites, remarque \ref{sites-remark-cartesian-cocontinuous}.
On peut alors construire un site $\mathcal{C}_{total}$ comme suit.
\begin{enumerate}
\item Un objet de $\mathcal{C}_{total}$ est un objet $U$ de
$\mathcal{C}_n$ pour un certain $n$,
\item un morphisme $(\varphi, f) : U \to V$ de $\mathcal{C}_{total}$
est donné par une application $\varphi : [m] \to [n]$ avec
$U \in \Ob(\mathcal{C}_n)$, $V \in \Ob(\mathcal{C}_m)$,
et un morphisme $f : u_\varphi(U) \to V$ de $\mathcal{C}_m$, et
\item un recouvrement $\{(\text{id}, f_i) :  U_i \to U\}$ dans $\mathcal{C}_{total}$
est donné par un entier $n$ et un recouvrement $\{f_i : U_i \to U\}$
de $\mathcal{C}_n$.
\end{enumerate}
\end{lemma}

\begin{proof}
La composée de $(\varphi, f) : U \to V$ et de $(\psi, g) : V \to W$
est donnée par $(\varphi \circ \psi, g \circ u_\psi(f))$.
On utilise ici l'égalité $u_\psi \circ u_\varphi = u_{\varphi \circ \psi}$.

\medskip\noindent
Soit $\{(\text{id}, f_i) :  U_i \to U\}$ un recouvrement comme en (3),
et soit $(\varphi, g) : W \to U$ un morphisme avec
$W \in \Ob(\mathcal{C}_m)$. Nous affirmons que
$$
W \times_{(\varphi, g), U, (\text{id}, f_i)} U_i =
W \times_{g, U, f_i} U_i
$$
dans la catégorie $\mathcal{C}_{total}$, où le membre de droite
est l'objet de $\mathcal{C}_m$ défini dans
Sites, remarque \ref{sites-remark-cartesian-cocontinuous},
qui existe d'après la propriété $P$. La compatibilité de ce type de produit fibré
avec la composition des foncteurs démontre l'affirmation (détails omis).
Comme la famille $\{W \times_{g, U, f_i} U_i \to W\}$ est un
recouvrement de $\mathcal{C}_m$ d'après la propriété $P$, on voit que
la famille des recouvrements est stable par changement
de base. Les autres axiomes d'un site sont immédiats.
\end{proof}

\begin{situation}
\label{situation-simplicial-site}
On se trouve ici dans l'un des deux cas suivants :
\begin{enumerate}
\item[(A)] $\mathcal{C}$ est un objet simplicial de la catégorie dont les
objets sont les sites et dont les morphismes sont les morphismes de sites. Pour tout
morphisme $\varphi : [m] \to [n]$ de $\Delta$, on dispose d'un morphisme de sites
$f_\varphi : \mathcal{C}_n \to \mathcal{C}_m$ donné par un foncteur continu
$u_\varphi : \mathcal{C}_m \to \mathcal{C}_n$.
\item[(B)] $\mathcal{C}$ est un objet simplicial de la catégorie dont les
objets sont les sites et dont les morphismes sont des foncteurs cocontinus possédant la
propriété $P$ de Sites, remarque \ref{sites-remark-cartesian-cocontinuous}.
Pour tout morphisme $\varphi : [m] \to [n]$ de $\Delta$, on dispose d'un foncteur cocontinu
$u_\varphi : \mathcal{C}_n \to \mathcal{C}_m$ qui induit un
morphisme de topos $f_\varphi : \Sh(\mathcal{C}_n) \to \Sh(\mathcal{C}_m)$.
\end{enumerate}
Comme d'habitude, on notera $f_\varphi^{-1}$ et $f_{\varphi, *}$ les foncteurs
image inverse et image directe. Notons $\mathcal{C}_{total}$ le
site défini au
lemme \ref{lemma-simplicial-site-site} (cas A) ou au
lemme \ref{lemma-simplicial-cocontinuous-site} (cas B).
\end{situation}

\noindent
Soit $\mathcal{C}$ comme dans la situation \ref{situation-simplicial-site}.
Soit $\mathcal{F}$ un faisceau sur $\mathcal{C}_{total}$.
Il résulte clairement de la définition des recouvrements que la restriction
de $\mathcal{F}$ aux objets de $\mathcal{C}_n$ définit un faisceau
$\mathcal{F}_n$ sur le site $\mathcal{C}_n$. Pour tout
$\varphi : [m] \to [n]$, les applications de restriction de $\mathcal{F}$
le long des morphismes $(\varphi, f) : U \to V$ avec
$U \in \Ob(\mathcal{C}_n)$ et $V \in \Ob(\mathcal{C}_m)$
définissent un élément $\mathcal{F}(\varphi)$ de
$$
\Mor_{\Sh(\mathcal{C}_m)}(\mathcal{F}_m, f_{\varphi, *}\mathcal{F}_n) =
\Mor_{\Sh(\mathcal{C}_n)}(f_\varphi^{-1}\mathcal{F}_m, \mathcal{F}_n)
$$
De plus, étant donnés $\varphi : [m] \to [n]$ et $\psi : [l] \to [m]$,
les diagrammes
$$
\vcenter{
\xymatrix{
\mathcal{F}_l \ar[rr]_{\mathcal{F}(\varphi \circ \psi)}
\ar[rd]_{\mathcal{F}(\psi)}
& & f_{\varphi \circ \psi, *} \mathcal{F}_n \\
& f_{\psi, *}\mathcal{F}_m \ar[ur]_{f_{\psi, *}\mathcal{F}(\varphi)}
}
}
\quad\text{et}\quad
\vcenter{
\xymatrix{
f_{\varphi \circ \psi}^{-1}\mathcal{F}_l
\ar[rr]_{\mathcal{F}(\varphi \circ \psi)}
\ar[rd]_{f_\varphi^{-1}\mathcal{F}(\psi)}
& & \mathcal{F}_n \\
& f_\varphi^{-1}\mathcal{F}_m \ar[ur]_{\mathcal{F}(\varphi)}
}
}
$$
sont commutatifs. La réciproque est manifestement vraie elle aussi : si l'on a un système
$(\{\mathcal{F}_n\}_{n \geq 0},
\{\mathcal{F}(\varphi)\}_{\varphi \in \text{Flèches}(\Delta)})$
satisfaisant les conditions de commutativité ci-dessus,
on obtient alors un faisceau sur $\mathcal{C}_{total}$.

\begin{lemma}
\label{lemma-describe-sheaves-simplicial-site-site}
Dans la situation \ref{situation-simplicial-site}, il existe une équivalence de
catégories entre
\begin{enumerate}
\item $\Sh(\mathcal{C}_{total})$, et
\item la catégorie des systèmes $(\mathcal{F}_n, \mathcal{F}(\varphi))$
décrits ci-dessus.
\end{enumerate}
En particulier, le topos $\Sh(\mathcal{C}_{total})$ ne dépend que des
topos $\Sh(\mathcal{C}_n)$ et des morphismes de topos $f_\varphi$.
\end{lemma}

\begin{proof}
Voir la discussion ci-dessus.
\end{proof}

\begin{lemma}
\label{lemma-restriction-to-components-site}
Dans la situation \ref{situation-simplicial-site}, le foncteur
$\mathcal{C}_n \to \mathcal{C}_{total}$, $U \mapsto U$, est continu
et cocontinu. Le morphisme de
topos associé $g_n : \Sh(\mathcal{C}_n) \to \Sh(\mathcal{C}_{total})$ vérifie
\begin{enumerate}
\item $g_n^{-1}$ associe au faisceau $\mathcal{F}$ sur $\mathcal{C}_{total}$
le faisceau $\mathcal{F}_n$ sur $\mathcal{C}_n$,
\item $g_n^{-1} : \Sh(\mathcal{C}_{total}) \to \Sh(\mathcal{C}_n)$
admet un adjoint à gauche $g^{Sh}_{n!}$,
\item pour $\mathcal{G}$ dans $\Sh(\mathcal{C}_n)$, la restriction de
$g_{n!}^{Sh}\mathcal{G}$ à $\mathcal{C}_m$ est
$\coprod\nolimits_{\varphi : [n] \to [m]} f_\varphi^{-1}\mathcal{G}$,
\item $g_{n!}^{Sh}$ commute aux limites finies connexes,
\item $g_n^{-1} : \textit{Ab}(\mathcal{C}_{total}) \to
\textit{Ab}(\mathcal{C}_n)$ admet un adjoint à gauche $g_{n!}$,
\item pour $\mathcal{G}$ dans $\textit{Ab}(\mathcal{C}_n)$, la restriction de
$g_{n!}\mathcal{G}$ à $\mathcal{C}_m$ est
$\bigoplus\nolimits_{\varphi : [n] \to [m]} f_\varphi^{-1}\mathcal{G}$, et
\item $g_{n!}$ est exact.
\end{enumerate}
\end{lemma}

\begin{proof}
Cas A. Si $\{U_i \to U\}_{i \in I}$ est un recouvrement dans $\mathcal{C}_n$,
alors son image $\{U_i \to U\}_{i \in I}$ est un recouvrement dans $\mathcal{C}_{total}$
par définition (lemme \ref{lemma-simplicial-site-site}). Pour un morphisme
$V \to U$ de $\mathcal{C}_n$, le produit fibré
$V \times_U U_i$ dans $\mathcal{C}_n$ est aussi
le produit fibré dans $\mathcal{C}_{total}$ (d'après l'affirmation de la
démonstration du lemme \ref{lemma-simplicial-site-site}).
Notre foncteur est donc continu. D'autre part, il
définit une bijection entre les recouvrements de $U$ dans $\mathcal{C}_n$
et les recouvrements de $U$ dans $\mathcal{C}_{total}$. Il est donc
assurément cocontinu.

\medskip\noindent
Cas B. Si $\{U_i \to U\}_{i \in I}$ est un recouvrement dans $\mathcal{C}_n$,
alors son image $\{U_i \to U\}_{i \in I}$ est un recouvrement dans $\mathcal{C}_{total}$
par définition (lemme \ref{lemma-simplicial-cocontinuous-site}). Pour un morphisme
$V \to U$ de $\mathcal{C}_n$, le produit fibré
$V \times_U U_i$ dans $\mathcal{C}_n$ est aussi
le produit fibré dans $\mathcal{C}_{total}$ (d'après l'affirmation de la
démonstration du lemme \ref{lemma-simplicial-cocontinuous-site}).
Notre foncteur est donc continu. D'autre part, il
définit une bijection entre les recouvrements de $U$ dans $\mathcal{C}_n$
et les recouvrements de $U$ dans $\mathcal{C}_{total}$. Il est donc
assurément cocontinu.

\medskip\noindent
À ce stade, la partie (1) et l'existence de $g^{Sh}_{n!}$ et de $g_{n!}$
dans les cas A et B résultent de
Sites, lemmes \ref{sites-lemma-cocontinuous-morphism-topoi} et
\ref{sites-lemma-when-shriek},
ainsi que de
Modules sur les sites, lemme \ref{sites-modules-lemma-g-shriek-adjoint}.

\medskip\noindent
Démonstration de (3). Soit $\mathcal{G}$ un faisceau sur $\mathcal{C}_n$.
Considérons le faisceau $\mathcal{H}$ sur $\mathcal{C}_{total}$
dont le terme de degré $m$ est le faisceau
$$
\mathcal{H}_m = \coprod\nolimits_{\varphi : [n] \to [m]}
f_\varphi^{-1}\mathcal{G}
$$
donné dans la partie (3) de l'énoncé du lemme.
Étant donnée une application $\psi : [m] \to [m']$, l'application
$\mathcal{H}(\psi) : f_\psi^{-1}\mathcal{H}_m \to \mathcal{H}_{m'}$
est donnée sur les composantes par les identifications
$$
f_\psi^{-1} f_\varphi^{-1} \mathcal{G} \to
f_{\psi \circ \varphi}^{-1}\mathcal{G}
$$
Observons qu'une application $\alpha : \mathcal{H} \to \mathcal{F}$
de faisceaux sur $\mathcal{C}_{total}$ fournit une application
$\mathcal{G} \to \mathcal{F}_n$
qui correspond à la restriction de $\alpha_n$ à la composante
$\mathcal{G}$ de $\mathcal{H}_n$. Réciproquement, étant donnée une application
$\beta : \mathcal{G} \to \mathcal{F}_n$ de faisceaux sur $\mathcal{C}_n$,
on peut définir
$\alpha : \mathcal{H} \to \mathcal{F}$ en prenant pour $\alpha_m$
l'application qui, sur les composantes
$$
f_\varphi^{-1}\mathcal{G} \to \mathcal{F}_m
$$
utilise les applications adjointes à $\mathcal{F}(\varphi) \circ f_\varphi^{-1}\beta$.
Nous omettons la vérification que ces deux constructions donnent
des applications inverses l'une de l'autre
$$
\Mor_{\Sh(\mathcal{C}_n)}(\mathcal{G}, \mathcal{F}_n) =
\Mor_{\Sh(\mathcal{C}_{total})}(\mathcal{H}, \mathcal{F})
$$
Ainsi, $\mathcal{H} = g^{Sh}_{n!}\mathcal{G}$, comme voulu.

\medskip\noindent
Démonstration de (4). Si $\mathcal{G}$ est un faisceau abélien sur $\mathcal{C}_n$,
on procède exactement de la même manière que ci-dessus, à ceci près que
l'on définit $\mathcal{H}$ comme le faisceau abélien sur $\mathcal{C}_{total}$
dont le terme de degré $m$ est le faisceau
$$
\bigoplus\nolimits_{\varphi : [n] \to [m]} f_\varphi^{-1}\mathcal{G}
$$
avec des applications de transition définies exactement comme ci-dessus. La bijection
$$
\Mor_{\textit{Ab}(\mathcal{C}_n)}(\mathcal{G}, \mathcal{F}_n) =
\Mor_{\textit{Ab}(\mathcal{C}_{total})}(\mathcal{H}, \mathcal{F})
$$
se démontre exactement comme ci-dessus.
Ainsi, $\mathcal{H} = g_{n!}\mathcal{G}$, comme voulu.

\medskip\noindent
Les propriétés d'exactitude de $g^{Sh}_{n!}$ et de $g_{n!}$ résultent
des formules données pour ces foncteurs.
\end{proof}

\begin{lemma}
\label{lemma-restriction-injective-to-component-site}
\begin{slogan}
Un faisceau abélien injectif sur un site simplicial est injectif sur chaque composante
\end{slogan}
Dans la situation \ref{situation-simplicial-site},
si $\mathcal{I}$ est injectif dans $\textit{Ab}(\mathcal{C}_{total})$,
alors $\mathcal{I}_n$ est injectif dans $\textit{Ab}(\mathcal{C}_n)$.
Si $\mathcal{I}^\bullet$ est un complexe K-injectif dans
$\textit{Ab}(\mathcal{C}_{total})$,
alors $\mathcal{I}_n^\bullet$ est K-injectif dans $\textit{Ab}(\mathcal{C}_n)$.
\end{lemma}

\begin{proof}
La première assertion résulte de
Homologie, lemme \ref{homology-lemma-adjoint-preserve-injectives}
et du
lemme \ref{lemma-restriction-to-components-site}.
La seconde résulte de
Catégories dérivées, lemme \ref{derived-lemma-adjoint-preserve-K-injectives}
et du
lemme \ref{lemma-restriction-to-components-site}.
\end{proof}







\section{Augmentations de sites simpliciaux}
\label{section-augmentation-simplicial-sites}

\noindent
Nous poursuivons selon la méthode décrite dans la
section \ref{section-simplicial-sites},
en explicitant le sens des augmentations dans les cas A et B
traités dans cette section.

\begin{remark}
\label{remark-augmentation-site}
Dans la situation \ref{situation-simplicial-site}, une
{\it augmentation $a_0$ vers un site $\mathcal{D}$} désignera
\begin{enumerate}
\item[(A)] un morphisme de sites $a_0 : \mathcal{C}_0 \to \mathcal{D}$
donné par un foncteur continu $u_0 : \mathcal{D} \to \mathcal{C}_0$
tel que, pour tous $\varphi, \psi : [0] \to [n]$, on ait
$u_\varphi \circ u_0 = u_\psi \circ u_0$.
\item[(B)] un morphisme $a_0 : \Sh(\mathcal{C}_0) \to \Sh(\mathcal{D})$
de topos donné par un foncteur cocontinu $u_0 : \mathcal{C}_0 \to \mathcal{D}$
tel que, pour tous $\varphi, \psi : [0] \to [n]$, on ait
$u_0 \circ u_\varphi = u_0 \circ u_\psi$.
\end{enumerate}
\end{remark}

\begin{lemma}
\label{lemma-augmentation-site}
Dans la situation \ref{situation-simplicial-site}, soit $a_0$ une
augmentation vers un site $\mathcal{D}$ au sens de la
remarque \ref{remark-augmentation-site}. Alors $a_0$ induit
\begin{enumerate}
\item un morphisme de topos $a_n : \Sh(\mathcal{C}_n) \to \Sh(\mathcal{D})$
pour tout $n \geq 0$,
\item un morphisme de topos $a : \Sh(\mathcal{C}_{total}) \to \Sh(\mathcal{D})$
\end{enumerate}
tel que
\begin{enumerate}
\item pour tout $\varphi : [m] \to [n]$, on ait $a_m \circ f_\varphi = a_n$,
\item si $g_n : \Sh(\mathcal{C}_n) \to \Sh(\mathcal{C}_{total})$
est défini comme au lemme \ref{lemma-restriction-to-components-site}, alors
$a \circ g_n = a_n$, et
\item $a_*\mathcal{F}$, pour $\mathcal{F} \in \Sh(\mathcal{C}_{total})$,
est l'égalisateur des deux applications
$a_{0, *}\mathcal{F}_0 \to a_{1, *}\mathcal{F}_1$.
\end{enumerate}
\end{lemma}

\begin{proof}
Cas A. Soit $u_n : \mathcal{D} \to \mathcal{C}_n$ la valeur commune
des foncteurs $u_\varphi \circ u_0$ pour $\varphi : [0] \to [n]$.
Alors $u_n$ correspond à un morphisme de sites
$a_n : \mathcal{C}_n \to \mathcal{D}$, voir
Sites, lemme \ref{sites-lemma-composition-morphisms-sites}.
Le même lemme montre que, pour tout $\varphi : [m] \to [n]$, on a
$a_m \circ f_\varphi = a_n$.

\medskip\noindent
Cas B. Soit $u_n : \mathcal{C}_n \to \mathcal{D}$ la valeur commune
des foncteurs $u_0 \circ u_\varphi$ pour $\varphi : [0] \to [n]$.
Alors $u_n$ est cocontinu et définit donc un morphisme de topos
$a_n : \Sh(\mathcal{C}_n) \to \Sh(\mathcal{D)}$, voir
Sites, lemme \ref{sites-lemma-composition-cocontinuous}.
Le même lemme montre que, pour tout $\varphi : [m] \to [n]$, on a
$a_m \circ f_\varphi = a_n$.

\medskip\noindent
Considérons le foncteur $a^{-1} : \Sh(\mathcal{D}) \to \Sh(\mathcal{C}_{total})$
qui, à un faisceau d'ensembles $\mathcal{G}$, associe le faisceau
$\mathcal{F} = a^{-1}\mathcal{G}$ dont les composantes sont $a_n^{-1}\mathcal{G}$
et dont les applications de transition $\mathcal{F}(\varphi)$ sont les identifications
$$
f_\varphi^{-1}\mathcal{F}_m =
f_\varphi^{-1} a_m^{-1}\mathcal{G} =
a_n^{-1}\mathcal{G} =
\mathcal{F}_n
$$
pour $\varphi : [m] \to [n]$, voir la description de
$\Sh(\mathcal{C}_{total})$ dans le
lemme \ref{lemma-describe-sheaves-simplicial-site-site}.
Puisque les foncteurs $a_n^{-1}$ sont exacts, $a^{-1}$ est un foncteur exact.
Enfin, pour $a_* : \Sh(\mathcal{C}_{total}) \to \Sh(\mathcal{D})$,
on prend le foncteur qui, à un faisceau $\mathcal{F}$ sur $\Sh(\mathcal{D})$,
associe
$$
\xymatrix{
a_*\mathcal{F} \ar@{=}[r] &
\text{Égalisateur}(a_{0, *}\mathcal{F}_0
\ar@<1ex>[r] \ar@<-1ex>[r] &
a_{1, *}\mathcal{F}_1)
}
$$
Les deux applications proviennent ici des deux applications $\varphi : [0] \to [1]$
par
$$
a_{0, *}\mathcal{F}_0 \to
a_{0, *}f_{\varphi, *} f_\varphi^{-1}\mathcal{F}_0
\xrightarrow{\mathcal{F}(\varphi)}
a_{0, *}f_{\varphi, *} \mathcal{F}_1 = a_{1, *}\mathcal{F}_1
$$
où la première flèche provient de $1 \to f_{\varphi, *} f_\varphi^{-1}$.
Notons $\mathcal{G}_\bullet$ le faisceau cosimplicial constant
de valeur $\mathcal{G}$, et $a_{\bullet, *}\mathcal{F}$
le faisceau cosimplicial ayant $a_{n, *}\mathcal{F}_n$ en degré $n$.
Par l'adjonction usuelle pour les morphismes de topos $a_n$, on voit qu'une
application $a^{-1}\mathcal{G} \to \mathcal{F}$
revient à se donner une application
$$
\mathcal{G}_\bullet \longrightarrow a_{\bullet, *}\mathcal{F}
$$
de faisceaux cosimpliciaux.
Par dualité avec Méthodes simpliciales, lemme \ref{simplicial-lemma-augmentation-howto},
cela revient à se donner une application $\mathcal{G} \to a_*\mathcal{F}$.
Ainsi, $a^{-1}$ et $a_*$ sont des foncteurs adjoints, et l'on obtient
notre morphisme de topos $a$\footnote{Dans le cas B, le morphisme $a$
correspond au foncteur cocontinu
$\mathcal{C}_{total} \to \mathcal{D}$ qui envoie
$U$ dans $\mathcal{C}_n$ sur $u_n(U)$.}. Les égalités
$a \circ g_n = f_n$ résultent immédiatement des définitions.
\end{proof}



\section{Morphismes de sites simpliciaux}
\label{section-morphism-simplicial-sites}

\noindent
Nous poursuivons de la manière décrite dans la
section \ref{section-simplicial-sites},
en explicitant ce que signifient les morphismes de sites simpliciaux
dans les cas A et B traités dans cette section.

\begin{remark}
\label{remark-morphism-simplicial-sites}
Soient $\mathcal{C}_n, f_\varphi, u_\varphi$ et
$\mathcal{C}'_n, f'_\varphi, u'_\varphi$ comme dans la
situation \ref{situation-simplicial-site}. Par
{\it morphisme $h$ entre sites simpliciaux}, on entendra
\begin{enumerate}
\item[(A)] Des morphismes de sites
$h_n : \mathcal{C}_n \to \mathcal{C}'_n$
tels que $f'_\varphi \circ h_n = h_m \circ f_\varphi$
comme morphismes de sites pour tout $\varphi : [m] \to [n]$.
\item[(B)] Des foncteurs cocontinus
$v_n : \mathcal{C}_n \to \mathcal{C}'_n$
induisant des morphismes de topos $h_n : \Sh(\mathcal{C}_n) \to \Sh(\mathcal{C}'_n)$
tels que $u'_\varphi \circ v_n = v_m \circ u_\varphi$
comme foncteurs pour tout $\varphi : [m] \to [n]$.
\end{enumerate}
Dans les deux cas, on a
$f'_\varphi \circ h_n = h_m \circ f_\varphi$
comme morphismes de topos ; voir
Sites, lemme \ref{sites-lemma-composition-cocontinuous}
pour le cas B et Sites,
définition \ref{sites-definition-composition-morphisms-sites}
pour le cas A.
\end{remark}

\begin{lemma}
\label{lemma-morphism-simplicial-sites}
Soient $\mathcal{C}_n, f_\varphi, u_\varphi$ et
$\mathcal{C}'_n, f'_\varphi, u'_\varphi$ comme dans la
situation \ref{situation-simplicial-site}.
Soit $h$ un morphisme entre sites simpliciaux comme dans la
remarque \ref{remark-morphism-simplicial-sites}.
Alors on obtient un morphisme de topos
$$
h_{total} : \Sh(\mathcal{C}_{total}) \to \Sh(\mathcal{C}'_{total})
$$
et des diagrammes commutatifs
$$
\xymatrix{
\Sh(\mathcal{C}_n) \ar[d]_{g_n} \ar[r]_{h_n} &
\Sh(\mathcal{C}'_n) \ar[d]^{g'_n} \\
\Sh(\mathcal{C}_{total}) \ar[r]^{h_{total}} &
\Sh(\mathcal{C}'_{total})
}
$$
De plus, on a $(g'_n)^{-1} \circ h_{total, *} = h_{n, *} \circ g_n^{-1}$.
\end{lemma}

\begin{proof}
Cas A. Disons que $h_n$ correspond au foncteur continu
$v_n : \mathcal{C}'_n \to \mathcal{C}_n$. On peut alors définir
un foncteur $v_{total} : \mathcal{C}'_{total} \to \mathcal{C}_{total}$
en utilisant $v_n$ en degré $n$. Il s'agit clairement d'un foncteur continu
(voir la définition des recouvrements dans le lemme \ref{lemma-simplicial-site-site}).
Soient
$h_{total}^{-1} = v_{total, s} :
\Sh(\mathcal{C}'_{total}) \to \Sh(\mathcal{C}_{total})$ et
$h_{total, *} = v_{total}^s = v_{total}^p :
\Sh(\mathcal{C}_{total}) \to \Sh(\mathcal{C}'_{total})$
le couple de foncteurs adjoints construit et étudié dans
Sites, sections \ref{sites-section-continuous-functors} et
\ref{sites-section-morphism-sites}.
Pour voir que $h_{total}$ est un morphisme de topos,
il reste à vérifier que $h_{total}^{-1}$ est exact.
On observe d'abord que
$(g'_n)^{-1} \circ h_{total, *} = h_{n, *} \circ g_n^{-1}$ ;
cela résulte immédiatement du calcul des sections sur un objet $U$
de $\mathcal{C}'_n$. Ainsi, si l'on considère un faisceau $\mathcal{F}$
sur $\mathcal{C}_{total}$ comme un système $(\mathcal{F}_n, \mathcal{F}(\varphi))$
comme dans le lemme \ref{lemma-describe-sheaves-simplicial-site-site}, alors
$h_{total, *}\mathcal{F}$ correspond
au système $(h_{n, *}\mathcal{F}_n, h_{n, *}\mathcal{F}(\varphi))$.
Manifestement, le foncteur
$(\mathcal{F}'_n, \mathcal{F}'(\varphi)) \to
(h_n^{-1}\mathcal{F}'_n, h_n^{-1}\mathcal{F}'(\varphi))$
est son adjoint à gauche. Par unicité des adjoints, on conclut que
$h_{total}^{-1}$ est donné sur les systèmes par cette règle. En particulier,
$h_{total}^{-1}$ est exact (d'après la description des faisceaux sur
$\mathcal{C}_{total}$ donnée dans le lemme et l'exactitude des
foncteurs $h_n^{-1}$), et l'on obtient notre morphisme de topos.
Enfin, on obtient également $g_n^{-1} \circ h_{total}^{-1} =
h_n^{-1} \circ (g'_n)^{-1}$, ce qui prouve que le
diagramme affiché dans le lemme est commutatif.

\medskip\noindent
Cas B. On dispose ici d'un foncteur
$v_{total} : \mathcal{C}_{total} \to \mathcal{C}'_{total}$
obtenu en utilisant $v_n$ en degré $n$. Il s'agit clairement d'un foncteur cocontinu
(voir la définition des recouvrements dans le lemme \ref{lemma-simplicial-cocontinuous-site}).
Soit $h_{total}$ le morphisme de topos associé à $v_{total}$.
La commutativité du diagramme affiché dans le lemme résulte
immédiatement de Sites, lemme \ref{sites-lemma-composition-cocontinuous}.
En prenant les adjoints à gauche, la dernière égalité du lemme devient
$$
h_{total}^{-1} \circ (g'_n)^{Sh}_! = g^{Sh}_{n!} \circ h_n^{-1}
$$
Cela résulte immédiatement de la description explicite des foncteurs
$(g'_n)^{Sh}_!$ et $g^{Sh}_{n!}$ donnée dans le
lemme \ref{lemma-restriction-to-components-site},
du fait que $h_n^{-1} \circ (f'_\varphi)^{-1} =
f_\varphi^{-1} \circ h_m^{-1}$ pour $\varphi : [m] \to [n]$, et
du fait que l'on sait déjà que $h_{total}^{-1}$ commute
aux restrictions aux composantes de degré $n$ des sites simpliciaux.
\end{proof}

\begin{lemma}
\label{lemma-direct-image-morphism-simplicial-sites}
Avec les notations et les hypothèses du lemme \ref{lemma-morphism-simplicial-sites},
pour $K \in D(\mathcal{C}_{total})$, on a
$(g'_n)^{-1}Rh_{total, *}K = Rh_{n, *}g_n^{-1}K$.
\end{lemma}

\begin{proof}
Soit $\mathcal{I}^\bullet$ un complexe K-injectif sur $\mathcal{C}_{total}$
représentant $K$. Alors $g_n^{-1}K$ est représenté par
$g_n^{-1}\mathcal{I}^\bullet = \mathcal{I}_n^\bullet$,
qui est K-injectif d'après le
lemme \ref{lemma-restriction-injective-to-component-site}.
On a $(g'_n)^{-1}h_{total, *}\mathcal{I}^\bullet =
h_{n, *}g_n^{-1}\mathcal{I}_n^\bullet$ d'après le
lemme \ref{lemma-morphism-simplicial-sites},
ce qui donne l'égalité voulue.
\end{proof}

\begin{remark}
\label{remark-morphism-augmentation-simplicial-sites}
Soient $\mathcal{C}_n, f_\varphi, u_\varphi$ et
$\mathcal{C}'_n, f'_\varphi, u'_\varphi$ comme dans la
situation \ref{situation-simplicial-site}.
Soit $a_0$, resp.\ $a'_0$, une augmentation
vers un site $\mathcal{D}$, resp.\ $\mathcal{D}'$,
comme dans la remarque \ref{remark-augmentation-site}.
Soit $h$ un morphisme entre sites simpliciaux comme dans la
remarque \ref{remark-morphism-simplicial-sites}.
On dit qu'un morphisme de topos $h_{-1} : \Sh(\mathcal{D}) \to \Sh(\mathcal{D}')$
est {\it compatible avec $h$, $a_0$, $a'_0$} si
\begin{enumerate}
\item[(A)] $h_{-1}$ provient d'un morphisme de sites
$h_{-1} : \mathcal{D} \to \mathcal{D}'$
tel que $a'_0 \circ h_0 = h_{-1} \circ a_0$
comme morphismes de sites.
\item[(B)] $h_{-1}$ provient d'un foncteur cocontinu
$v_{-1} : \mathcal{D} \to \mathcal{D}'$
tel que $u'_0 \circ v_0 = v_{-1} \circ u_0$
comme foncteurs.
\end{enumerate}
Dans les deux cas, on a $a'_0 \circ h_0 = h_{-1} \circ a_0$
comme morphismes de topos ; voir
Sites, lemme \ref{sites-lemma-composition-cocontinuous}
pour le cas B et Sites,
définition \ref{sites-definition-composition-morphisms-sites}
pour le cas A.
\end{remark}

\begin{lemma}
\label{lemma-morphism-augmentation-simplicial-sites}
Soient $\mathcal{C}_n, f_\varphi, u_\varphi, \mathcal{D}, a_0$,
$\mathcal{C}'_n, f'_\varphi, u'_\varphi, \mathcal{D}', a'_0$, et
$h_n$, $n \geq -1$, comme dans la
remarque \ref{remark-morphism-augmentation-simplicial-sites}.
On obtient alors un diagramme commutatif
$$
\xymatrix{
\Sh(\mathcal{C}_{total}) \ar[d]_a \ar[r]_{h_{total}} &
\Sh(\mathcal{C}'_{total}) \ar[d]^{a'} \\
\Sh(\mathcal{D}) \ar[r]^{h_{-1}} &
\Sh(\mathcal{D}')
}
$$
\end{lemma}

\begin{proof}
Le morphisme $h$ est défini dans le lemme \ref{lemma-morphism-simplicial-sites}.
Les morphismes $a$ et $a'$ sont définis dans le lemme \ref{lemma-augmentation-site}.
Il reste donc seulement à prouver la commutativité du diagramme.
Pour cela, on prouve que
$a^{-1} \circ h_{-1}^{-1} = h_{total}^{-1} \circ (a')^{-1}$.
D'après les diagrammes commutatifs du
lemme \ref{lemma-morphism-simplicial-sites}
et la description de $\Sh(\mathcal{C}_{total})$
et $\Sh(\mathcal{C}'_{total})$ en termes de composantes
donnée dans le lemme \ref{lemma-describe-sheaves-simplicial-site-site},
il suffit de montrer que
$$
\xymatrix{
\Sh(\mathcal{C}_n) \ar[d]_{a_n} \ar[r]_{h_n} &
\Sh(\mathcal{C}'_n) \ar[d]^{a'_n} \\
\Sh(\mathcal{D}) \ar[r]^{h_{-1}} &
\Sh(\mathcal{D}')
}
$$
est commutatif pour tout $n$. Cela résulte du cas $n = 0$
(qui est une hypothèse de la
remarque \ref{remark-morphism-augmentation-simplicial-sites}) ;
pour $n > 0$, on choisit $\varphi : [0] \to [n]$,
et la commutativité requise résulte alors du cas $n = 0$,
des relations $a_n = a_0 \circ f_\varphi$
et $a'_n = a'_0 \circ f'_\varphi$,
ainsi que des relations de commutation
$f'_\varphi \circ h_n = h_0 \circ f_\varphi$.
\end{proof}




\section{Sites simpliciaux annelés}
\label{section-simplicial-sites-modules}

\noindent
Munissons notre topos simplicial d'un faisceau d'anneaux.

\begin{lemma}
\label{lemma-restriction-module-to-components-site}
Dans la situation \ref{situation-simplicial-site}, soit $\mathcal{O}$
un faisceau d'anneaux sur $\mathcal{C}_{total}$.
Il existe un morphisme canonique de topos annelés
$g_n : (\Sh(\mathcal{C}_n), \mathcal{O}_n) \to
(\Sh(\mathcal{C}_{total}), \mathcal{O})$
qui coïncide, sur les topos sous-jacents, avec le morphisme $g_n$ du
lemme \ref{lemma-restriction-to-components-site}.
Le foncteur
$g_n^* : \textit{Mod}(\mathcal{O}) \to \textit{Mod}(\mathcal{O}_n)$
admet un adjoint à gauche $g_{n!}$.
Pour $\mathcal{G}$ dans $\textit{Mod}(\mathcal{O}_n)$, la
restriction de $g_{n!}\mathcal{G}$ à $\mathcal{C}_m$ est
$$
\bigoplus\nolimits_{\varphi : [n] \to [m]} f_\varphi^*\mathcal{G}
$$
où $f_\varphi : (\Sh(\mathcal{C}_m), \mathcal{O}_m) \to
(\Sh(\mathcal{C}_n), \mathcal{O}_n)$ est le morphisme de topos annelés
qui coïncide sur les topos avec le morphisme $f_\varphi$ défini précédemment et
qui utilise l'application
$\mathcal{O}(\varphi) : f_\varphi^{-1}\mathcal{O}_n \to \mathcal{O}_m$
sur les faisceaux d'anneaux.
\end{lemma}

\begin{proof}
D'après le lemme \ref{lemma-restriction-to-components-site}, on a
$g_n^{-1}\mathcal{O} = \mathcal{O}_n$, d'où notre
morphisme de topos annelés. D'après Modules sur les sites, lemme
\ref{sites-modules-lemma-lower-shriek-modules},
on obtient l'adjoint $g_{n!}$. Pour prouver la formule donnant $g_{n!}$,
on définit d'abord un faisceau de $\mathcal{O}$-modules $\mathcal{H}$
sur $\mathcal{C}_{total}$ dont la composante de degré $m$ est
le $\mathcal{O}_m$-module
$$
\mathcal{H}_m =
\bigoplus\nolimits_{\varphi : [n] \to [m]} f_\varphi^*\mathcal{G}
$$
Étant donnée une application $\psi : [m] \to [m']$, l'application
$\mathcal{H}(\psi) : f_\psi^{-1}\mathcal{H}_m \to \mathcal{H}_{m'}$
est donnée sur les composantes par
$$
f_\psi^{-1} f_\varphi^*\mathcal{G} \to
f_\psi^* f_\varphi^*\mathcal{G} \to
f_{\psi \circ \varphi}^*\mathcal{G}
$$
Puisque cette application $f_\psi^{-1}\mathcal{H}_m \to \mathcal{H}_{m'}$ est
semi-linéaire relativement à $\mathcal{O}(\psi) : f_\psi^{-1}\mathcal{O}_m \to \mathcal{O}_{m'}$,
elle définit bien un $\mathcal{O}$-module
(utiliser le lemme \ref{lemma-describe-sheaves-simplicial-site-site}).
On prouve alors directement que
$$
\Mor_{\mathcal{O}_n}(\mathcal{G}, \mathcal{F}_n) =
\Mor_{\mathcal{O}}(\mathcal{H}, \mathcal{F})
$$
en procédant comme dans la preuve du lemme \ref{lemma-restriction-to-components-site}.
Ainsi, $\mathcal{H} = g_{n!}\mathcal{G}$, comme voulu.
\end{proof}

\begin{lemma}
\label{lemma-restriction-injective-to-component-limp}
Dans la situation \ref{situation-simplicial-site},
soit $\mathcal{O}$ un faisceau d'anneaux sur $\mathcal{C}_{total}$.
Si $\mathcal{I}$ est injectif dans $\textit{Mod}(\mathcal{O})$, alors
$\mathcal{I}_n$ est un faisceau totalement acyclique sur $\mathcal{C}_n$.
\end{lemma}

\begin{proof}
Cela résulte de
Cohomologie sur les sites, lemme \ref{sites-cohomology-lemma-pullback-injective-limp},
appliqué au foncteur d'inclusion $\mathcal{C}_n \to \mathcal{C}_{total}$
et aux propriétés de celui-ci prouvées dans le lemme \ref{lemma-restriction-to-components-site}.
\end{proof}

\begin{lemma}
\label{lemma-exactness-g-shriek-modules}
Sous les hypothèses du
lemme \ref{lemma-restriction-module-to-components-site}, le foncteur
$g_{n!} : \textit{Mod}(\mathcal{O}_n) \to \textit{Mod}(\mathcal{O})$
est exact si les applications $f_\varphi^{-1}\mathcal{O}_n \to \mathcal{O}_m$
sont plates pour tout $\varphi : [n] \to [m]$.
\end{lemma}

\begin{proof}
Rappelons que $g_{n!}\mathcal{G}$ est le $\mathcal{O}$-module
dont la composante de degré $m$ est le $\mathcal{O}_m$-module
$$
\bigoplus\nolimits_{\varphi : [n] \to [m]} f_\varphi^*\mathcal{G}
$$
Ici, le morphisme de topos annelés
$f_\varphi : (\Sh(\mathcal{C}_m), \mathcal{O}_m) \to
(\Sh(\mathcal{C}_n), \mathcal{O}_n)$ utilise l'application
$f_\varphi^{-1}\mathcal{O}_n \to \mathcal{O}_m$ de l'
énoncé du lemme. Si ces applications sont plates, alors
$f_\varphi^*$ est exact
(Modules sur les sites, lemme \ref{sites-modules-lemma-flat-pullback-exact}).
Par définition du site $\mathcal{C}_{total}$, on voit que ces
foncteurs possèdent les propriétés d'exactitude voulues, et l'on conclut.
\end{proof}

\begin{lemma}
\label{lemma-restriction-injective-to-component-site-module}
Dans la situation \ref{situation-simplicial-site},
soit $\mathcal{O}$ un faisceau d'anneaux sur $\mathcal{C}_{total}$
tel que $f_\varphi^{-1}\mathcal{O}_n \to \mathcal{O}_m$
soit plat pour tout $\varphi : [n] \to [m]$.
Si $\mathcal{I}$ est injectif dans $\textit{Mod}(\mathcal{O})$, alors
$\mathcal{I}_n$ est injectif dans $\textit{Mod}(\mathcal{O}_n)$.
\end{lemma}

\begin{proof}
Cela résulte du
lemme \ref{homology-lemma-adjoint-preserve-injectives} de Homologie
et du
lemme \ref{lemma-exactness-g-shriek-modules}.
\end{proof}







\section{Morphismes de sites simpliciaux annelés}
\label{section-morphism-simplicial-sites-modules}

\noindent
Nous poursuivons la discussion de la section \ref{section-morphism-simplicial-sites}.

\begin{remark}
\label{remark-morphism-simplicial-sites-modules}
Soient $\mathcal{C}_n, f_\varphi, u_\varphi$ et
$\mathcal{C}'_n, f'_\varphi, u'_\varphi$ comme dans la
situation \ref{situation-simplicial-site}.
Soient $\mathcal{O}$ et $\mathcal{O}'$
des faisceaux d'anneaux sur $\mathcal{C}_{total}$ et $\mathcal{C}'_{total}$, respectivement.
On dira que $(h, h^\sharp)$ est un
{\it morphisme entre sites simpliciaux annelés}
si $h$ est un morphisme entre sites simpliciaux comme dans la
remarque \ref{remark-morphism-simplicial-sites}
et si $h^\sharp : h_{total}^{-1}\mathcal{O}' \to \mathcal{O}$,
ou, de manière équivalente, $h^\sharp : \mathcal{O}' \to h_{total, *}\mathcal{O}$,
est un homomorphisme de faisceaux d'anneaux.
\end{remark}

\begin{lemma}
\label{lemma-morphism-simplicial-sites-modules}
Soient $\mathcal{C}_n, f_\varphi, u_\varphi$ et
$\mathcal{C}'_n, f'_\varphi, u'_\varphi$ comme dans la
situation \ref{situation-simplicial-site}.
Soient $\mathcal{O}$ et $\mathcal{O}'$
des faisceaux d'anneaux sur $\mathcal{C}_{total}$ et $\mathcal{C}'_{total}$, respectivement.
Soit $(h, h^\sharp)$ un morphisme entre sites simpliciaux annelés comme dans la
remarque \ref{remark-morphism-simplicial-sites-modules}.
Alors on obtient un morphisme de topos annelés
$$
h_{total} :
(\Sh(\mathcal{C}_{total}), \mathcal{O})
\to
(\Sh(\mathcal{C}'_{total}), \mathcal{O}')
$$
et des diagrammes commutatifs
$$
\xymatrix{
(\Sh(\mathcal{C}_n), \mathcal{O}_n) \ar[d]_{g_n} \ar[r]_{h_n} &
(\Sh(\mathcal{C}'_n), \mathcal{O}'_n) \ar[d]^{g'_n} \\
(\Sh(\mathcal{C}_{total}), \mathcal{O}) \ar[r]^{h_{total}} &
(\Sh(\mathcal{C}'_{total}), \mathcal{O}')
}
$$
de topos annelés, où $g_n$ et $g'_n$ sont comme dans le
lemme \ref{lemma-restriction-module-to-components-site}.
De plus, on a
$(g'_n)^* \circ h_{total, *} = h_{n, *} \circ g_n^*$
comme foncteurs $\textit{Mod}(\mathcal{O}) \to \textit{Mod}(\mathcal{O}'_n)$.
\end{lemma}

\begin{proof}
Cela résulte des
lemmes \ref{lemma-morphism-simplicial-sites} et
\ref{lemma-restriction-module-to-components-site},
en tenant compte des faisceaux d'anneaux.
Un petit point est que, pour définir $h_n$ comme morphisme
de topos annelés, on pose
$h_n^\sharp = g_n^{-1}h^\sharp :
g_n^{-1}h_{total}^{-1}\mathcal{O}' \to g_n^{-1}\mathcal{O}$,
ce qui a un sens puisque
$g_n^{-1}h_{total}^{-1}\mathcal{O}' = h_n^{-1}(g'_n)^{-1}\mathcal{O}' =
h_n^{-1}\mathcal{O}'_n$ et $g_n^{-1}\mathcal{O} = \mathcal{O}_n$.
Notons que $g_n^*\mathcal{F} = g_n^{-1}\mathcal{F}$
pour un faisceau de $\mathcal{O}$-modules $\mathcal{F}$,
et de même pour $g'_n$ ; cela aide à comprendre pourquoi
$(g'_n)^* \circ h_{total, *} = h_{n, *} \circ g_n^*$
résulte de l'énoncé correspondant du
lemme \ref{lemma-morphism-simplicial-sites}.
\end{proof}

\begin{lemma}
\label{lemma-direct-image-morphism-simplicial-sites-modules}
Avec les notations et les hypothèses du
lemme \ref{lemma-morphism-simplicial-sites-modules},
pour $K \in D(\mathcal{O})$, on a
$(g'_n)^*Rh_{total, *}K = Rh_{n, *}g_n^*K$.
\end{lemma}

\begin{proof}
Rappelons que $g_n^* = g_n^{-1}$ parce que $g_n^{-1}\mathcal{O} = \mathcal{O}_n$
par la construction du lemme \ref{lemma-restriction-module-to-components-site}.
En particulier, $g_n^*$ est exact, et $Lg_n^*$ s'obtient en appliquant $g_n^*$
à n'importe quel complexe de modules représentant l'objet considéré. De même pour $g'_n$.
Il existe un morphisme canonique de changement de base
$(g'_n)^*Rh_{total, *}K \to Rh_{n, *}g_n^*K$ ; voir
Cohomologie sur les sites, remarque \ref{sites-cohomology-remark-base-change}.
D'après Cohomologie sur les sites, lemme
\ref{sites-cohomology-lemma-modules-abelian-unbounded},
l'image de ce morphisme dans $D(\mathcal{C}'_n)$ est le morphisme
$(g'_n)^{-1}Rh_{total, *}K_{ab} \to Rh_{n, *}g_n^{-1}K_{ab}$,
où $K_{ab}$ est l'image de $K$ dans $D(\mathcal{C}_{total})$.
On a prouvé que celui-ci est un isomorphisme dans le
lemme \ref{lemma-direct-image-morphism-simplicial-sites},
et le résultat en découle.
\end{proof}








\section{Cohomologie sur les sites simpliciaux}
\label{section-cohomology-simplicial-sites}

\noindent
Soit $\mathcal{C}$ comme dans la situation \ref{situation-simplicial-site}.
Dans l'énoncé des lemmes suivants, on notera
$g_n : \Sh(\mathcal{C}_n) \to \Sh(\mathcal{C}_{total})$ le
morphisme de topos du
lemme \ref{lemma-restriction-to-components-site}. Si $\varphi : [m] \to [n]$
est un morphisme de $\Delta$, alors le diagramme de topos
$$
\xymatrix{
\Sh(\mathcal{C}_n) \ar[rd]_{g_n} \ar[rr]_{f_\varphi} & &
\Sh(\mathcal{C}_m) \ar[ld]^{g_m} \\
& \Sh(\mathcal{C}_{total})
}
$$
n'est pas commutatif, mais il existe un $2$-morphisme $g_n \to g_m \circ f_\varphi$
provenant des applications
$\mathcal{F}(\varphi) : f_\varphi^{-1}\mathcal{F}_m \to \mathcal{F}_n$.
Voir Sites, section \ref{sites-section-2-category}.

\begin{lemma}
\label{lemma-simplicial-resolution-Z-site}
Dans la situation \ref{situation-simplicial-site} et avec les notations ci-dessus,
il existe un complexe
$$
\ldots \to g_{2!}\mathbf{Z} \to g_{1!}\mathbf{Z} \to g_{0!}\mathbf{Z}
$$
de faisceaux abéliens sur $\mathcal{C}_{total}$ qui constitue une résolution du
faisceau constant de valeur $\mathbf{Z}$ sur $\mathcal{C}_{total}$.
\end{lemma}

\begin{proof}
On utilisera sans autre mention la description des foncteurs $g_{n!}$ donnée dans le
lemme \ref{lemma-restriction-to-components-site}. Comme applications du complexe,
on prend $\sum (-1)^i d^n_i$, où
$d^n_i : g_{n!}\mathbf{Z} \to g_{n - 1!}\mathbf{Z}$ est l'
adjoint de l'application $\mathbf{Z} \to
\bigoplus_{[n - 1] \to [n]} \mathbf{Z} = g_n^{-1}g_{n - 1!}\mathbf{Z}$
correspondant au facteur étiqueté par $\delta^n_i : [n - 1] \to [n]$.
En appliquant $g_m^{-1}$ au complexe, on obtient le complexe
$$
\ldots \to
\bigoplus\nolimits_{\alpha \in \Mor_\Delta([2], [m])]} \mathbf{Z} \to
\bigoplus\nolimits_{\alpha \in \Mor_\Delta([1], [m])]} \mathbf{Z} \to
\bigoplus\nolimits_{\alpha \in \Mor_\Delta([0], [m])]} \mathbf{Z}
$$
sur $\mathcal{C}_m$.
Autrement dit, c'est le complexe associé au
faisceau abélien libre sur l'ensemble simplicial $\Delta[m]$ ; voir
Méthodes simpliciales, exemple \ref{simplicial-example-simplex-simplicial-set}.
Puisque $\Delta[m]$ est homotopiquement équivalent à $\Delta[0]$ ; voir
Méthodes simpliciales, exemple \ref{simplicial-example-simplex-contractible},
et puisque « prendre le faisceau abélien libre sur » est un foncteur,
on voit que le complexe ci-dessus est homotopiquement équivalent au
faisceau abélien libre sur $\Delta[0]$
(Méthodes simpliciales, remarque \ref{simplicial-remark-homotopy-better} et
lemme \ref{simplicial-lemma-homotopy-equivalence-s-N}).
Ce complexe est acyclique en degrés strictement positifs
et égal à $\mathbf{Z}$ en degré $0$.
\end{proof}

\begin{lemma}
\label{lemma-cech-complex}
Dans la situation \ref{situation-simplicial-site}, soit $\mathcal{F}$ un faisceau
abélien sur $\mathcal{C}_{total}$. Il existe un complexe canonique
$$
0 \to \Gamma(\mathcal{C}_{total}, \mathcal{F}) \to
\Gamma(\mathcal{C}_0, \mathcal{F}_0) \to
\Gamma(\mathcal{C}_1, \mathcal{F}_1) \to
\Gamma(\mathcal{C}_2, \mathcal{F}_2) \to \ldots
$$
qui est exact en degrés $-1, 0$ et partout exact
si $\mathcal{F}$ est injectif.
\end{lemma}

\begin{proof}
Observons que
$\Hom(\mathbf{Z}, \mathcal{F}) = \Gamma(\mathcal{C}_{total}, \mathcal{F})$
et que
$\Hom(g_{n!}\mathbf{Z}, \mathcal{F}) = \Gamma(\mathcal{C}_n, \mathcal{F}_n)$.
Ce lemme résulte donc immédiatement du
lemme \ref{lemma-simplicial-resolution-Z-site}
et du fait que $\Hom(-, \mathcal{F})$ est exact si
$\mathcal{F}$ est injectif.
\end{proof}

\begin{lemma}
\label{lemma-simplicial-sheaf-cohomology-site}
Dans la situation \ref{situation-simplicial-site}, pour $K$ dans
$D^+(\mathcal{C}_{total})$, il existe une suite spectrale
$(E_r, d_r)_{r \geq 0}$ avec
$$
E_1^{p, q} = H^q(\mathcal{C}_p, K_p),\quad
d_1^{p, q} : E_1^{p, q} \to E_1^{p + 1, q}
$$
qui aboutit à $H^{p + q}(\mathcal{C}_{total}, K)$.
Cette suite spectrale est fonctorielle en $K$.
\end{lemma}

\begin{proof}
Soit $\mathcal{I}^\bullet$ un complexe borné inférieurement d'injectifs
représentant $K$. Considérons le complexe double dont les termes sont
$$
A^{p, q} = \Gamma(\mathcal{C}_p, \mathcal{I}^q_p)
$$
où les flèches horizontales proviennent du lemme \ref{lemma-cech-complex}
et les flèches verticales des différentielles du
complexe $\mathcal{I}^\bullet$. Les lignes du complexe double sont exactes
en degrés strictement positifs et valent
$\Gamma(\mathcal{C}_{total}, \mathcal{I}^q)$ en degré $0$.
D'autre part, puisque la restriction à $\mathcal{C}_p$ est exacte
(lemme \ref{lemma-restriction-to-components-site}),
le complexe $\mathcal{I}_p^\bullet$ représente $K_p$ dans
$D(\mathcal{C}_p)$. Les faisceaux $\mathcal{I}_p^q$ sont des
faisceaux abéliens injectifs sur $\mathcal{C}_p$
(lemme \ref{lemma-restriction-injective-to-component-site}).
La cohomologie des colonnes calcule donc les groupes
$H^q(\mathcal{C}_p, K_p)$. On conclut en appliquant les
lemmes \ref{homology-lemma-first-quadrant-ss} et
\ref{homology-lemma-double-complex-gives-resolution} de Homologie.
\end{proof}

\begin{remark}
\label{remark-simplicial-cohomology-site}
On fait les mêmes hypothèses et on emploie les mêmes notations que dans le
lemme \ref{lemma-simplicial-sheaf-cohomology-site},
sauf que l'on ne suppose pas que $K$ dans $D(\mathcal{C}_{total})$
soit borné inférieurement. On affirme qu'il existe également dans ce cas une
suite spectrale naturelle. En effet, supposons que
$\mathcal{I}^\bullet$ soit un complexe K-injectif de faisceaux sur
$\mathcal{C}_{total}$, à termes injectifs, qui représente $K$.
On a
\begin{align*}
R\Gamma(\mathcal{C}_{total}, K)
& =
R\Hom(\mathbf{Z}, K) \\
& =
R\Hom(
\ldots \to g_{2!}\mathbf{Z} \to g_{1!}\mathbf{Z} \to g_{0!}\mathbf{Z}, K) \\
& =
\Gamma(\mathcal{C}_{total},
\SheafHom^\bullet(
\ldots \to g_{2!}\mathbf{Z} \to g_{1!}\mathbf{Z} \to g_{0!}\mathbf{Z},
\mathcal{I}^\bullet)) \\
& =
\text{Tot}_\pi(A^{\bullet, \bullet})
\end{align*}
où $A^{\bullet, \bullet}$ est le complexe double dont les termes sont
$A^{p, q} = \Gamma(\mathcal{C}_p, \mathcal{I}^q_p)$, et où $\text{Tot}_\pi$
désigne la totalisation par produits de ce complexe double.
En effet, la première égalité vaut dans tout site. La deuxième égalité
résulte du lemme \ref{lemma-simplicial-resolution-Z-site}.
La troisième égalité vaut parce que $\mathcal{I}^\bullet$ est K-injectif ; voir
Cohomologie sur les sites, sections \ref{sites-cohomology-section-hom-complexes}
et \ref{sites-cohomology-section-internal-hom}.
La dernière égalité résulte de la construction de $\SheafHom^\bullet$
et du fait que $\Hom(g_{p!}\mathbf{Z}, \mathcal{I}^q) =
\Gamma(\mathcal{C}_p, \mathcal{I}^q_p)$.
On obtient alors notre suite spectrale en considérant
$\text{Tot}_\pi(A^{\bullet, \bullet})$ comme un
complexe filtré avec $F^i\text{Tot}^n_\pi(A^{\bullet, \bullet}) =
\prod_{p + q = n,\ p \geq i} A^{p, q}$. La suite spectrale
ainsi obtenue se comporte comme la suite spectrale $({}'E_r, {}'d_r)_{r \geq 0}$ de
Homologie, section \ref{homology-section-double-complex}
(où est traité le cas de la totalisation par sommes directes),
si ce n'est pour les questions de régularité, de bornitude, de convergence et
d'aboutissement. En particulier, on obtient $E_1^{p, q} = H^q(\mathcal{C}_p, K_p)$
comme dans le lemme \ref{lemma-simplicial-sheaf-cohomology-site}.
\end{remark}

\begin{lemma}
\label{lemma-simplicial-sheaf-cohomology-site-zero}
Dans la situation \ref{situation-simplicial-site}, soit $K$ un objet de
$D(\mathcal{C}_{total})$.
\begin{enumerate}
\item Si $H^{-p}(\mathcal{C}_p, K_p) = 0$ pour tout $p \geq 0$, alors
$H^0(\mathcal{C}_{total}, K) = 0$.
\item Si $R\Gamma(\mathcal{C}_p, K_p) = 0$
pour tout $p \geq 0$, alors $R\Gamma(\mathcal{C}_{total}, K) = 0$.
\end{enumerate}
\end{lemma}

\begin{proof}
Avec les notations de la remarque \ref{remark-simplicial-cohomology-site},
on voit que $R\Gamma(\mathcal{C}_{total}, K)$ est
représenté par $\text{Tot}_\pi(A^{\bullet, \bullet})$.
L'hypothèse de (1) nous dit que $H^{-p}(A^{p, \bullet}) = 0$.
L'annulation de (1) résulte donc du
lemme \ref{more-algebra-lemma-vanishing-coh-prod-totalization} de Compléments d'algèbre.
L'assertion (2) résulte de l'assertion (1) en prenant des décalés.
\end{proof}

\begin{lemma}
\label{lemma-sanity-check}
Soit $\mathcal{C}$ comme dans la situation \ref{situation-simplicial-site}.
Soit $U \in \Ob(\mathcal{C}_n)$. Soit
$\mathcal{F} \in \textit{Ab}(\mathcal{C}_{total})$.
Alors $H^p(U, \mathcal{F}) = H^p(U, g_n^{-1}\mathcal{F})$,
où, dans le membre de gauche, $U$ est considéré comme un objet de $\mathcal{C}_{total}$.
\end{lemma}

\begin{proof}
Notons que l'expression « $U$ considéré comme objet de $\mathcal{C}_{total}$ »
est expliquée par la construction de $\mathcal{C}_{total}$ dans le
lemme \ref{lemma-simplicial-site-site} dans le cas (A) et le
lemme \ref{lemma-simplicial-cocontinuous-site} dans le cas (B).
L'égalité résulte alors du
lemme \ref{lemma-restriction-injective-to-component-site}
et de la définition de la cohomologie.
\end{proof}






\section{Cohomologie et augmentations de sites simpliciaux}
\label{section-cohomology-augmentation-simplicial-sites}

\noindent
Considérons un site simplicial $\mathcal{C}$ comme dans la situation
\ref{situation-simplicial-site}.
Soit $a_0$ une augmentation vers un site $\mathcal{D}$ comme dans la
remarque \ref{remark-augmentation-site}.
Par le lemme \ref{lemma-augmentation-site}, on obtient un morphisme de topoi
$$
a : \Sh(\mathcal{C}_{total}) \longrightarrow \Sh(\mathcal{D})
$$
ainsi que des morphismes de topoi
$g_n : \Sh(\mathcal{C}_n) \to \Sh(\mathcal{C}_{total})$
comme dans le lemme \ref{lemma-restriction-to-components-site}.
Les compositions $a \circ g_n$ sont notées
$a_n : \Sh(\mathcal{C}_n) \to \Sh(\mathcal{D})$.
De plus, la structure simpliciale donne des morphismes
de topoi
$f_\varphi : \Sh(\mathcal{C}_n) \to \Sh(\mathcal{C}_m)$
tels que $a_n \circ f_\varphi = a_m$ pour tout $\varphi : [m] \to [n]$.

\begin{lemma}
\label{lemma-simplicial-resolution-augmentation}
Dans la situation \ref{situation-simplicial-site}, soit
$a_0$ une augmentation vers un site $\mathcal{D}$
comme dans la remarque \ref{remark-augmentation-site}.
Pour tout faisceau abélien $\mathcal{G}$ sur $\mathcal{D}$,
il existe un complexe exact
$$
\ldots \to
g_{2!}(a_2^{-1}\mathcal{G}) \to
g_{1!}(a_1^{-1}\mathcal{G}) \to
g_{0!}(a_0^{-1}\mathcal{G}) \to
a^{-1}\mathcal{G} \to 0
$$
de faisceaux abéliens sur $\mathcal{C}_{total}$.
\end{lemma}

\begin{proof}
Nous conseillons au lecteur de lire d'abord la preuve du
lemme \ref{lemma-simplicial-resolution-Z-site}.
Nous utiliserons le lemme \ref{lemma-augmentation-site} et
la description des foncteurs $g_{n!}$ donnée dans le
lemme \ref{lemma-restriction-to-components-site}, sans autre mention.
En particulier, $g_{n!}(a_n^{-1}\mathcal{G})$ est le
faisceau sur $\mathcal{C}_{total}$ dont la restriction à $\mathcal{C}_m$
est le faisceau
$$
\bigoplus\nolimits_{\varphi : [n] \to [m]} f_\varphi^{-1}a_n^{-1}\mathcal{G} =
\bigoplus\nolimits_{\varphi : [n] \to [m]} a_m^{-1}\mathcal{G}
$$
Comme morphismes du complexe, prenons $\sum (-1)^i d^n_i$, où
$d^n_i : g_{n!}(a_n^{-1}\mathcal{G}) \to g_{n - 1!}(a_{n - 1}^{-1}\mathcal{G})$
est adjoint au morphisme
$a_n^{-1}\mathcal{G} \to \bigoplus_{[n - 1] \to [n]} a_n^{-1}\mathcal{G} =
g_n^{-1}g_{n - 1!}(a_{n - 1}^{-1}\mathcal{G})$
correspondant au facteur étiqueté $\delta^n_i : [n - 1] \to [n]$.
Le morphisme $g_{0!}(a_0^{-1}\mathcal{G}) \to a^{-1}\mathcal{G}$ est adjoint
au morphisme identité de $a_0^{-1}\mathcal{G}$.
L'application de $g_m^{-1}$ au complexe de chaînes en degrés
$\ldots, 2, 1, 0$ donne le complexe
$$
\ldots \to
\bigoplus\nolimits_{\alpha \in \Mor_\Delta([2], [m])]} a_m^{-1}\mathcal{G} \to
\bigoplus\nolimits_{\alpha \in \Mor_\Delta([1], [m])]} a_m^{-1}\mathcal{G} \to
\bigoplus\nolimits_{\alpha \in \Mor_\Delta([0], [m])]} a_m^{-1}\mathcal{G}
$$
sur $\mathcal{C}_m$. Il est égal à $a_m^{-1}\mathcal{G}$
tensorisé sur le faisceau constant $\mathbf{Z}$ avec le complexe
$$
\ldots \to
\bigoplus\nolimits_{\alpha \in \Mor_\Delta([2], [m])]} \mathbf{Z} \to
\bigoplus\nolimits_{\alpha \in \Mor_\Delta([1], [m])]} \mathbf{Z} \to
\bigoplus\nolimits_{\alpha \in \Mor_\Delta([0], [m])]} \mathbf{Z}
$$
discuté dans la preuve du lemme \ref{lemma-simplicial-resolution-Z-site}.
Nous y avons vu que ce complexe est homotopiquement équivalent à
$\mathbf{Z}$ placé en degré $0$, ce qui termine la preuve.
\end{proof}

\begin{lemma}
\label{lemma-augmentation-cech-complex}
Dans la situation \ref{situation-simplicial-site}, soit
$a_0$ une augmentation vers un site $\mathcal{D}$
comme dans la remarque \ref{remark-augmentation-site}.
Pour un faisceau abélien $\mathcal{F}$ sur $\mathcal{C}_{total}$
il existe un complexe canonique
$$
0 \to a_*\mathcal{F} \to a_{0, *}\mathcal{F}_0 \to a_{1, *}\mathcal{F}_1 \to
a_{2, *}\mathcal{F}_2 \to \ldots
$$
sur $\mathcal{D}$ qui est exact en degrés $-1, 0$ et
exact partout si $\mathcal{F}$ est injectif.
\end{lemma}

\begin{proof}
Pour construire le complexe, il suffit, par le lemme de Yoneda, de construire pour tout
faisceau abélien $\mathcal{G}$ sur $\mathcal{D}$ un complexe
$$
0 \to \Hom(\mathcal{G}, a_*\mathcal{F}) \to
\Hom(\mathcal{G}, a_{0, *}\mathcal{F}_0) \to
\Hom(\mathcal{G}, a_{1, *}\mathcal{F}_1) \to \ldots
$$
fonctoriellement en $\mathcal{G}$. Pour cela, appliquons $\Hom(-, \mathcal{F})$
au complexe exact du lemme \ref{lemma-simplicial-resolution-augmentation}
et utilisons l'adjonction entre image inverse et image directe.
Les propriétés d'exactitude en degrés $-1, 0$ découlent de
la construction, puisque $\Hom(-, \mathcal{F})$ est exact à gauche.
Si $\mathcal{F}$ est un faisceau abélien injectif, alors le complexe
est exact car $\Hom(-, \mathcal{F})$ est exact.
\end{proof}

\begin{lemma}
\label{lemma-augmentation-spectral-sequence}
Dans la situation \ref{situation-simplicial-site}, soit
$a_0$ une augmentation vers un site $\mathcal{D}$
comme dans la remarque \ref{remark-augmentation-site}.
Pour tout $K$ dans $D^+(\mathcal{C}_{total})$, il existe une suite
spectrale
$(E_r, d_r)_{r \geq 0}$ avec
$$
E_1^{p, q} = R^qa_{p, *} K_p,\quad
d_1^{p, q} : E_1^{p, q} \to E_1^{p + 1, q}
$$
qui aboutit à $R^{p + q}a_*K$. Cette suite spectrale est fonctorielle en $K$.
\end{lemma}

\begin{proof}
Soit $\mathcal{I}^\bullet$ un complexe borné inférieurement d'injectifs
représentant $K$. Considérons le complexe double dont les termes sont
$$
A^{p, q} = a_{p, *}\mathcal{I}^q_p
$$
où les flèches horizontales proviennent du
lemme \ref{lemma-augmentation-cech-complex}
et les flèches verticales des différentiels du complexe
$\mathcal{I}^\bullet$. Les lignes du double complexe sont exactes
en degrés positifs et s'évaluent à $a_*\mathcal{I}^q$ en degré $0$.
D'autre part, puisque la restriction à $\mathcal{C}_p$ est exacte
(lemme \ref{lemma-restriction-to-components-site})
le complexe $\mathcal{I}_p^\bullet$ représente $K_p$ dans
$D(\mathcal{C}_p)$. Les faisceaux $\mathcal{I}_p^q$ sont des faisceaux abéliens injectifs
sur $\mathcal{C}_p$
(lemme \ref{lemma-restriction-injective-to-component-site}).
La cohomologie des colonnes calcule donc $R^qa_{p, *}K_p$.
On conclut en appliquant les
lemmes \ref{homology-lemma-first-quadrant-ss} et
\ref{homology-lemma-double-complex-gives-resolution}.
\end{proof}





\section{Cohomologie sur les sites simpliciaux annelés}
\label{section-cohomology-simplicial-sites-modules}

\noindent
Cette section est l'analogue de
la section \ref{section-cohomology-simplicial-sites}
pour les faisceaux de modules.

\medskip\noindent
Dans la situation \ref{situation-simplicial-site}, soit $\mathcal{O}$
un faisceau d'anneaux sur $\mathcal{C}_{total}$.
Dans les énoncés des lemmes suivants, nous noterons
$g_n : (\Sh(\mathcal{C}_n), \mathcal{O}_n) \to
(\Sh(\mathcal{C}_{total}), \mathcal{O})$
le morphisme de topos annelés du lemme
\ref{lemma-restriction-module-to-components-site}.
Si $\varphi : [m] \to [n]$ est un morphisme de $\Delta$, alors le diagramme
de topoi annelés
$$
\xymatrix{
(\Sh(\mathcal{C}_n), \mathcal{O}_n) \ar[rd]_{g_n} \ar[rr]_{f_\varphi} & &
(\Sh(\mathcal{C}_m), \mathcal{O}_m) \ar[ld]^{g_m} \\
& (\Sh(\mathcal{C}_{total}), \mathcal{O})
}
$$
n'est pas commutatif, mais il existe un $2$-morphisme $g_n \to g_m \circ f_\varphi$
provenant des morphismes
$\mathcal{F}(\varphi) : f_\varphi^{-1}\mathcal{F}_m \to \mathcal{F}_n$.
Voir Sites, section \ref{sites-section-2-category}.

\begin{lemma}
\label{lemma-simplicial-resolution-ringed}
Dans la situation \ref{situation-simplicial-site}, soit $\mathcal{O}$
un faisceau d'anneaux sur $\mathcal{C}_{total}$. Il existe un complexe
$$
\ldots \to g_{2!}\mathcal{O}_2 \to g_{1!}\mathcal{O}_1 \to g_{0!}\mathcal{O}_0
$$
de $\mathcal{O}$-modules qui est une résolution de
$\mathcal{O}$.
Ici, $g_{n!}$ est comme dans le lemme \ref{lemma-restriction-module-to-components-site}.
\end{lemma}

\begin{proof}
Nous utiliserons la description de $g_{n!}$ donnée dans le
lemme \ref{lemma-restriction-to-components-site}.
Comme morphismes du complexe, prenons $\sum (-1)^i d^n_i$, où
$d^n_i : g_{n!}\mathcal{O}_n \to g_{n - 1!}\mathcal{O}_{n - 1}$
est adjoint au morphisme
$\mathcal{O}_n \to \bigoplus_{[n - 1] \to [n]} \mathcal{O}_n =
g_n^*g_{n - 1!}\mathcal{O}_{n - 1}$
correspondant au facteur étiqueté $\delta^n_i : [n - 1] \to [n]$.
Alors $g_m^{-1}$ appliqué au complexe donne le complexe
$$
\ldots \to
\bigoplus\nolimits_{\alpha \in \Mor_\Delta([2], [m])]} \mathcal{O}_m \to
\bigoplus\nolimits_{\alpha \in \Mor_\Delta([1], [m])]} \mathcal{O}_m \to
\bigoplus\nolimits_{\alpha \in \Mor_\Delta([0], [m])]} \mathcal{O}_m
$$
sur $\mathcal{C}_m$.
Autrement dit, il s'agit du complexe associé au
$\mathcal{O}_m$-module libre engendré par l'ensemble simplicial $\Delta[m]$; voir
Simplicial, exemple \ref{simplicial-example-simplex-simplicial-set}.
Puisque $\Delta[m]$ est homotopiquement équivalent à $\Delta[0]$, voir
Simplicial, exemple \ref{simplicial-example-simplex-contractible},
et puisque « prendre le faisceau abélien libre engendré par » est un foncteur,
on voit que le complexe ci-dessus est homotopiquement équivalent au
faisceau abélien libre engendré par $\Delta[0]$
(Simplicial, remarque \ref{simplicial-remark-homotopy-better} et
lemme \ref{simplicial-lemma-homotopy-equivalence-s-N}).
Ce complexe est acyclique en degrés positifs
et égal à $\mathcal{O}_m$ en degré $0$.
\end{proof}

\begin{lemma}
\label{lemma-cech-complex-modules}
Dans la situation \ref{situation-simplicial-site}, soit $\mathcal{O}$
un faisceau d'anneaux. Soit $\mathcal{F}$ un faisceau
de $\mathcal{O}$-modules. Il existe un complexe canonique
$$
0 \to \Gamma(\mathcal{C}_{total}, \mathcal{F}) \to
\Gamma(\mathcal{C}_0, \mathcal{F}_0) \to
\Gamma(\mathcal{C}_1, \mathcal{F}_1) \to
\Gamma(\mathcal{C}_2, \mathcal{F}_2) \to \ldots
$$
qui est exact en degrés $-1, 0$ et exact partout
si $\mathcal{F}$ est un $\mathcal{O}$-module injectif.
\end{lemma}

\begin{proof}
Observons que
$\Hom(\mathcal{O}, \mathcal{F}) = \Gamma(\mathcal{C}_{total}, \mathcal{F})$
et
$\Hom(g_{n!}\mathcal{O}_n, \mathcal{F}) = \Gamma(\mathcal{C}_n, \mathcal{F}_n)$.
Ce lemme est donc une conséquence immédiate du lemme
\ref{lemma-simplicial-resolution-ringed}
et du fait que $\Hom(-, \mathcal{F})$ est exact si
$\mathcal{F}$ est injectif.
\end{proof}

\begin{lemma}
\label{lemma-simplicial-module-cohomology-site}
Dans la situation \ref{situation-simplicial-site}, soit $\mathcal{O}$
un faisceau d'anneaux. Pour $K$ dans $D^+(\mathcal{O})$
il existe une suite spectrale $(E_r, d_r)_{r \geq 0}$ telle que
$$
E_1^{p, q} = H^q(\mathcal{C}_p, K_p),\quad
d_1^{p, q} : E_1^{p, q} \to E_1^{p + 1, q}
$$
qui aboutit à $H^{p + q}(\mathcal{C}_{total}, K)$.
Cette suite spectrale est fonctorielle en $K$.
\end{lemma}

\begin{proof}
Soit $\mathcal{I}^\bullet$ un complexe borné inférieurement de
$\mathcal{O}$-modules injectifs représentant $K$. Considérons le complexe double de termes
$$
A^{p, q} = \Gamma(\mathcal{C}_p, \mathcal{I}^q_p)
$$
où les flèches horizontales proviennent du
lemme \ref{lemma-cech-complex-modules}
et les flèches verticales des différentiels du complexe
$\mathcal{I}^\bullet$. Notons que
$\Gamma(\mathcal{D}, -) =
\Hom_{\mathcal{O}_\mathcal{D}}(\mathcal{O}_\mathcal{D}, -)$
sur $\textit{Mod}(\mathcal{O}_\mathcal{D})$. Le lemme
affirme donc que les lignes du complexe double sont exactes
en degrés positifs et donnent
$\Gamma(\mathcal{C}_{total}, \mathcal{I}^q)$ en degré $0$.
Ainsi le complexe total associé au double complexe
calcule $R\Gamma(\mathcal{C}_{total}, K)$ d'après
Homologie, lemme \ref{homology-lemma-double-complex-gives-resolution}.
D'un autre côté, puisque la restriction à $\mathcal{C}_p$ est exacte,
(lemme \ref{lemma-restriction-to-components-site})
le complexe $\mathcal{I}_p^\bullet$ représente $K_p$ dans
$D(\mathcal{C}_p)$. Les faisceaux $\mathcal{I}_p^q$
sont totalement acycliques sur $\mathcal{C}_p$
(lemme \ref{lemma-restriction-injective-to-component-limp}).
La cohomologie des colonnes calcule donc les groupes
$H^q(\mathcal{C}_p, K_p)$ par le lemme d'acyclicité de Leray
(Catégories dérivées, lemme \ref{derived-lemma-leray-acyclicity})
et
Cohomologie sur les sites, lemme \ref{sites-cohomology-lemma-limp-acyclic}.
On conclut en appliquant
Homologie, lemme \ref{homology-lemma-first-quadrant-ss}.
\end{proof}

\begin{lemma}
\label{lemma-sanity-check-modules}
Dans la situation \ref{situation-simplicial-site}, soit $\mathcal{O}$
un faisceau d'anneaux. Soient $U \in \Ob(\mathcal{C}_n)$ et
$\mathcal{F} \in \textit{Mod}(\mathcal{O})$.
Alors $H^p(U, \mathcal{F}) = H^p(U, g_n^*\mathcal{F})$
où, au membre de gauche, $U$ est considéré comme un objet de
$\mathcal{C}_{total}$.
\end{lemma}

\begin{proof}
Observons que « $U$ considéré comme objet de $\mathcal{C}_{total}$ »
est expliqué par la construction de $\mathcal{C}_{total}$ dans le
lemme \ref{lemma-simplicial-site-site} dans le cas (A), et dans le
lemme \ref{lemma-simplicial-cocontinuous-site} dans le cas (B).
Dans les deux cas, le foncteur $\mathcal{C}_n \to \mathcal{C}$
est continu et cocontinu, voir
le lemme \ref{lemma-restriction-to-components-site}, et
$g_n^{-1}\mathcal{O} = \mathcal{O}_n$ par définition.
Le résultat est donc un cas particulier du lemme
de Cohomologie sur les sites
\ref{sites-cohomology-lemma-pullback-same-cohomology}.
\end{proof}






\section{Cohomologie et augmentations des sites simpliciaux annelés}
\label{section-cohomology-augmentation-ringed-simplicial-sites}

\noindent
Cette section est l'analogue de
la section \ref{section-cohomology-augmentation-simplicial-sites}
pour les faisceaux de modules.

\medskip\noindent
Considérons un site simplicial $\mathcal{C}$ comme dans
la situation \ref{situation-simplicial-site}.
Soit $a_0$ une augmentation vers un site $\mathcal{D}$ comme dans
la remarque \ref{remark-augmentation-site}.
Soit $\mathcal{O}$ un faisceau d'anneaux sur $\mathcal{C}_{total}$.
Soit $\mathcal{O}_\mathcal{D}$ un faisceau d'anneaux sur $\mathcal{D}$.
Supposons qu’on nous donne un morphisme
$$
a^\sharp : \mathcal{O}_\mathcal{D} \longrightarrow a_*\mathcal{O}
$$
où $a$ est comme dans le lemme \ref{lemma-augmentation-site}.
Par conséquent, on obtient un morphisme de topoi annelés
$$
a :
(\Sh(\mathcal{C}_{total}), \mathcal{O})
\longrightarrow
(\Sh(\mathcal{D}), \mathcal{O}_\mathcal{D})
$$
Considérons le morphisme de topos annelés $g_n : (\Sh(\mathcal{C}_n), \mathcal{O}_n) \to
(\Sh(\mathcal{C}_{total}), \mathcal{O})$
du
lemme \ref{lemma-restriction-module-to-components-site}. En formant
la composition $a_n = a \circ g_n$
(lemme \ref{lemma-augmentation-site})
dans la catégorie des topoi annelés, on obtient
$$
a_n :
(\Sh(\mathcal{C}_n), \mathcal{O}_n)
\longrightarrow
(\Sh(\mathcal{D}), \mathcal{O}_\mathcal{D})
$$
À l'aide des morphismes de transition $f_\varphi^{-1}\mathcal{O}_m \to \mathcal{O}_n$
on obtient des morphismes de topoi annelés
$$
f_\varphi : (\Sh(\mathcal{C}_n), \mathcal{O}_n) \to
(\Sh(\mathcal{C}_m), \mathcal{O}_m)
$$
tels que $a_m \circ f_\varphi = a_n$ comme morphismes de topoi annelés
pour tout $\varphi : [m] \to [n]$.

\begin{lemma}
\label{lemma-flat-augmentation-modules}
Avec les notations ci-dessus, le morphisme
$a : (\Sh(\mathcal{C}_{total}), \mathcal{O}) \to
(\Sh(\mathcal{D}), \mathcal{O}_\mathcal{D})$
est plat si et seulement si
$a_n : (\Sh(\mathcal{C}_n), \mathcal{O}_n) \to
(\Sh(\mathcal{D}), \mathcal{O}_\mathcal{D})$
est plat pour $n \geq 0$.
\end{lemma}

\begin{proof}
Puisque $g_n : (\Sh(\mathcal{C}_n), \mathcal{O}_n) \to
(\Sh(\mathcal{C}_{total}), \mathcal{O})$ est plat, on voit
que, si $a$ est plat, alors $a_n = a \circ g_n$ est plat comme
composé. Réciproquement, supposons $a_n$ plat pour tout $n$.
Il faut vérifier que $\mathcal{O}$ est plat comme faisceau de
$a^{-1}\mathcal{O}_\mathcal{D}$-modules. Soit $\mathcal{F} \to \mathcal{G}$
un morphisme injectif de $a^{-1}\mathcal{O}_\mathcal{D}$-modules.
Nous devons montrer que
$$
\mathcal{F} \otimes_{a^{-1}\mathcal{O}_\mathcal{D}} \mathcal{O}
\to
\mathcal{G} \otimes_{a^{-1}\mathcal{O}_\mathcal{D}} \mathcal{O}
$$
est injectif. On peut le vérifier sur $\mathcal{C}_n$, c'est-à-dire après
application de $g_n^{-1}$. Puisque $g_n^* = g_n^{-1}$ et que
$g_n^{-1}\mathcal{O} = \mathcal{O}_n$, on obtient
$$
g_n^{-1}\mathcal{F} \otimes_{g_n^{-1}a^{-1}\mathcal{O}_\mathcal{D}}
\mathcal{O}_n
\to
g_n^{-1}\mathcal{G} \otimes_{g_n^{-1}a^{-1}\mathcal{O}_\mathcal{D}}
\mathcal{O}_n
$$
qui est injectif car
$g_n^{-1}a^{-1}\mathcal{O}_\mathcal{D} = a_n^{-1}\mathcal{O}_\mathcal{D}$
et, par hypothèse, $a_n$ est plat.
\end{proof}

\begin{lemma}
\label{lemma-simplicial-resolution-augmentation-modules}
Avec les notations ci-dessus, pour tout $\mathcal{O}_\mathcal{D}$-module $\mathcal{G}$,
il existe un complexe exact
$$
\ldots \to
g_{2!}(a_2^*\mathcal{G}) \to
g_{1!}(a_1^*\mathcal{G}) \to
g_{0!}(a_0^*\mathcal{G}) \to
a^*\mathcal{G} \to 0
$$
de faisceaux de $\mathcal{O}$-modules sur $\mathcal{C}_{total}$.
Ici, $g_{n!}$ est comme dans le lemme \ref{lemma-restriction-module-to-components-site}.
\end{lemma}

\begin{proof}
Observons que $a^*\mathcal{G}$ est le $\mathcal{O}$-module sur
$\mathcal{C}_{total}$ dont la restriction à $\mathcal{C}_m$
est le $\mathcal{O}_m$-module $a_m^*\mathcal{G}$.
La description des foncteurs $g_{n!}$ sur les modules
dans le lemme \ref{lemma-restriction-module-to-components-site}
montre que $g_{n!}(a_n^*\mathcal{G})$ est le
$\mathcal{O}$-module sur $\mathcal{C}_{total}$
dont la restriction à $\mathcal{C}_m$ est le $\mathcal{O}_m$-module
$$
\bigoplus\nolimits_{\varphi : [n] \to [m]} f_\varphi^*a_n^*\mathcal{G} =
\bigoplus\nolimits_{\varphi : [n] \to [m]} a_m^*\mathcal{G}
$$
Le reste de la démonstration est identique à celui du lemme
\ref{lemma-simplicial-resolution-augmentation},
après remplacement de $a_m^{-1}\mathcal{G}$ par $a_m^*\mathcal{G}$.
\end{proof}

\begin{lemma}
\label{lemma-augmentation-cech-complex-modules}
Avec les notations ci-dessus.
Pour un $\mathcal{O}$-module $\mathcal{F}$ sur $\mathcal{C}_{total}$
il existe un complexe canonique
$$
0 \to a_*\mathcal{F} \to a_{0, *}\mathcal{F}_0 \to a_{1, *}\mathcal{F}_1 \to
a_{2, *}\mathcal{F}_2 \to \ldots
$$
de $\mathcal{O}_\mathcal{D}$-modules, exact en degrés $-1, 0$.
Si $\mathcal{F}$ est un $\mathcal{O}$-module injectif, alors le complexe
est exact à tous les degrés et reste exact en appliquant le foncteur
$\Hom_{\mathcal{O}_\mathcal{D}}(\mathcal{G}, -)$ pour tout
$\mathcal{O}_\mathcal{D}$-module $\mathcal{G}$.
\end{lemma}

\begin{proof}
Pour construire le complexe, il suffit, par le lemme de Yoneda, de construire pour tout
$\mathcal{O}_\mathcal{D}$-module $\mathcal{G}$ sur $\mathcal{D}$
un complexe
$$
0 \to \Hom_{\mathcal{O}_\mathcal{D}}(\mathcal{G}, a_*\mathcal{F}) \to
\Hom_{\mathcal{O}_\mathcal{D}}(\mathcal{G}, a_{0, *}\mathcal{F}_0) \to
\Hom_{\mathcal{O}_\mathcal{D}}(\mathcal{G}, a_{1, *}\mathcal{F}_1) \to \ldots
$$
fonctoriellement en $\mathcal{G}$. Pour cela, appliquons
$\Hom_\mathcal{O}(-, \mathcal{F})$
au complexe exact du
lemme \ref{lemma-simplicial-resolution-augmentation-modules}
et utilisons l'adjonction entre image inverse et image directe.
Les propriétés d'exactitude en degrés $-1, 0$ découlent de
la construction, puisque $\Hom_\mathcal{O}(-, \mathcal{F})$ est exact à gauche.
Si $\mathcal{F}$ est un $\mathcal{O}$-module injectif, alors le complexe
est exact car $\Hom_\mathcal{O}(-, \mathcal{F})$ est exact.
\end{proof}

\begin{lemma}
\label{lemma-augmentation-spectral-sequence-modules}
Avec les notations ci-dessus, pour tout $K$ dans $D^+(\mathcal{O})$, il existe une suite
spectrale $(E_r, d_r)_{r \geq 0}$ dans $\textit{Mod}(\mathcal{O}_\mathcal{D})$
avec
$$
E_1^{p, q} = R^qa_{p, *} K_p
$$
qui aboutit à $R^{p + q}a_*K$. Cette suite spectrale est fonctorielle en $K$.
\end{lemma}

\begin{proof}
Soit $\mathcal{I}^\bullet$ un complexe borné inférieurement de
$\mathcal{O}$-modules injectifs représentant $K$. Considérons le complexe double de termes
$$
A^{p, q} = a_{p, *}\mathcal{I}^q_p
$$
où les flèches horizontales proviennent du
lemme \ref{lemma-augmentation-cech-complex-modules}
et les flèches verticales des différentiels du complexe
$\mathcal{I}^\bullet$. Le lemme
affirme que les lignes du complexe double sont exactes
en degrés positifs et donnent
$a_*\mathcal{I}^q$ en degré $0$.
Ainsi le complexe total associé au double complexe
calcule $Ra_*K$ d'après
Homologie, lemme \ref{homology-lemma-double-complex-gives-resolution}.
D'autre part, puisque la restriction à $\mathcal{C}_p$ est exacte
(lemme \ref{lemma-restriction-to-components-site})
le complexe $\mathcal{I}_p^\bullet$ représente $K_p$ dans
$D(\mathcal{C}_p)$. Les faisceaux $\mathcal{I}_p^q$
sont totalement acycliques sur $\mathcal{C}_p$
(lemme \ref{lemma-restriction-injective-to-component-limp}).
La cohomologie des colonnes est donc donnée par les faisceaux
$R^qa_{p, *}K_p$ par le lemme d'acyclicité de Leray
(Catégories dérivées, lemme \ref{derived-lemma-leray-acyclicity})
et
Cohomologie sur les sites, lemme \ref{sites-cohomology-lemma-limp-acyclic}.
On conclut en appliquant
Homologie, lemme \ref{homology-lemma-first-quadrant-ss}.
\end{proof}





\section{Faisceaux et modules cartésiens}
\label{section-cartesian}

\noindent
Voici la définition.

\begin{definition}
\label{definition-cartesian-sheaf}
Dans la situation \ref{situation-simplicial-site}.
\begin{enumerate}
\item Un faisceau $\mathcal{F}$ d'ensembles ou de groupes abéliens sur
$\mathcal{C}_{total}$ est {\it cartésien} si les morphismes
$\mathcal{F}(\varphi) : f_\varphi^{-1}\mathcal{F}_m \to \mathcal{F}_n$
sont des isomorphismes pour tout $\varphi : [m] \to [n]$.
\item Si $\mathcal{O}$ est un faisceau d'anneaux sur $\mathcal{C}_{total}$,
alors un faisceau $\mathcal{F}$ de $\mathcal{O}$-modules est
{\it cartésien} si les morphismes $f_\varphi^*\mathcal{F}_m \to \mathcal{F}_n$
sont des isomorphismes pour tout $\varphi : [m] \to [n]$.
\item Un objet $K$ de $D(\mathcal{C}_{total})$ est {\it cartésien} si les morphismes
$f_\varphi^{-1}K_m \to K_n$
sont des isomorphismes pour tout $\varphi : [m] \to [n]$.
\item Si $\mathcal{O}$ est un faisceau d'anneaux sur $\mathcal{C}_{total}$, alors
un objet $K$ de $D(\mathcal{O})$ est {\it cartésien} si les morphismes
$Lf_\varphi^*K_m \to K_n$
sont des isomorphismes pour tout $\varphi : [m] \to [n]$.
\end{enumerate}
\end{definition}

\noindent
Bien entendu, il existe une notion générale de section cartésienne d'une
catégorie fibrée, dont les notions ci-dessus ne sont que des exemples.
La propriété sur les images inverses doit uniquement être vérifiée pour les dégénérescences.

\begin{lemma}
\label{lemma-check-cartesian-module}
Dans la situation \ref{situation-simplicial-site}.
\begin{enumerate}
\item Un faisceau $\mathcal{F}$ d'ensembles ou de groupes abéliens est cartésien
si et seulement si les morphismes
$(f_{\delta^n_j})^{-1}\mathcal{F}_{n - 1} \to \mathcal{F}_n$
sont des isomorphismes.
\item Un objet $K$ de $D(\mathcal{C}_{total})$ est cartésien
si et seulement si les morphismes
$(f_{\delta^n_j})^{-1}K_{n - 1} \to K_n$
sont des isomorphismes.
\item Si $\mathcal{O}$ est un faisceau d'anneaux sur $\mathcal{C}_{total}$,
un faisceau $\mathcal{F}$ de $\mathcal{O}$-modules est cartésien
si et seulement si les morphismes
$(f_{\delta^n_j})^*\mathcal{F}_{n - 1} \to \mathcal{F}_n$
sont des isomorphismes.
\item Si $\mathcal{O}$ est un faisceau d'anneaux sur $\mathcal{C}_{total}$,
un objet $K$ de $D(\mathcal{O})$ est cartésien
si et seulement si les morphismes
$L(f_{\delta^n_j})^*K_{n - 1} \to K_n$
sont des isomorphismes.
\item Ajouter d'autres cas ici.
\end{enumerate}
\end{lemma}

\begin{proof}
Dans chaque cas, le point essentiel est que les foncteurs image inverse
se composent en un foncteur du même type; pour la partie (4), voir
Cohomologie sur les sites, lemme
\ref{sites-cohomology-lemma-derived-pullback-composition}.
Nous montrons comment fonctionne l'argument dans le cas (1) et omettons la preuve
dans les autres cas.
La catégorie $\Delta$ est engendrée par les morphismes
$\delta^n_j$ et $\sigma^n_j$; voir
Simplicial, lemme \ref{simplicial-lemma-face-degeneracy}.
Il suffit donc de vérifier que les morphismes
$(f_{\delta^n_j})^{-1}\mathcal{F}_{n - 1} \to \mathcal{F}_n$
et $(f_{\sigma^n_j})^{-1}\mathcal{F}_{n + 1} \to \mathcal{F}_n$ sont
des isomorphismes; voir
Simplicial, lemme \ref{simplicial-lemma-characterize-simplicial-object}
pour les notations. Puisque $\sigma^n_j \circ \delta_j^{n + 1} = \text{id}_{[n]}$
la composition
$$
\mathcal{F}_n =
(f_{\sigma^n_j})^{-1}
(f_{\delta_j^{n + 1}})^{-1}
\mathcal{F}_n \to
(f_{\sigma^n_j})^{-1}
\mathcal{F}_{n + 1} \to
\mathcal{F}_n
$$
est l'identité. Ainsi, le résultat pour $\delta^{n + 1}_j$ implique le résultat
pour $\sigma^n_j$.
\end{proof}

\begin{lemma}
\label{lemma-augmentation-cartesian-module}
Dans la situation \ref{situation-simplicial-site}, soit
$a_0$ une augmentation vers un site $\mathcal{D}$ comme dans la
remarque \ref{remark-augmentation-site}.
\begin{enumerate}
\item L'image inverse $a^{-1}\mathcal{G}$ d'un faisceau d'ensembles ou de groupes abéliens
sur $\mathcal{D}$ est cartésienne.
\item L'image inverse $a^{-1}K$ d'un objet $K$ de $D(\mathcal{D})$
est cartésienne.
\end{enumerate}
Soit $\mathcal{O}$ un faisceau d'anneaux sur $\mathcal{C}_{total}$ et
$\mathcal{O}_\mathcal{D}$ un faisceau d'anneaux sur $\mathcal{D}$
et $a^\sharp : \mathcal{O}_\mathcal{D} \to a_*\mathcal{O}$ un
morphisme comme dans la section
\ref{section-cohomology-augmentation-ringed-simplicial-sites}.
\begin{enumerate}
\item[(3)] L'image inverse $a^*\mathcal{F}$ d'un faisceau de
$\mathcal{O}_\mathcal{D}$-modules est cartésienne.
\item[(4)] L'image inverse dérivée $La^*K$ d'un objet
$K$ de $D(\mathcal{O}_\mathcal{D})$ est cartésienne.
\end{enumerate}
\end{lemma}

\begin{proof}
Cela découle immédiatement des identités
$a_m \circ f_\varphi = a_n$ pour tout $\varphi : [m] \to [n]$.
Voir le lemme \ref{lemma-augmentation-site} et la discussion dans la section
\ref{section-cohomology-augmentation-ringed-simplicial-sites}.
\end{proof}

\begin{lemma}
\label{lemma-characterize-cartesian}
Dans la situation \ref{situation-simplicial-site}.
La catégorie des faisceaux cartésiens d'ensembles (resp.\ de groupes abéliens)
est équivalente à la catégorie des couples $(\mathcal{F}, \alpha)$
où $\mathcal{F}$ est un faisceau d'ensembles (resp.\ de groupes abéliens)
sur $\mathcal{C}_0$ et
$$
\alpha :
(f_{\delta_1^1})^{-1}\mathcal{F}
\longrightarrow (f_{\delta_0^1})^{-1}\mathcal{F}
$$
est un isomorphisme de faisceaux d'ensembles (resp.\ de groupes abéliens)
sur $\mathcal{C}_1$ tel que
$(f_{\delta^2_1})^{-1}\alpha =
(f_{\delta^2_0})^{-1}\alpha \circ (f_{\delta^2_2})^{-1}\alpha$
comme morphismes de faisceaux sur $\mathcal{C}_2$.
\end{lemma}

\begin{proof}
Nous abrégeons
$d^n_j = f_{\delta^n_j} : \Sh(\mathcal{C}_n) \to \Sh(\mathcal{C}_{n - 1})$.
La condition imposée à $\alpha$ dans l'énoncé du lemme est bien définie, car
$$
d^1_1 \circ d^2_2 = d^1_1 \circ d^2_1, \quad
d^1_1 \circ d^2_0 = d^1_0 \circ d^2_2, \quad
d^1_0 \circ d^2_0 = d^1_0 \circ d^2_1
$$
comme morphismes de topoi $\Sh(\mathcal{C}_2) \to \Sh(\mathcal{C}_0)$, voir
Simplicial, remarque \ref{simplicial-remark-relations}. On peut donc
représenter ces morphismes comme suit
$$
\xymatrix{
& (d^2_0)^{-1}(d^1_1)^{-1}\mathcal{F} \ar[r]_-{(d^2_0)^{-1}\alpha} &
(d^2_0)^{-1}(d^1_0)^{-1}\mathcal{F} \ar@{=}[rd] & \\
(d^2_2)^{-1}(d^1_0)^{-1}\mathcal{F} \ar@{=}[ru] & & &
(d^2_1)^{-1}(d^1_0)^{-1}\mathcal{F} \\
& (d^2_2)^{-1}(d^1_1)^{-1}\mathcal{F} \ar[lu]^{(d^2_2)^{-1}\alpha} \ar@{=}[r] &
(d^2_1)^{-1}(d^1_1)^{-1}\mathcal{F} \ar[ru]_{(d^2_1)^{-1}\alpha}
}
$$
et la condition signifie que le diagramme est commutatif. Or,
étant donné un faisceau cartésien $\mathcal{G}$ d'ensembles (resp.\ de groupes abéliens)
sur $\mathcal{C}_{total}$,
on peut poser $\mathcal{F} = \mathcal{G}_0$ et prendre pour $\alpha$ le composé
$$
(d_1^1)^{-1}\mathcal{G}_0 \to \mathcal{G}_1
\leftarrow (d_1^0)^{-1}\mathcal{G}_0
$$
où les flèches sont inversibles puisque $\mathcal{G}$ est cartésien.
Pour prouver que ce foncteur
est une équivalence, construisons-en un quasi-inverse. La construction du
quasi-inverse est analogue à celle décrite dans
Descente, section \ref{descent-section-descent-modules}, à laquelle nous empruntons
les notations $\tau^n_i : [0] \to [n]$, $0 \mapsto i$ et
$\tau^n_{ij} : [1] \to [n]$, $0 \mapsto i$, $1 \mapsto j$.
Plus précisément, pour un couple $(\mathcal{F}, \alpha)$
comme dans le lemme, posons $\mathcal{G}_n = (f_{\tau^n_n})^{-1}\mathcal{F}$.
Pour $\varphi : [n] \to [m]$, définissons
$\mathcal{G}(\varphi) : (f_\varphi)^{-1}\mathcal{G}_n \to \mathcal{G}_m$
en utilisant
$$
\xymatrix{
(f_\varphi)^{-1}\mathcal{G}_n \ar@{=}[r] &
(f_\varphi)^{-1}(f_{\tau^n_n})^{-1}\mathcal{F} \ar@{=}[r] &
(f_{\tau^m_{\varphi(n)}})^{-1}\mathcal{F} \ar@{=}[r] &
(f_{\tau^m_{\varphi(n)m}})^{-1}(d^1_1)^{-1}\mathcal{F}
\ar[d]^{(f_{\tau^m_{\varphi(n)m}})^{-1}\alpha} \\
&
\mathcal{G}_m \ar@{=}[r] &
(f_{\tau^m_m})^{-1}\mathcal{F} \ar@{=}[r] &
(f_{\tau^m_{\varphi(n)m}})^{-1}(d^1_0)^{-1}\mathcal{F}
}
$$
Nous laissons au lecteur la vérification que la commutativité du diagramme
ci-dessus entraîne la compatibilité des morphismes avec la composition et définit donc un
faisceau sur $\mathcal{C}_{total}$; voir
le lemme \ref{lemma-describe-sheaves-simplicial-site-site}.
Nous laissons également au lecteur le soin de vérifier
que les deux foncteurs sont quasi inverses l'un de l'autre.
\end{proof}

\begin{lemma}
\label{lemma-characterize-cartesian-modules}
Dans la situation \ref{situation-simplicial-site},
soit $\mathcal{O}$ un faisceau d'anneaux sur $\mathcal{C}_{total}$.
La catégorie des $\mathcal{O}$-modules cartésiens
est équivalente à la catégorie des couples $(\mathcal{F}, \alpha)$
où $\mathcal{F}$ est un $\mathcal{O}_0$-module
et
$$
\alpha :
(f_{\delta_1^1})^*\mathcal{F}
\longrightarrow (f_{\delta_0^1})^*\mathcal{F}
$$
est un isomorphisme de $\mathcal{O}_1$-modules tel que
$(f_{\delta^2_1})^*\alpha =
(f_{\delta^2_0})^*\alpha \circ (f_{\delta^2_2})^*\alpha$
comme morphismes de $\mathcal{O}_2$-modules.
\end{lemma}

\begin{proof}
La démonstration est identique à celle du
lemme \ref{lemma-characterize-cartesian},
l'image inverse des faisceaux de groupes abéliens étant remplacée
par l'image inverse des modules.
\end{proof}

\begin{lemma}
\label{lemma-Serre-subcat-cartesian-modules}
Dans la situation \ref{situation-simplicial-site}.
\begin{enumerate}
\item La sous-catégorie pleine des faisceaux abéliens cartésiens est une
sous-catégorie de Serre faible de $\textit{Ab}(\mathcal{C}_{total})$.
Les colimites de systèmes de faisceaux abéliens cartésiens sont cartésiennes.
\item Soit $\mathcal{O}$ un faisceau d'anneaux sur $\mathcal{C}_{total}$
tel que les morphismes
$$
f_{\delta^n_j} : (\Sh(\mathcal{C}_n), \mathcal{O}_n)
\to (\Sh(\mathcal{C}_{n - 1}), \mathcal{O}_{n - 1})
$$
sont plats. La sous-catégorie pleine des $\mathcal{O}$-modules cartésiens est une
sous-catégorie de Serre faible de $\textit{Mod}(\mathcal{O})$.
Les colimites de systèmes de $\mathcal{O}$-modules cartésiens sont cartésiennes.
\end{enumerate}
\end{lemma}

\begin{proof}
Pour voir qu'on obtient en (1) une sous-catégorie de Serre faible,
vérifions les conditions énumérées dans
Homologie, lemme \ref{homology-lemma-characterize-weak-serre-subcategory}.
D'abord, si $\varphi : \mathcal{F} \to \mathcal{G}$ est un morphisme
entre faisceaux abéliens cartésiens, alors
$\Ker(\varphi)$ et $\Coker(\varphi)$ sont également cartésiens
car les foncteurs de restriction
$\Sh(\mathcal{C}_{total}) \to \Sh(\mathcal{C}_n)$
et les foncteurs $f_\varphi^{-1}$ sont exacts.
De même, si
$$
0 \to \mathcal{F} \to \mathcal{H} \to \mathcal{G} \to 0
$$
est une suite exacte courte de faisceaux abéliens sur $\mathcal{C}_{total}$
avec $\mathcal{F}$ et $\mathcal{G}$ cartésiens, il s'ensuit que
$\mathcal{H}$ est cartésien par le lemme des cinq. Pour les colimites,
utilisons le fait qu'elles commutent aux images inverses, puisque les foncteurs image inverse sont adjoints
à gauche. Dans le cas des modules,
on raisonne de même en utilisant l'exactitude de l'image inverse par un morphisme plat
(Modules sur les sites, lemme \ref{sites-modules-lemma-flat-pullback-exact})
et le fait qu'il suffit de vérifier la condition
pour $f_{\delta^n_j}$; voir le lemme \ref{lemma-check-cartesian-module}.
\end{proof}

\begin{remark}[Avertissement]
\label{remark-warning-cartesian-modules}
Malgré le lemme \ref{lemma-Serre-subcat-cartesian-modules}, il
peut arriver que la catégorie des $\mathcal{O}$-modules cartésiens soit abélienne
sans être une sous-catégorie de Serre de $\textit{Mod}(\mathcal{O})$.
Plus précisément, supposons que nous sachions seulement que
$f_{\delta_1^1}$ et $f_{\delta_0^1}$ sont plats.
Il résulte alors facilement du lemme
\ref{lemma-characterize-cartesian-modules}
que la catégorie des $\mathcal{O}$-modules cartésiens est abélienne.
Mais si $f_{\delta_0^2}$ n'est pas plat (par exemple),
il n'y a aucune raison pour que le foncteur d'inclusion
de la catégorie des $\mathcal{O}$-modules cartésiens
dans la catégorie de tous les $\mathcal{O}$-modules soit exact.
\end{remark}

\begin{lemma}
\label{lemma-derived-cartesian-modules}
Dans la situation \ref{situation-simplicial-site}.
\begin{enumerate}
\item Un objet $K$ de $D(\mathcal{C}_{total})$ est cartésien si et seulement
si $H^q(K)$ est un faisceau abélien cartésien pour tout $q$.
\item Soit $\mathcal{O}$ un faisceau
d'anneaux sur $\mathcal{C}_{total}$ tel que les morphismes
$f_{\delta^n_j} : (\Sh(\mathcal{C}_n), \mathcal{O}_n)
\to (\Sh(\mathcal{C}_{n - 1}), \mathcal{O}_{n - 1})$ sont plats.
Alors un objet $K$ de $D(\mathcal{O})$ est cartésien si et seulement
si $H^q(K)$ est un $\mathcal{O}$-module cartésien pour tout $q$.
\end{enumerate}
\end{lemma}

\begin{proof}
La partie (1) résulte de ce que les foncteurs image inverse $(f_\varphi)^{-1}$
sont exacts. La partie (2) découle de la caractérisation dans le lemme
\ref{lemma-check-cartesian-module}
et du fait que $L(f_{\delta^n_j})^* = (f_{\delta^n_j})^*$
par planéité.
\end{proof}

\begin{lemma}
\label{lemma-derived-cartesian-shriek}
Dans la situation \ref{situation-simplicial-site}.
\begin{enumerate}
\item Un objet $K$ de $D(\mathcal{C}_{total})$ est cartésien si et seulement
le morphisme canonique
$$
g_{n!}K_n \longrightarrow
g_{n!}\mathbf{Z} \otimes^\mathbf{L}_\mathbf{Z} K
$$
est un isomorphisme pour tout $n$.
\item Soit $\mathcal{O}$ un faisceau d'anneaux sur $\mathcal{C}_{total}$
tel que les morphismes $f_\varphi^{-1}\mathcal{O}_n \to \mathcal{O}_m$
sont plats pour tout $\varphi : [n] \to [m]$. Alors un objet $K$ de
$D(\mathcal{O})$ est cartésien si et seulement si le morphisme canonique
$$
g_{n!}K_n \longrightarrow
g_{n!}\mathcal{O}_n \otimes^\mathbf{L}_\mathcal{O} K
$$
est un isomorphisme pour tout $n$.
\end{enumerate}
\end{lemma}

\begin{proof}
Pour (1). Puisque $g_{n!}$ est exact, il induit un foncteur
sur les catégories dérivées adjoint à $g_n^{-1}$.
Le morphisme est adjoint au morphisme
$K_n \to (g_n^{-1}g_{n!}\mathbf{Z}) \otimes^\mathbf{L}_\mathbf{Z} K_n$
correspondant à $\mathbf{Z} \to g_n^{-1}g_{n!}\mathbf{Z}$
qui à son tour est adjoint à
$\text{id} : g_{n!}\mathbf{Z} \to g_{n!}\mathbf{Z}$.
En utilisant la description de $g_{n!}$
donnée dans le lemme \ref{lemma-restriction-to-components-site}
on voit que la restriction à $\mathcal{C}_m$ de ce morphisme
est
$$
\bigoplus\nolimits_{\varphi : [n] \to [m]} f_\varphi^{-1}K_n
\longrightarrow
\bigoplus\nolimits_{\varphi : [n] \to [m]} K_m
$$
L'assertion est donc claire.

\medskip\noindent
Pour (2). Puisque $g_{n!}$ est exact
(lemme \ref{lemma-exactness-g-shriek-modules}), il induit un foncteur
sur les catégories dérivées adjoint à $g_n^*$ (également exact).
Le morphisme est adjoint au morphisme
$K_n \to (g_n^*g_{n!}\mathcal{O}_n) \otimes^\mathbf{L}_{\mathcal{O}_n} K_n$
correspondant à $\mathcal{O}_n \to g_n^*g_{n!}\mathcal{O}_n$
qui à son tour est adjoint à
$\text{id} : g_{n!}\mathcal{O}_n \to g_{n!}\mathcal{O}_n$.
En utilisant la description de $g_{n!}$
donnée dans le lemme \ref{lemma-restriction-module-to-components-site}
on voit que la restriction à $\mathcal{C}_m$ de ce morphisme
est
$$
\bigoplus\nolimits_{\varphi : [n] \to [m]} f_\varphi^*K_n
\longrightarrow
\bigoplus\nolimits_{\varphi : [n] \to [m]}
f_\varphi^*\mathcal{O}_n \otimes_{\mathcal{O}_m} K_m =
\bigoplus\nolimits_{\varphi : [n] \to [m]} K_m
$$
L'assertion est donc claire.
\end{proof}

\begin{lemma}
\label{lemma-quasi-coherent-sheaf}
Dans la situation \ref{situation-simplicial-site},
soit $\mathcal{O}$ un faisceau d'anneaux sur $\mathcal{C}_{total}$.
Soit $\mathcal{F}$ un faisceau de $\mathcal{O}$-modules.
Alors $\mathcal{F}$ est quasi-cohérent au sens de
Modules sur les sites, définition \ref{sites-modules-definition-site-local},
si et seulement si $\mathcal{F}$ est cartésien
et si $\mathcal{F}_n$ est un $\mathcal{O}_n$-module quasi-cohérent pour tout $n$.
\end{lemma}

\begin{proof}
Supposons que $\mathcal{F}$ soit quasi-cohérent. Puisque les images inverses des modules quasi-cohérents
sont quasi-cohérents
(Modules sur les sites, lemme \ref{sites-modules-lemma-local-pullback}),
nous voyons que $\mathcal{F}_n$ est un module $\mathcal{O}_n$ quasi-cohérent
pour tout $n$. Pour montrer que $\mathcal{F}$ est cartésien, soit $U$
un objet de $\mathcal{C}_n$ pour un certain $n$. Regardons $U$
comme un objet de $\mathcal{C}_{total}$. Comme $\mathcal{F}$
est quasi-cohérent, il existe un recouvrement $\{U_i \to U\}$
et pour chaque $i$ une présentation
$$
\bigoplus\nolimits_{j \in J_i} \mathcal{O}_{\mathcal{C}_{total}/U_i} \to
\bigoplus\nolimits_{k \in K_i} \mathcal{O}_{\mathcal{C}_{total}/U_i} \to
\mathcal{F}|_{\mathcal{C}_{total}/U_i} \to 0
$$
Notons que $\{U_i \to U\}$ est un recouvrement de $\mathcal{C}_n$ par
la construction du site $\mathcal{C}_{total}$.
Soit ensuite $V$ un objet de $\mathcal{C}_m$ pour un certain $m$, et soit
$V \to U$ un morphisme de $\mathcal{C}_{total}$ au-dessus de
$\varphi : [n] \to [m]$. Les produits fibrés $V_i = V \times_U U_i$
existent et donnent un recouvrement induit $\{V_i \to V\}$ dans $\mathcal{C}_m$.
En restreignant la présentation ci-dessus aux sites
$\mathcal{C}_n/U_i$ et $\mathcal{C}_m/V_i$, on obtient
les présentations
$$
\bigoplus\nolimits_{j \in J_i} \mathcal{O}_{\mathcal{C}_m/U_i} \to
\bigoplus\nolimits_{k \in K_i} \mathcal{O}_{\mathcal{C}_m/U_i} \to
\mathcal{F}_n|_{\mathcal{C}_n/U_i} \to 0
$$
et
$$
\bigoplus\nolimits_{j \in J_i} \mathcal{O}_{\mathcal{C}_m/V_i} \to
\bigoplus\nolimits_{k \in K_i} \mathcal{O}_{\mathcal{C}_m/V_i} \to
\mathcal{F}_m|_{\mathcal{C}_m/V_i} \to 0
$$
Ces présentations sont compatibles avec le morphisme
$\mathcal{F}(\varphi) : f_\varphi^*\mathcal{F}_n \to \mathcal{F}_m$
(car ce morphisme est défini à l'aide des morphismes de restriction de $\mathcal{F}$
le long des morphismes de $\mathcal{C}_{total}$ recouvrant $\varphi$).
On en déduit que $\mathcal{F}(\varphi)|_{\mathcal{C}_m/V_i}$
est un isomorphisme. Comme $\{V_i \to V\}$ est un recouvrement,
$\mathcal{F}(\varphi)|_{\mathcal{C}_m/V}$ est donc un isomorphisme.
L'arbitraire de $V$ et de $U$ montre que $\mathcal{F}$ est cartésien.
(Dans le cas (A), on utilise Sites, lemme \ref{sites-lemma-morphism-of-sites-covering}.)

\medskip\noindent
Réciproquement, supposons $\mathcal{F}_n$
quasi-cohérent pour tout $n$ et que $\mathcal{F}$ est cartésien.
Alors, pour tout $n$ et tout objet $U$ de $\mathcal{C}_n$, nous
pouvons choisir un recouvrement $\{U_i \to U\}$ de $\mathcal{C}_n$
et pour chacun $i$ une présentation
$$
\bigoplus\nolimits_{j \in J_i} \mathcal{O}_{\mathcal{C}_m/U_i} \to
\bigoplus\nolimits_{k \in K_i} \mathcal{O}_{\mathcal{C}_m/U_i} \to
\mathcal{F}_n|_{\mathcal{C}_n/U_i} \to 0
$$
En prenant l'image inverse sur $\mathcal{C}_{total}/U_i$, on obtient les complexes
$$
\bigoplus\nolimits_{j \in J_i} \mathcal{O}_{\mathcal{C}_{total}/U_i} \to
\bigoplus\nolimits_{k \in K_i} \mathcal{O}_{\mathcal{C}_{total}/U_i} \to
\mathcal{F}|_{\mathcal{C}_{total}/U_i} \to 0
$$
de modules sur $\mathcal{C}_{total}/U_i$. La cartésianité de
$\mathcal{F}$ implique alors l'exactitude de ce complexe.
Nous omettons les détails.
\end{proof}











\section{Systèmes simpliciaux de la catégorie dérivée}
\label{section-glueing}

\noindent
Dans cette section, nous allons prouver un cas particulier de
\cite[Proposition 3.2.9]{BBD} dans le cadre des catégories dérivées
de faisceaux abéliens. Le cas des modules
est abordé dans la section \ref{section-glueing-modules}.

\begin{definition}
\label{definition-cartesian-derived}
Dans la situation \ref{situation-simplicial-site}. Un
{\it système simplicial de la catégorie dérivée}
se compose des données suivantes
\begin{enumerate}
\item pour tout $n$, un objet $K_n$ de $D(\mathcal{C}_n)$,
\item pour tout $\varphi : [m] \to [n]$, un morphisme
$K_\varphi : f_\varphi^{-1}K_m \to K_n$ dans $D(\mathcal{C}_n)$
\end{enumerate}
soumises à la condition
$$
K_{\varphi \circ \psi} = K_\varphi \circ f_\varphi^{-1}K_\psi :
f_{\varphi \circ \psi}^{-1}K_l = f_\varphi^{-1} f_\psi^{-1}K_l
\longrightarrow
K_n
$$
pour tout couple de morphismes $\varphi : [m] \to [n]$ et $\psi : [l] \to [m]$ de $\Delta$.
On dit que le système simplicial est {\it cartésien} si les morphismes $K_\varphi$
sont des isomorphismes pour tout $\varphi$.
Étant donné deux systèmes simpliciaux de la catégorie dérivée,
il existe une notion évidente de
{\it morphisme de systèmes simpliciaux de la catégorie dérivée}.
\end{definition}

\noindent
Nous avons intentionnellement donné à cette notion un nom ridiculement long.
Le but est de montrer qu'un système simplicial de la catégorie dérivée
provient d'un objet de $D(\mathcal{C}_{total})$ sous certaines
hypothèses.

\begin{lemma}
\label{lemma-cartesian-objects-derived}
Dans la situation \ref{situation-simplicial-site}.
Si $K \in D(\mathcal{C}_{total})$ est un objet,
alors $(K_n, K(\varphi))$ est un système simplicial de la catégorie dérivée.
Si $K$ est cartésien, le système l’est également.
\end{lemma}

\begin{proof}
C'est évident.
\end{proof}

\begin{lemma}
\label{lemma-construct-cartesian-objects-derived}
Dans la situation \ref{situation-simplicial-site},
supposons donnés $K_0 \in D(\mathcal{C}_0)$ et un isomorphisme
$$
\alpha :
f_{\delta_1^1}^{-1}K_0
\longrightarrow
f_{\delta_0^1}^{-1}K_0
$$
satisfaisant à la condition de cocycle. Posons
$\tau^n_i : [0] \to [n]$, $0 \mapsto i$ et
puis posons $K_n = f_{\tau^n_n}^{-1}K_0$. Les $K_n$
forment un système simplicial cartésien de la catégorie dérivée.
\end{lemma}

\begin{proof}
Comparer avec le lemme \ref{lemma-characterize-cartesian}
et sa preuve (où la condition de cocycle est aussi explicitée).
La construction est analogue à celle décrite dans
Descente, section \ref{descent-section-descent-modules}, à laquelle
nous empruntons les notations $\tau^n_i : [0] \to [n]$, $0 \mapsto i$ et
$\tau^n_{ij} : [1] \to [n]$, $0 \mapsto i$, $1 \mapsto j$.
Pour $\varphi : [n] \to [m]$, définissons
$K_\varphi : f_\varphi^{-1}K_n \to K_m$
en utilisant
$$
\xymatrix{
f_\varphi^{-1}K_n \ar@{=}[r] &
f_\varphi^{-1} f_{\tau^n_n}^{-1}K_0 \ar@{=}[r] &
f_{\tau^m_{\varphi(n)}}^{-1}K_0 \ar@{=}[r] &
f_{\tau^m_{\varphi(n)m}}^{-1}f_{\delta^1_1}^{-1}K_0
\ar[d]_{f_{\tau^m_{\varphi(n)m}}^{-1}\alpha} \\
&
K_m \ar@{=}[r] &
f_{\tau^m_m}^{-1}K_0 \ar@{=}[r] &
f_{\tau^m_{\varphi(n)m}}^{-1}f_{\delta^1_0}^{-1}K_0
}
$$
Nous laissons au lecteur le soin de vérifier que la condition de cocycle
entraîne la compatibilité des morphismes avec la composition (dans leurs catégories dérivées
respectives) et définit donc un système simplicial
dans la catégorie dérivée.
\end{proof}

\begin{lemma}
\label{lemma-abelian-postnikov}
Dans la situation \ref{situation-simplicial-site}. Soit $K$
un objet de $D(\mathcal{C}_{total})$. Posons
$$
X_n = (g_{n!}\mathbf{Z})
\otimes^\mathbf{L}_\mathbf{Z} K
\quad\text{et}\quad
Y_n =
(g_{n!}\mathbf{Z} \to \ldots \to g_{0!}\mathbf{Z})[-n]
\otimes^\mathbf{L}_\mathbf{Z} K
$$
comme objets de $D(\mathcal{C}_{total})$, les morphismes étant
comme dans le lemme \ref{lemma-simplicial-resolution-Z-site}.
Avec les morphismes canoniques évidents $Y_n \to X_n$ et
$Y_0 \to Y_1[1] \to Y_2[2] \to \ldots$, on a
\begin{enumerate}
\item les triangles distingués $Y_n \to X_n \to Y_{n - 1} \to Y_n[1]$
définissent un système de Postnikov
(Catégories dérivées, définition \ref{derived-definition-postnikov-system})
pour $\ldots \to X_2 \to X_1 \to X_0$,
\item $K = \text{hocolim} Y_n[n]$ dans $D(\mathcal{C}_{total})$.
\end{enumerate}
\end{lemma}

\begin{proof}
D'abord, si $K = \mathbf{Z}$, c'est la construction de
Catégories dérivées, exemple \ref{derived-example-key-postnikov},
appliquée au complexe
$$
\ldots \to
g_{2!}\mathbf{Z} \to
g_{1!}\mathbf{Z} \to
g_{0!}\mathbf{Z}
$$
dans $\textit{Ab}(\mathcal{C}_{total})$, compte tenu du fait que
ce complexe représente $K = \mathbf{Z}$ dans $D(\mathcal{C}_{total})$
par le lemme \ref{lemma-simplicial-resolution-Z-site}.
Le cas général en découle, puisque le foncteur exact
$- \otimes^\mathbf{L}_\mathbf{Z} K$ transforme les systèmes de Postnikov en
systèmes de Postnikov et
que $- \otimes^\mathbf{L}_\mathbf{Z} K$ commute aux colimites homotopiques.
\end{proof}

\begin{lemma}
\label{lemma-nullity-cartesian-objects-derived}
Dans la situation \ref{situation-simplicial-site}.
Soient $K, K' \in D(\mathcal{C}_{total})$.
Supposons que
\begin{enumerate}
\item $K$ est cartésien,
\item $\Hom(K_i[i], K'_i) = 0$ pour $i > 0$ et
\item $\Hom(K_i[i + 1], K'_i) = 0$ pour $i \geq 0$.
\end{enumerate}
Alors tout morphisme $K \to K'$ induisant le morphisme nul $K_0 \to K'_0$ est nul.
\end{lemma}

\begin{proof}
Considérons les objets $X_n$ et le système de Postnikov $Y_n$
associé à $K$ dans le lemme \ref{lemma-abelian-postnikov}.
Comme $K = \text{hocolim} Y_n[n]$, le morphisme $K \to K'$ fournit
une famille compatible de morphismes $Y_n[n] \to K'$.
D'après (1) et le lemme \ref{lemma-derived-cartesian-shriek}, on a
$X_n = g_{n!}K_n$. Puisque $Y_0 = X_0$, la nullité de
$K_0 \to K'_0$ entraîne celle de $Y_0 \to K'$.
Supposons par récurrence que le morphisme $Y_n[n] \to K'$ soit nul
pour un certain $n \geq 0$. Le triangle distingué
$$
Y_n[n] \to Y_{n + 1}[n + 1] \to X_{n + 1}[n + 1] \to Y_n[n + 1]
$$
on déduit une suite exacte
$$
\Hom(X_{n + 1}[n + 1], K') \to
\Hom(Y_{n + 1}[n + 1], K') \to
\Hom(Y_n[n], K')
$$
Comme $X_{n + 1}[n + 1] = g_{n + 1!}K_{n + 1}[n + 1]$, le premier groupe est égal à
$$
\Hom(K_{n + 1}[n + 1], K'_{n + 1})
$$
qui est nul par l'hypothèse (2). Par récurrence, tous les morphismes
$Y_n[n] \to K'$ sont nuls. Considérons le triangle distingué qui définit
$$
\bigoplus Y_n[n] \to
\bigoplus Y_n[n] \to
K \to
(\bigoplus Y_n[n])[1]
$$
la colimite homotopique. En raisonnant comme ci-dessus, on voit qu'il suffit
de montrer que
$$
\Hom((\bigoplus Y_n[n])[1], K') = \prod \Hom(Y_n[n + 1], K')
$$
soit nul pour tout $n \geq 0$. Pour le vérifier, en raisonnant comme ci-dessus,
il suffit de montrer que
$$
\Hom(K_n[n + 1], K'_n)  = 0
$$
pour tout $n \geq 0$, ce qui découle de la condition (3).
\end{proof}

\begin{lemma}
\label{lemma-hom-cartesian-objects-derived}
Dans la situation \ref{situation-simplicial-site},
soient $K, K' \in D(\mathcal{C}_{total})$.
Supposons que
\begin{enumerate}
\item $K$ est cartésien,
\item $\Hom(K_i[i - 1], K'_i) = 0$ pour $i > 1$.
\end{enumerate}
Alors tout morphisme $\{K_n \to K'_n\}$ entre les systèmes simpliciaux associés
à $K$ et $K'$ provient d'un morphisme $K \to K'$ dans $D(\mathcal{C}_{total})$.
\end{lemma}

\begin{proof}
Soit $\{K_n \to K'_n\}_{n \geq 0}$
un morphisme de systèmes simpliciaux de la catégorie dérivée.
Considérons les objets $X_n$ et le système de Postnikov $Y_n$
associés au $K$ du lemme \ref{lemma-abelian-postnikov}.
D'après (1) et le lemme \ref{lemma-derived-cartesian-shriek}, on a
$X_n = g_{n!}K_n$. En particulier, le morphisme $K_0 \to K'_0$
induit un morphisme $X_0 \to K'$. Puisque $\{K_n \to K'_n\}$
est un morphisme de systèmes, un calcul (omis) montre que
la composition
$$
X_1 \to X_0 \to K'
$$
est nulle. Comme $Y_0 = X_0$ et que $Y_1$ s'insère dans un
triangle distingué
$$
Y_1 \to X_1 \to Y_0 \to Y_1[1]
$$
on en déduit l'existence d'un morphisme $Y_1[1] \to K'$ dont la composition
avec $X_0 = Y_0 \to Y_1[1]$ est le morphisme $X_0 \to K'$
ci-dessus. Supposons donné un morphisme $Y_n[n] \to K'$ pour $n \geq 1$.
À partir du triangle distingué
$$
X_{n + 1}[n] \to Y_n[n] \to Y_{n + 1}[n + 1] \to X_{n + 1}[n + 1]
$$
on obtient une suite exacte
$$
\Hom(Y_{n + 1}[n + 1], K') \to
\Hom(Y_n[n], K') \to
\Hom(X_{n + 1}[n], K')
$$
Comme $X_{n + 1}[n] = g_{n + 1!}K_{n + 1}[n]$, le dernier groupe est égal à
$$
\Hom(K_{n + 1}[n], K'_{n + 1})
$$
qui est nul par l'hypothèse (2). Par récurrence, on construit un système de morphismes
$Y_n[n] \to K'$ compatible avec les morphismes de transition et se réduisant
au morphisme donné sur $Y_0$. Cela fournit un morphisme
$$
\gamma :
K = \text{hocolim} Y_n[n]
\longrightarrow
K'
$$
Ce morphisme rend en particulier commutatif le diagramme
$$
\xymatrix{
X_0 \ar[rd] \ar[r] &
K \ar[d]^\gamma \\
& K'
}
$$
est commutatif. Par restriction à
$\mathcal{C}_0$, on en déduit que le morphisme $\gamma_0 : K_0 \to K'_0$
est le même que le premier morphisme $K_0 \to K'_0$ du morphisme
de systèmes simpliciaux. Comme $K$ est cartésien, on en déduit aisément que
$\{\gamma_n\}$ est le morphisme de systèmes simpliciaux dont on était parti.
\end{proof}

\begin{lemma}
\label{lemma-cartesian-object-derived-from-simplicial}
Dans la situation \ref{situation-simplicial-site}. Soit
$(K_n, K_\varphi)$ un système simplicial de la catégorie dérivée.
Supposons que
\begin{enumerate}
\item $(K_n, K_\varphi)$ est cartésien,
\item $\Hom(K_i[t], K_i) = 0$ pour $i \geq 0$ et $t > 0$.
\end{enumerate}
Il existe alors un objet cartésien $K$ de $D(\mathcal{C}_{total})$
dont le système simplicial associé est isomorphe à $(K_n, K_\varphi)$.
\end{lemma}

\begin{proof}
Posons $X_n = g_{n!}K_n$ dans $D(\mathcal{C}_{total})$. Pour tout $n \geq 1$,
on a
$$
\Hom(X_n, X_{n - 1}) =
\Hom(K_n, g_n^{-1}g_{n - 1!}K_{n - 1}) =
\bigoplus\nolimits_{\varphi : [n - 1] \to [n]}
\Hom(K_n, f_\varphi^{-1}K_{n - 1})
$$
On obtient ainsi un morphisme $X_n \to X_{n - 1}$ correspondant à la somme alternée
des morphismes
$K_\varphi^{-1} : K_n \to f_\varphi^{-1}K_{n - 1}$
où $\varphi$ parcourt $\delta^n_0, \ldots, \delta^n_n$.
Nous pouvons le faire car $K_\varphi$ est inversible par l'hypothèse (1).
Comparer avec la définition des morphismes
dans la preuve du lemme \ref{lemma-simplicial-resolution-Z-site}.
On obtient un complexe
$$
\ldots \to X_2 \to X_1 \to X_0
$$
dans $D(\mathcal{C}_{total})$. Nous omettons le calcul qui montre
que les composés sont nuls. D'après
Catégories dérivées, lemme \ref{derived-lemma-existence-postnikov-system},
si l'on a
$$
\Hom(X_i[i - j - 2], X_j) = 0\text{ pour }i > j + 2
$$
on peut prolonger ce complexe en un système de Postnikov.
Le groupe est égal à
$$
\Hom(K_i[i - j - 2], g_i^{-1}g_{j!}K_j)
$$
En utilisant à nouveau la cartésianité de $(K_n, K_\varphi)$, on voit que
$g_i^{-1}g_{j!}K_j$ est isomorphe à une somme directe finie de copies de
$K_i$. Ce groupe s'annule donc par l'hypothèse (2).
Prenons le système de Postnikov défini par $Y_0 = X_0$ et les triangles distingués
$Y_n \to X_n \to Y_{n - 1} \to Y_n[1]$ pour $n \geq 1$.
Nous définissons
$$
K = \text{hocolim} Y_n[n]
$$
Pour terminer la preuve, nous devons montrer que $g_m^{-1}K$ est isomorphe
à $K_m$ pour tout $m$, de façon compatible avec les morphismes $K_\varphi$. Notons que
$$
g_m^{-1} K = \text{hocolim} g_m^{-1}Y_n[n]
$$
et que $g_m^{-1}Y_n[n]$ est un système Postnikov pour $g_m^{-1}X_n$.
Considérons les isomorphismes
$$
g_m^{-1}X_n =
\bigoplus\nolimits_{\varphi : [n] \to [m]} f_\varphi^{-1}K_n
\xrightarrow{\bigoplus K_\varphi}
\bigoplus\nolimits_{\varphi : [n] \to [m]} K_m
$$
Ces morphismes définissent un isomorphisme de complexes
$$
\xymatrix{
\ldots \ar[r] &
g_m^{-1}X_2 \ar[r] \ar[d] &
g_m^{-1}X_1 \ar[r] \ar[d] &
g_m^{-1}X_0 \ar[d] \\
\ldots \ar[r] &
\bigoplus\nolimits_{\varphi : [2] \to [m]} K_m \ar[r] &
\bigoplus\nolimits_{\varphi : [1] \to [m]} K_m \ar[r] &
\bigoplus\nolimits_{\varphi : [0] \to [m]} K_m
}
$$
dans $D(\mathcal{C}_m)$; les flèches de la rangée inférieure sont celles de
dans la preuve du lemme \ref{lemma-simplicial-resolution-Z-site}.
Les carrés sont commutatifs par le choix des flèches du complexe
$\ldots \to X_2 \to X_1 \to X_0$ ; nous omettons le calcul.
Le complexe formé par la rangée inférieure admet la tour de Postnikov
$$
Y'_{m, n} =
\left(\bigoplus\nolimits_{\varphi : [n] \to [m]} \mathbf{Z} \to
\ldots \to
\bigoplus\nolimits_{\varphi : [0] \to [m]} \mathbf{Z}\right)[-n]
\otimes^\mathbf{L}_\mathbf{Z} K_m
$$
et l'on a $\text{hocolim} Y'_{m, n} = K_m$
(comparer avec la preuve du lemme \ref{lemma-abelian-postnikov}
et Catégories dérivées, exemple \ref{derived-example-key-postnikov}).
La seconde partie de
Catégories dérivées, lemme \ref{derived-lemma-existence-postnikov-system},
permet de prolonger les morphismes verticaux du grand diagramme en un isomorphisme
de systèmes de Postnikov, pourvu que l'on ait
$$
\Hom(g_m^{-1}X_i[i - j - 1], \bigoplus\nolimits_{\varphi : [j] \to [m]} K_m)
= 0\text{ pour }i > j + 1
$$
Cette annulation découle de $\Hom(K_m[i - j - 1], K_m) = 0$ pour $i > j + 1$,
par l'hypothèse (2). Choisissons un isomorphisme
de systèmes de Postnikov $\gamma_{m, n} : g_m^{-1}Y_n \to Y'_{m, n}$
dans $D(\mathcal{C}_m)$. Par unicité des colimites homotopiques,
il existe un isomorphisme
$$
g_m^{-1} K = \text{hocolim} g_m^{-1}Y_n[n]
\xrightarrow{\gamma_m}
\text{hocolim} Y'_{m, n} = K_m
$$
compatible avec $\gamma_{m, n}$.

\medskip\noindent
Il reste à montrer que les morphismes $\gamma_m$ s'inscrivent dans les diagrammes commutatifs
$$
\xymatrix{
f_\varphi^{-1}g_m^{-1}K \ar[d]_{f_\varphi^{-1}\gamma_m} \ar[r]_{K(\varphi)} &
g_n^{-1}K \ar[d]^{\gamma_n} \\
f_\varphi^{-1}K_m \ar[r]^{K_\varphi} &
K_n
}
$$
pour tout $\varphi : [m] \to [n]$. Considérons le diagramme
$$
\xymatrix{
f_\varphi^{-1}(\bigoplus_{\psi : [0] \to [m]} f_\psi^{-1}K_0)
\ar@{=}[r] \ar[d]_{f_\varphi^{-1}(\bigoplus K_\psi)} &
f_\varphi^{-1}g_m^{-1}X_0 \ar[d] \ar[r]_{X_0(\varphi)} &
g_n^{-1}X_0 \ar[d] &
\bigoplus_{\chi : [0] \to [n]} f_\chi^{-1}K_0
\ar@{=}[l] \ar[d]^{\bigoplus K_\chi} \\
f_\varphi^{-1}(\bigoplus_{\psi : [0] \to [m]} K_m) \ar@{=}[d] &
f_\varphi^{-1}g_m^{-1}K \ar[d]_{f_\varphi^{-1}\gamma_m} \ar[r]_{K(\varphi)} &
g_n^{-1}K \ar[d]^{\gamma_n} &
\bigoplus_{\chi : [0] \to [n]} K_n \ar@{=}[d] \\
f_\varphi^{-1}Y'_{0, m} \ar[r] &
f_\varphi^{-1}K_m \ar[r]^{K_\varphi} &
K_n &
Y'_{0, n} \ar[l]
}
$$
Le carré du milieu supérieur est commutatif car $X_0 \to K$ est un morphisme
d'objets simpliciaux. Les rectangles de gauche et de droite sont
commutatifs, car $\gamma_m$ et $\gamma_n$ sont respectivement compatibles avec
$\gamma_{0, m}$ et $\gamma_{0, n}$, qui sont les flèches
$\bigoplus K_\psi$ et $\bigoplus K_\chi$ dans le diagramme.
Le rectangle extérieur du diagramme
est commutatif, car $(K_n, K_\varphi)$ est un système simplicial
et le morphisme $X_0(\varphi)$ est donné par les identifications évidentes
$f_\varphi^{-1}f_\psi^{-1}K_0 = f_{\varphi \circ \psi}^{-1}K_0$.
Notons que la flèche $\bigoplus_\psi K_m \to Y'_{0, m} \to K_m$
induit un isomorphisme sur chacun des facteurs directs
(en raison de notre construction explicite des systèmes de Postnikov
$Y'_{i, j}$ ci-dessus).
Ainsi, si l'on prend un facteur direct du
coin supérieur gauche, celui-ci s'envoie isomorphiquement sur
$f_\varphi^{-1}g_m^{-1}K$ car $\gamma_m$ est un isomorphisme.
En explicitant ce qui précède
et en ne considérant que ce facteur direct, on conclut que le carré
central inférieur est lui aussi commutatif. Cela termine la preuve.
\end{proof}









\section{Systèmes simpliciaux de la catégorie dérivée : modules}
\label{section-glueing-modules}

\noindent
Dans cette section, nous allons prouver un cas particulier de
\cite[Proposition 3.2.9]{BBD} dans le cadre des catégories
dérivées de $\mathcal{O}$-modules. Le cas (un peu) plus simple des faisceaux abéliens
est discuté dans la section \ref{section-glueing}.

\begin{definition}
\label{definition-cartesian-derived-modules}
Dans la situation \ref{situation-simplicial-site}. Soit $\mathcal{O}$
un faisceau d'anneaux sur $\mathcal{C}_{total}$. Un
{\it système simplicial de la catégorie dérivée de modules}
se compose des données suivantes
\begin{enumerate}
\item pour tout $n$, un objet $K_n$ de $D(\mathcal{O}_n)$,
\item pour tout $\varphi : [m] \to [n]$, un morphisme
$K_\varphi : Lf_\varphi^*K_m \to K_n$ dans $D(\mathcal{O}_n)$
\end{enumerate}
soumises à la condition
$$
K_{\varphi \circ \psi} = K_\varphi \circ Lf_\varphi^*K_\psi :
Lf_{\varphi \circ \psi}^*K_l = Lf_\varphi^* Lf_\psi^*K_l
\longrightarrow
K_n
$$
pour tout couple de morphismes $\varphi : [m] \to [n]$ et $\psi : [l] \to [m]$ de $\Delta$.
On dit que le système simplicial est {\it cartésien} si les morphismes $K_\varphi$
sont des isomorphismes pour tout $\varphi$.
Étant donnés deux systèmes simpliciaux de la catégorie dérivée,
il existe une notion évidente de
{\it morphisme de systèmes simpliciaux de la catégorie dérivée de modules}.
\end{definition}

\noindent
Nous avons intentionnellement donné à cette notion un nom ridiculement long.
Le but est de montrer qu'un système simplicial de la catégorie dérivée
de modules provient d'un objet de $D(\mathcal{O})$ sous certaines
hypothèses.

\begin{lemma}
\label{lemma-cartesian-objects-derived-modules}
Dans la situation \ref{situation-simplicial-site}, soit $\mathcal{O}$ un faisceau d'anneaux
sur $\mathcal{C}_{total}$.
Si $K \in D(\mathcal{O})$ est un objet, alors $(K_n, K(\varphi))$
est un système simplicial de la catégorie dérivée de modules.
Si $K$ est cartésien, le système l’est également.
\end{lemma}

\begin{proof}
Cela résulte immédiatement des définitions.
\end{proof}

\begin{lemma}
\label{lemma-construct-cartesian-objects-derived-modules}
Dans la situation \ref{situation-simplicial-site}, soit $\mathcal{O}$ un faisceau d'anneaux
sur $\mathcal{C}_{total}$.
Supposons donnés $K_0 \in D(\mathcal{O}_0)$ et un isomorphisme
$$
\alpha :
L(f_{\delta_1^1})^*K_0
\longrightarrow
L(f_{\delta_0^1})^*K_0
$$
satisfaisant à la condition de cocycle. Posons $\tau^n_i : [0] \to [n]$, $0 \mapsto i$
et posons $K_n = Lf_{\tau^n_n}^*K_0$. Les objets $K_n$ sont les termes d'un
système simplicial cartésien de la catégorie dérivée de modules.
\end{lemma}

\begin{proof}
Comparer avec les lemmes \ref{lemma-construct-cartesian-objects-derived}
et \ref{lemma-characterize-cartesian} et sa preuve
(également pour voir la condition du cocycle énoncée).
La construction est analogue à la construction discutée dans
Descente, section \ref{descent-section-descent-modules}, à laquelle nous empruntons
la notation $\tau^n_i : [0] \to [n]$, $0 \mapsto i$ et
$\tau^n_{ij} : [1] \to [n]$, $0 \mapsto i$, $1 \mapsto j$.
Pour $\varphi : [n] \to [m]$, définissons
$K_\varphi : L(f_\varphi)^*K_n \to K_m$
à l'aide du diagramme
$$
\xymatrix{
L(f_\varphi)^*K_n \ar@{=}[r] &
L(f_\varphi)^* L(f_{\tau^n_n})^*K_0 \ar@{=}[r] &
L(f_{\tau^m_{\varphi(n)}})^*K_0 \ar@{=}[r] &
L(f_{\tau^m_{\varphi(n)m}})^* L(f_{\delta^1_1})^*K_0
\ar[d]_{L(f_{\tau^m_{\varphi(n)m}})^*\alpha} \\
&
K_m \ar@{=}[r] &
L(f_{\tau^m_m})^*K_0 \ar@{=}[r] &
L(f_{\tau^m_{\varphi(n)m}})^* L(f_{\delta^1_0})^*K_0
}
$$
Nous laissons au lecteur le soin de vérifier que la condition de cocycle
entraîne la compatibilité des morphismes avec la composition (dans leurs catégories dérivées
respectives) et définit donc un système simplicial
de la catégorie dérivée de modules.
\end{proof}

\begin{lemma}
\label{lemma-modules-postnikov}
Dans la situation \ref{situation-simplicial-site}, soit $\mathcal{O}$
un faisceau d'anneaux sur $\mathcal{C}_{total}$. Soit $K$
un objet de $D(\mathcal{C}_{total})$. Posons
$$
X_n = (g_{n!}\mathcal{O}_n)
\otimes^\mathbf{L}_\mathcal{O} K
\quad\text{et}\quad
Y_n =
(g_{n!}\mathcal{O}_n \to \ldots \to g_{0!}\mathcal{O}_0)[-n]
\otimes^\mathbf{L}_\mathcal{O} K
$$
comme objets de $D(\mathcal{O})$, les morphismes étant
comme dans le lemme \ref{lemma-simplicial-resolution-Z-site}.
Avec les morphismes canoniques évidents $Y_n \to X_n$ et
$Y_0 \to Y_1[1] \to Y_2[2] \to \ldots$, on a
\begin{enumerate}
\item les triangles distingués $Y_n \to X_n \to Y_{n - 1} \to Y_n[1]$
définissent un système de Postnikov
(Catégories dérivées, définition \ref{derived-definition-postnikov-system})
pour $\ldots \to X_2 \to X_1 \to X_0$,
\item $K = \text{hocolim} Y_n[n]$ dans $D(\mathcal{O})$.
\end{enumerate}
\end{lemma}

\begin{proof}
D'abord, si $K = \mathcal{O}$, c'est la construction de
Catégories dérivées, exemple \ref{derived-example-key-postnikov},
appliquée au complexe
$$
\ldots \to
g_{2!}\mathcal{O}_2 \to
g_{1!}\mathcal{O}_1 \to
g_{0!}\mathcal{O}_0
$$
dans $\textit{Ab}(\mathcal{C}_{total})$ combiné au fait que
ce complexe représente $K = \mathcal{O}$ dans $D(\mathcal{C}_{total})$
par le lemme \ref{lemma-simplicial-resolution-ringed}.
Le cas général en découle, puisque le foncteur exact
$- \otimes^\mathbf{L}_\mathcal{O} K$ transforme les systèmes de Postnikov en
systèmes de Postnikov et
que $- \otimes^\mathbf{L}_\mathcal{O} K$ commute aux colimites homotopiques.
\end{proof}

\begin{lemma}
\label{lemma-nullity-cartesian-modules-derived}
Dans la situation \ref{situation-simplicial-site}, soit $\mathcal{O}$
un faisceau d'anneaux sur $\mathcal{C}_{total}$.
Soient $K, K' \in D(\mathcal{O})$.
Supposons que
\begin{enumerate}
\item $f_\varphi^{-1}\mathcal{O}_n \to \mathcal{O}_m$ est plat pour
$\varphi : [m] \to [n]$,
\item $K$ est cartésien,
\item $\Hom(K_i[i], K'_i) = 0$ pour $i > 0$ et
\item $\Hom(K_i[i + 1], K'_i) = 0$ pour $i \geq 0$.
\end{enumerate}
Alors tout morphisme $K \to K'$ induisant le morphisme nul $K_0 \to K'_0$ est nul.
\end{lemma}

\begin{proof}
La démonstration est identique à celle du
lemme \ref{lemma-nullity-cartesian-objects-derived}, à ceci près qu'elle utilise le
lemme \ref{lemma-modules-postnikov} au lieu du
lemme \ref{lemma-abelian-postnikov}.
\end{proof}

\begin{lemma}
\label{lemma-hom-cartesian-modules-derived}
Dans la situation \ref{situation-simplicial-site}, soit $\mathcal{O}$
un faisceau d'anneaux sur $\mathcal{C}_{total}$.
Soient $K, K' \in D(\mathcal{O})$.
Supposons que
\begin{enumerate}
\item $f_\varphi^{-1}\mathcal{O}_n \to \mathcal{O}_m$ est plat pour
$\varphi : [m] \to [n]$,
\item $K$ est cartésien,
\item $\Hom(K_i[i - 1], K'_i) = 0$ pour $i > 1$.
\end{enumerate}
Alors tout morphisme $\{K_n \to K'_n\}$ entre les systèmes simpliciaux associés
à $K$ et $K'$ provient d'un morphisme $K \to K'$ dans $D(\mathcal{O})$.
\end{lemma}

\begin{proof}
La démonstration est identique à celle du
lemme \ref{lemma-hom-cartesian-objects-derived}, à ceci près qu'elle utilise le
lemme \ref{lemma-modules-postnikov} au lieu du
lemme \ref{lemma-abelian-postnikov}.
\end{proof}

\begin{lemma}
\label{lemma-cartesian-module-derived-from-simplicial}
Dans la situation \ref{situation-simplicial-site}, soit $\mathcal{O}$
un faisceau d'anneaux sur $\mathcal{C}_{total}$. Soit
$(K_n, K_\varphi)$ un système simplicial de la catégorie dérivée
de modules. Supposons que
\begin{enumerate}
\item $f_\varphi^{-1}\mathcal{O}_n \to \mathcal{O}_m$ est plat pour
$\varphi : [m] \to [n]$,
\item $(K_n, K_\varphi)$ est cartésien,
\item $\Hom(K_i[t], K_i) = 0$ pour $i \geq 0$ et $t > 0$.
\end{enumerate}
Il existe alors un objet cartésien $K$ de $D(\mathcal{O})$
dont le système simplicial associé est isomorphe à $(K_n, K_\varphi)$.
\end{lemma}

\begin{proof}
La démonstration est identique à celle du
lemme \ref{lemma-cartesian-object-derived-from-simplicial},
avec les modifications suivantes
\begin{enumerate}
\item utiliser $g_n^* = Lg_n^*$ partout au lieu de $g_n^{-1}$,
\item utiliser $f_\varphi^* = Lf_\varphi^*$ partout au lieu de $f_\varphi^{-1}$,
\item se référer au lemme \ref{lemma-simplicial-resolution-ringed}
au lieu du lemme \ref{lemma-simplicial-resolution-Z-site},
\item dans la construction de $Y'_{m, n}$ utiliser
$\mathcal{O}_m$ au lieu de $\mathbf{Z}$,
\item comparer avec la preuve du lemme \ref{lemma-modules-postnikov}
plutôt que la preuve du lemme \ref{lemma-abelian-postnikov}.
\end{enumerate}
Ceci termine la preuve.
\end{proof}







\section{Le site associé à un objet semi-représentable}
\label{section-semi-representable}

\noindent
Soit $\mathcal{C}$ un site. Rappelons qu'un {\it objet semi-représentable}
de $\mathcal{C}$ est simplement une famille $\{U_i\}_{i \in I}$
d'objets de $\mathcal{C}$. Un
{\it morphisme $\{U_i\}_{i \in I} \to \{V_j\}_{j \in J}$ d'objets
semi-représentables} est donné par une application $\alpha : I \to J$
et pour tout $i \in I$ un morphisme $f_i : U_i \to V_{\alpha(i)}$
de $\mathcal{C}$.
La catégorie d'objets semi-représentables de $\mathcal{C}$
est notée $\text{SR}(\mathcal{C})$.
Voir Hyperrecouvrements, définition \ref{hypercovering-definition-SR},
et la section qui la contient pour davantage de précisions.

\medskip\noindent
Pour un objet semi-représentable $K = \{U_i\}_{i \in I}$ de $\mathcal{C}$,
on pose
$$
\mathcal{C}/K = \coprod\nolimits_{i \in I} \mathcal{C}/U_i
$$
l'union disjointe des localisations de $\mathcal{C}$ en $U_i$.
Cette catégorie porte une structure naturelle de site, dont les recouvrements sont
hérités de ceux des localisations $\mathcal{C}/U_i$.
Le site $\mathcal{C}/K$ s'appelle
{\it localisation de $\mathcal{C}$ en $K$}.
Observons qu'un faisceau sur $\mathcal{C}/K$ est la même chose qu'
une famille de faisceaux $\mathcal{F}_i$ sur $\mathcal{C}/U_i$, c'est-à-dire
$$
\Sh(\mathcal{C}/K) = \prod\nolimits_{i \in I} \Sh(\mathcal{C}/U_i)
$$
Cette observation est parfois utile pour comprendre la situation.

\medskip\noindent
Soient $\mathcal{C}$ un site et $K = \{U_i\}_{i \in I}$ un objet de
$\text{SR}(\mathcal{C})$. Il existe un foncteur de localisation
continu et cocontinu $j : \mathcal{C}/K \to \mathcal{C}$ qui est
le produit des foncteurs de localisation
$j_i : \mathcal{C}/V_i \to \mathcal{C}$.
On obtient les foncteurs $j_!$, $j^{-1}$, $j_*$ exactement
comme dans Sites, section \ref{sites-section-localize}.
En termes de la décomposition en produit
$\Sh(\mathcal{C}/K) = \prod\nolimits_{i \in I} \Sh(\mathcal{C}/U_i)$
nous avons
$$
\begin{matrix}
j_! & : &
(\mathcal{F}_i)_{i \in I} &
\longmapsto &
\coprod j_{i, !}\mathcal{F}_i \\
j^{-1} & : &
\mathcal{G} &
\longmapsto &
(j_i^{-1}\mathcal{G})_{i \in I} \\
j_* & : &
(\mathcal{F}_i)_{i \in I} &
\longmapsto &
\prod j_{i, *}\mathcal{F}_i
\end{matrix}
$$
ce que le lecteur vérifie aisément.

\medskip\noindent
Soit $f : K \to L$ un morphisme de $\text{SR}(\mathcal{C})$.
On obtient alors un foncteur continu et cocontinu
$$
v : \mathcal{C}/K \longrightarrow \mathcal{C}/L
$$
en appliquant aux composantes la construction de Sites, lemme \ref{sites-lemma-relocalize}.
aux composantes. Plus précisément, supposons $f = (\alpha, f_i)$
où $K = \{U_i\}_{i \in I}$, $L = \{V_j\}_{j \in J}$, $\alpha : I \to J$,
et $f_i : U_i \to V_{\alpha(i)}$. Le foncteur $v$ envoie la composante
$\mathcal{C}/U_i$ dans la composante $\mathcal{C}/V_{\alpha(i)}$
par la construction du lemme précité. En particulier,
on obtient un morphisme
$$
f : \Sh(\mathcal{C}/K) \to \Sh(\mathcal{C}/L)
$$
de topoi. En termes des décompositions en produits
$\Sh(\mathcal{C}/K) = \prod\nolimits_{i \in I} \Sh(\mathcal{C}/U_i)$ et
$\Sh(\mathcal{C}/L) = \prod\nolimits_{j \in J} \Sh(\mathcal{C}/V_j)$
le lecteur vérifiera que
$$
\begin{matrix}
f_! & : &
(\mathcal{F}_i)_{i \in I} &
\longmapsto &
(\coprod\nolimits_{i \in I, \alpha(i) = j} f_{i, !}\mathcal{F}_i)_{j \in J} \\
f^{-1} & : &
(\mathcal{G}_j)_{j \in J} &
\longmapsto &
(f_i^{-1}\mathcal{G}_{\alpha(i)})_{i \in I} \\
f_* & : &
(\mathcal{F}_i)_{i \in I} &
\longmapsto &
(\prod\nolimits_{i \in I, \alpha(i) = j} f_{i, *}\mathcal{F}_i)_{j \in J}
\end{matrix}
$$
où $f_i : \Sh(\mathcal{C}/U_i) \to \Sh(\mathcal{C}/V_{\alpha(i)})$
est le morphisme associé au foncteur de localisation
$\mathcal{C}/U_i \to \mathcal{C}/V_{\alpha(i)}$ correspondant à
$f_i : U_i \to V_{\alpha(i)}$.

\begin{lemma}
\label{lemma-has-P}
Soit $\mathcal{C}$ un site.
\begin{enumerate}
\item Pour $K$ dans $\text{SR}(\mathcal{C})$, le foncteur
$j : \mathcal{C}/K \to \mathcal{C}$ est continu,
cocontinu et possède la propriété P de
Sites, remarque \ref{sites-remark-cartesian-cocontinuous}.
\item Pour $f : K \to L$ dans $\text{SR}(\mathcal{C})$,
le foncteur $v : \mathcal{C}/K \to \mathcal{C}/L$ (voir ci-dessus)
est continu, cocontinu et possède la propriété P de
Sites, remarque \ref{sites-remark-cartesian-cocontinuous}.
\end{enumerate}
\end{lemma}

\begin{proof}
Pour (2). Avec les notations de la discussion précédant le lemme,
les foncteurs de localisation $\mathcal{C}/U_i \to \mathcal{C}/V_{\alpha(i)}$
sont continus et cocontinus par
Sites, section \ref{sites-section-localize}
et possèdent la propriété $P$ d'après
Sites, remarque \ref{sites-remark-localization-cartesian-cocontinuous}.
On en déduit formellement que $v$ est continu, cocontinu et possède la propriété $P$.
Nous laissons au lecteur les détails, ainsi que la preuve de (1).
\end{proof}

\begin{lemma}
\label{lemma-push-pull-localization}
Soit $\mathcal{C}$ un site et soit $K$ dans $\text{SR}(\mathcal{C})$.
Pour $\mathcal{F}$ dans $\Sh(\mathcal{C})$, on a
$$
j_*j^{-1}\mathcal{F} = \SheafHom(F(K)^\#, \mathcal{F})
$$
où $F$ est comme dans
Hyperrecouvrements, définition \ref{hypercovering-definition-SR-F}.
\end{lemma}

\begin{proof}
Écrivons $K = \{U_i\}_{i \in I}$.
En utilisant la description des foncteurs $j^{-1}$ et $j_*$
donnée ci-dessus, nous voyons que
$$
j_*j^{-1}\mathcal{F} =
\prod\nolimits_{i \in I} j_{i, *}(\mathcal{F}|_{\mathcal{C}/U_i}) =
\prod\nolimits_{i \in I} \SheafHom(h_{U_i}^\#, \mathcal{F})
$$
La seconde égalité résulte de Sites, lemme \ref{sites-lemma-hom-sheaf-hU}.
Puisque $F(K) = \coprod h_{U_i}$ dans $\textit{PSh}(\mathcal{C}$,
nous avons $F(K)^\# = \coprod h_{U_i}^\#$ dans $\Sh(\mathcal{C})$
et puisque $\SheafHom(-, \mathcal{F})$ transforme les coproduits en
produits (cela résulte immédiatement de la construction de
Sites, section \ref{sites-section-glueing-sheaves}), on conclut.
\end{proof}

\begin{lemma}
\label{lemma-localize-compare}
Soit $\mathcal{C}$ un site.
\begin{enumerate}
\item Pour $K$ dans $\text{SR}(\mathcal{C})$, le foncteur $j_!$
donne une équivalence $\Sh(\mathcal{C}/K) \to \Sh(\mathcal{C})/F(K)^\#$
où $F$ est comme dans
Hyperrecouvrements, définition \ref{hypercovering-definition-SR-F}.
\item Le foncteur $j^{-1} : \Sh(\mathcal{C}) \to \Sh(\mathcal{C}/K)$
correspond, via l'identification de (1), à
$\mathcal{F} \mapsto (\mathcal{F} \times F(K)^\# \to F(K)^\#)$.
\item Pour $f : K \to L$ dans $\text{SR}(\mathcal{C})$, le foncteur
$f^{-1}$ correspond, via les identifications de (1), au foncteur
$\Sh(\mathcal{C})/F(L)^\# \to \Sh(\mathcal{C})/F(K)^\#$,
$(\mathcal{G} \to F(L)^\#) \mapsto
(\mathcal{G} \times_{F(L)^\#} F(K)^\# \to F(K)^\#)$.
\end{enumerate}
\end{lemma}

\begin{proof}
Observons que, si $K = \{U_i\}_{i \in I}$, la catégorie
$\Sh(\mathcal{C}/K)$ se décompose en produit des catégories
$\Sh(\mathcal{C}/U_i)$. Observons que
$F(K)^\# = \coprod_{i \in I} h_{U_i}^\#$ (coproduit en faisceaux).
$\Sh(\mathcal{C})/F(K)^\#$ est donc le produit des catégories
$\Sh(\mathcal{C})/h_{U_i}^\#$.
Ainsi, (1) et (2) résultent des assertions
correspondantes pour chaque $i$; voir
Sites, lemmes \ref{sites-lemma-essential-image-j-shriek} et
\ref{sites-lemma-compute-j-shriek-restrict}.
De même, si $L = \{V_j\}_{j \in J}$ et si $f$ est donné
par $\alpha : I \to J$ et $f_i : U_i \to V_{\alpha(i)}$,
alors on peut appliquer
Sites, lemme \ref{sites-lemma-relocalize-explicit}
à chacun des morphismes de relocalisation
$\mathcal{C}/U_i \to \mathcal{C}/V_{\alpha(i)}$
pour obtenir (3).
\end{proof}

\begin{lemma}
\label{lemma-localize-injective}
Soit $\mathcal{C}$ un site et soit $K$ dans $\text{SR}(\mathcal{C})$.
Le foncteur $j^{-1}$ envoie les faisceaux abéliens injectifs sur des faisceaux
abéliens injectifs. De même, le foncteur $j^{-1}$ envoie les complexes K-injectifs
de faisceaux abéliens sur des complexes K-injectifs de faisceaux
abéliens.
\end{lemma}

\begin{proof}
La première assertion est la généralisation naturelle du lemme
de Cohomologie sur les sites
\ref{sites-cohomology-lemma-cohomology-of-open}
aux objets semi-représentables.
Elle découle en fait de ce lemme
par la décomposition en produit de $\Sh(\mathcal{C}/K)$
et la description du foncteur $j^{-1}$ donnée ci-dessus.
La seconde assertion est la généralisation naturelle du lemme
de Cohomologie sur les sites
\ref{sites-cohomology-lemma-restrict-K-injective-to-open}
et s'en déduit par la décomposition du topos en produit.

\medskip\noindent
Autre démonstration. Puisque $j$ induit une localisation de topos d'après la
partie (1) du lemme \ref{lemma-localize-compare},
on le déduit aussi immédiatement de
Cohomologie sur les sites, lemmes \ref{sites-cohomology-lemma-cohomology-of-open}
et \ref{sites-cohomology-lemma-restrict-K-injective-to-open}
après agrandissement du site; comparer avec la preuve de
Cohomologie sur les sites, lemme
\ref{sites-cohomology-lemma-cohomology-on-sheaf-sets}
dans le cas des faisceaux injectifs.
\end{proof}

\begin{remark}[Variante au-dessus d'un objet]
\label{remark-semi-representable-over-object}
Soit $\mathcal{C}$ un site et soit $X \in \Ob(\mathcal{C})$.
La catégorie $\text{SR}(\mathcal{C}, X)$
des {\it objets semi-représentables au-dessus de $X$}
est définie par la formule
$\text{SR}(\mathcal{C}, X) = \text{SR}(\mathcal{C}/X)$.
Voir Hyperrecouvrements, définition \ref{hypercovering-definition-SR}.
Nous pouvons donc appliquer la discussion ci-dessus au site
$\mathcal{C}/X$. En résumé, les constructions ci-dessus donnent
\begin{enumerate}
\item un site $\mathcal{C}/K$ pour $K$ dans $\text{SR}(\mathcal{C}, X)$,
\item une décomposition
$\Sh(\mathcal{C}/K) = \prod \Sh(\mathcal{C}/U_i)$ si $K = \{U_i/X\}$,
\item un foncteur de localisation $j : \mathcal{C}/K \to \mathcal{C}/X$,
\item un morphisme $f : \Sh(\mathcal{C}/K) \to \Sh(\mathcal{C}/L)$
pour $f : K \to L$ dans $\text{SR}(\mathcal{C}, X)$.
\end{enumerate}
Tous les résultats de cette section sont valables dans cette situation en remplaçant partout
$\mathcal{C}$ par $\mathcal{C}/X$.
\end{remark}

\begin{remark}[Variante annelée]
\label{remark-semi-representable-ringed}
Soit $\mathcal{C}$ un site et soit $\mathcal{O}_\mathcal{C}$
un faisceau d'anneaux sur $\mathcal{C}$. Alors, pour tout
objet semi-représentable $K$ de $\mathcal{C}$, le site
$\mathcal{C}/K$ est annelé et muni du faisceau d'anneaux
$\mathcal{O}_K = j^{-1}\mathcal{O}_\mathcal{C}$.
Les constructions ci-dessus donnent
\begin{enumerate}
\item un site annelé $(\mathcal{C}/K, \mathcal{O}_K)$
pour $K$ dans $\text{SR}(\mathcal{C})$,
\item une décomposition
$\textit{Mod}(\mathcal{O}_K) =
\prod \textit{Mod}(\mathcal{O}_{U_i})$ si $K = \{U_i\}$,
\item un morphisme de localisation
$j : (\Sh(\mathcal{C}/K), \mathcal{O}_K) \to
(\Sh(\mathcal{C}), \mathcal{O}_\mathcal{C})$
de topoi annelés,
\item un morphisme
$f : (\Sh(\mathcal{C}/K), \mathcal{O}_K) \to
(\Sh(\mathcal{C}/L), \mathcal{O}_L)$ de topoi annelés
pour $f : K \to L$ dans $\text{SR}(\mathcal{C})$.
\end{enumerate}
La plupart des résultats ci-dessus sont valables dans ce cadre. Par exemple, le foncteur
$j^*$ a un adjoint exact à gauche
$$
j_! : \textit{Mod}(\mathcal{O}_K) \to \textit{Mod}(\mathcal{O}_\mathcal{C}),
$$
qui, en termes de la décomposition en produit donnée en (2), envoie
$(\mathcal{F}_i)_{i \in I}$ à $\bigoplus j_{i, !}\mathcal{F}_i$.
De même, étant donné $f : K \to L$ comme ci-dessus, le foncteur $f^*$ a
un adjoint exact à gauche
$f_! : \textit{Mod}(\mathcal{O}_K) \to \textit{Mod}(\mathcal{O}_L)$.
Ainsi, les foncteurs $j^*$ et $f^*$ sont exacts, c'est-à-dire que
$j$ et $f$ sont des morphismes plats de topos annelés (ce qui découle aussi
des égalités $\mathcal{O}_K = j^{-1}\mathcal{O}_\mathcal{C}$
et $\mathcal{O}_K = f^{-1}\mathcal{O}_L$).
\end{remark}

\begin{remark}[Variante annelée au-dessus d'un objet]
\label{remark-semi-representable-ringed-over-object}
Soit $\mathcal{C}$ un site. Soit $\mathcal{O}_\mathcal{C}$
un faisceau d'anneaux sur $\mathcal{C}$. Soit $X \in \Ob(\mathcal{C})$
et notons $\mathcal{O}_X = \mathcal{O}_\mathcal{C}|_{\mathcal{C}/U}$.
On peut alors combiner les constructions données dans les
remarques \ref{remark-semi-representable-over-object}
et \ref{remark-semi-representable-ringed} pour obtenir
\begin{enumerate}
\item un site annelé $(\mathcal{C}/K, \mathcal{O}_K)$
pour $K$ dans $\text{SR}(\mathcal{C}, X)$,
\item une décomposition
$\textit{Mod}(\mathcal{O}_K) =
\prod \textit{Mod}(\mathcal{O}_{U_i})$ si $K = \{U_i\}$,
\item un morphisme de localisation
$j : (\Sh(\mathcal{C}/K), \mathcal{O}_K) \to
(\Sh(\mathcal{C}/X), \mathcal{O}_X)$
de topoi annelés,
\item un morphisme
$f : (\Sh(\mathcal{C}/K), \mathcal{O}_K) \to
(\Sh(\mathcal{C}/L), \mathcal{O}_L)$ de topoi annelés
pour $f : K \to L$ dans $\text{SR}(\mathcal{C}, X)$.
\end{enumerate}
Bien entendu, tous les résultats mentionnés dans
la remarque \ref{remark-semi-representable-ringed}
sont également valables dans ce cadre.
\end{remark}





\section{Le site associé à un objet semi-représentable simplicial}
\label{section-simplicial-semi-representable}

\noindent
Soient $\mathcal{C}$ un site et $K$ un objet simplicial de
$\text{SR}(\mathcal{C})$. Comme d'habitude, posons $K_n = K([n])$ et notons
$K(\varphi) : K_n \to K_m$ le morphisme associé à $\varphi : [m] \to [n]$.
D'après la construction de la
section \ref{section-semi-representable}, on obtient un objet simplicial
$n \mapsto \mathcal{C}/K_n$ dans la catégorie dont les objets sont des sites et
dont les morphismes sont des foncteurs cocontinus. Autrement dit, on obtient
un diagramme comme dans le cas (B) de la section \ref{section-simplicial-sites}.
Les foncteurs satisfont la propriété P par le lemme \ref{lemma-has-P}.
Nous pouvons donc appliquer le lemme
\ref{lemma-simplicial-cocontinuous-site}
pour obtenir un site $(\mathcal{C}/K)_{total}$.

\medskip\noindent
Nous pouvons décrire explicitement le site $(\mathcal{C}/K)_{total}$ comme suit.
Écrivons $K_n = \{U_{n, i}\}_{i \in I_n}$. Pour $\varphi : [m] \to [n]$
le morphisme $K(\varphi) : K_n \to K_m$ est donné par une application
$\alpha(\varphi) : I_n \to I_m$ et les morphismes
$f_{\varphi, i} : U_{n, i} \to U_{m, \alpha(\varphi)(i)}$ pour $i \in I_n$.
On a alors
\begin{enumerate}
\item un objet de $(\mathcal{C}/K)_{total}$
correspond à un objet $(U/U_{n, i})$
de $\mathcal{C}/U_{n, i}$ pour un certain $n$ et un certain $i \in I_n$;
\item un morphisme entre $U/U_{n, i}$ et $V/U_{m, j}$
est une paire $(\varphi, f)$ où $\varphi : [m] \to [n]$,
$j = \alpha(\varphi)(i)$ et $f : U \to V$ est un morphisme de
$\mathcal{C}$ tel que
$$
\vcenter{
\xymatrix{
U \ar[r]_f \ar[d] & V \ar[d] \\
U_{n, i} \ar[r]^-{f_{\varphi, i}} &
U_{m, j}
}
}
$$
est commutatif, et
\item les recouvrements de l'objet $U/U_{n, i}$ sont construits
à partir d'un recouvrement $\{f_j : U_j \to U\}$ dans $\mathcal{C}$, en déclarant
$\{(\text{id}, f_j) : U_j/U_{n, i} \to U/U_{n, i}\}$
recouvrement dans $(\mathcal{C}/K)_{total}$.
\end{enumerate}
Toute la théorie générale développée pour les sites simpliciaux s'applique à
$(\mathcal{C}/K)_{total}$. Observons que le foncteur oublieux évident
$$
j_{total} : (\mathcal{C}/K)_{total} \longrightarrow \mathcal{C}
$$
est continu et cocontinu. Le morphisme de topoi
associé provient d'une augmentation évidente.

\begin{lemma}
\label{lemma-augmentation-simplicial-semi-representable}
Soient $\mathcal{C}$ un site et $K$ un objet simplicial de
$\text{SR}(\mathcal{C})$. Le foncteur de localisation
$j_0 : \mathcal{C}/K_0 \to \mathcal{C}$ définit une augmentation
$a_0 : \Sh(\mathcal{C}/K_0) \to \Sh(\mathcal{C})$, comme dans le cas (B) de
la remarque \ref{remark-augmentation-site}.
Les morphismes de topoi correspondants
$$
a_n : \Sh(\mathcal{C}/K_n) \longrightarrow \Sh(\mathcal{C}),\quad
a : \Sh((\mathcal{C}/K)_{total}) \longrightarrow \Sh(\mathcal{C})
$$
du lemme \ref{lemma-augmentation-site}
sont égaux aux morphismes de topos associés aux
foncteurs de localisation continus et cocontinus
$j_n : \mathcal{C}/K_n \to \mathcal{C}$ et
$j_{total} : (\mathcal{C}/K)_{total} \to \mathcal{C}$.
\end{lemma}

\begin{proof}
Cela résulte immédiatement des définitions.
Voir notamment la note de bas de page dans la preuve du lemme
\ref{lemma-augmentation-site}
pour la relation entre $a$ et $j_{total}$.
\end{proof}

\begin{lemma}
\label{lemma-comparison}
Avec les hypothèses et notations du
lemme \ref{lemma-augmentation-simplicial-semi-representable},
nous avons les propriétés suivantes :
\begin{enumerate}
\item il existe un foncteur
$a^{Sh}_! : \Sh((\mathcal{C}/K)_{total}) \to \Sh(\mathcal{C})$
adjoint à gauche de $a^{-1} : \Sh(\mathcal{C}) \to \Sh((\mathcal{C}/K)_{total})$;
\item il existe un foncteur
$a_! : \textit{Ab}((\mathcal{C}/K)_{total}) \to \textit{Ab}(\mathcal{C})$
adjoint à gauche de
$a^{-1} : \textit{Ab}(\mathcal{C}) \to \textit{Ab}((\mathcal{C}/K)_{total})$,
\item le foncteur $a^{-1}$ associe à un faisceau
$\mathcal{F}$ dans $\Sh(\mathcal{C})$ le faisceau sur $(\mathcal{C}/K)_{total}$
qui en degré $n$ est égal à $a_n^{-1}\mathcal{F}$,
\item le foncteur $a_*$ associe à un faisceau abélien $\mathcal{G}$ sur
$\textit{Ab}((\mathcal{C}/K)_{total})$ l'égaliseur des deux morphismes
$j_{0, *}\mathcal{G}_0 \to j_{1, *}\mathcal{G}_1$,
\end{enumerate}
\end{lemma}

\begin{proof}
Les parties (3) et (4) valent pour toute augmentation d'un
site simplicial; voir le lemme \ref{lemma-augmentation-site}.
Les parties (1) et (2) suivent de ce que $j_{total}$ est continu et cocontinu.
Le foncteur $a^{Sh}_!$ est construit dans
Sites, lemme \ref{sites-lemma-when-shriek},
et le foncteur $a_!$ est construit dans Modules
sur les sites, lemme
\ref{sites-modules-lemma-g-shriek-adjoint}.
\end{proof}

\begin{lemma}
\label{lemma-sanity-check-simplicial-semi-representable}
Soient $\mathcal{C}$ un site et $K$ un objet simplicial de
$\text{SR}(\mathcal{C})$. Soit $U/U_{n, i}$ un objet de
$\mathcal{C}/K_n$. Soit
$\mathcal{F} \in \textit{Ab}((\mathcal{C}/K)_{total})$.
Alors
$$
H^p(U, \mathcal{F}) = H^p(U, \mathcal{F}_{n, i})
$$
où
\begin{enumerate}
\item au membre de gauche, $U$ est considéré comme un objet de
$\mathcal{C}_{total}$ et
\item au membre de droite, $\mathcal{F}_{n, i}$ est la $i$-ième composante
du faisceau $\mathcal{F}_n$ sur $\mathcal{C}/K_n$
dans la décomposition $\Sh(\mathcal{C}/K_n) = \prod \Sh(\mathcal{C}/U_{n, i})$
de la section \ref{section-semi-representable}.
\end{enumerate}
\end{lemma}

\begin{proof}
Cela découle immédiatement du lemme \ref{lemma-sanity-check}
et des décompositions en produits de la section \ref{section-semi-representable}.
\end{proof}

\begin{remark}[Variante au-dessus d'un objet]
\label{remark-augmentation-over-object}
Soient $\mathcal{C}$ un site et $X \in \Ob(\mathcal{C})$.
Rappelons que nous avons une catégorie
$\text{SR}(\mathcal{C}, X) = \text{SR}(\mathcal{C}/X)$
d'objets semi-représentables au-dessus de $X$;
voir la remarque \ref{remark-semi-representable-over-object}.
Nous pouvons appliquer la discussion ci-dessus au site
$\mathcal{C}/X$. En bref, les constructions ci-dessus donnent
\begin{enumerate}
\item un site $(\mathcal{C}/K)_{total}$ pour un objet simplicial $K$
de $\text{SR}(\mathcal{C}, X)$,
\item un foncteur de localisation
$j_{total} : (\mathcal{C}/K)_{total} \to \mathcal{C}/X$,
\item des foncteurs de localisation $j_n : \mathcal{C}/K_n \to \mathcal{C}/X$,
\item un morphisme de topoi
$a : \Sh((\mathcal{C}/K)_{total}) \to \Sh(\mathcal{C}/X)$,
\item des morphismes de topoi
$a_n : \Sh(\mathcal{C}/K_n) \to \Sh(\mathcal{C}/X)$,
\item un foncteur
$a^{Sh}_! : \Sh((\mathcal{C}/K)_{total}) \to \Sh(\mathcal{C}/X)$
adjoint à gauche de $a^{-1}$, et
\item un foncteur
$a_! : \textit{Ab}((\mathcal{C}/K)_{total}) \to \textit{Ab}(\mathcal{C}/X)$
adjoint à gauche de $a^{-1}$.
\end{enumerate}
Tous les résultats de cette section valent dans ce cadre.
Pour le démontrer, on remplace partout
le site $\mathcal{C}$ par $\mathcal{C}/X$.
\end{remark}

\begin{remark}[Variante annelée]
\label{remark-augmentation-ringed}
Soient $\mathcal{C}$ un site et $\mathcal{O}_\mathcal{C}$ un faisceau d'anneaux.
Étant donné un objet semi-représentable simplicial $K$ de $\mathcal{C}$,
nous définissons $\mathcal{O} = a^{-1}\mathcal{O}_\mathcal{C}$, où $a$
est comme dans les lemmes \ref{lemma-augmentation-simplicial-semi-representable} et
\ref{lemma-comparison}.
Les constructions ci-dessus, en tenant compte des faisceaux d'anneaux
comme dans la remarque \ref{remark-semi-representable-ringed}, donnent
\begin{enumerate}
\item un site annelé $((\mathcal{C}/K)_{total}, \mathcal{O})$
pour un objet $K$ simplicial de $\text{SR}(\mathcal{C})$,
\item un morphisme de topoi annelés
$a : (\Sh((\mathcal{C}/K)_{total}), \mathcal{O}) \to
(\Sh(\mathcal{C}), \mathcal{O}_\mathcal{C})$,
\item des morphismes de topoi annelés
$a_n : (\Sh(\mathcal{C}/K_n), \mathcal{O}_n) \to
(\Sh(\mathcal{C}), \mathcal{O}_\mathcal{C})$,
\item un foncteur
$a_! : \textit{Mod}(\mathcal{O}) \to \textit{Mod}(\mathcal{O}_\mathcal{C})$
adjoint à gauche de $a^*$.
\end{enumerate}
Le foncteur $a_!$ existe (mais en général il n'est pas exact)
car $a^{-1}\mathcal{O}_\mathcal{C} = \mathcal{O}$
et l'on peut remplacer l'emploi du
lemme \ref{sites-modules-lemma-g-shriek-adjoint} de Modules sur les sites
dans la preuve du lemme \ref{lemma-comparison}
par Modules sur les sites, lemme \ref{sites-modules-lemma-lower-shriek-modules}.
Comme indiqué dans la remarque \ref{remark-semi-representable-ringed}
il existe des foncteurs exacts
$a_{n!} : \textit{Mod}(\mathcal{O}_n) \to
\textit{Mod}(\mathcal{O}_\mathcal{C})$
adjoints à gauche de $a_n^*$. Par conséquent, les morphismes $a$ et $a_n$ sont plats.
La remarque \ref{remark-semi-representable-ringed}
implique que le morphisme de topoi annelés
$f_\varphi : (\Sh(\mathcal{C}/K_n), \mathcal{O}_n) \to
(\Sh(\mathcal{C}/K_m), \mathcal{O}_m)$
pour $\varphi : [m] \to [n]$ est plat, et il existe un foncteur exact
$f_{\varphi !} : \textit{Mod}(\mathcal{O}_n) \to \textit{Mod}(\mathcal{O}_m)$
adjoint à gauche de $f_\varphi^*$. Il en résulte que, pour
le morphisme plat de topos annelés
$g_n : (\Sh(\mathcal{C}/K_n), \mathcal{O}_n) \to
(\Sh((\mathcal{C}/K)_{total}), \mathcal{O})$
le foncteur $g_{n!} : \textit{Mod}(\mathcal{O}_n) \to
\textit{Mod}(\mathcal{O})$, adjoint à gauche de $g_n^*$, est exact; voir
le lemme \ref{lemma-exactness-g-shriek-modules}.
\end{remark}

\begin{remark}[Variante annelée au-dessus d'un objet]
\label{remark-augmentation-ringed-over-object}
Soient $\mathcal{C}$ un site et $\mathcal{O}_\mathcal{C}$ un faisceau d'anneaux.
Soit $X \in \Ob(\mathcal{C})$ et désignons
$\mathcal{O}_X = \mathcal{O}_\mathcal{C}|_{\mathcal{C}/X}$.
On peut alors combiner les constructions données dans les
remarques \ref{remark-augmentation-over-object} et
\ref{remark-augmentation-ringed}
pour obtenir
\begin{enumerate}
\item un site annelé $((\mathcal{C}/K)_{total}, \mathcal{O})$
pour un objet simplicial $K$ de $\text{SR}(\mathcal{C}, X)$,
\item un morphisme de topoi annelés
$a : (\Sh((\mathcal{C}/K)_{total}), \mathcal{O}) \to
(\Sh(\mathcal{C}/X), \mathcal{O}_X)$,
\item des morphismes de topoi annelés
$a_n : (\Sh(\mathcal{C}/K_n), \mathcal{O}_n) \to
(\Sh(\mathcal{C}/X), \mathcal{O}_X)$,
\item un foncteur
$a_! : \textit{Mod}(\mathcal{O}) \to \textit{Mod}(\mathcal{O}_X)$
adjoint à gauche de $a^*$.
\end{enumerate}
Bien entendu, tous les résultats mentionnés dans la remarque
\ref{remark-augmentation-ringed}
valent aussi dans ce cadre.
\end{remark}






\section{Descente cohomologique pour les hyperrecouvrements}
\label{section-cohomological-descent-hypercoverings}

\noindent
Soit $\mathcal{C}$ un site. Dans cette section, nous supposons que $\mathcal{C}$
admet des égalisateurs et des produits fibrés. Soit $K$ un hyperrecouvrement
au sens d'Hyperrecouvrements, définition
\ref{hypercovering-definition-hypercovering-variant}. Nous étudierons
l'augmentation
$$
a : \Sh((\mathcal{C}/K)_{total}) \longrightarrow \Sh(\mathcal{C})
$$
de la section \ref{section-simplicial-semi-representable}.

\begin{lemma}
\label{lemma-hypercovering-descent-sheaves}
Soit $\mathcal{C}$ un site admettant des égalisateurs et des produits fibrés.
Soit $K$ un hyperrecouvrement. Alors
\begin{enumerate}
\item $a^{-1} : \Sh(\mathcal{C}) \to \Sh((\mathcal{C}/K)_{total})$
est pleinement fidèle, d'image essentielle les faisceaux cartésiens d'ensembles,
\item $a^{-1} : \textit{Ab}(\mathcal{C}) \to
\textit{Ab}((\mathcal{C}/K)_{total})$
est pleinement fidèle, d'image essentielle les faisceaux cartésiens
de groupes abéliens.
\end{enumerate}
Dans les deux cas, $a_*$ est un foncteur quasi-inverse.
\end{lemma}

\begin{proof}
Le cas des faisceaux abéliens découle immédiatement du cas
de faisceaux d'ensembles, puisque le foncteur $a^{-1}$ commute aux produits.
Dans la suite de la preuve, nous travaillons avec des faisceaux d'ensembles.
Observons que $a^{-1}\mathcal{F}$ est cartésien pour
tout $\mathcal{F}$ dans $\Sh(\mathcal{C})$, d'après le
lemme \ref{lemma-augmentation-cartesian-module}.
Il suffit de montrer que le morphisme d'adjonction
$\mathcal{F} \to a_*a^{-1}\mathcal{F}$
est un isomorphisme pour tout $\mathcal{F}$ dans $\Sh(\mathcal{C})$
et que, pour tout faisceau cartésien
$\mathcal{G}$ sur $(\mathcal{C}/K)_{total}$,
le morphisme d'adjonction
$a^{-1}a_*\mathcal{G} \to \mathcal{G}$ est un isomorphisme.

\medskip\noindent
Soit $\mathcal{F}$ un faisceau sur $\mathcal{C}$.
Rappelons que $a_*a^{-1}\mathcal{F}$ est l'égaliseur
des deux morphismes $a_{0, *}a_0^{-1}\mathcal{F} \to a_{1, *}a_1^{-1}\mathcal{F}$;
voir le lemme \ref{lemma-comparison}.
D'après le lemme \ref{lemma-push-pull-localization},
$$
a_{0, *}a_0^{-1}\mathcal{F} = \SheafHom(F(K_0)^\#, \mathcal{F})
\quad\text{et}\quad
a_{1, *}a_1^{-1}\mathcal{F} = \SheafHom(F(K_1)^\#, \mathcal{F})
$$
D'autre part, on sait que
$$
\xymatrix{
F(K_1)^\# \ar@<1ex>[r] \ar@<-1ex>[r] &
F(K_0)^\# \ar[r] & \text{objet final }*\text{ de }\Sh(\mathcal{C})
}
$$
est un diagramme coégalisateur dans les faisceaux d'ensembles, par définition d'
un hyperrecouvrement. Il suffit donc de prouver
que $\SheafHom(-, \mathcal{F})$ transforme les coégaliseurs
en égaliseurs, ce qui résulte immédiatement de la construction
dans Sites, section \ref{sites-section-glueing-sheaves}.

\medskip\noindent
Soit $\mathcal{G}$ un faisceau cartésien sur $(\mathcal{C}/K)_{total}$.
Montrons que $\mathcal{G} = a^{-1}\mathcal{F}$ pour un certain faisceau
$\mathcal{F}$ sur $\mathcal{C}$. Cela achèvera la preuve, car
on aura alors $a^{-1}a_*\mathcal{G} = a^{-1}a_*a^{-1}\mathcal{F} =
a^{-1}\mathcal{F} = \mathcal{G}$ par le résultat du paragraphe précédent.
Posons $\mathcal{K}_n = F(K_n)^\#$ pour $n \geq 0$. On dispose alors de morphismes de faisceaux
$$
\xymatrix{
\mathcal{K}_2
\ar@<1ex>[r]
\ar@<0ex>[r]
\ar@<-1ex>[r]
&
\mathcal{K}_1
\ar@<0.5ex>[r]
\ar@<-0.5ex>[r]
&
\mathcal{K}_0
}
$$
provenant du fait que $K$ est un objet semi-représentable simplicial.
Le fait que $K$ soit un hyperrecouvrement signifie que
$$
\mathcal{K}_1 \to \mathcal{K}_0 \times \mathcal{K}_0
\quad\text{et}\quad
\mathcal{K}_2 \to
\left(\text{cosk}_1(
\xymatrix{
\mathcal{K}_1
\ar@<0.5ex>[r]
\ar@<-0.5ex>[r]
&
\mathcal{K}_0 \ar[l]
})\right)_2
$$
sont des morphismes surjectifs de faisceaux. En utilisant la description des faisceaux cartésiens sur
$(\mathcal{C}/K)_{total}$ donnée dans le lemme \ref{lemma-characterize-cartesian}
et en utilisant la description de $\Sh(\mathcal{C}/K_n)$ dans
le lemme \ref{lemma-localize-compare},
on voit que notre problème peut se formuler entièrement\footnote{Bien que
la formulation précise n'a pas d'importance, nous l'expliquons :
le problème consiste à montrer qu'étant donné un objet
$\mathcal{G}_0/\mathcal{K}_0$ de $\Sh(\mathcal{C})/\mathcal{K}_0$
et un isomorphisme
$$
\alpha :
\mathcal{G}_0 \times_{\mathcal{K}_0, \mathcal{K}(\delta^1_1)} \mathcal{K}_1 \to
\mathcal{G}_0 \times_{\mathcal{K}_0, \mathcal{K}(\delta^1_0)} \mathcal{K}_1
$$
sur $\mathcal{K}_1$ satisfaisant une condition de cocycle dans
$\Sh(\mathcal{C})/\mathcal{K}_2$, il existe
$\mathcal{F}$ dans $\Sh(\mathcal{C})$ et un isomorphisme
$\mathcal{F} \times \mathcal{K}_0 \to \mathcal{G}_0$ sur $\mathcal{K}_0$
compatible avec $\alpha$.} en termes de
\begin{enumerate}
\item le topos $\Sh(\mathcal{C})$, et
\item l'objet simplicial $\mathcal{K}$ de $\Sh(\mathcal{C})$
dont les termes sont $\mathcal{K}_n$.
\end{enumerate}
Ainsi, après avoir remplacé $\mathcal{C}$ par un autre site $\mathcal{C}'$
comme dans Sites, lemme \ref{sites-lemma-topos-good-site}, on peut supposer
$\mathcal{C}$ admet toutes les limites finies,
la topologie sur $\mathcal{C}$ est sous-canonique,
une famille $\{V_j \to V\}$ de morphismes de $\mathcal{C}$
est un recouvrement si et seulement si $\coprod h_{V_j} \to V$ est surjectif, et
il existe un objet simplicial $U$ de $\mathcal{C}$
tel que $\mathcal{K}_n = h_{U_n}$ comme faisceaux simpliciaux.
En transportant la situation par ces équivalences, on peut supposer
$K_n = \{U_n\}$ pour tout $n$.

\medskip\noindent
Soit $X$ l'objet final de $\mathcal{C}$.
Alors $\{U_0 \to X\}$ est un recouvrement,
$\{U_1 \to U_0 \times U_0\}$ est un recouvrement et
$\{U_2 \to (\text{cosk}_1 \text{sk}_1 U)_2\}$ est un recouvrement.
Notons $d^n_i : U_n \to U_{n - 1}$ et
$s^n_j : U_n \to U_{n + 1}$ les morphismes correspondant
à $\delta^n_i$ et $\sigma^n_j$, comme dans
Simplicial, définition \ref{simplicial-definition-face-degeneracy}.
Par abus de notation, pour un morphisme
$c : V \to W$ de $\mathcal{C}$, on note le morphisme de topos
$c : \Sh(\mathcal{C}/V) \to \Sh(\mathcal{C}/W)$ de la même lettre.
Le faisceau $\mathcal{G}$ est alors déterminé par un faisceau $\mathcal{G}_0$
sur $\mathcal{C}/U_0$ et un isomorphisme
$\alpha : (d^1_1)^{-1}\mathcal{G}_0 \to (d^1_0)^{-1}\mathcal{G}_0$
satisfaisant à la condition de cocycle sur $\mathcal{C}/U_2$
formulée dans le lemme \ref{lemma-characterize-cartesian}.
Puisque $\{U_2 \to (\text{cosk}_1 \text{sk}_1 U)_2\}$
est un recouvrement, le foncteur image inverse correspondant
sur les faisceaux est fidèle (nous omettons cette vérification élémentaire).
On peut donc remplacer $U$ par $\text{cosk}_1 \text{sk}_1 U$, car
cela remplace $U_2$ par $(\text{cosk}_1 \text{sk}_1 U)_2$ sans modifier
$U_1$ ni $U_0$. Alors
$$
(d^2_0, d^2_1, d^2_2) : U_2 \to U_1 \times U_1 \times U_1
$$
est un monomorphisme dont l'image sur les points à valeurs dans $T$ est
décrite dans Simplicial, lemme \ref{simplicial-lemma-work-out}.
En particulier, il existe un morphisme $c$ tel que le diagramme suivant soit commutatif :
$$
\xymatrix{
U_1 \times_{(d^1_1, d^1_0), U_0 \times U_0, (d^1_1, d^1_0)} U_1
\ar[d] \ar[rr]_c & & U_2 \ar[d] \\
U_1 \times U_1
\ar[rr]^{(\text{pr}_1, \text{pr}_2, s^0_0 \circ d^1_1 \circ \text{pr}_1)} & &
U_1 \times U_1 \times U_1
}
$$
car le parcours dans l'autre sens définit un point de $U_2$.
L'image inverse de la condition de cocycle vérifiée par $\alpha$ sur $U_2$
exprime que les images inverses de $\alpha$
par les deux projections sur
$U_1 \times_{(d^1_1, d^1_0), U_0 \times U_0, (d^1_1, d^1_0)} U_1$
coïncident, tandis que l'image inverse de $\alpha$ par
$s^0_0 \circ d^1_1 \circ \text{pr}_1$ est le morphisme identité
(en effet, l'image inverse de $\alpha$ par $s^0_0$ est l'identité).
D'après Sites, lemme \ref{sites-lemma-glue-maps},
on en déduit que $\alpha$ provient d'un isomorphisme
$$
\alpha' : \text{pr}_1^{-1}\mathcal{G}_0 \to \text{pr}_2^{-1}\mathcal{G}_0
$$
de faisceaux sur $\mathcal{C}/U_0 \times U_0$.
Enfin, le morphisme $U_2 \to U_0 \times U_0 \times U_0$
est surjectif sur les faisceaux associés, comme on le voit en utilisant la surjectivité
de $U_1 \to U_0 \times U_0$
et la description de $U_2$ indiquée ci-dessus. Par conséquent, $\alpha'$
satisfait la condition de cocycle sur $U_0 \times U_0 \times U_0$.
La preuve s'achève en appliquant
Sites, lemme \ref{sites-lemma-mapping-property-glue}
au recouvrement $\{U_0 \to X\}$.
\end{proof}

\begin{lemma}
\label{lemma-hypercovering-cech-complex}
Soit $\mathcal{C}$ un site admettant des égalisateurs et des produits fibrés.
Soit $K$ un hyperrecouvrement. Le complexe de Čech
du lemme \ref{lemma-augmentation-cech-complex}, associé à
$a^{-1}\mathcal{F}$
$$
a_{0, *}a_0^{-1}\mathcal{F} \to a_{1, *}a_1^{-1}\mathcal{F} \to
a_{2, *}a_2^{-1}\mathcal{F} \to \ldots
$$
s'identifie au complexe $\SheafHom(s(\mathbf{Z}_{F(K)}^\#), \mathcal{F})$.
Ici, $s(\mathbf{Z}_{F(K)}^\#)$ est le complexe défini dans
Hyperrecouvrements, définition \ref{hypercovering-definition-homology}.
\end{lemma}

\begin{proof}
D'après le lemme \ref{lemma-push-pull-localization}, on a
$$
a_{n, *}a_n^{-1}\mathcal{F} = \SheafHom'(F(K_n)^\#, \mathcal{F})
$$
où $\SheafHom'$ est comme dans Sites, section \ref{sites-section-glueing-sheaves}.
Les morphismes de bord du complexe du
lemme \ref{lemma-augmentation-cech-complex}
proviennent de la structure simpliciale.
L'égalité des complexes résulte donc
des identifications canoniques
$\SheafHom'(\mathcal{G}, \mathcal{F}) =
\SheafHom(\mathbf{Z}_\mathcal{G}, \mathcal{F})$ pour
$\mathcal{G}$ dans $\Sh(\mathcal{C})$.
\end{proof}

\begin{lemma}
\label{lemma-hypercovering-descent-bounded-abelian}
Soit $\mathcal{C}$ un site admettant des égalisateurs et des produits fibrés.
Soit $K$ un hyperrecouvrement. Pour
$E \in D(\mathcal{C})$, le morphisme
$$
E \longrightarrow Ra_*a^{-1}E
$$
est un isomorphisme.
\end{lemma}

\begin{proof}
D'abord, soit $\mathcal{I}$ un faisceau abélien injectif sur $\mathcal{C}$.
La suite spectrale du lemme
\ref{lemma-augmentation-spectral-sequence},
appliquée au faisceau $a^{-1}\mathcal{I}$, dégénère puisque
$(a^{-1}\mathcal{I})_p = a_p^{-1}\mathcal{I}$
est injectif par le lemme \ref{lemma-localize-injective}.
Ainsi le complexe
$$
a_{0, *}a_0^{-1}\mathcal{I} \to
a_{1, *}a_1^{-1}\mathcal{I} \to
a_{2, *}a_2^{-1}\mathcal{I} \to\ldots
$$
calcule $Ra_*a^{-1}\mathcal{I}$. Par
le lemme \ref{lemma-hypercovering-cech-complex},
ce complexe s'identifie à
$\SheafHom(s(\mathbf{Z}_{F(K)}^\#), \mathcal{I})$.
Puisque $K$ est un hyperrecouvrement,
$s(\mathbf{Z}_{F(K)}^\#)$ est exact en degrés $> 0$ par
Hyperrecouvrements, lemme \ref{hypercovering-lemma-acyclic-hypercover-sheaves}
appliqué au préfaisceau simplicial $F(K)$.
Puisque $\mathcal{I}$ est injectif, le foncteur $\SheafHom(-, \mathcal{I})$
est exact; on en déduit que
$\SheafHom(s(\mathbf{Z}_{F(K)}^\#), \mathcal{I})$
est exact en degrés positifs. Ainsi,
$R^pa_*a^{-1}\mathcal{I} = 0$ pour $p > 0$.
D'autre part, $\mathcal{I} = a_*a^{-1}\mathcal{I}$ d'après le
du lemme \ref{lemma-hypercovering-descent-sheaves}.

\medskip\noindent
Cas borné. Soit $E \in D^+(\mathcal{C})$.
Choisissons un complexe borné inférieurement $\mathcal{I}^\bullet$ d'injectifs
représentant $E$. D'après le premier paragraphe et le
lemme d'acyclicité de Leray
(Catégories dérivées, lemme \ref{derived-lemma-leray-acyclicity}),
$Ra_*a^{-1}\mathcal{I}^\bullet$
est calculé par le complexe
$a_*a^{-1}\mathcal{I}^\bullet = \mathcal{I}^\bullet$
ce qui prouve le lemme dans ce cas.

\medskip\noindent
Cas non borné. Nous invitons le lecteur à sauter ce passage, puisque l'argument
est le même que ci-dessus, sauf que nous utilisons la représentation explicite
par doubles complexes pour contourner les problèmes de convergence.
Soit $E \in D(\mathcal{C})$.
Pour montrer que $E \to Ra_*a^{-1}E$ est un isomorphisme,
il suffit de montrer, pour tout objet $U$ de $\mathcal{C}$, que
$$
R\Gamma(U, E) = R\Gamma(U, Ra_*a^{-1}E)
$$
Calculons les deux membres et montrons que le morphisme $E \to Ra_*a^{-1}E$
induit un isomorphisme. Choisissons un complexe K-injectif
$\mathcal{I}^\bullet$ représentant $E$, puis un quasi-isomorphisme
$a^{-1}\mathcal{I}^\bullet \to \mathcal{J}^\bullet$
pour un complexe K-injectif $\mathcal{J}^\bullet$ sur
$(\mathcal{C}/K)_{total}$. Nous avons
$$
R\Gamma(U, E) = R\Hom(\mathbf{Z}_U^\#, E)
$$
et
$$
R\Gamma(U, Ra_*a^{-1}E) = R\Hom(\mathbf{Z}_U^\#, Ra_*a^{-1}E) =
R\Hom(a^{-1}\mathbf{Z}_U^\#, a^{-1}E)
$$
D'après le lemme \ref{lemma-simplicial-resolution-augmentation},
nous avons un quasi-isomorphisme
$$
\Big(\ldots \to
g_{2!}(a_2^{-1}\mathbf{Z}_U^\#) \to
g_{1!}(a_1^{-1}\mathbf{Z}_U^\#) \to
g_{0!}(a_0^{-1}\mathbf{Z}_U^\#)\Big)
\longrightarrow
a^{-1}\mathbf{Z}_U^\#
$$
Par conséquent, $R\Hom(a^{-1}\mathbf{Z}_U^\#, a^{-1}E)$ est égal à
$$
R\Gamma((\mathcal{C}/K)_{total},
R\SheafHom(
\ldots \to
g_{2!}(a_2^{-1}\mathbf{Z}_U^\#) \to
g_{1!}(a_1^{-1}\mathbf{Z}_U^\#) \to
g_{0!}(a_0^{-1}\mathbf{Z}_U^\#),
\mathcal{J}^\bullet))
$$
D'après la construction de Cohomologie sur les sites, section
\ref{sites-cohomology-section-internal-hom}
et puisque $\mathcal{J}^\bullet$ est K-injectif, on voit que
l'objet précédent est représenté par le complexe de groupes abéliens dont les termes sont
$$
\prod\nolimits_{p + q = n}
\Hom(g_{p!}(a_p^{-1}\mathbf{Z}_U^\#), \mathcal{J}^q) =
\prod\nolimits_{p + q = n}
\Hom(a_p^{-1}\mathbf{Z}_U^\#, g_p^{-1}\mathcal{J}^q)
$$
Voir Cohomologie sur les sites, lemmes
\ref{sites-cohomology-lemma-RHom-into-K-injective} et
\ref{sites-cohomology-lemma-section-RHom-over-U} pour plus d’informations.
Ainsi, $R\Gamma(U, Ra_*a^{-1}E)$ est calculé par
le complexe total produit $\text{Tot}_\pi(B^{\bullet, \bullet})$
avec $B^{p, q} = \Hom(a_p^{-1}\mathbf{Z}_U^\#, g_p^{-1}\mathcal{J}^q)$.
Pour l'autre membre, raisonnons de la même manière. Observons d'abord que
$$
s(\mathbf{Z}_{F(K)}^\#) \longrightarrow \mathbf{Z}
$$
est un quasi-isomorphisme de complexes sur $\mathcal{C}$
d'après Hyperrecouvrements, lemme \ref{hypercovering-lemma-acyclic-hypercover-sheaves}.
Puisque $\mathbf{Z}_U^\#$ est un faisceau plat de $\mathbf{Z}$-modules,
on voit que
$$
s(\mathbf{Z}_{F(K)}^\#) \otimes_\mathbf{Z} \mathbf{Z}_U^\#
\longrightarrow
\mathbf{Z}_U^\#
$$
est un quasi-isomorphisme. Par conséquent,
$R\Hom(\mathbf{Z}_U^\#, E)$ est égal à
$$
R\Gamma(\mathcal{C}, R\SheafHom(
s(\mathbf{Z}_{F(K)}^\#) \otimes_\mathbf{Z} \mathbf{Z}_U^\#,
\mathcal{I}^\bullet))
$$
Par construction de $R\SheafHom$, et puisque $\mathcal{I}^\bullet$
est K-injectif, l'objet précédent est représenté par le complexe de groupes abéliens
dont les termes sont
$$
\prod\nolimits_{p + q = n}
\Hom(\mathbf{Z}^\#_{K_p} \otimes_\mathbf{Z} \mathbf{Z}_U^\#, \mathcal{I}^q)
=
\prod\nolimits_{p + q = n}
\Hom(a_p^{-1}\mathbf{Z}_U^\#, a_p^{-1}\mathcal{I}^q)
$$
L'égalité des termes découle du fait que
$\mathbf{Z}^\#_{K_p} \otimes_\mathbf{Z} \mathbf{Z}_U^\# =
a_{p!}a_p^{-1}\mathbf{Z}_U^\#$ d'après Modules sur les sites,
remarque \ref{sites-modules-remark-j-shriek-tensor}.
Ainsi, $R\Gamma(U, E)$ est calculé par
le complexe total produit $\text{Tot}_\pi(A^{\bullet, \bullet})$
avec $A^{p, q} = \Hom(a_p^{-1}\mathbf{Z}_U^\#, a_p^{-1}\mathcal{I}^q)$.

\medskip\noindent
Puisque $\mathcal{I}^\bullet$ est K-injectif, nous voyons que
$a_p^{-1}\mathcal{I}^\bullet$ est K-injectif, voir
le lemme \ref{lemma-localize-injective}.
Puisque $\mathcal{J}^\bullet$ est K-injectif, nous voyons que
$g_p^{-1}\mathcal{J}^\bullet$ est K-injectif, voir
le lemme \ref{lemma-restriction-injective-to-component-site}.
Ces deux complexes représentent l'objet $a_p^{-1}E$.
Ainsi, pour tout $p \geq 0$, le morphisme de complexes
$$
A^{p, \bullet} = \Hom(a_p^{-1}\mathbf{Z}_U^\#, a_p^{-1}\mathcal{I}^\bullet)
\longrightarrow
\Hom(a_p^{-1}\mathbf{Z}_U^\#, g_p^{-1}\mathcal{J}^\bullet) = B^{p, \bullet}
$$
obtenu en appliquant $g_p^{-1}$ au morphisme donné
$a^{-1}\mathcal{I}^\bullet \to \mathcal{J}^\bullet$
est un quasi-isomorphisme, puisque ces deux complexes calculent
$$
R\Hom(a_p^{-1}\mathbf{Z}_U^\#, a_p^{-1}E)
$$
D'après Compléments d'algèbre, lemme \ref{more-algebra-lemma-prod-qis-gives-qis},
nous concluons que la flèche verticale droite dans le diagramme commutatif
$$
\xymatrix{
R\Gamma(U, E) \ar[r] \ar[d] &
\text{Tot}_\pi(A^{\bullet, \bullet}) \ar[d] \\
R\Gamma(U, Ra_*a^{-1}E) \ar[r] &
\text{Tot}_\pi(B^{\bullet, \bullet})
}
$$
est un quasi-isomorphisme. Puisque nous avons vu plus haut que les flèches horizontales
sont des quasi-isomorphismes, la flèche verticale gauche l'est également.
\end{proof}

\begin{lemma}
\label{lemma-compare-cohomology-hypercovering}
Soit $\mathcal{C}$ un site admettant des égalisateurs et des produits fibrés.
Soit $K$ un hyperrecouvrement.
Nous avons alors un isomorphisme canonique
$$
R\Gamma(\mathcal{C}, E) =
R\Gamma((\mathcal{C}/K)_{total}, a^{-1}E)
$$
pour $E \in D(\mathcal{C})$.
\end{lemma}

\begin{proof}
Cela découle du lemme \ref{lemma-hypercovering-descent-bounded-abelian}
car $R\Gamma((\mathcal{C}/K)_{total}, -) =
R\Gamma(\mathcal{C}, -) \circ Ra_*$ par
Cohomologie sur les sites, remarque \ref{sites-cohomology-remark-before-Leray}.
\end{proof}

\begin{lemma}
\label{lemma-hypercovering-equivalence-bounded}
Soit $\mathcal{C}$ un site admettant des égalisateurs et des produits fibrés.
Soit $K$ un hyperrecouvrement.
Soit $\mathcal{A} \subset \textit{Ab}((\mathcal{C}/K)_{total})$
la sous-catégorie de Serre faible des faisceaux abéliens cartésiens.
Alors le foncteur $a^{-1}$ définit une équivalence
$$
D^+(\mathcal{C}) \longrightarrow D_\mathcal{A}^+((\mathcal{C}/K)_{total})
$$
de quasi-inverse $Ra_*$.
\end{lemma}

\begin{proof}
Observons que $\mathcal{A}$ est une sous-catégorie de Serre faible d'après
le lemme \ref{lemma-Serre-subcat-cartesian-modules}.
L'équivalence est une conséquence formelle des
résultats obtenus jusqu'ici. On utilise les lemmes
\ref{lemma-hypercovering-descent-sheaves} et
\ref{lemma-hypercovering-descent-bounded-abelian}, ainsi que
Cohomologie sur les sites, lemme \ref{sites-cohomology-lemma-equivalence-bounded}.
\end{proof}

\noindent
Nous invitons le lecteur à ignorer la remarque suivante.

\begin{remark}
\label{remark-compare-cohomology-hypercovering-presheaf}
Soient $\mathcal{C}$ un site et $\mathcal{G}$ un préfaisceau d'ensembles sur
$\mathcal{C}$. Si $\mathcal{C}$ admet des égalisateurs et des produits fibrés,
nous avons défini la notion d'hyperrecouvrement de $\mathcal{G}$ dans
Hyperrecouvrements, définition \ref{hypercovering-definition-hypercovering-variant}.
Tous les résultats de cette section se transposent à
ce cadre.
Pour le voir,
considérons la localisation $\mathcal{C}/\mathcal{G}$
de $\mathcal{C}$ en $\mathcal{G}$, définie exactement comme dans
Sites, lemme \ref{sites-lemma-localize-topos-site}
(qui n'est énoncé que pour les faisceaux; le topos
$\Sh(\mathcal{C}/\mathcal{G})$ est égal à la localisation
du topos $\Sh(\mathcal{C})$ au faisceau $\mathcal{G}^\#$).
Le lecteur vérifie alors aisément que le site
$\mathcal{C}/\mathcal{G}$ admet des produits fibrés et des égalisateurs,
et qu'un hyperrecouvrement de $\mathcal{G}$ dans $\mathcal{C}$
équivaut à un hyperrecouvrement pour le site $\mathcal{C}/\mathcal{G}$.
Ainsi, en remplaçant le site $\mathcal{C}$ par $\mathcal{C}/\mathcal{G}$
dans les lemmes ci-dessus sur les hyperrecouvrements, on obtient les démonstrations des
résultats correspondants pour les hyperrecouvrements de $\mathcal{G}$.
Exemple : pour un hyperrecouvrement $K$ de $\mathcal{G}$, on a
$$
R\Gamma(\mathcal{C}/\mathcal{G}, E) =
R\Gamma((\mathcal{C}/K)_{total}, a^{-1}E)
$$
pour $E \in D^+(\mathcal{C}/\mathcal{G})$, où
$a : \Sh((\mathcal{C}/K)_{total}) \to \Sh(\mathcal{C}/\mathcal{G})$
est l'augmentation canonique. C'est un cas particulier du lemme
\ref{lemma-compare-cohomology-hypercovering}.
Notons $R\Gamma(\mathcal{G}, -) : D(\mathcal{C}) \to D(\textit{Ab})$ le foncteur
défini comme le foncteur dérivé du foncteur
$H^0(\mathcal{G}, -) = H^0(\mathcal{G}^\#, -)$
étudié dans Hyperrecouvrements, section
\ref{hypercovering-section-hypercoverings-verdier}, et dans
Cohomologie sur les sites, section \ref{sites-cohomology-section-limp}.
Nous avons
$$
R\Gamma(\mathcal{G}, E) = R\Gamma(\mathcal{C}/\mathcal{G}, j^{-1}E)
$$
d'après l'analogue du lemme de Cohomologie sur les sites
\ref{sites-cohomology-lemma-cohomology-of-open}
pour le foncteur de localisation $j : \mathcal{C}/\mathcal{G} \to \mathcal{C}$.
En combinant ces égalités, on obtient
$$
R\Gamma(\mathcal{G}, E) =
R\Gamma((\mathcal{C}/K)_{total}, a^{-1}j^{-1}E) =
R\Gamma((\mathcal{C}/K)_{total}, g^{-1}E)
$$
pour $E \in D^+(\mathcal{C})$, où
$g : \Sh((\mathcal{C}/K)_{total}) \to \Sh(\mathcal{C})$
est la composition de $a$ et $j$.
\end{remark}







\section{Descente cohomologique pour les hyperrecouvrements : modules}
\label{section-cohomological-descent-hypercoverings-modules}

\noindent
Soit $\mathcal{C}$ un site et soit $\mathcal{O}_\mathcal{C}$
un faisceau d'anneaux. Supposons que $\mathcal{C}$
admette des égalisateurs et des produits fibrés, et soit $K$ un hyperrecouvrement
au sens d'Hyperrecouvrements, définition
\ref{hypercovering-definition-hypercovering-variant}. Nous étudierons la descente cohomologique
pour l'augmentation
$$
a :
(\Sh((\mathcal{C}/K)_{total}), \mathcal{O})
\longrightarrow
(\Sh(\mathcal{C}), \mathcal{O}_\mathcal{C})
$$
de la remarque \ref{remark-augmentation-ringed}.

\begin{lemma}
\label{lemma-hypercovering-descent-modules}
Soit $\mathcal{C}$ un site admettant des égalisateurs et des produits fibrés.
Soit $\mathcal{O}_\mathcal{C}$ un faisceau d'anneaux.
Soit $K$ un hyperrecouvrement. Avec les notations ci-dessus,
$$
a^* : \textit{Mod}(\mathcal{O}_\mathcal{C}) \to \textit{Mod}(\mathcal{O})
$$
est pleinement fidèle, d'image essentielle les $\mathcal{O}$-modules cartésiens.
Le foncteur $a_*$ est un quasi-inverse.
\end{lemma}

\begin{proof}
Puisque $a^{-1}\mathcal{O}_\mathcal{C} = \mathcal{O}$, on a
$a^* = a^{-1}$. Le lemme résulte donc immédiatement
du lemme \ref{lemma-hypercovering-descent-sheaves}.
\end{proof}

\begin{lemma}
\label{lemma-hypercovering-descent-bounded-modules}
Soit $\mathcal{C}$ un site admettant des égalisateurs et des produits fibrés.
Soit $\mathcal{O}_\mathcal{C}$ un faisceau d'anneaux.
Soit $K$ un hyperrecouvrement. Pour
$E \in D(\mathcal{O}_\mathcal{C})$, le morphisme
$$
E \longrightarrow Ra_*La^*E
$$
est un isomorphisme.
\end{lemma}

\begin{proof}
Puisque $a^{-1}\mathcal{O}_\mathcal{C} = \mathcal{O}$, on a
$La^* = a^* = a^{-1}$. De plus, $Ra_*$ coïncide avec
$Ra_*$ sur les faisceaux abéliens; voir Cohomologie
sur les sites, lemme
\ref{sites-cohomology-lemma-modules-abelian-unbounded}.
Le lemme en résulte immédiatement,
du lemme \ref{lemma-hypercovering-descent-bounded-abelian}.
\end{proof}

\begin{lemma}
\label{lemma-compare-cohomology-hypercovering-modules}
Soit $\mathcal{C}$ un site admettant des égalisateurs et des produits fibrés.
Soit $\mathcal{O}_\mathcal{C}$ un faisceau d'anneaux.
Soit $K$ un hyperrecouvrement.
Nous avons alors un isomorphisme canonique
$$
R\Gamma(\mathcal{C}, E) =
R\Gamma((\mathcal{C}/K)_{total}, La^*E)
$$
pour $E \in D(\mathcal{O}_\mathcal{C})$.
\end{lemma}

\begin{proof}
Cela découle du lemme \ref{lemma-hypercovering-descent-bounded-modules}
car $R\Gamma((\mathcal{C}/K)_{total}, -) =
R\Gamma(\mathcal{C}, -) \circ Ra_*$ par
Cohomologie sur les sites, remarque \ref{sites-cohomology-remark-before-Leray},
ou par
Cohomologie sur les sites, lemme \ref{sites-cohomology-lemma-Leray-unbounded}.
\end{proof}

\begin{lemma}
\label{lemma-hypercovering-equivalence-bounded-modules}
Soit $\mathcal{C}$ un site admettant des égalisateurs et des produits fibrés.
Soit $\mathcal{O}_\mathcal{C}$ un faisceau d'anneaux.
Soit $K$ un hyperrecouvrement.
Soit $\mathcal{A} \subset \textit{Mod}(\mathcal{O})$
la sous-catégorie de Serre faible des $\mathcal{O}$-modules cartésiens.
Alors le foncteur $La^*$ définit une équivalence
$$
D^+(\mathcal{O}_\mathcal{C}) \longrightarrow D_\mathcal{A}^+(\mathcal{O})
$$
de quasi-inverse $Ra_*$.
\end{lemma}

\begin{proof}
Observons que $\mathcal{A}$ est une sous-catégorie de Serre faible d'après le
lemme \ref{lemma-Serre-subcat-cartesian-modules}
(les hypothèses requises découlent de la discussion de la
remarque \ref{remark-augmentation-ringed}).
L'équivalence est une conséquence formelle
des résultats obtenus jusqu'ici. On utilise les lemmes
\ref{lemma-hypercovering-descent-modules},
\ref{lemma-hypercovering-descent-bounded-modules}, ainsi que Cohomologie
sur les sites, lemme \ref{sites-cohomology-lemma-equivalence-bounded}.
\end{proof}






\section{Descente cohomologique pour les hyperrecouvrements d'un objet}
\label{section-cohomological-descent-hypercoverings-X}

\noindent
Dans cette section, nous supposons que $\mathcal{C}$ admet des produits fibrés
et nous fixons $X \in \Ob(\mathcal{C})$. Soit $K$ un hyperrecouvrement de $X$
au sens de
Hyperrecouvrements, définition \ref{hypercovering-definition-hypercovering}.
Nous étudierons l'augmentation
$$
a : \Sh((\mathcal{C}/K)_{total}) \longrightarrow \Sh(\mathcal{C}/X)
$$
de la remarque \ref{remark-augmentation-over-object}.
Observons que $\mathcal{C}/X$ est un site qui admet des égalisateurs
et des produits fibrés, et que $K$ est un hyperrecouvrement
du site $\mathcal{C}/X$\footnote{La réciproque peut être fausse :
si $K$ est un objet simplicial de
$\text{SR}(\mathcal{C}, X) = \text{SR}(\mathcal{C}/X)$
qui est un hyperrecouvrement du site $\mathcal{C}/X$ au sens de
Hyperrecouvrements, définition \ref{hypercovering-definition-hypercovering-variant},
alors $K$ ne définit pas nécessairement un hyperrecouvrement de $X$.
Par exemple, si $K_0 = \{U_{0, i}\}_{i \in I_0}$,
cette dernière condition garantit que
$\{U_{0, i} \to X\}$ est un recouvrement de $\mathcal{C}$
tandis que la première exige seulement que
$\coprod h_{U_{0, i}}^\# \to h_X^\#$ soit un morphisme surjectif
de faisceaux;} voir Hyperrecouvrements, lemme
\ref{hypercovering-lemma-hypercovering-F}.
Ainsi, tout résultat démontré pour les hyperrecouvrements
dans la section \ref{section-cohomological-descent-hypercoverings}
se transpose immédiatement à la situation de la présente section.

\begin{lemma}
\label{lemma-hypercovering-X-descent-sheaves}
Soit $\mathcal{C}$ un site avec des produits fibrés et $X \in \Ob(\mathcal{C})$.
Soit $K$ un hyperrecouvrement de $X$. Alors
\begin{enumerate}
\item $a^{-1} : \Sh(\mathcal{C}/X) \to \Sh((\mathcal{C}/K)_{total})$
est pleinement fidèle, d'image essentielle les faisceaux cartésiens d'ensembles,
\item $a^{-1} : \textit{Ab}(\mathcal{C}/X) \to
\textit{Ab}((\mathcal{C}/K)_{total})$
est pleinement fidèle, d'image essentielle les faisceaux cartésiens
de groupes abéliens.
\end{enumerate}
Dans les deux cas, $a_*$ est un foncteur quasi-inverse.
\end{lemma}

\begin{proof}
D'après les remarques \ref{remark-semi-representable-over-object} et
\ref{remark-augmentation-over-object}, ainsi que la discussion de
l'introduction de cette section,
le résultat découle du lemme \ref{lemma-hypercovering-descent-sheaves}.
\end{proof}

\begin{lemma}
\label{lemma-hypercovering-X-descent-bounded-abelian}
Soit $\mathcal{C}$ un site admettant des produits fibrés et soit $X \in \Ob(\mathcal{C})$.
Soit $K$ un hyperrecouvrement de $X$. Pour
$E \in D(\mathcal{C}/X)$, le morphisme
$$
E \longrightarrow Ra_*a^{-1}E
$$
est un isomorphisme.
\end{lemma}

\begin{proof}
D'après les remarques \ref{remark-semi-representable-over-object} et
\ref{remark-augmentation-over-object}, ainsi que la discussion de
l'introduction de cette section,
le résultat découle du lemme \ref{lemma-hypercovering-descent-bounded-abelian}.
\end{proof}

\begin{lemma}
\label{lemma-compare-cohomology-hypercovering-X}
Soit $\mathcal{C}$ un site avec des produits fibrés et $X \in \Ob(\mathcal{C})$.
Soit $K$ un hyperrecouvrement de $X$.
Nous avons alors un isomorphisme canonique
$$
R\Gamma(X, E) = R\Gamma((\mathcal{C}/K)_{total}, a^{-1}E)
$$
pour $E \in D(\mathcal{C}/X)$.
\end{lemma}

\begin{proof}
D'après les remarques \ref{remark-semi-representable-over-object} et
\ref{remark-augmentation-over-object},
le résultat découle du lemme \ref{lemma-compare-cohomology-hypercovering}.
\end{proof}

\begin{lemma}
\label{lemma-hypercovering-X-equivalence-bounded}
Soit $\mathcal{C}$ un site avec des produits fibrés et $X \in \Ob(\mathcal{C})$.
Soit $K$ un hyperrecouvrement de $X$.
Soit $\mathcal{A} \subset \textit{Ab}((\mathcal{C}/K)_{total})$
la sous-catégorie de Serre faible des faisceaux abéliens cartésiens.
Alors le foncteur $a^{-1}$ définit une équivalence
$$
D^+(\mathcal{C}/X) \longrightarrow D_\mathcal{A}^+((\mathcal{C}/K)_{total})
$$
de quasi-inverse $Ra_*$.
\end{lemma}

\begin{proof}
D'après les remarques \ref{remark-semi-representable-over-object} et
\ref{remark-augmentation-over-object},
le résultat découle du lemme \ref{lemma-hypercovering-equivalence-bounded}.
\end{proof}








\section{Descente cohomologique pour les hyperrecouvrements d'un objet : modules}
\label{section-cohomological-descent-hypercoverings-X-modules}

\noindent
Dans cette section, nous supposons que $\mathcal{C}$ admet des produits fibrés
et nous fixons $X \in \Ob(\mathcal{C})$. Soit $K$ un hyperrecouvrement de $X$
au sens de
Hyperrecouvrements, définition \ref{hypercovering-definition-hypercovering}.
Soit $\mathcal{O}_\mathcal{C}$ un faisceau d'anneaux sur $\mathcal{C}$.
Posons $\mathcal{O}_X = \mathcal{O}_\mathcal{C}|_{\mathcal{C}/X}$.
Nous étudierons l'augmentation
$$
a :
(\Sh((\mathcal{C}/K)_{total}), \mathcal{O})
\longrightarrow
(\Sh(\mathcal{C}/X), \mathcal{O}_X)
$$
de la remarque \ref{remark-augmentation-ringed-over-object}.
Observons que $\mathcal{C}/X$ est un site qui admet des égalisateurs
et des produits fibrés, et que $K$ est un hyperrecouvrement
du site $\mathcal{C}/X$.
Par conséquent, les résultats de cette section découlent immédiatement
des résultats correspondants dans la section
\ref{section-cohomological-descent-hypercoverings-modules}.

\begin{lemma}
\label{lemma-hypercovering-X-descent-modules}
Soient $\mathcal{C}$ un site admettant des produits fibrés et $X \in \Ob(\mathcal{C})$.
Soit $\mathcal{O}_\mathcal{C}$ un faisceau d'anneaux.
Soit $K$ un hyperrecouvrement de $X$. Avec les notations ci-dessus,
$$
a^* : \textit{Mod}(\mathcal{O}_X) \to \textit{Mod}(\mathcal{O})
$$
est pleinement fidèle, d'image essentielle les $\mathcal{O}$-modules cartésiens.
Le foncteur $a_*$ est un quasi-inverse.
\end{lemma}

\begin{proof}
D'après les remarques \ref{remark-semi-representable-ringed-over-object} et
\ref{remark-augmentation-ringed-over-object}, ainsi que la discussion de
l'introduction de cette section,
le résultat découle du lemme \ref{lemma-hypercovering-descent-modules}.
\end{proof}

\begin{lemma}
\label{lemma-hypercovering-X-descent-bounded-modules}
Soient $\mathcal{C}$ un site admettant des produits fibrés et $X \in \Ob(\mathcal{C})$.
Soit $\mathcal{O}_\mathcal{C}$ un faisceau d'anneaux.
Soit $K$ un hyperrecouvrement de $X$. Pour
$E \in D(\mathcal{O}_X)$, le morphisme
$$
E \longrightarrow Ra_*La^*E
$$
est un isomorphisme.
\end{lemma}

\begin{proof}
D'après les remarques \ref{remark-semi-representable-ringed-over-object} et
\ref{remark-augmentation-ringed-over-object}, ainsi que la discussion de
l'introduction de cette section,
le résultat découle du lemme \ref{lemma-hypercovering-descent-bounded-modules}.
\end{proof}

\begin{lemma}
\label{lemma-compare-cohomology-hypercovering-X-modules}
Soient $\mathcal{C}$ un site admettant des produits fibrés et $X \in \Ob(\mathcal{C})$.
Soit $\mathcal{O}_\mathcal{C}$ un faisceau d'anneaux.
Soit $K$ un hyperrecouvrement de $X$.
Nous avons alors un isomorphisme canonique
$$
R\Gamma(X, E) = R\Gamma((\mathcal{C}/K)_{total}, La^*E)
$$
pour $E \in D(\mathcal{O}_\mathcal{C})$.
\end{lemma}

\begin{proof}
D'après les remarques \ref{remark-semi-representable-ringed-over-object} et
\ref{remark-augmentation-ringed-over-object}, ainsi que la discussion de
l'introduction de cette section,
le résultat découle du lemme \ref{lemma-compare-cohomology-hypercovering-modules}.
\end{proof}

\begin{lemma}
\label{lemma-hypercovering-X-equivalence-bounded-modules}
Soient $\mathcal{C}$ un site admettant des produits fibrés et $X \in \Ob(\mathcal{C})$.
Soit $\mathcal{O}_\mathcal{C}$ un faisceau d'anneaux.
Soit $K$ un hyperrecouvrement de $X$.
Soit $\mathcal{A} \subset \textit{Mod}(\mathcal{O})$
la sous-catégorie de Serre faible des $\mathcal{O}$-modules cartésiens.
Alors le foncteur $La^*$ définit une équivalence
$$
D^+(\mathcal{O}_X) \longrightarrow D_\mathcal{A}^+(\mathcal{O})
$$
de quasi-inverse $Ra_*$.
\end{lemma}

\begin{proof}
D'après les remarques \ref{remark-semi-representable-ringed-over-object} et
\ref{remark-augmentation-ringed-over-object}, ainsi que la discussion de
l'introduction de cette section,
le résultat découle du lemme \ref{lemma-hypercovering-equivalence-bounded-modules}.
\end{proof}










\section{Hyperrecouvrement par un objet simplicial du site}
\label{section-hypercovering}

\noindent
Soit $\mathcal{C}$ un site admettant des produits fibrés et
soit $X \in \Ob(\mathcal{C})$.
Dans cette section, nous expliquons les résultats de la section
\ref{section-cohomological-descent-hypercoverings-X}
dans le cas où notre hyperrecouvrement est donné par
un objet simplicial du site.
Soit $U$ un objet simplicial de $\mathcal{C}$.
Comme d'habitude, on note $U_n = U([n])$ et $f_\varphi : U_n \to U_m$
le morphisme $f_\varphi = U(\varphi)$ correspondant à
$\varphi : [m] \to [n]$.
Supposons que nous ayons une augmentation
$$
a : U \to X
$$
On en déduit un site simplicial $(\mathcal{C}/U)_{total}$
et un morphisme d'augmentation
$$
a : \Sh((\mathcal{C}/U)_{total}) \longrightarrow \Sh(\mathcal{C}/X)
$$
Plus précisément, $U$ détermine un objet simplicial $K$ de
$\text{SR}(\mathcal{C}, X)$ dont le terme de degré $n$ est
$K_n = \{U_n \to X\}$; on peut donc appliquer les constructions de la
remarque \ref{remark-augmentation-over-object}.
Plus précisément, un objet du site $(\mathcal{C}/U)_{total}$ est de la forme
$V/U_n$, et un morphisme $(\varphi, f) : V/U_n \to W/U_m$ est donné
par un morphisme $\varphi : [m] \to [n]$ dans $\Delta$ et un morphisme
$f : V \to W$ tel que le diagramme
$$
\xymatrix{
V \ar[r]_f \ar[d] & W \ar[d] \\
U_n \ar[r]^{f_\varphi} & U_m
}
$$
est commutatif. Le morphisme de topoi $a$ est défini par le foncteur cocontinu
$V/U_n \mapsto V/X$. Voilà toute la construction.

\medskip\noindent
Dans cette section, nous dirons que l'augmentation $a : U \to X$ est un
{\it hyperrecouvrement de $X$ dans $\mathcal{C}$} si les conditions suivantes sont satisfaites :
\begin{enumerate}
\item $\{U_0 \to X\}$ est un recouvrement de $\mathcal{C}$,
\item $\{U_1 \to U_0 \times_X U_0\}$ est un recouvrement de $\mathcal{C}$,
\item $\{U_{n + 1} \to (\text{cosk}_n\text{sk}_n U)_{n + 1}\}$
est un recouvrement de $\mathcal{C}$ pour $n \geq 1$.
\end{enumerate}
Cela équivaut à dire que $K$ (comme ci-dessus) est un hyperrecouvrement
de $X$; voir
Hyperrecouvrements, exemple \ref{hypercovering-example-hypercovering-in-C}.

\begin{lemma}
\label{lemma-hypercovering-X-simple-descent-sheaves}
Soient $\mathcal{C}$ un site admettant des produits fibrés et $X \in \Ob(\mathcal{C})$.
Soit $a : U \to X$ un hyperrecouvrement de $X$ dans $\mathcal{C}$ comme défini ci-dessus.
Alors
\begin{enumerate}
\item $a^{-1} : \Sh(\mathcal{C}/X) \to \Sh((\mathcal{C}/U)_{total})$
est pleinement fidèle, d'image essentielle les faisceaux cartésiens d'ensembles,
\item $a^{-1} : \textit{Ab}(\mathcal{C}/X) \to
\textit{Ab}((\mathcal{C}/U)_{total})$
est pleinement fidèle, d'image essentielle les faisceaux cartésiens
de groupes abéliens.
\end{enumerate}
Dans les deux cas, $a_*$ est un foncteur quasi-inverse.
\end{lemma}

\begin{proof}
Il s'agit d'un cas particulier du lemme
\ref{lemma-hypercovering-X-descent-sheaves}.
\end{proof}

\begin{lemma}
\label{lemma-hypercovering-X-simple-descent-bounded-abelian}
Soient $\mathcal{C}$ un site admettant des produits fibrés et $X \in \Ob(\mathcal{C})$.
Soit $a : U \to X$ un hyperrecouvrement de $X$ dans $\mathcal{C}$ comme défini ci-dessus.
Pour $E \in D(\mathcal{C}/X)$, le morphisme
$$
E \longrightarrow Ra_*a^{-1}E
$$
est un isomorphisme.
\end{lemma}

\begin{proof}
Il s'agit d'un cas particulier du lemme
\ref{lemma-hypercovering-X-descent-bounded-abelian}.
\end{proof}

\begin{lemma}
\label{lemma-compare-cohomology-hypercovering-X-simple}
Soient $\mathcal{C}$ un site admettant des produits fibrés et $X \in \Ob(\mathcal{C})$.
Soit $a : U \to X$ un hyperrecouvrement de $X$ dans $\mathcal{C}$ comme défini ci-dessus.
Nous avons alors un isomorphisme canonique
$$
R\Gamma(X, E) = R\Gamma((\mathcal{C}/U)_{total}, a^{-1}E)
$$
pour $E \in D(\mathcal{C}/X)$.
\end{lemma}

\begin{proof}
Il s'agit d'un cas particulier du lemme
\ref{lemma-compare-cohomology-hypercovering-X}.
\end{proof}

\begin{lemma}
\label{lemma-hypercovering-X-simple-equivalence-bounded}
Soient $\mathcal{C}$ un site admettant des produits fibrés et $X \in \Ob(\mathcal{C})$.
Soit $a : U \to X$ un hyperrecouvrement de $X$ dans $\mathcal{C}$ comme défini ci-dessus.
Soit $\mathcal{A} \subset \textit{Ab}((\mathcal{C}/U)_{total})$
la sous-catégorie de Serre faible des faisceaux abéliens cartésiens.
Alors le foncteur $a^{-1}$ définit une équivalence
$$
D^+(\mathcal{C}/X) \longrightarrow D_\mathcal{A}^+((\mathcal{C}/U)_{total})
$$
de quasi-inverse $Ra_*$.
\end{lemma}

\begin{proof}
C'est un cas particulier du
lemme \ref{lemma-hypercovering-X-equivalence-bounded}.
\end{proof}

\begin{lemma}
\label{lemma-sr-when-fibre-products}
Soit $U$ un objet simplicial d'un site $\mathcal{C}$
admettant des produits fibrés.
\begin{enumerate}
\item $\mathcal{C}/U$ définit un objet simplicial
dans la catégorie dont les objets sont des sites et
dont les morphismes sont des morphismes de sites,
\item la construction du lemme \ref{lemma-simplicial-site-site}
appliquée à la structure de (1)
reproduit le site $(\mathcal{C}/U)_{total}$ ci-dessus,
\item si $a : U \to X$ est une augmentation, alors
$a_0 : \mathcal{C}/U_0 \to \mathcal{C}/X$ est une augmentation
comme dans la partie (A) de la remarque \ref{remark-augmentation-site} et donne le
même morphisme de topoi
$a : \Sh((\mathcal{C}/U)_{total}) \to \Sh(\mathcal{C}/X)$
identique à celui défini ci-dessus.
\end{enumerate}
\end{lemma}

\begin{proof}
Pour un morphisme d'objets $V \to W$ de $\mathcal{C}$, le morphisme de localisation
$j : \mathcal{C}/V \to \mathcal{C}/W$ est adjoint à gauche au
foncteur de changement de base $\mathcal{C}/W \to \mathcal{C}/V$.
Le foncteur de changement de base est continu et induit le même morphisme de topoi
que $j$. Voir
Sites, lemme \ref{sites-lemma-relocalize-given-fibre-products}.
Cela prouve (1).

\medskip\noindent
La partie (2) résulte de ce qu'un morphisme $V/U_n \to W/U_m$
de la catégorie construite
dans le lemme \ref{lemma-simplicial-site-site}
est un morphisme $V \to W \times_{U_m, f_\varphi} U_n$ sur $U_n$
équivaut à un morphisme $f : V \to W$
sur le morphisme $f_\varphi : U_n \to U_m$, c'est-à-dire
la même chose qu'un morphisme dans la catégorie $(\mathcal{C}/U)_{total}$
défini ci-dessus. L'égalité des ensembles de recouvrements résulte
immédiatement de la définition.

\medskip\noindent
Nous omettons la preuve de (3).
\end{proof}







\section{Hyperrecouvrement par un objet simplicial du site : modules}
\label{section-hypercovering-modules}

\noindent
Soient $\mathcal{C}$ un site admettant des produits fibrés et $X \in \Ob(\mathcal{C})$.
Soit $\mathcal{O}_\mathcal{C}$ un faisceau d'anneaux sur $\mathcal{C}$.
Soit $U \to X$ un hyperrecouvrement de $X$ dans $\mathcal{C}$, au sens de la
section \ref{section-hypercovering}. Nous étudions ici l'augmentation
suivante :
$$
a :
(\Sh((\mathcal{C}/U)_{total}), \mathcal{O})
\longrightarrow
(\Sh(\mathcal{C}/X), \mathcal{O}_X)
$$
que l'on obtient en considérant $U$ comme un objet semi-représentable simplicial
de $\mathcal{C}/X$, dont le terme de degré $n$ est le singleton
$\{U_n/X\}$, puis en appliquant les constructions de la
remarque \ref{remark-augmentation-ringed-over-object}.
Ainsi tous les résultats de cette section sont des conséquences immédiates
des résultats correspondants dans la section
\ref{section-cohomological-descent-hypercoverings-X-modules}.

\begin{lemma}
\label{lemma-hypercovering-X-simple-descent-modules}
Soient $\mathcal{C}$ un site admettant des produits fibrés et $X \in \Ob(\mathcal{C})$.
Soit $\mathcal{O}_\mathcal{C}$ un faisceau d'anneaux.
Soit $U$ un hyperrecouvrement de $X$ dans $\mathcal{C}$. Avec les notations ci-dessus,
$$
a^* : \textit{Mod}(\mathcal{O}_X) \to \textit{Mod}(\mathcal{O})
$$
est pleinement fidèle, d'image essentielle les $\mathcal{O}$-modules cartésiens.
Le foncteur $a_*$ est un quasi-inverse.
\end{lemma}

\begin{proof}
Il s'agit d'un cas particulier du lemme
\ref{lemma-hypercovering-X-descent-modules}.
\end{proof}

\begin{lemma}
\label{lemma-hypercovering-X-simple-descent-bounded-modules}
Soient $\mathcal{C}$ un site admettant des produits fibrés et $X \in \Ob(\mathcal{C})$.
Soit $\mathcal{O}_\mathcal{C}$ un faisceau d'anneaux.
Soit $U$ un hyperrecouvrement de $X$ dans $\mathcal{C}$. Pour
$E \in D(\mathcal{O}_X)$, le morphisme
$$
E \longrightarrow Ra_*La^*E
$$
est un isomorphisme.
\end{lemma}

\begin{proof}
Il s'agit d'un cas particulier du lemme
\ref{lemma-hypercovering-X-descent-bounded-modules}.
\end{proof}

\begin{lemma}
\label{lemma-compare-cohomology-hypercovering-X-simple-modules}
Soient $\mathcal{C}$ un site admettant des produits fibrés et $X \in \Ob(\mathcal{C})$.
Soit $\mathcal{O}_\mathcal{C}$ un faisceau d'anneaux.
Soit $U$ un hyperrecouvrement de $X$ dans $\mathcal{C}$.
Nous avons alors un isomorphisme canonique
$$
R\Gamma(X, E) = R\Gamma((\mathcal{C}/U)_{total}, La^*E)
$$
pour $E \in D(\mathcal{O}_\mathcal{C})$.
\end{lemma}

\begin{proof}
Il s'agit d'un cas particulier du lemme
\ref{lemma-compare-cohomology-hypercovering-X-modules}.
\end{proof}

\begin{lemma}
\label{lemma-hypercovering-X-simple-equivalence-bounded-modules}
Soient $\mathcal{C}$ un site admettant des produits fibrés et $X \in \Ob(\mathcal{C})$.
Soit $\mathcal{O}_\mathcal{C}$ un faisceau d'anneaux.
Soit $U$ un hyperrecouvrement de $X$ dans $\mathcal{C}$.
Soit $\mathcal{A} \subset \textit{Mod}(\mathcal{O})$
la sous-catégorie de Serre faible des $\mathcal{O}$-modules cartésiens.
Alors le foncteur $La^*$ définit une équivalence
$$
D^+(\mathcal{O}_X) \longrightarrow D_\mathcal{A}^+(\mathcal{O})
$$
de quasi-inverse $Ra_*$.
\end{lemma}

\begin{proof}
Il s'agit d'un cas particulier du lemme
\ref{lemma-hypercovering-X-equivalence-bounded-modules}.
\end{proof}







\section{Descente cohomologique non bornée pour les hyperrecouvrements}
\label{section-unbounded-cohomological-descent}

\noindent
Dans cette section, nous étudions la descente cohomologique non bornée.
Les résultats eux-mêmes seront des conséquences immédiates de
nos résultats sur la descente cohomologique bornée dans les
sections précédentes et des lemmes de Cohomologie sur les sites
\ref{sites-cohomology-lemma-equivalence-unbounded-one} et/ou
\ref{sites-cohomology-lemma-equivalence-unbounded-two} ; le vrai travail réside
dans la mise en place de la notation et le choix des hypothèses appropriées.
Cette discussion est motivée par celle de \cite{six-I},
bien que les détails diffèrent sensiblement.

\medskip\noindent
Soit $(\mathcal{C}, \mathcal{O}_\mathcal{C})$ un site annelé.
Supposons donnée, pour tout objet $U$ de $\mathcal{C}$,
une sous-catégorie de Serre faible $\mathcal{A}_U \subset \textit{Mod}(\mathcal{O}_U)$
satisfaisant aux propriétés suivantes :
\begin{enumerate}
\item
\label{item-restriction}
pour tout morphisme $U \to V$ de $\mathcal{C}$, le foncteur de restriction
$\textit{Mod}(\mathcal{O}_V) \to \textit{Mod}(\mathcal{O}_U)$
envoie $\mathcal{A}_V$ dans $\mathcal{A}_U$,
\item
\label{item-local}
pour tout recouvrement $\{U_i \to U\}_{i \in I}$ de $\mathcal{C}$,
un objet $\mathcal{F}$ de $\textit{Mod}(\mathcal{O}_U)$
appartient à $\mathcal{A}_U$ si et seulement si la restriction de
$\mathcal{F}$ à $\mathcal{C}/U_i$ appartient à $\mathcal{A}_{U_i}$
pour tout $i \in I$.
\item
\label{item-bounded-dimension}
il existe un sous-ensemble $\mathcal{B} \subset \Ob(\mathcal{C})$
tel que
\begin{enumerate}
\item tout objet de $\mathcal{C}$ admet un recouvrement dont les membres
appartiennent à $\mathcal{B}$, et
\item pour tout $V \in \mathcal{B}$, il existe un entier $d_V$
et un système cofinal $\text{Cov}_V$ de recouvrements de $V$ tel
que
$$
H^p(V_i, \mathcal{F}) = 0 \text{ pour }
\{V_i \to V\} \in \text{Cov}_V,\ p > d_V, \text{ et }
\mathcal{F} \in \Ob(\mathcal{A}_V)
$$
\end{enumerate}
\end{enumerate}
Remarquons que cette propriété doit être satisfaite pour $\mathcal{F}$ dans
$\mathcal{A}_V$ et pas seulement pour les objets « globaux »
(elle est donc plus forte que la condition de
Cohomologie sur les sites, situation \ref{sites-cohomology-situation-olsson-laszlo}).
Dans cette situation, il existe une sous-catégorie de Serre faible
$\mathcal{A} \subset \textit{Mod}(\mathcal{O}_\mathcal{C})$
constituée des objets dont la restriction à $\mathcal{C}/U$
appartient à $\mathcal{A}_U$ pour tout $U \in \Ob(\mathcal{C})$.
De plus, il existe des catégories dérivées
$D_\mathcal{A}(\mathcal{O}_\mathcal{C})$ et
$D_{\mathcal{A}_U}(\mathcal{O}_U)$, et les foncteurs de restriction
envoient ces catégories les unes dans les autres.

\begin{example}
\label{example-quasi-coherent-spaces-etale}
Soit $S$ un schéma et $X$ un espace algébrique sur $S$.
Soit $\mathcal{C} = X_{spaces, \etale}$ le site \'etale
sur la catégorie des espaces algébriques étales sur $X$; voir
Propriétés des espaces, définition
\ref{spaces-properties-definition-spaces-etale-site}.
Notons $\mathcal{O}_\mathcal{C}$ le faisceau structural, c'est-à-dire le faisceau
donné par la règle $U \mapsto \Gamma(U, \mathcal{O}_U)$.
Notons $\mathcal{A}_U$ la catégorie des $\mathcal{O}_U$-modules quasi-cohérents.
Soit $\mathcal{B} = \Ob(\mathcal{C})$ et, pour $V \in \mathcal{B}$,
posons $d_V = 0$ et notons $\text{Cov}_V$ l'ensemble des
recouvrements $\{V_i \to V\}$ tels que $V_i$ soit affine pour tout $i$.
Alors les hypothèses (1), (2) et (3) sont satisfaites.
Voir Propriétés des espaces, lemmes
\ref{spaces-properties-lemma-pullback-quasi-coherent} et
\ref{spaces-properties-lemma-properties-quasi-coherent}
pour les propriétés (1) et (2); l'annulation en (3) découle de
Cohomologie des schémas, lemme
\ref{coherent-lemma-quasi-coherent-affine-cohomology-zero}
et de la discussion de Cohomologie des espaces, section
\ref{spaces-cohomology-section-higher-direct-image}.
\end{example}

\begin{example}
\label{example-etale}
Soit $S$ l'un des types de schémas suivants :
\begin{enumerate}
\item le spectre d'un corps fini,
\item le spectre d'un corps séparablement clos,
\item le spectre d'un anneau local noethérien strictement hensélien,
\item le spectre d'un anneau local noethérien hensélien
à corps résiduel fini,
\item ajouter d'autres exemples ici.
\end{enumerate}
Soit $\Lambda$ un anneau fini dont l'ordre est inversible sur $S$.
Soit $\mathcal{C} \subset (\Sch/S)_\etale$
la sous-catégorie pleine formée des schémas localement de type fini
sur $S$, munis de la topologie étale.
Soit $\mathcal{O}_\mathcal{C} = \underline{\Lambda}$ le
faisceau constant. Posons $\mathcal{A}_U = \textit{Mod}(\mathcal{O}_U)$;
autrement dit, on considère tous les faisceaux étales de $\Lambda$-modules.
Soit $\mathcal{B} \subset \Ob(\mathcal{C})$
l'ensemble des objets quasi-compacts. Pour $V \in \mathcal{B}$, posons
$$
d_V = 1 + 2\dim(S) +
\sup\nolimits_{v \in V}(\text{trdeg}_{\kappa(s)}(\kappa(v)) +
2 \dim \mathcal{O}_{V, v})
$$
et notons $\text{Cov}_V$ l'ensemble des recouvrements étales $\{V_i \to V\}$
tels que $V_i$ soit quasi-compact pour tout $i$.
Le choix de la borne $d_V$ provient du théorème de Gabber
sur la dimension cohomologique. Pour vérifier que la condition (3)
est satisfaite avec ce choix, on utilise
\cite[Expos\'e VIII-A, Corollary 1.2 and Lemma 2.2]{Traveaux}
ainsi que des arguments élémentaires sur les dimensions cohomologiques des corps.
Nous ajoutons $1$ dans la formule parce que notre liste contient des cas où $S$
peut avoir un corps résiduel fini.
Nous reviendrons sur cet exemple plus tard (insérer la référence future).
\end{example}

\noindent
Soit $(\mathcal{C}, \mathcal{O}_\mathcal{C})$ un site annelé.
Supposons données des sous-catégories de Serre faibles
$\mathcal{A}_U \subset \textit{Mod}(\mathcal{O}_U)$
satisfassent à la condition (\ref{item-restriction}).
Alors
\begin{enumerate}
\item pour tout objet semi-représentable $K = \{U_i\}_{i \in I}$,
nous obtenons une sous-catégorie de Serre faible
$\mathcal{A}_K \subset \textit{Mod}(\mathcal{O}_K)$
en prenant $\prod \mathcal{A}_{U_i} \subset 
\prod \textit{Mod}(\mathcal{O}_{U_i}) = \textit{Mod}(\mathcal{O}_K)$, et
\item étant donné un morphisme d'objets semi-représentables
$f : K \to L$, le foncteur image inverse
$f^* : \textit{Mod}(\mathcal{O}_L) \to \textit{Mod}(\mathcal{O}_L)$
envoie $\mathcal{A}_L$ dans $\mathcal{A}_K$.
\end{enumerate}
Voir la remarque \ref{remark-semi-representable-ringed} pour les notations et les
explications. En particulier, pour un objet semi-représentable simplicial $K$,
l'expression « un objet $\mathcal{F}$ de
$\textit{Mod}(\mathcal{O})$, au sens de la remarque \ref{remark-augmentation-ringed},
dont les restrictions $\mathcal{F}_n$ appartiennent à
$\mathcal{A}_{K_n}$ pour tout $n$ » est sans ambiguïté.

\begin{lemma}
\label{lemma-hypercovering-equivalence-modules}
Soit $(\mathcal{C}, \mathcal{O}_\mathcal{C})$ un site annelé.
Supposons données des sous-catégories de Serre faibles
$\mathcal{A}_U \subset \textit{Mod}(\mathcal{O}_U)$
satisfaisant aux conditions (\ref{item-restriction}),
(\ref{item-local}) et (\ref{item-bounded-dimension}) ci-dessus.
Supposons que $\mathcal{C}$ admette des égalisateurs et des produits fibrés, et
soit $K$ un hyperrecouvrement.
Considérons $((\mathcal{C}/K)_{total}, \mathcal{O})$ comme dans
la remarque \ref{remark-augmentation-ringed}.
Soit $\mathcal{A}_{total} \subset \textit{Mod}(\mathcal{O})$
la sous-catégorie de Serre faible formée des $\mathcal{O}$-modules cartésiens
$\mathcal{F}$ tels que la restriction $\mathcal{F}_n$ appartienne à
$\mathcal{A}_{K_n}$ pour tout $n$ (au sens défini ci-dessus).
Alors le foncteur $La^*$ définit une équivalence
$$
D_\mathcal{A}(\mathcal{O}_\mathcal{C})
\longrightarrow
D_{\mathcal{A}_{total}}(\mathcal{O})
$$
de quasi-inverse $Ra_*$.
\end{lemma}

\begin{proof}
Les $\mathcal{O}$-modules cartésiens forment une sous-catégorie de Serre faible d'après le
lemme \ref{lemma-Serre-subcat-cartesian-modules}
(les hypothèses requises découlent de la discussion de la
remarque \ref{remark-augmentation-ringed}).
Puisque les foncteurs de restriction
$g_n^* : \textit{Mod}(\mathcal{O}) \to \textit{Mod}(\mathcal{O}_n)$
sont exacts, il s'ensuit que $\mathcal{A}_{total}$ est une
sous-catégorie de Serre faible.

\medskip\noindent
Montrons que $a^* : \mathcal{A} \to \mathcal{A}_{total}$
est une équivalence de catégories, de quasi-inverse $La_*$.
On sait déjà que $La_*a^*\mathcal{F} = \mathcal{F}$ par la
version bornée
(lemme \ref{lemma-hypercovering-equivalence-bounded-modules}).
Il est clair que $a^*\mathcal{F}$ appartient à $\mathcal{A}_{total}$
pour tout $\mathcal{F}$ dans $\mathcal{A}$. Réciproquement, soit
$\mathcal{G} \in \mathcal{A}_{total}$. Comme $\mathcal{G}$
est cartésien, on a $\mathcal{G} = a^*\mathcal{F}$
pour un certain $\mathcal{O}_\mathcal{C}$-module $\mathcal{F}$, d'après le
lemme \ref{lemma-hypercovering-descent-modules}.
Nous voulons montrer que $\mathcal{F}$ appartient à $\mathcal{A}$.
Fixons $U \in \Ob(\mathcal{C})$. Il faut montrer que la restriction
de $\mathcal{F}$ à $\mathcal{C}/U$ appartient à $\mathcal{A}_U$.
Comme d'habitude, écrivons $K_0 = \{U_{0, i}\}_{i \in I_0}$.
Puisque $K$ est un hyperrecouvrement, le morphisme $\coprod_{i \in I_0} h_{U_{0, i}} \to *$
devient surjectif après passage au faisceau associé. Il existe donc
un recouvrement $\{U_j \to U\}_{j \in J}$ et une application $\tau : J \to I_0$
et pour chaque $j \in J$ un morphisme $\varphi_j : U_j \to U_{0, \tau(j)}$.
Puisque $\mathcal{G}_0 = a_0^*\mathcal{F}$, on voit
que la restriction de $\mathcal{F}$ à $\mathcal{C}/U_j$
est égale à la restriction de la composante $\tau(j)$ de
$\mathcal{G}_0$ à $\mathcal{C}/U_j$ via le morphisme
$\varphi_j : U_j \to U_{0, \tau(i)}$. Par conséquent,
la condition (\ref{item-restriction}) donne que $\mathcal{F}|_{\mathcal{C}/U_j}$
appartient à $\mathcal{A}_{U_j}$; la condition
(\ref{item-local}) donne alors que $\mathcal{F}|_{\mathcal{C}/U}$
appartient à $\mathcal{A}_U$.

\medskip\noindent
En particulier, l'énoncé du lemme est bien défini.
Le lemme découle alors de Cohomologie sur les sites,
lemme \ref{sites-cohomology-lemma-equivalence-unbounded-one}.
L'hypothèse (1) est claire (voir la remarque \ref{remark-augmentation-ringed}).
Les hypothèses (2) et (3) ont été démontrées au paragraphe précédent.
L'hypothèse (4) découle immédiatement de (\ref{item-bounded-dimension}).
Pour l'hypothèse (5), soit $\mathcal{B}_{total}$ l'ensemble des objets
$U/U_{n, i}$ du site $(\mathcal{C}/K)_{total}$
tels que $U \in \mathcal{B}$, où $\mathcal{B}$ est comme dans
(\ref{item-bounded-dimension}). Nous utilisons ici la description de
la description du site $(\mathcal{C}/K)_{total}$ donnée dans la section
\ref{section-simplicial-semi-representable}.
De plus, posons $\text{Cov}_{U/U_{n, i}}$ égal à $\text{Cov}_U$
et $d_{U/U_{n, i}}$ égal à $d_U$, où $\text{Cov}_U$ et $d_U$
sont fournis par (\ref{item-bounded-dimension}).
Avec ces choix, la condition (5) est satisfaite.
Cela découle immédiatement du lemme
\ref{lemma-sanity-check-simplicial-semi-representable}
et du fait que $\mathcal{F} \in \mathcal{A}_{total}$
implique $\mathcal{F}_n \in \mathcal{A}_{K_n}$ et donc
$\mathcal{F}_{n, i} \in \mathcal{A}_{U_{n, i}}$.
(Le lecteur préoccupé par la différence entre la cohomologie
des faisceaux abéliens et la cohomologie
des faisceaux de modules peut consulter Cohomologie sur les sites, lemme
\ref{sites-cohomology-lemma-cohomology-modules-abelian-agree}.)
\end{proof}










\section{Recollement de complexes}
\label{section-glueing-complexes}

\noindent
Cette section prolonge
Cohomologie, section \ref{cohomology-section-glueing-complexes}.
Le but est de prouver une légère généralisation de \cite[Theorem 3.2.4]{BBD}.
Notre méthode différera légèrement, en ce qu'elle utilise
les résultats des sections \ref{section-glueing} et
\ref{section-glueing-modules}. Nous redémontrerons également la
version non bornée en suivant la démonstration de \cite{six-I}.

\medskip\noindent
Conseil au lecteur : nous suggérons de lire d'abord l'énoncé du
lemme \ref{lemma-bbd-glueing-cech}, ainsi que sa seconde démonstration.

\medskip\noindent
Voici la situation qui nous intéresse.

\begin{situation}
\label{situation-locally-given}
Soit $(\mathcal{C}, \mathcal{O}_\mathcal{C})$ un site annelé. On nous donne
\begin{enumerate}
\item une catégorie $\mathcal{B}$ et un foncteur
$u : \mathcal{B} \to \mathcal{C}$,
\item un objet $E_U$ dans $D(\mathcal{O}_{u(U)})$ pour $U \in \Ob(\mathcal{B})$,
\item un isomorphisme $\rho_a : E_U|_{\mathcal{C}/u(V)} \to E_V$ dans
$D(\mathcal{O}_{u(V)})$ pour $a : V \to U$ dans $\mathcal{B}$
\end{enumerate}
avec la compatibilité suivante : pour toutes flèches composables
$b : W \to V$ et $a : V \to U$ de $\mathcal{B}$, on a
$\rho_{a \circ b} = \rho_b \circ \rho_a|_{\mathcal{C}/u(W)}$.
\end{situation}

\noindent
Nous ne pourrons rien démontrer dans ce cadre sans faire davantage
d'hypothèses. Un cas intéressant est celui où $\mathcal{B}$ est une sous-catégorie pleine
telle que chaque objet de $\mathcal{C}$ a un recouvrement
dont les membres sont des objets de $\mathcal{B}$ (c'est le cas étudié
dans \cite{BBD}). Il importe ici d'autoriser aussi les situations où
ce n'est pas le cas; la principale variante est celle où l'on dispose d'un morphisme
de sites $f : \mathcal{C} \to \mathcal{D}$ et $\mathcal{B}$
est une sous-catégorie pleine de $\mathcal{D}$ telle que chaque objet de
$\mathcal{D}$ admet un recouvrement dont les membres sont des objets de $\mathcal{B}$.

\medskip\noindent
Dans la situation \ref{situation-locally-given}, une {\it solution}
sera une paire $(E, \rho_U)$ où $E$ est un objet de
$D(\mathcal{O}_\mathcal{C})$
et où les morphismes $\rho_U : E|_{\mathcal{C}/u(U)} \to E_U$
pour $U \in \Ob(\mathcal{B})$
sont des isomorphismes tels que
l'on a $\rho_a \circ \rho_U|_{\mathcal{C}/u(V)} = \rho_V$
pour $a : V \to U$ dans $\mathcal{B}$.

\begin{lemma}
\label{lemma-prepare-bbd-glueing}
Dans la situation \ref{situation-locally-given}.
Supposons que les auto-Ext négatifs de $E_U$ dans $D(\mathcal{O}_{u(U)})$ soient nuls.
Soit $L$ un objet simplicial de $\text{SR}(\mathcal{B})$.
Considérons l'objet simplicial $K = u(L)$ de $\text{SR}(\mathcal{C})$
et considérons $((\mathcal{C}/K)_{total}, \mathcal{O})$ comme dans la
remarque \ref{remark-augmentation-ringed}.
Il existe un objet cartésien $E$ de $D(\mathcal{O})$
tel que, si l'on écrit $L_n = \{U_{n, i}\}_{i \in I_n}$,
l'image de $E$ dans $D(\mathcal{O}_{\mathcal{C}/u(U_{n, i})})$
s'identifie de manière compatible à $E_{U_{n, i}}$ (voir la preuve).
De plus, $E$ est unique à isomorphisme unique près.
\end{lemma}

\begin{proof}
Rappelons que
$\Sh(\mathcal{C}/K_n) = \prod_{i \in I_n} \Sh(\mathcal{C}/u(U_{n, i}))$
et de même pour les catégories de modules. Cette décomposition en produit
se transporte également aux catégories dérivées de faisceaux de modules.
De plus, cette décomposition en produit est compatible avec
les morphismes de l'objet semi-représentable simplicial $K$.
Voir la section \ref{section-semi-representable}.
Nous pouvons donc définir $E_n = \prod_{i \in I_n} E_{U_{n, i}}$
(produit « formel ») dans $D(\mathcal{O}_n)$.
En prenant les produits (formels) des morphismes $\rho_a$ de
la situation \ref{situation-locally-given},
nous obtenons des isomorphismes $E_\varphi : f_\varphi^*E_n \to E_m$.
L'hypothèse sur les composés des morphismes $\rho_a$
implique immédiatement que $(E_n, E_\varphi)$
définit un système simplicial de la catégorie dérivée de modules
comme dans la définition \ref{definition-cartesian-derived-modules}.
L'annulation des auto-Ext négatifs supposée dans le lemme entraîne que
$\Hom(E_n[t], E_n) = 0$ pour $n \geq 0$ et $t > 0$.
Par le
lemme \ref{lemma-cartesian-module-derived-from-simplicial},
on obtient $E$.
L'unicité jusqu'à l'isomorphisme unique découle des lemmes
\ref{lemma-nullity-cartesian-modules-derived} et
\ref{lemma-hom-cartesian-modules-derived}.
\end{proof}

\begin{lemma}[Lemme de recollement BBD]
\label{lemma-bbd-glueing}
Dans la situation \ref{situation-locally-given}. Supposons que
\begin{enumerate}
\item $\mathcal{C}$ admet des égalisateurs et des produits fibrés,
\item il existe un morphisme de sites $f : \mathcal{C} \to \mathcal{D}$
donné par un foncteur continu $u : \mathcal{D} \to \mathcal{C}$
tel que
\begin{enumerate}
\item $\mathcal{D}$ admet des égalisateurs et des produits fibrés, et $u$
commute à leur formation,
\item $\mathcal{B}$ est une sous-catégorie pleine de $\mathcal{D}$
et $u : \mathcal{B} \to \mathcal{C}$ est la restriction de $u$,
\item tout objet de $\mathcal{D}$ admet un recouvrement dont les membres
sont des objets de $\mathcal{B}$,
\end{enumerate}
\item pour tout $U$ dans $\mathcal{B}$, tous les auto-Ext négatifs de $E_U$
dans $D(\mathcal{O}_{u(U)})$ sont nuls, et
\item il existe un $t \in \mathbf{Z}$ tel que $H^i(E_U) = 0$ pour $i < t$
et tout $U \in \Ob(\mathcal{B})$.
\end{enumerate}
Il existe alors une solution unique à isomorphisme unique près.
\end{lemma}

\begin{proof}
D'après Hyperrecouvrements, lemme \ref{hypercovering-lemma-hypercovering-site},
il existe un hyperrecouvrement $L$ pour le site $\mathcal{D}$ tel que
$L_n = \{U_{n, i}\}_{i \in I_n}$ avec $U_{i, n} \in \Ob(\mathcal{B})$.
Posons $K = u(L)$. Le lemme \ref{lemma-prepare-bbd-glueing} donne
un objet cartésien $E$ de $D(\mathcal{O})$ sur le site
$(\mathcal{C}/K)_{total}$, dont la restriction est $E_{U_{n, i}}$ sur
$\mathcal{C}/u(U_{n, i})$, de façon compatible.
L'hypothèse sur $t$ implique que $E \in D^+(\mathcal{O})$.
D'après Hyperrecouvrements, lemme \ref{hypercovering-lemma-hypercovering-morphism-sites},
$K$ est lui aussi un hyperrecouvrement.
Le lemme \ref{lemma-hypercovering-equivalence-bounded-modules}
donne $E = a^*F$ pour un certain $F$ dans $D^+(\mathcal{O}_\mathcal{C})$.

\medskip\noindent
Pour prouver que $F$ est une solution, nous utiliserons la construction de
$L_0$ et $L_1$ donnée dans la preuve de
Hyperrecouvrements, lemme \ref{hypercovering-lemma-hypercovering-site}.
(C'est un peu inélégant, mais il ne semble pas y avoir de moyen
tout à fait direct de l'éviter.)

\medskip\noindent
Plus précisément, $I_0 = \Ob(\mathcal{B})$ et donc
$L_0 = \{U\}_{U \in \Ob(\mathcal{B})}$.
La restriction de l'isomorphisme $a^*F \to E$ aux composantes
$\mathcal{C}/u(U)$ de $\mathcal{C}/K_0$ définit les isomorphismes
$\rho_U : F|_{\mathcal{C}/u(U)} \to E_U$ pour $U \in \Ob(\mathcal{B})$
par construction de $E$.

\medskip\noindent
Pour vérifier que les $\rho_U$ satisfont à la condition de compatibilité
avec les morphismes $\rho_a$ de la situation \ref{situation-locally-given},
on utilise le fait que $I_1$ contient l'ensemble
$$
\Omega =
\{(U, V, W, a, b) \mid U, V, W \in \mathcal{B}, a : U \to V, b : U \to W\}
$$
et que, pour $i = (U, V, W, a, b)$ dans $\Omega$, on a
$U_{1, i} = U$. De plus, les composantes $f_{\delta^1_0, i}$ et
$f_{\delta^1_1, i}$ des deux morphismes $K_1 \to K_0$ sont les morphismes
$$
a : U \to V \quad\text{et}\quad b : U \to V
$$
La compatibilité du
lemme \ref{lemma-prepare-bbd-glueing} donne donc
$$
\rho_a \circ \rho_V|_{\mathcal{C}/u(U)} = \rho_U
\quad\text{et}\quad
\rho_b \circ \rho_W|_{\mathcal{C}/u(U)} = \rho_U
$$
En prenant par exemple $i = (U, V, U, a, \text{id}_U) \in \Omega$, on obtient
la compatibilité voulue. L'unicité de $F$ découle
de celle de $E$ dans le lemme précédent (nous omettons ce détail).
\end{proof}

\begin{lemma}[Lemme de recollement BBD non borné]
\label{lemma-bbd-unbounded-glueing}
Dans la situation \ref{situation-locally-given}. Supposons que
\begin{enumerate}
\item $\mathcal{C}$ admet des égalisateurs et des produits fibrés,
\item il existe un morphisme de sites $f : \mathcal{C} \to \mathcal{D}$
donné par un foncteur continu $u : \mathcal{D} \to \mathcal{C}$
tel que
\begin{enumerate}
\item $\mathcal{D}$ admet des égalisateurs et des produits fibrés, et $u$
commute à leur formation,
\item $\mathcal{B}$ est une sous-catégorie pleine de $\mathcal{D}$
et $u : \mathcal{B} \to \mathcal{C}$ est la restriction de $u$,
\item chaque objet de $\mathcal{D}$ possède un recouvrement dont les membres
sont des objets de $\mathcal{B}$,
\end{enumerate}
\item tous les auto-Ext négatifs de $E_U$ dans $D(\mathcal{O}_{u(U)})$ sont nuls, et
\item il existe des sous-catégories de Serre faibles
$\mathcal{A}_U \subset \textit{Mod}(\mathcal{O}_U)$ pour tout
$U \in \Ob(\mathcal{C})$, satisfaisant aux conditions (\ref{item-restriction}),
(\ref{item-local}) et (\ref{item-bounded-dimension}),
\item $E_U \in D_{\mathcal{A}_U}(\mathcal{O}_U)$.
\end{enumerate}
Il existe alors une solution unique à isomorphisme unique près.
\end{lemma}

\begin{proof}
La démonstration est {\bf exactement} la même que celle du
lemme \ref{lemma-bbd-glueing}. La seule différence est que
$E$ est un objet de $D_{\mathcal{A}_{total}}(\mathcal{O})$
et l'on utilise donc le lemme \ref{lemma-hypercovering-equivalence-modules}
pour obtenir $F$ avec $E = a^*F$
au lieu du lemme \ref{lemma-hypercovering-equivalence-bounded-modules}.
\end{proof}

\noindent
Voici un exemple d’application de la théorie générale ci-dessus.

\begin{lemma}
\label{lemma-bbd-glueing-cech}
\begin{reference}
Courriel de Martin Olsson daté du 9 septembre 2021.
\end{reference}
Soit $(\mathcal{C}, \mathcal{O}_\mathcal{C})$ un site annelé.
Supposons que $\mathcal{C}$ admette des produits fibrés.
Soit $\{U_i \to X\}_{i \in I}$ un recouvrement dans $\mathcal{C}$.
Pour $i \in I$, soit $E_i$ un objet de $D(\mathcal{O}_{U_i})$, et
pour $i, j \in I$, soit
$$
\rho_{ij} : E_i|_{\mathcal{C}/U_{ij}} \longrightarrow E_j|_{\mathcal{C}/U_{ij}}
$$
est un isomorphisme dans $D(\mathcal{O}_{U_{ij}})$, où
$U_{ij} = U_i \times_X U_j$. Supposons que
\begin{enumerate}
\item les $\rho_{ij}$ satisfassent à la condition de cocycle sur
$U_i \times_X U_j \times_X U_k$ pour tout $i, j, k \in I$,
\item $\SheafExt^p_{\mathcal{O}_{U_i}}(E_i, E_i) = 0$
pour tout $p < 0$ et tout $i \in I$, et
\item il existe un $t \in \mathbf{Z}$ tel que
$H^p(E_i) = 0$ pour $p < t$ et tout $i \in I$.
\end{enumerate}
Alors il existe un unique couple $(E, \rho_i)$, où $E$ est un objet de
$D(\mathcal{O}_X)$ et les $\rho_i : E|_{U_i} \to E_i$ sont des isomorphismes dans
$D(\mathcal{O}_{U_i})$ compatibles avec les $\rho_{ij}$.
\end{lemma}

\begin{proof}[Première démonstration]
Dans cette preuve, nous déduisons le résultat du lemme général
\ref{lemma-bbd-glueing}. Nous invitons le lecteur à consulter plutôt
la seconde démonstration.

\medskip\noindent
Nous pouvons remplacer $\mathcal{C}$ par $\mathcal{C}/X$.
On peut donc supposer que $X$ est l'objet final de $\mathcal{C}$
et que $\mathcal{C}$ admet toutes les limites finies.

\medskip\noindent
Soit $\mathcal{B}$ la sous-catégorie pleine de $\mathcal{C}$
constituée des objets $U \in \Ob(\mathcal{C})$ pour lesquels il existe
un $i(U) \in I$ et un morphisme $a_U : U \to U_{i(U)}$.
Notons $E_U = a_U^*E_{i(U)}$ dans $D(\mathcal{O}_U)$,
image inverse (restriction) de $E_i$ par $a_U$.
Pour un morphisme $a : U \to U'$ de $\mathcal{B}$,
on obtient un morphisme
$(a_{U'} \circ a, a_U) : U \to U_{i(U')} \times_X U_{i(U)} = U_{i(U')i(U)}$
et donc un isomorphisme
$$
\rho_a :
a^*E_{U'} = a^*a_{U'}^*E_{i(U')}
\xrightarrow{(a_{U'} \circ a, a_U)^*\rho_{i(U')i(U)}}
a_{U}^*E_{i(U)} = E_{U}
$$
dans $D(\mathcal{O}_U)$. Les données $\mathcal{B}, E_U, \rho_a$
sont celles de la situation \ref{situation-locally-given}; les
isomorphismes $\rho_a$ satisfont la condition de cocycle
précisément grâce à la condition (1) de l'énoncé du lemme
(détails omis).

\medskip\noindent
Nous allons appliquer le lemme \ref{lemma-bbd-glueing} avec $\mathcal{B}$,
$E_U$, $\rho_a$ comme ci-dessus et avec $\mathcal{D} = \mathcal{C}$ et
$f : \mathcal{C} \to \mathcal{D}$ le morphisme identité.
Les hypothèses (1) et (2)(a) du lemme \ref{lemma-bbd-glueing} ont été vérifiées ci-dessus.
L’hypothèse (2)(b) du lemme \ref{lemma-bbd-glueing} est claire.
L'hypothèse (2)(c) du lemme \ref{lemma-bbd-glueing} est satisfaite, car
$\{U_i \to X\}$ est un recouvrement\footnote{En fait, il suffirait
que le morphisme $\coprod_{i \in I} h_{U_i} \to h_X$ devienne
surjectif après passage au faisceau associé; le lemme resterait vrai dans ce cas
avec la même preuve.}.
L'hypothèse (3) du lemme \ref{lemma-bbd-glueing} est satisfaite,
car nous avons supposé l'annulation de tous les faisceaux d'Ext négatifs
de $E_i$; cela implique, pour tout objet $U$ au-dessus de $U_i$, que
les auto-Ext négatifs de $E_i|_U$ sont nuls.
L'hypothèse (4) du lemme \ref{lemma-bbd-glueing} est satisfaite, car
nous avons supposé que les faisceaux de cohomologie de chaque $E_i$ sont nuls
en degrés strictement inférieurs à $t$.

\medskip\noindent
Nous obtenons une solution unique $(E, \rho_U)$.
En posant $\rho_i = \rho_{U_i}$, on obtient le lemme.
\end{proof}

\begin{proof}[Seconde démonstration]
Nous esquissons une preuve plus directe.
Notons $K$ l'hyperrecouvrement de Čech de $X$ associé au
recouvrement $\{U_i \to X\}_{i \in I}$; voir
Hyperrecouvrements, exemple \ref{hypercovering-example-cech}.
Ainsi par exemple $K_0 = \{U_i \to X\}_{i \in I}$ et
$K_1 = \{U_i \times_X U_j \to X\}_{i, j \in I}$ et ainsi de suite.
Considérons $((\mathcal{C}/K)_{total}, \mathcal{O})$, $a$ et $a_n$ comme dans la
remarque \ref{remark-augmentation-ringed-over-object}.
Les objets $E_i$ déterminent un objet $M_0$ dans
$D(\mathcal{O}_0) = \prod D(\mathcal{O}_{U_i})$.
De même, les isomorphismes $\rho_{ij}$ déterminent un isomorphisme
$$
\alpha :
L(f_{\delta_1^1})^*M_0
\longrightarrow
L(f_{\delta_0^1})^*M_0
$$
dans $D(\mathcal{O}_1)$ satisfaisant la condition de cocycle.
Par le lemme \ref{lemma-construct-cartesian-objects-derived-modules}
on obtient un système simplicial cartésien $(M_n)$ de la catégorie dérivée.
L'annulation supposée des faisceaux d'Ext négatifs montre que
les auto-Ext négatifs des objets $M_n$ sont nuls.
On obtient donc un objet cartésien $M$ de $D(\mathcal{O})$
dont le système simplicial associé est isomorphe à $(M_n)$ d'après le
lemme \ref{lemma-cartesian-module-derived-from-simplicial}.
Puisque les faisceaux de cohomologie de $M$ sont nuls en degrés $< t$,
le lemme \ref{lemma-hypercovering-X-equivalence-bounded-modules}
donne $M = La^*E$ pour un certain $E$ dans $D(\mathcal{O}_X)$.
L'isomorphisme $La^*E \to M$ restreint à $\mathcal{C}/U_i$
induit les isomorphismes $\rho_i$. Nous omettons la vérification
de la compatibilité avec les isomorphismes $\rho_{ij}$.
\end{proof}












\section{Hyperrecouvrements propres en topologie}
\label{section-proper-hypercovering}

\noindent
Travaillons dans la catégorie $\textit{LC}$ des espaces topologiques séparés et localement
quasi compacts, munie des applications continues; voir
Cohomologie sur les sites, section \ref{sites-cohomology-section-cohomology-LC}.
Soient $X$ un objet de $\textit{LC}$ et $U$ un objet
simplicial de $\textit{LC}$, muni d'une augmentation
$$
a : U \to X
$$
On dit que $U$ est un {\it hyperrecouvrement propre} de $X$ si
\begin{enumerate}
\item $U_0 \to X$ est un morphisme propre surjectif,
\item $U_1 \to U_0 \times_X U_0$ est un morphisme propre surjectif,
\item $U_{n + 1} \to (\text{cosk}_n\text{sk}_n U)_{n + 1}$
est un morphisme propre surjectif pour $n \geq 1$.
\end{enumerate}
La catégorie $\textit{LC}$ admet toutes les limites finies; les cosquelettes
utilisés dans la formulation ci-dessus existent donc.
$$
\fbox{Principe : les hyperrecouvrements propres permettent de calculer la cohomologie.}
$$
Une idée essentielle de la démonstration consiste à munir
$\textit{LC}$ d'une topologie plus fine que la topologie usuelle, telle que :
(a) tout morphisme propre surjectif définisse un recouvrement;
(b) la cohomologie, pour cette topologie plus fine, des faisceaux usuels
coïncide avec leur cohomologie usuelle. Les propriétés
(a) et (b) sont valables pour la topologie qc, voir Cohomologie
sur les sites, section \ref{sites-cohomology-section-cohomology-LC}.
Une fois (a) et (b) établis, le principe découle
des résultats précédents de ce chapitre.

\begin{lemma}
\label{lemma-compare-simplicial-objects}
Soit $U$ un objet simplicial de $\textit{LC}$ et soit $a : U \to X$
une augmentation. Il existe un diagramme commutatif
$$
\xymatrix{
\Sh((\textit{LC}_{qc}/U)_{total}) \ar[r]_-h \ar[d]_{a_{qc}} &
\Sh(U_{Zar}) \ar[d]^a \\
\Sh(\textit{LC}_{qc}/X) \ar[r]^-{h_{-1}} &
\Sh(X)
}
$$
où la flèche verticale gauche est définie dans la section
\ref{section-hypercovering}
et la flèche verticale droite est définie dans le lemme
\ref{lemma-augmentation}.
\end{lemma}

\begin{proof}
Écrivons $\Sh(X) = \Sh(X_{Zar})$. Observons que
$(\textit{LC}_{qc}/U)_{total}$ et $U_{Zar}$ relèvent tous deux
du cas (A) de la situation \ref{situation-simplicial-site}.
Cela résulte immédiatement de la construction de
$U_{Zar}$ dans la section \ref{section-simplicial-top}
et découle du lemme \ref{lemma-sr-when-fibre-products}
pour $(\textit{LC}_{qc}/U)_{total}$.
Considérons alors les foncteurs
$U_{n, Zar} \to \textit{LC}_{qc}/U_n$, $U \mapsto U/U_n$
et
$X_{Zar} \to \textit{LC}_{qc}/X$, $U \mapsto U/X$.
Ces foncteurs définissent des morphismes de sites, comme on l'a vu
dans Cohomologie sur les sites, section \ref{sites-cohomology-section-cohomology-LC}.
On obtient ainsi un morphisme de sites simpliciaux compatible avec les augmentations
comme dans la remarque \ref{remark-morphism-augmentation-simplicial-sites}
et on peut appliquer le lemme
\ref{lemma-morphism-augmentation-simplicial-sites} pour conclure.
\end{proof}

\begin{lemma}
\label{lemma-descent-sheaves-for-proper-hypercovering}
Soit $U$ un objet simplicial de $\textit{LC}$, muni d'une augmentation $a : U \to X$.
Si $a : U \to X$ définit un hyperrecouvrement propre de $X$,
alors
$$
a^{-1} : \Sh(X) \to \Sh(U_{Zar})
\quad\text{et}\quad
a^{-1} : \textit{Ab}(X) \to \textit{Ab}(U_{Zar})
$$
sont pleinement fidèles, d'image essentielle les faisceaux cartésiens, et
admettent $a_*$ pour quasi-inverse. Ici, $a : \Sh(U_{Zar}) \to \Sh(X)$ est comme dans le lemme
\ref{lemma-augmentation}.
\end{lemma}

\begin{proof}
Nous démontrerons l'énoncé pour les faisceaux d'ensembles. Ce sera une
conséquence presque formelle des résultats déjà établis.
Considérons le diagramme du lemme \ref{lemma-compare-simplicial-objects}.
D'après Cohomologie sur les sites, lemme \ref{sites-cohomology-lemma-describe-pullback-pi},
le foncteur $(h_{-1})^{-1}$ est pleinement fidèle, de quasi-inverse $h_{-1, *}$.
Il en va de même pour les composantes $h_n$ de $h$.
Par la description des foncteurs $h^{-1}$ et $h_*$ du lemme
\ref{lemma-morphism-simplicial-sites}
on conclut que $h^{-1}$ est pleinement fidèle, de quasi-inverse $h_*$.
Observons que $U$ est un hyperrecouvrement de $X$ dans $\textit{LC}_{qc}$
(au sens de la section \ref{section-hypercovering}), d'après
Cohomologie sur les sites, lemme
\ref{sites-cohomology-lemma-proper-surjective-is-qc-covering}.
Par le lemme \ref{lemma-hypercovering-X-simple-descent-sheaves}
on voit que $a_{qc}^{-1}$ est pleinement fidèle, de quasi-inverse $a_{qc, *}$
et d'image essentielle les faisceaux cartésiens sur
$(\textit{LC}_{qc}/U)_{total}$.
Une chasse formelle dans le diagramme montre alors que
$a^{-1}$ est pleinement fidèle.

\medskip\noindent
Enfin, supposons que $\mathcal{G}$ soit un faisceau cartésien sur $U_{Zar}$.
Alors $h^{-1}\mathcal{G}$ est un faisceau cartésien sur $\textit{LC}_{qc}/U$.
Ainsi, $h^{-1}\mathcal{G} = a_{qc}^{-1}\mathcal{H}$ pour un certain faisceau
$\mathcal{H}$ sur $\textit{LC}_{qc}/X$.
Calculons :
\begin{align*}
(h_{-1})^{-1}(a_*\mathcal{G})
& =
(h_{-1})^{-1}
\text{Eq}(
\xymatrix{
a_{0, *}\mathcal{G}_0
\ar@<1ex>[r] \ar@<-1ex>[r] &
a_{1, *}\mathcal{G}_1
}
) \\
& =
\text{Eq}(
\xymatrix{
(h_{-1})^{-1}a_{0, *}\mathcal{G}_0
\ar@<1ex>[r] \ar@<-1ex>[r] &
(h_{-1})^{-1}a_{1, *}\mathcal{G}_1
}
) \\
& =
\text{Eq}(
\xymatrix{
a_{qc, 0, *}h_0^{-1}\mathcal{G}_0
\ar@<1ex>[r] \ar@<-1ex>[r] &
a_{qc, 1, *}h_1^{-1}\mathcal{G}_1
}
) \\
& =
\text{Eq}(
\xymatrix{
a_{qc, 0, *}a_{qc, 0}^{-1}\mathcal{H}
\ar@<1ex>[r] \ar@<-1ex>[r] &
a_{qc, 1, *}a_{qc, 1}^{-1}\mathcal{H}
}
) \\
& =
a_{qc, *}a_{qc}^{-1}\mathcal{H} \\
& =
\mathcal{H}
\end{align*}
La première égalité découle du lemme \ref{lemma-augmentation};
la deuxième, de l'exactitude du foncteur $(h_{-1})^{-1}$;
la troisième, de
Cohomologie sur les sites, lemme \ref{sites-cohomology-lemma-push-pull-LC}
(ici nous utilisons que $a_0 : U_0 \to X$ et $a_1: U_1 \to X$ sont propres),
la quatrième, de $a_{qc}^{-1}\mathcal{H} = h^{-1}\mathcal{G}$;
la cinquième, du lemme \ref{lemma-augmentation-site}; enfin la sixième
a été établie ci-dessus. Puisque $a_{qc}^{-1}\mathcal{H} = h^{-1}\mathcal{G}$,
on en déduit que $h^{-1}\mathcal{G} \cong h^{-1}a^{-1}a_*\mathcal{G}$,
ce qui termine la preuve par la pleine fidélité de $h^{-1}$.
\end{proof}

\begin{lemma}
\label{lemma-cohomological-descent-for-proper-hypercovering}
Soit $U$ un objet simplicial de $\textit{LC}$ et soit $a : U \to X$
une augmentation. Si $a : U \to X$ est un hyperrecouvrement propre de $X$,
alors pour $K \in D^+(X)$
$$
K \to Ra_*(a^{-1}K)
$$
est un isomorphisme, où $a : \Sh(U_{Zar}) \to \Sh(X)$ est comme dans le lemme
\ref{lemma-augmentation}.
\end{lemma}

\begin{proof}
Considérons le diagramme du lemme \ref{lemma-compare-simplicial-objects}.
Observons que $Rh_{n, *}h_n^{-1}$ est le foncteur identité
sur $D^+(U_n)$ d'après Cohomologie sur les sites, lemme
\ref{sites-cohomology-lemma-cohomological-descent-LC}.
Par conséquent, $Rh_*h^{-1}$ est le foncteur identité sur
$D^+(U_{Zar})$ par
le lemme \ref{lemma-direct-image-morphism-simplicial-sites}.
Nous avons
\begin{align*}
Ra_*(a^{-1}K)
& =
Ra_*Rh_*h^{-1}a^{-1}K \\
& =
Rh_{-1, *}Ra_{qc, *}a_{qc}^{-1}(h_{-1})^{-1}K \\
& =
Rh_{-1, *}(h_{-1})^{-1}K \\
& =
K
\end{align*}
La première égalité découle de la discussion ci-dessus; la deuxième,
de la commutativité du diagramme du lemme
\ref{lemma-compare-simplicial-objects}; la troisième, du lemme
\ref{lemma-hypercovering-X-simple-descent-bounded-abelian}
($U$ est un hyperrecouvrement de $X$ dans $\textit{LC}_{qc}$ d'après Cohomologie
sur les sites, lemme
\ref{sites-cohomology-lemma-proper-surjective-is-qc-covering}),
et la dernière, du lemme déjà utilisé de Cohomologie sur les sites
\ref{sites-cohomology-lemma-cohomological-descent-LC}.
\end{proof}

\begin{lemma}
\label{lemma-compute-via-proper-hypercovering}
Soit $U$ un objet simplicial de $\textit{LC}$ et soit $a : U \to X$
une augmentation. Si $U$ est un hyperrecouvrement propre de $X$, alors
$$
R\Gamma(X, K) = R\Gamma(U_{Zar}, a^{-1}K)
$$
pour $K \in D^+(X)$, où $a : \Sh(U_{Zar}) \to \Sh(X)$
est comme dans le lemme \ref{lemma-augmentation}.
\end{lemma}

\begin{proof}
Cela découle du lemme
\ref{lemma-cohomological-descent-for-proper-hypercovering}
car $R\Gamma(U_{Zar}, -) = R\Gamma(X, -) \circ Ra_*$ par
Cohomologie sur les sites, remarque \ref{sites-cohomology-remark-before-Leray}.
\end{proof}

\begin{lemma}
\label{lemma-proper-hypercovering-equivalence-bounded}
Soit $U$ un objet simplicial de $\textit{LC}$ et soit $a : U \to X$
une augmentation.
Soit $\mathcal{A} \subset \textit{Ab}(U_{Zar})$
la sous-catégorie de Serre faible des faisceaux abéliens cartésiens.
Si $U$ est un hyperrecouvrement propre de $X$, alors
le foncteur $a^{-1}$ définit une équivalence
$$
D^+(X) \longrightarrow D_\mathcal{A}^+(U_{Zar})
$$
de quasi-inverse $Ra_*$, où $a : \Sh(U_{Zar}) \to \Sh(X)$
est comme dans le lemme \ref{lemma-augmentation}.
\end{lemma}

\begin{proof}
Observons que $\mathcal{A}$ est une sous-catégorie de Serre faible d'après le
lemme \ref{lemma-Serre-subcat-cartesian-modules}.
L'équivalence est une conséquence formelle
des résultats obtenus jusqu'ici. On utilise les
lemmes \ref{lemma-descent-sheaves-for-proper-hypercovering},
\ref{lemma-cohomological-descent-for-proper-hypercovering}, ainsi que
Cohomologie sur les sites, lemme \ref{sites-cohomology-lemma-equivalence-bounded}.
\end{proof}

\begin{lemma}
\label{lemma-spectral-sequence-proper-hypercovering}
Soit $U$ un objet simplicial de $\textit{LC}$ et soit
$a : U \to X$ une augmentation. Soit $\mathcal{F}$ un faisceau abélien
sur $X$. Notons $\mathcal{F}_n$ son image inverse sur $U_n$.
Si $U$ est un hyperrecouvrement propre de $X$, alors
il existe une suite spectrale canonique
$$
E_1^{p, q} = H^q(U_p, \mathcal{F}_p)
$$
convergeant vers $H^{p + q}(X, \mathcal{F})$.
\end{lemma}

\begin{proof}
Conséquence immédiate des lemmes \ref{lemma-compute-via-proper-hypercovering}
et \ref{lemma-simplicial-sheaf-cohomology}.
\end{proof}




\section{Schémas simpliciaux}
\label{section-simplicial}

\noindent
Un {\it schéma simplicial} est un objet simplicial dans la catégorie des schémas,
voir Simplicial, définition \ref{simplicial-definition-simplicial-object}.
Rappelons qu'un schéma simplicial ressemble à
$$
\xymatrix{
X_2
\ar@<2ex>[r]
\ar@<0ex>[r]
\ar@<-2ex>[r]
&
X_1
\ar@<1ex>[r]
\ar@<-1ex>[r]
\ar@<1ex>[l]
\ar@<-1ex>[l]
&
X_0
\ar@<0ex>[l]
}
$$
Il y a ici deux morphismes $d^1_0, d^1_1 : X_1 \to X_0$
et un seul morphisme $s^0_0 : X_0 \to X_1$, etc.
Ces morphismes satisfont notamment aux relations
$d^1_0 \circ s^0_0 = \text{id}_{X_0} = d^1_1 \circ s^0_0$, voir
Simplicial, lemme \ref{simplicial-lemma-characterize-simplicial-object}.
Il est utile de considérer $d^n_i : X_n \to X_{n - 1}$
comme la « projection qui oublie la $i$-ième coordonnée », et
$s^n_j : X_n \to X_{n + 1}$ comme le « morphisme diagonal qui répète
la $j$-ième coordonnée ».

\medskip\noindent
Un {\it morphisme de schémas simpliciaux} $h : X \to Y$ est la même chose
qu'un morphisme d'objets simpliciaux dans la catégorie des schémas,
voir Simplicial, définition \ref{simplicial-definition-simplicial-object}.
Ainsi, $h$ est constitué de morphismes de schémas $h_n : X_n \to Y_n$
tels que $h_{n - 1} \circ d^n_j = d^n_j \circ h_n$ et
$h_{n + 1} \circ s^n_j = s^n_j \circ h_n$ chaque fois que ces expressions sont définies.

\medskip\noindent
Une {\it augmentation} d'un schéma simplicial $X$ est un morphisme
de schémas $a_0 : X_0 \to S$ tel que $a_0 \circ d^1_0 = a_0 \circ d^1_1$.
Voir Simplicial, section \ref{simplicial-section-augmentation}.

\medskip\noindent
Soit $X$ un schéma simplicial. La construction de la section
\ref{section-simplicial-top} appliquée à l'espace topologique simplicial
sous-jacent donne un site $X_{Zar}$.
En revanche, pour tout $n$, nous avons le petit site de Zariski
$X_{n, Zar}$ (Topologies, définition
\ref{topologies-definition-big-small-Zariski})
et pour tout morphisme $\varphi : [m] \to [n]$
on a un morphisme de sites
$f_\varphi = X(\varphi)_{small} : X_{n, Zar} \to X_{m, Zar}$,
associé au morphisme de schémas
$X(\varphi) : X_n \to X_m$ (Topologies, lemme
\ref{topologies-lemma-morphism-big-small}).
Cela donne un objet simplicial $\mathcal{C}$ dans la catégorie des sites.
Dans le lemme \ref{lemma-simplicial-site-site}, on a construit le site
associé $\mathcal{C}_{total}$. Le foncteur qui associe son image à une immersion ouverte
définit une équivalence $\mathcal{C}_{total} \to X_{Zar}$ qui
identifie les faisceaux, c'est-à-dire $\Sh(\mathcal{C}_{total}) = \Sh(X_{Zar})$.
La différence entre $\mathcal{C}_{total}$ et $X_{Zar}$
est similaire à la différence entre le petit site Zariski $S_{Zar}$
et l'espace topologique sous-jacent de $S$.
Nous identifierons tacitement ces sites dans la suite.

\medskip\noindent
Soit $X_{Zar}$ le site associé à un schéma simplicial $X$.
Il existe un faisceau d'anneaux $\mathcal{O}$ sur $X_{Zar}$ dont la restriction
à $X_n$ est le faisceau structural $\mathcal{O}_{X_n}$. Cela résulte
du lemme \ref{lemma-describe-sheaves-simplicial-site} ou du lemme
\ref{lemma-describe-sheaves-simplicial-site-site}. Nous dirons que
{\it $\mathcal{O}$ est le faisceau structural du schéma simplicial $X$}.
À ce stade, tout ce qui a été développé pour les sites simpliciaux (annelés)
s'applique, voir les sections \ref{section-simplicial-sites},
\ref{section-augmentation-simplicial-sites},
\ref{section-morphism-simplicial-sites},
\ref{section-simplicial-sites-modules},
\ref{section-cohomology-simplicial-sites},
\ref{section-cohomology-augmentation-simplicial-sites},
\ref{section-cohomology-simplicial-sites-modules},
\ref{section-cohomology-augmentation-ringed-simplicial-sites},
\ref{section-cartesian},
\ref{section-glueing} et
\ref{section-glueing-modules}.

\medskip\noindent
Soit $X$ un schéma simplicial muni de son faisceau structural $\mathcal{O}$.
Comme sur tout topos annelé, il existe une notion
de {\it $\mathcal{O}$-module quasi-cohérent sur $X_{Zar}$}; voir Modules
sur les sites, définition \ref{sites-modules-definition-site-local}.
Cependant, un $\mathcal{O}$-module quasi-cohérent sur $X_{Zar}$ n'est
autre qu'un $\mathcal{O}$-module cartésien $\mathcal{F}$ dont les restrictions
$\mathcal{F}_n$ sont quasi-cohérentes sur $X_n$; voir le
lemme \ref{lemma-quasi-coherent-sheaf}.

\medskip\noindent
Soit $h : X \to Y$ un morphisme de schémas simpliciaux. Le
lemme \ref{lemma-simplicial-space-site-functorial} ou
(la preuve du) lemme \ref{lemma-morphism-simplicial-sites}
fournit un morphisme de sites $h_{Zar} : X_{Zar} \to Y_{Zar}$.
Rappelons que $h_{Zar}^{-1}$ et $h_{Zar, *}$ s'expriment simplement
en termes des composantes; voir le lemme
\ref{lemma-describe-functoriality} ou le lemme
\ref{lemma-morphism-simplicial-sites}.
Notons $\mathcal{O}_X$ et $\mathcal{O}_Y$ les faisceaux structuraux
respectifs de $X$ et de $Y$. On définit
$h_{Zar}^\sharp : h_{Zar, *}\mathcal{O}_X \to \mathcal{O}_Y$
comme le morphisme de faisceaux d'anneaux sur $Y_{Zar}$ donné par
$h_n^\sharp : h_{n, *}\mathcal{O}_{X_n} \to \mathcal{O}_{Y_n}$ sur $Y_n$.
On obtient un morphisme de sites annelés
$$
h_{Zar} : (X_{Zar}, \mathcal{O}_X) \longrightarrow (Y_{Zar}, \mathcal{O}_Y)
$$

\medskip\noindent
Soit $X$ un schéma simplicial muni de son faisceau structural $\mathcal{O}$.
Soit $S$ un schéma et $a_0 : X_0 \to S$ une augmentation de $X$.
Le
lemme \ref{lemma-augmentation} ou le
lemme \ref{lemma-augmentation-site}
fournit un morphisme correspondant de topoi $a : \Sh(X_{Zar}) \to \Sh(S)$.
Observons que $a^{-1}\mathcal{G}$ est le faisceau sur $X_{Zar}$ de composantes
$a_n^{-1}\mathcal{G}$. On peut donc utiliser les morphismes
$a_n^\sharp : a_n^{-1}\mathcal{O}_S \to \mathcal{O}_{X_n}$ pour définir
un morphisme $a^\sharp : a^{-1}\mathcal{O}_S \to \mathcal{O}$ ou, de manière équivalente,
par adjonction, un morphisme $a^\sharp : \mathcal{O}_S \to a_*\mathcal{O}$
(qui, comme d'habitude, porte le même nom). Cela nous place dans la situation
discutée dans la section
\ref{section-cohomology-augmentation-ringed-simplicial-sites}.
On obtient donc un morphisme de topoi annelés
$$
a : (\Sh(X_{Zar}), \mathcal{O}) \longrightarrow (\Sh(S), \mathcal{O}_S)
$$

\medskip\noindent
Enfin, supposons donné un morphisme
$h : X \to Y$ entre des schémas simpliciaux $X$ et $Y$, munis de leurs faisceaux structuraux
$\mathcal{O}_X$, $\mathcal{O}_Y$, d'augmentations
$a_0 : X_0 \to X_{-1}$, $b_0 : Y_0 \to Y_{-1}$ et un morphisme
$h_{-1} : X_{-1} \to Y_{-1}$ tel que le diagramme
$$
\xymatrix{
X_0 \ar[r]_{h_0} \ar[d]_{a_0} & Y_0 \ar[d]^{b_0} \\
X_{-1} \ar[r]^{h_{-1}} & Y_{-1}
}
$$
soit commutatif. Les constructions décrites ci-dessus
fournissent alors le diagramme commutatif suivant de morphismes de topoi annelés :
$$
\xymatrix{
(\Sh(X_{Zar}), \mathcal{O}_X) \ar[r]_{h_{Zar}} \ar[d]_a &
(\Sh(Y_{Zar}), \mathcal{O}_Y) \ar[d]^b \\
(\Sh(X_{-1}), \mathcal{O}_{X_{-1}}) \ar[r]^{h_{-1}} &
(\Sh(Y_{-1}), \mathcal{O}_{Y_{-1}})
}
$$








\section{Descente en termes de schémas simpliciaux}
\label{section-simplicial-descent}

\noindent
Les morphismes cartésiens sont définis comme suit.

\begin{definition}
\label{definition-cartesian-morphism}
Soit $a : Y \to X$ un morphisme de schémas simpliciaux.
On dit que $a$ est {\it cartésien}, ou que {\it $Y$ est cartésien au-dessus de $X$},
si pour tout morphisme $\varphi : [n] \to [m]$ de $\Delta$ le diagramme
correspondant
$$
\xymatrix{
Y_m \ar[r]_a \ar[d]_{Y(\varphi)} & X_m \ar[d]^{X(\varphi)}\\
Y_n \ar[r]^{a} & X_n
}
$$
est un carré cartésien dans la catégorie des schémas.
\end{definition}

\noindent
Les morphismes cartésiens sont liés aux données de descente. Commençons par un lemme général
décrivant, pour un schéma simplicial fixé, la catégorie des schémas simpliciaux cartésiens
au-dessus de celui-ci. Dans ce lemme, $f^* : \Sch/X \to \Sch/Y$ désigne
le foncteur de changement de base associé à un morphisme de schémas $f : Y \to X$.

\begin{lemma}
\label{lemma-characterize-cartesian-schemes}
Soit $X$ un schéma simplicial. La catégorie des schémas simpliciaux cartésiens
sur $X$ est équivalente à la catégorie des couples $(V, \varphi)$
où $V$ est un schéma sur $X_0$ et
$$
\varphi :
V \times_{X_0, d^1_1} X_1
\longrightarrow
X_1 \times_{d^1_0, X_0} V
$$
est un isomorphisme sur $X_1$ tel que
$(s_0^0)^*\varphi = \text{id}_V$ et tel que
$$
(d^2_1)^*\varphi = (d^2_0)^*\varphi \circ (d^2_2)^*\varphi
$$
comme morphismes de schémas sur $X_2$.
\end{lemma}

\begin{proof}
L'égalité affichée a un sens puisque
$d^1_1 \circ d^2_2 = d^1_1 \circ d^2_1$,
$d^1_1 \circ d^2_0 = d^1_0 \circ d^2_2$ et
$d^1_0 \circ d^2_0 = d^1_0 \circ d^2_1$ comme morphismes $X_2 \to X_0$, voir
Simplicial, remarque \ref{simplicial-remark-relations}; on peut donc
représenter ces morphismes comme suit
$$
\xymatrix{
&
X_2 \times_{d^1_1 \circ d^2_0, X_0} V
\ar[r]_-{(d^2_0)^*\varphi} &
X_2 \times_{d^1_0 \circ d^2_0, X_0} V
\ar@{=}[rd] & \\
X_2 \times_{d^1_0 \circ d^2_2, X_0} V
\ar@{=}[ru] & & &
X_2 \times_{d^1_0 \circ d^2_1, X_0} V \\
&
X_2 \times_{d^1_1 \circ d^2_2, X_0} V
\ar[lu]^{(d^2_2)^*\varphi} \ar@{=}[r] &
X_2 \times_{d^1_1 \circ d^2_1, X_0} V
\ar[ru]_{(d^2_1)^*\varphi}
}
$$
et la condition signifie que le diagramme est commutatif. Or,
à tout schéma simplicial cartésien $Y$ au-dessus de $X$, on peut associer
$V = Y_0$ et prendre pour $\varphi$ le composé
$$
V \times_{X_0, d^1_1} X_1 =
Y_0 \times_{X_0, d^1_1} X_1 = Y_1 =
X_1 \times_{X_0, d^1_0} Y_0 =
X_1 \times_{X_0, d^1_0} V
$$
des identifications données par la structure cartésienne. Pour montrer que ce foncteur
est une équivalence, construisons-en un quasi-inverse. La construction du
quasi-inverse est analogue à celle décrite dans
Descente, section \ref{descent-section-descent-modules}, à laquelle nous empruntons
les notations $\tau^n_i : [0] \to [n]$, $0 \mapsto i$ et
$\tau^n_{ij} : [1] \to [n]$, $0 \mapsto i$, $1 \mapsto j$.
Plus précisément, pour un couple $(V, \varphi)$
comme dans le lemme, posons $Y_n = X_n \times_{X(\tau^n_n), X_0} V$.
Pour $\beta : [n] \to [m]$, définissons alors
$V(\beta) : Y_m \to Y_n$ comme le changement de base par $X(\tau^m_{\beta(n)m})$
du morphisme $\varphi$ composé avec la projection
$X_m \times_{X(\beta), X_n} Y_n \to Y_n$. Cette définition a un sens, car
$$
X_m \times_{X(\tau^m_{\beta(n)m}), X_1} X_1 \times_{d^1_1, X_0} V
=
X_m \times_{X(\tau^m_m), X_0} V = Y_m
$$
et
$$
X_m \times_{X(\tau^m_{\beta(n)m}), X_1} X_1 \times_{d^1_0, X_0} V =
X_m \times_{X(\tau^m_{\beta(n)}), X_0} V =
X_m \times_{X(\beta), X_n} Y_n.
$$
Nous laissons au lecteur la vérification que la commutativité
du diagramme ci-dessus
entraîne la compatibilité des morphismes avec la composition, ainsi que la vérification
que les deux foncteurs sont quasi inverses l'un de l'autre.
\end{proof}

\begin{definition}
\label{definition-fibre-products-simplicial-scheme}
Soit $f : X \to S$ un morphisme de schémas. Le {\it schéma simplicial
associé à $f$}, noté $(X/S)_\bullet$, est le foncteur
$\Delta^{opp} \to \Sch$, $[n] \mapsto X \times_S \ldots \times_S X$
décrit dans
Simplicial, exemple \ref{simplicial-example-fibre-products-simplicial-object}.
\end{definition}

\noindent
Ainsi, $(X/S)_n$ est le produit fibré de $(n + 1)$ copies de $X$ sur $S$.
Le morphisme $d^1_0 : X \times_S X \to X$ est la projection
$(x_0, x_1) \mapsto x_1$, et le morphisme $d^1_1$ est l'autre
projection. Le morphisme $s^0_0$ est le morphisme diagonal
$X \to X \times_S X$.

\begin{lemma}
\label{lemma-cartesian-over}
Soit $f : X \to S$ un morphisme de schémas.
Soit $\pi : Y \to (X/S)_\bullet$ un morphisme cartésien
de schémas simpliciaux.
Posons $V = Y_0$ considéré comme un schéma sur $X$.
Les morphismes $d^1_0, d^1_1 : Y_1 \to Y_0$ et le morphisme
$\pi_1 : Y_1 \to X \times_S X$ induisent des isomorphismes
$$
\xymatrix{
V \times_S X & &
Y_1 \ar[ll]_-{(d^1_1, \text{pr}_1 \circ \pi_1)}
\ar[rr]^-{(\text{pr}_0 \circ \pi_1, d^1_0)} & &
X \times_S V.
}
$$
Notons $\varphi : V \times_S X \to X \times_S V$
l'isomorphisme ainsi obtenu.
Alors le couple $(V, \varphi)$ est une donnée de descente relative
à $X \to S$.
\end{lemma}

\begin{proof}
C'est un cas particulier d'une des assertions du
lemme \ref{lemma-characterize-cartesian-schemes},
car l'équation affichée de ce lemme est
équivalente à la condition de cocycle de
Descente, définition \ref{descent-definition-descent-datum}.
\end{proof}

\begin{lemma}
\label{lemma-cartesian-equivalent-descent-datum}
Soit $f : X \to S$ un morphisme de schémas. La construction
$$
\begin{matrix}
\text{catégorie des schémas cartésiens} \\
\text{au-dessus de } (X/S)_\bullet
\end{matrix}
\longrightarrow
\begin{matrix}
\text{catégorie des données de descente} \\
\text{relatives à } X/S
\end{matrix}
$$
du lemme \ref{lemma-cartesian-over}
est une équivalence de catégories.
\end{lemma}

\begin{proof}
Le foncteur de gauche à droite est donné dans le lemme
\ref{lemma-cartesian-over}.
Il s’agit donc d’un cas particulier du lemme
\ref{lemma-characterize-cartesian-schemes}.
\end{proof}

\noindent
Nous pouvons reformuler le changement de base du
lemme \ref{descent-lemma-pullback} de Descente comme suit.
Supposons donnés un morphisme de schémas simpliciaux $f : X' \to X$ et un morphisme cartésien
de schémas simpliciaux $Y \to X$. Alors
le produit fibré (considéré comme une « image inverse »)
$$
f^*Y = Y \times_X X'
$$
est un schéma simplicial cartésien sur $X'$.
Supposons qu'on donne un diagramme commutatif de morphismes de schémas
$$
\xymatrix{
X' \ar[r]_f \ar[d] & X \ar[d] \\
S' \ar[r] & S.
}
$$
On obtient ainsi un morphisme de schémas simpliciaux
$$
f_\bullet : (X'/S')_\bullet \longrightarrow (X/S)_\bullet.
$$
Nous affirmons que le « changement de base » $f_\bullet^*$ associé au morphisme
$f_\bullet : (X'/S')_\bullet \to (X/S)_\bullet$ correspond, via
le lemme \ref{lemma-cartesian-equivalent-descent-datum},
au changement de base défini, en termes de données de descente, dans
Descente, lemme
\ref{descent-lemma-pullback}.







\section{Modules quasi-cohérents sur les schémas simpliciaux}
\label{section-modules-simplicial}

\begin{lemma}
\label{lemma-pullback-cartesian-module}
Soit $f : V \to U$ un morphisme de schémas simpliciaux. Pour un
module quasi-cohérent $\mathcal{F}$ sur $U_{Zar}$, l'image inverse
$f^*\mathcal{F}$ est un module quasi-cohérent sur $V_{Zar}$.
\end{lemma}

\begin{proof}
Rappelons que $\mathcal{F}$ est cartésien avec $\mathcal{F}_n$
quasi-cohérent, voir le lemme \ref{lemma-quasi-coherent-sheaf}.
D'après le lemme \ref{lemma-describe-functoriality}, nous voyons que
$(f^*\mathcal{F})_n = f_n^*\mathcal{F}_n$ (quelques détails omis).
$(f^*\mathcal{F})_n$ est donc quasi-cohérent.
Le même fait et la propriété cartésienne pour $\mathcal{F}$
impliquent la propriété cartésienne pour $f^*\mathcal{F}$.
Ainsi $\mathcal{F}$ est quasi-cohérent par
le lemme \ref{lemma-quasi-coherent-sheaf}, une fois encore.
\end{proof}

\begin{lemma}
\label{lemma-pushforward-cartesian-module}
Soit $f : V \to U$ un morphisme cartésien de schémas simpliciaux.
Supposons que les morphismes $d^n_j : U_n \to U_{n - 1}$ soient
plats et que les morphismes $V_n \to U_n$ soient quasi-compacts et quasi-séparés.
Pour un module quasi-cohérent $\mathcal{G}$ sur $V_{Zar}$
l'image directe $f_*\mathcal{G}$ est un module quasi-cohérent sur $U_{Zar}$.
\end{lemma}

\begin{proof}
Si $\mathcal{F} = f_* \mathcal{G}$, alors
$\mathcal{F}_n = f_{n , *}\mathcal{G}_n$ par
le lemme \ref{lemma-describe-functoriality}.
Les morphismes $\mathcal{F}(\varphi)$ sont définis à l'aide des morphismes de changement de base; voir
Cohomologie, section \ref{cohomology-section-base-change-map}.
Les faisceaux $\mathcal{F}_n$ sont quasi-cohérents par
Schémas, lemme \ref{schemes-lemma-push-forward-quasi-coherent}
et le fait que $\mathcal{G}_n$ est quasi-cohérent
d'après le lemme \ref{lemma-quasi-coherent-sheaf}.
Les morphismes de changement de base le long des dégénérescences
$d^n_j$ sont des isomorphismes d'après Cohomologie des schémas, lemme
\ref{coherent-lemma-flat-base-change-cohomology}
et le fait que $\mathcal{G}$ est cartésien
d'après le lemme \ref{lemma-quasi-coherent-sheaf}.
Ainsi, $\mathcal{F}$ est cartésien d'après le
lemme \ref{lemma-check-cartesian-module}.
Par conséquent, $\mathcal{F}$ est quasi-cohérent d'après le
lemme \ref{lemma-quasi-coherent-sheaf}.
\end{proof}

\begin{lemma}
\label{lemma-adjoint-functors-cartesian-modules}
Soit $f : V \to U$ un morphisme cartésien de schémas
simpliciaux. Supposons que les morphismes $d^n_j : U_n \to U_{n - 1}$ soient
plats et que les morphismes $V_n \to U_n$ soient quasi-compacts et quasi-séparés.
Alors $f^*$ et $f_*$ forment une paire de foncteurs adjoints
entre les catégories de modules quasi-cohérents sur $U_{Zar}$ et $V_{Zar}$.
\end{lemma}

\begin{proof}
Nous avons vu dans les lemmes \ref{lemma-pullback-cartesian-module} et
\ref{lemma-pushforward-cartesian-module}
que l'assertion a un sens. La propriété d'adjonction résulte
du fait que chaque $f_n^*$ est adjoint à $f_{n, *}$.
\end{proof}

\begin{lemma}
\label{lemma-cartesian-modules-with-section}
Soit $f : X \to S$ un morphisme de schémas qui admet une section
\footnote{En fait, il suffirait de supposer que $f$
admet fpqc-localement sur $S$ une section, puisque les modules quasi-cohérents
satisfont la descente d'après Descente,
section \ref{descent-section-fpqc-descent-quasi-coherent}.}.
Soit $(X/S)_\bullet$ le schéma simplicial
associé à $X \to S$, voir
définition \ref{definition-fibre-products-simplicial-scheme}.
Alors l'image inverse définit une équivalence entre la catégorie des
$\mathcal{O}_S$-modules quasi-cohérents et la catégorie des modules quasi-cohérents
sur $((X/S)_\bullet)_{Zar}$.
\end{lemma}

\begin{proof}
Soit $\sigma : S \to X$ une section de $f$. Soit $(\mathcal{F}, \alpha)$
une paire comme dans le lemme \ref{lemma-characterize-cartesian-modules}.
Posons $\mathcal{G} = \sigma^*\mathcal{F}$. Considérons le diagramme
$$
\xymatrix{
X \ar[r]_-{(\sigma \circ f, 1)} \ar[d]_f &
X \times_S X \ar[d]^{\text{pr}_0} \ar[r]_-{\text{pr}_1} & X \\
S \ar[r]^\sigma & X
}
$$
Notons que $\text{pr}_0 = d^1_1$ et $\text{pr}_1 = d^1_0$. Nous voyons donc que
$(\sigma \circ f, 1)^*\alpha$ définit un isomorphisme
$$
f^*\mathcal{G} = (\sigma \circ f, 1)^*\text{pr}_0^*\mathcal{F}
\longrightarrow
(\sigma \circ f, 1)^*\text{pr}_1^*\mathcal{F} = \mathcal{F}
$$
Nous omettons de vérifier que cet isomorphisme est compatible
avec $\alpha$ et l'isomorphisme canonique
$\text{pr}_0^*f^*\mathcal{G} \to \text{pr}_1^*f^*\mathcal{G}$.
\end{proof}






\section{Groupoïdes et schémas simpliciaux}
\label{section-groupoids-simplicial}

\noindent
Étant donné un groupoïde en schémas, nous pouvons construire un schéma simplicial.
Il s'avère que la catégorie des faisceaux quasi-cohérents sur un groupoïde
est équivalente à la catégorie des faisceaux cartésiens quasi-cohérents
sur le schéma simplicial associé.

\begin{lemma}
\label{lemma-groupoid-simplicial}
Soit $(U, R, s, t, c, e, i)$ un groupoïde en schémas sur $S$.
Il existe un schéma simplicial $X$ sur $S$
avec les propriétés suivantes
\begin{enumerate}
\item $X_0 = U$, $X_1 = R$, $X_2 = R \times_{s, U, t} R$,
\item $s_0^0 = e : X_0 \to X_1$,
\item $d^1_0 = s : X_1 \to X_0$, $d^1_1 = t : X_1 \to X_0$,
\item $s_0^1 = (e \circ t, 1) : X_1 \to X_2$,
$s_1^1 = (1, e \circ t) : X_1 \to X_2$,
\item $d^2_0 = \text{pr}_1 : X_2 \to X_1$,
$d^2_1 = c : X_2 \to X_1$,
$d^2_2 = \text{pr}_0$ et
\item $X = \text{cosk}_2 \text{sk}_2 X$.
\end{enumerate}
Pour tout $n$, on a $X_n = R \times_{s, U, t} \ldots \times_{s, U, t} R$
à $n$ facteurs. Le morphisme $d^n_j : X_n \to X_{n - 1}$ est donné sur les
foncteurs de points par
$$
(r_1, \ldots, r_n) \longmapsto (r_1, \ldots, c(r_j, r_{j + 1}), \ldots, r_n)
$$
pour $1 \leq j \leq n - 1$ alors que
$d^n_0(r_1, \ldots, r_n) = (r_2, \ldots, r_n)$ et
$d^n_n(r_1, \ldots, r_n) = (r_1, \ldots, r_{n - 1})$.
\end{lemma}

\begin{proof}
Il suffit de vérifier que les règles prescrites en (1), (2), (3), (4), (5)
définissent un schéma simplicial $2$-tronqué $U'$ sur $S$, puisque (6)
nous permet de définir $X = \text{cosk}_2 U'$, voir
Simplicial, lemme \ref{simplicial-lemma-existence-cosk}.
En utilisant l'approche du foncteur de points, il suffit de vérifier que
si $(\text{Ob}, \text{Arrows}, s, t, c, e, i)$ est un groupoïde, alors
$$
\xymatrix{
\text{Arrows} \times_{s, \text{Ob}, t} \text{Arrows}
\ar@<8ex>[d]^{\text{pr}_0}
\ar@<0ex>[d]_c
\ar@<-8ex>[d]_{\text{pr}_1}
\\
\text{Arrows}
\ar@<4ex>[d]^t
\ar@<-4ex>[d]_s
\ar@<4ex>[u]^{1, e}
\ar@<-4ex>[u]_{e, 1}
\\
\text{Ob}
\ar@<0ex>[u]_e
}
$$
est un ensemble simplicial $2$-tronqué. Nous omettons les détails.

\medskip\noindent
Enfin, la description de $X_n$ pour $n > 2$ suit par induction de
la description de $X_0$, $X_1$, $X_2$ et
Simplicial, remarque \ref{simplicial-remark-inductive-coskeleton}, et du
lemme \ref{simplicial-lemma-work-out}. On peut aussi montrer que
l'application de $\text{cosk}_2$ à l'ensemble simplicial $2$-tronqué ci-dessus
donne un ensemble simplicial dont le terme de degré $n$ est égal à
$\text{Arrows} \times_{s, \text{Ob}, t} \ldots \times_{s, \text{Ob}, t}
\text{Arrows}$ à $n$ facteurs et les morphismes de dégénérescence décrits dans le lemme.
Certains détails ont été omis.
\end{proof}

\begin{lemma}
\label{lemma-quasi-coherent-groupoid-simplicial}
Soient $S$ un schéma et $(U, R, s, t, c)$ un groupoïde en schémas
sur $S$. Soit $X$ le schéma simplicial sur $S$ construit
dans le lemme \ref{lemma-groupoid-simplicial}.
Alors la catégorie des modules quasi-cohérents sur $(U, R, s, t, c)$
est équivalente à la catégorie des modules quasi-cohérents sur $X_{Zar}$.
\end{lemma}

\begin{proof}
Cela ressort clairement des lemmes
\ref{lemma-quasi-coherent-sheaf} et
\ref{lemma-characterize-cartesian-modules}
et Groupoïdes, définition \ref{groupoids-definition-groupoid-module}.
\end{proof}

\noindent
Dans le lemme suivant, nous utiliserons la notion de morphisme cartésien
$V \to U$ de schémas simpliciaux, telle qu'elle est définie dans la définition
\ref{definition-cartesian-morphism}.

\begin{lemma}
\label{lemma-quasi-coherent-groupoid-R-cartesian}
Soit $(U, R, s, t, c)$ un groupoïde en schémas sur un schéma $S$.
Soit $X$ le schéma simplicial sur $S$ construit
dans le lemme \ref{lemma-groupoid-simplicial}.
Soit $(R/U)_\bullet$ le schéma simplicial
associé à $s : R \to U$, voir
définition \ref{definition-fibre-products-simplicial-scheme}.
Il existe un morphisme cartésien $t_\bullet : (R/U)_\bullet \to X$
de schémas simpliciaux, dont les morphismes en petits degrés sont donnés par
$$
\xymatrix{
R \times_{s, U, s} R \times_{s, U, s} R
\ar@<3ex>[r]_-{\text{pr}_{12}}
\ar@<0ex>[r]_-{\text{pr}_{02}}
\ar@<-3ex>[r]_-{\text{pr}_{01}}
\ar[dd]_{(r_0, r_1, r_2) \mapsto (r_0 \circ r_1^{-1}, r_1 \circ r_2^{-1})} &
R \times_{s, U, s} R
\ar@<1ex>[r]_-{\text{pr}_1} \ar@<-2ex>[r]_-{\text{pr}_0}
\ar[dd]_{(r_0, r_1) \mapsto r_0 \circ r_1^{-1}} &
R \ar[dd]^t
\\
\\
R \times_{s, U, t} R
\ar@<3ex>[r]_{\text{pr}_1}
\ar@<0ex>[r]_c
\ar@<-3ex>[r]_{\text{pr}_0} &
R \ar@<1ex>[r]_s \ar@<-2ex>[r]_t &
U
}
$$
\end{lemma}

\begin{proof}
Pour tout $n$, nous définissons $(R/U)_\bullet \to X_n$ par la règle
$$
(r_0, \ldots, r_n)
\longrightarrow
(r_0 \circ r_1^{-1}, \ldots, r_{n - 1} \circ r_n^{-1})
$$
La compatibilité avec les morphismes de dégénérescence résulte de la description de ceux-ci
dans le lemme \ref{lemma-groupoid-simplicial}.
On omet de vérifier que les morphismes respectent les $s^n_j$.
Le lemme \ref{groupoids-lemma-diagram-pull} de Groupoïdes
(avec les rôles de $s$ et $t$ inversés)
montre que les deux carrés de droite sont cartésiens. Exactement de la même manière,
on montre que tous les autres carrés sont également cartésiens. Le
morphisme est donc cartésien.
\end{proof}




\section{Les données de descente définissent des relations d'équivalence}
\label{section-equivalence-relation}

\noindent
Dans la section \ref{section-simplicial-descent}, nous avons vu comment les données de descente relatives à
$X \to S$ peuvent se formuler en termes de schémas simpliciaux cartésiens
au-dessus de $(X/S)_\bullet$. Nous les relions ici aux relations d'équivalence
comme suit.

\begin{lemma}
\label{lemma-equivalence-relation}
Soit $f : X \to S$ un morphisme de schémas.
Soit $\pi : Y \to (X/S)_\bullet$ un morphisme cartésien de schémas
simpliciaux; voir les définitions \ref{definition-cartesian-morphism} et
\ref{definition-fibre-products-simplicial-scheme}.
Alors le morphisme
$$
j = (d^1_1, d^1_0) : Y_1 \to Y_0 \times_S Y_0
$$
définit une relation d'équivalence sur $Y_0$ au-dessus de $S$;
voir Groupoïdes, définition \ref{groupoids-definition-equivalence-relation}.
\end{lemma}

\begin{proof}
Notons que $j$ est un monomorphisme. En effet, le composé
$Y_1 \to Y_0 \times_S Y_0 \to Y_0 \times_S X$
est un isomorphisme car $\pi$ est cartésien.

\medskip\noindent
Considérons le morphisme
$$
(d^2_2, d^2_0) : Y_2 \to Y_1 \times_{d^1_0, Y_0, d^1_1} Y_1.
$$
Ce morphisme est bien défini car $d_0 \circ d_2 = d_1 \circ d_0$;
voir Simplicial, remarque \ref{simplicial-remark-relations}.
Il s'agit en outre d'un morphisme au-dessus de $(X/S)_2$. C'est un isomorphisme
car $Y \to (X/S)_\bullet$ est cartésien. Notons par exemple que le membre
de droite est isomorphe à
$Y_0 \times_{\pi_0, X, \text{pr}_1} (X \times_S X \times_S X) =
X \times_S Y_0 \times_S X$
car $\pi$ est cartésien. Détails omis.

\medskip\noindent
Comme dans Groupoïdes, définition \ref{groupoids-definition-equivalence-relation},
on note $t = \text{pr}_0 \circ j = d^1_1$ et
$s = \text{pr}_1 \circ j = d^1_0$.
L'isomorphisme ci-dessus, combiné au morphisme
$d^2_1 : Y_2 \to Y_1$, nous donne un morphisme de composition
$$
c : Y_1 \times_{s, Y_0, t} Y_1 \longrightarrow Y_1
$$
sur $Y_0 \times_S Y_0$. Cela implique immédiatement
que pour tout schéma $T/S$ la relation
$Y_1(T) \subset Y_0(T) \times Y_0(T)$ est transitive.

\medskip\noindent
La réflexivité découle du fait que la restriction
du morphisme $j$ à la diagonale
$\Delta : X \to X \times_S X$ est un isomorphisme
(utilisez à nouveau la propriété cartésienne de $\pi$).

\medskip\noindent
Pour voir la symétrie, nous considérons le morphisme
$$
(d^2_2, d^2_1) : Y_2 \to Y_1 \times_{d^1_1, Y_0, d^1_1} Y_1.
$$
Ce morphisme est bien défini car $d_1 \circ d_2 = d_1 \circ d_1$;
voir Simplicial, remarque \ref{simplicial-remark-relations}.
C'est un isomorphisme
car $Y \to (X/S)_\bullet$ est cartésien.
Notons par exemple que le membre
de droite est isomorphe à
$Y_0 \times_{\pi_0, X, \text{pr}_0} (X \times_S X \times_S X) =
Y_0 \times_S X \times_S X$
car $\pi$ est cartésien. Détails omis.

\medskip\noindent
Soit $T/S$ un schéma. Convenons que $a \sim b$, pour $a, b \in Y_0(T)$,
signifie $(a, b) \in Y_1(T)$.
L'isomorphisme $(d^2_2, d^2_1)$ au-dessus de
implique que si $a \sim b$ et $a \sim c$, alors $b \sim c$.
Combiné avec la réflexivité, cela montre que $\sim$ est
une relation d'équivalence.
\end{proof}







\section{Un exemple}
\label{section-example}

\noindent
Dans cette section, on montre que les réunions disjointes de spectres
d'anneaux artiniens peuvent être descendues le long d'un morphisme de schémas
quasi compact, surjectif et plat.

\begin{lemma}
\label{lemma-equivalence-classes-points}
Soit $X \to S$ un morphisme de schémas. Supposons que $Y \to (X/S)_\bullet$
soit un morphisme cartésien de schémas simpliciaux. Pour un point $y \in Y_0$, posons
$$
T_y = \{y' \in Y_0 \mid \exists\ y_1 \in Y_1:
d^1_1(y_1) = y, d^1_0(y_1) = y'\}
$$
comme sous-ensemble de $Y_0$. Alors $y \in T_y$ et
$T_y \cap T_{y'} \not = \emptyset \Rightarrow T_y = T_{y'}$.
\end{lemma}

\begin{proof}
Il suffit de combiner le lemme \ref{lemma-equivalence-relation} et le
lemme de Groupoïdes
\ref{groupoids-lemma-pre-equivalence-equivalence-relation-points}.
\end{proof}

\begin{lemma}
\label{lemma-quasi-compact}
Soit $X \to S$ un morphisme de schémas.
Supposons que $Y \to (X/S)_\bullet$ soit un morphisme cartésien de schémas simpliciaux.
Soit $y \in Y_0$ un point. Si $X \to S$ est quasi-compact, alors
$$
T_y = \{y' \in Y_0 \mid \exists\ y_1 \in Y_1:
d^1_1(y_1) = y, d^1_0(y_1) = y'\}
$$
est un sous-ensemble quasi-compact de $Y_0$.
\end{lemma}

\begin{proof}
Soit $F_y$ la fibre schématique de $d^1_1 : Y_1 \to Y_0$
à $y$. On voit alors que $T_y$ est l'image du morphisme
$$
\xymatrix{
F_y \ar[r] \ar[d] &
Y_1 \ar[r]^{d^1_0} \ar[d]^{d^1_1} &
Y_0 \\
y \ar[r] &
Y_0 &
}
$$
Notons que $F_y$ est quasi-compact. Cela prouve le lemme.
\end{proof}

\begin{lemma}
\label{lemma-descent-disjoint-union-Artinian-along-fields}
Soit $X \to S$ un morphisme plat surjectif quasi-compact.
Soit $(V, \varphi)$ une donnée de descente relative
à $X \to S$. Si $V$ est une réunion disjointe de spectres
d'anneaux artiniens, alors $(V, \varphi)$ est effective.
\end{lemma}

\begin{proof}
Soit $Y \to (X/S)_\bullet$ le morphisme cartésien de schémas simpliciaux
correspondant à $(V, \varphi)$ par
le lemme \ref{lemma-cartesian-equivalent-descent-datum}.
Notons que $Y_0 = V$.
Écrivons $V = \coprod_{i \in I} \Spec(A_i)$, où chaque $A_i$ est local
artinien. Notons en outre $v_i \in V$ l'unique point fermé de
$\Spec(A_i)$ pour tout $i \in I$. Posons $i \sim j$ si et seulement si
$v_i \in T_{v_j}$, avec les notations du
lemme \ref{lemma-equivalence-classes-points} ci-dessus.
D'après les lemmes \ref{lemma-equivalence-classes-points} et \ref{lemma-quasi-compact},
il s'agit d'une relation d'équivalence à classes d'équivalence finies.
Posons $\overline{I} = I/\sim$. On peut alors écrire
$V = \coprod_{\overline{i} \in \overline{I}} V_{\overline{i}}$
avec
$V_{\overline{i}} = \coprod_{i \in \overline{i}} \Spec(A_i)$.
Par construction, nous voyons que
$\varphi : V \times_S X \to X \times_S V$ envoie
les sous-espaces ouverts et fermés $V_{\overline{i}} \times_S X$
dans les sous-espaces ouverts et fermés $X \times_S V_{\overline{i}}$.
Autrement dit, nous obtenons des données de descente
$(V_{\overline{i}}, \varphi_{\overline{i}})$, et
$(V, \varphi)$ en est le coproduit dans la catégorie des
données de descente.
Puisque chacun des $V_{\overline{i}}$ est une union finie de spectres
d'anneaux locaux artiniens, le morphisme $V_{\overline{i}} \to X$
est affine, voir Morphismes, lemme \ref{morphisms-lemma-Artinian-affine}.
Puisque $\{X \to S\}$ est un recouvrement fpqc, toutes les
données de descente $(V_{\overline{i}}, \varphi_{\overline{i}})$ sont effectives
d'après Descente, lemme \ref{descent-lemma-affine}.
\end{proof}

\noindent
Certes, le lemme ci-dessus a un champ d'application très limité !









\section{Espaces algébriques simpliciaux}
\label{section-simplicial-algebraic-spaces}

\noindent
Soit $S$ un schéma. Un {\it espace algébrique simplicial}
est un objet simplicial dans la catégorie des espaces algébriques sur $S$,
voir Simplicial, définition \ref{simplicial-definition-simplicial-object}.
Rappelons qu'un espace algébrique simplicial ressemble à
$$
\xymatrix{
X_2
\ar@<2ex>[r]
\ar@<0ex>[r]
\ar@<-2ex>[r]
&
X_1
\ar@<1ex>[r]
\ar@<-1ex>[r]
\ar@<1ex>[l]
\ar@<-1ex>[l]
&
X_0
\ar@<0ex>[l]
}
$$
Ici, il y a deux morphismes $d^1_0, d^1_1 : X_1 \to X_0$
et un seul morphisme $s^0_0 : X_0 \to X_1$, etc.
Ces morphismes satisfont notamment aux relations
$d^1_0 \circ s^0_0 = \text{id}_{X_0} = d^1_1 \circ s^0_0$, voir
Simplicial, lemme \ref{simplicial-lemma-characterize-simplicial-object}.
Il est utile de considérer $d^n_i : X_n \to X_{n - 1}$
comme la « projection qui oublie la $i$-ième coordonnée », et
$s^n_j : X_n \to X_{n + 1}$ comme le « morphisme diagonal qui répète
la $j$-ième coordonnée ».

\medskip\noindent
Un {\it morphisme d'espaces algébriques simpliciaux} $h : X \to Y$ est la même
chose qu'un morphisme d'objets simpliciaux dans
la catégorie des espaces algébriques sur $S$,
voir Simplicial, définition \ref{simplicial-definition-simplicial-object}.
Ainsi, $h$ est constitué de morphismes d'espaces algébriques $h_n : X_n \to Y_n$
tels que $h_{n - 1} \circ d^n_j = d^n_j \circ h_n$ et
$h_{n + 1} \circ s^n_j = s^n_j \circ h_n$ chaque fois que ces expressions sont définies.

\medskip\noindent
Une {\it augmentation} $a : X \to X_{-1}$
d'un espace algébrique simplicial $X$ est donnée par un morphisme
d'espaces algébriques $a_0 : X_0 \to X_{-1}$
tel que $a_0 \circ d^1_0 = a_0 \circ d^1_1$.
Voir Simplicial, section \ref{simplicial-section-augmentation}.
Dans cette situation, on note toujours $a_n : X_n \to X_{-1}$ le morphisme induit
pour $n \geq 0$.

\medskip\noindent
Soit $X$ un espace algébrique simplicial. Pour chaque $n$, nous avons le site
$X_{n, spaces, \etale}$ (Propriétés des espaces, définition
\ref{spaces-properties-definition-spaces-etale-site})
et pour chaque morphisme $\varphi : [m] \to [n]$ nous avons un morphisme de sites
$$
f_\varphi = X(\varphi)_{spaces, \etale} :
X_{n, spaces, \etale} \to X_{m, spaces, \etale},
$$
associé au morphisme des espaces algébriques
$X(\varphi) : X_n \to X_m$ (Propriétés des espaces, lemme
\ref{spaces-properties-lemma-functoriality-etale-site}).
Cela donne un objet simplicial dans la catégorie des sites.
Dans le lemme \ref{lemma-simplicial-site-site}, nous avons construit un site
associé que nous notons $X_{spaces, \etale}$.
Un objet du site $X_{spaces, \etale}$ est
un espace algébrique $U$ \'etale sur $X_n$ pour un certain $n$;
et un morphisme $(\varphi, f) : U/X_n \to V/X_m$ est donné
par un morphisme $\varphi : [m] \to [n]$ dans $\Delta$ et un morphisme
$f : U \to V$ d'espaces algébriques tel que le diagramme
$$
\xymatrix{
U \ar[r]_f \ar[d] & V \ar[d] \\
X_n \ar[r]^{f_\varphi} & X_m
}
$$
est commutatif. Considérons les sous-catégories pleines
$$
X_{affine, \etale} \subset X_\etale \subset X_{spaces, \etale}
$$
dont les objets sont $U/X_n$ avec $U$ affine, respectivement un schéma.
Munies de leurs topologies naturelles
(voir
Propriétés des espaces, lemme \ref{spaces-properties-lemma-alternative},
définition \ref{spaces-properties-definition-etale-site} et
lemme \ref{spaces-properties-lemma-compare-etale-sites}),
ces catégories ont des foncteurs d'inclusion qui définissent des équivalences de topoi
$$
\Sh(X_{affine, \etale}) = \Sh(X_\etale) = \Sh(X_{spaces, \etale})
$$
Dans la suite, nous identifierons tacitement ces topos.
On dira que $X_\etale$ est le {\it petit site \'etale de $X$}
et son topos est le {\it petit topos \'etale de $X$}.

\medskip\noindent
Soit $X_\etale$ le petit site \'etale d'un espace algébrique simplicial $X$.
Il existe un faisceau d'anneaux $\mathcal{O}$ sur $X_\etale$ dont la restriction
à $X_n$ est le faisceau structural $\mathcal{O}_{X_n}$. Cela découle du lemme
\ref{lemma-describe-sheaves-simplicial-site-site}. Nous dirons
{\it $\mathcal{O}$ est le faisceau structural de l'espace algébrique simplicial
$X$}.
À ce stade, tout ce qui a été développé pour les sites simpliciaux (annelés)
s'applique, voir les sections \ref{section-simplicial-sites},
\ref{section-augmentation-simplicial-sites},
\ref{section-morphism-simplicial-sites},
\ref{section-simplicial-sites-modules},
\ref{section-cohomology-simplicial-sites},
\ref{section-cohomology-augmentation-simplicial-sites},
\ref{section-cohomology-simplicial-sites-modules},
\ref{section-cohomology-augmentation-ringed-simplicial-sites},
\ref{section-cartesian},
\ref{section-glueing} et
\ref{section-glueing-modules}.

\medskip\noindent
Soit $X$ un espace algébrique simplicial muni de son faisceau structural $\mathcal{O}$.
Comme sur tout topos annelé, il existe une notion
de {\it $\mathcal{O}$-module quasi-cohérent sur $X_\etale$}; voir Modules
sur les sites, définition \ref{sites-modules-definition-site-local}.
Cependant, un $\mathcal{O}$-module quasi-cohérent sur $X_\etale$ n'est
autre qu'un $\mathcal{O}$-module cartésien $\mathcal{F}$ dont les restrictions
$\mathcal{F}_n$ sont quasi-cohérentes sur $X_n$; voir le lemme
\ref{lemma-quasi-coherent-sheaf}.

\medskip\noindent
Soit $h : X \to Y$ un morphisme d'espaces algébriques simpliciaux sur $S$.
D'après le lemme \ref{lemma-morphism-simplicial-sites}, appliqué aux morphismes
des sites $(h_n)_{spaces, \etale} : X_{spaces, \etale} \to Y_{spaces, \etale}$
(Propriétés des espaces, lemme
\ref{spaces-properties-lemma-functoriality-etale-site})
on obtient un morphisme de petits topoi étales
$h_\etale : \Sh(X_\etale) \to \Sh(Y_\etale)$.
Rappelons que $h_\etale^{-1}$ et $h_{\etale, *}$ se décrivent simplement
en termes de leurs composantes; voir le
lemme \ref{lemma-morphism-simplicial-sites}.
Notons $\mathcal{O}_X$, resp.\ $\mathcal{O}_Y$, le faisceau structural
de $X$, resp.\ de $Y$. On définit
$h_\etale^\sharp : h_{\etale, *}\mathcal{O}_X \to \mathcal{O}_Y$
comme le morphisme de faisceaux d'anneaux sur $Y_\etale$ donné par
$h_n^\sharp : h_{n, *}\mathcal{O}_{X_n} \to \mathcal{O}_{Y_n}$ sur $Y_n$.
On obtient un morphisme de topoi annelés
$$
h_\etale :
(\Sh(X_\etale), \mathcal{O}_X)
\longrightarrow
(\Sh(Y_\etale), \mathcal{O}_Y)
$$

\medskip\noindent
Soit $X$ un espace algébrique simplicial muni de son faisceau structural $\mathcal{O}$.
Soit $X_{-1}$ un espace algébrique sur $S$ et soit $a_0 : X_0 \to X_{-1}$
une augmentation de $X$. Par
le lemme \ref{lemma-augmentation-site},
appliqué au morphisme des sites
$(a_0)_{spaces, \etale} :
X_{0, spaces, \etale} \to X_{-1, spaces, \etale}$
on obtient un morphisme correspondant de topoi
$a : \Sh(X_\etale) \to \Sh(X_{-1, \etale})$.
Observons que $a^{-1}\mathcal{G}$ est le faisceau sur
$X_\etale$ de composantes $a_n^{-1}\mathcal{G}$. On peut donc utiliser les morphismes
$a_n^\sharp : a_n^{-1}\mathcal{O}_{X_{-1}} \to \mathcal{O}_{X_n}$ pour définir
un morphisme $a^\sharp : a^{-1}\mathcal{O}_{X_{-1}} \to \mathcal{O}$ ou, de manière équivalente,
par adjonction, un morphisme $a^\sharp : \mathcal{O}_{X_{-1}} \to a_*\mathcal{O}$
(qui, comme d'habitude, porte le même nom). Cela nous place dans la situation
discutée dans la section
\ref{section-cohomology-augmentation-ringed-simplicial-sites}.
On obtient donc un morphisme de topoi annelés
$$
a :
(\Sh(X_\etale), \mathcal{O})
\longrightarrow
(\Sh(X_{-1}), \mathcal{O}_{X_{-1}})
$$

\medskip\noindent
Faisons une dernière observation. Supposons donné un morphisme
$h : X \to Y$ entre des espaces algébriques simpliciaux $X$ et $Y$, munis de leurs faisceaux structuraux
$\mathcal{O}_X$, $\mathcal{O}_Y$, d'augmentations
$a_0 : X_0 \to X_{-1}$, $b_0 : Y_0 \to Y_{-1}$ et d'un morphisme
$h_{-1} : X_{-1} \to Y_{-1}$ tel que le diagramme
$$
\xymatrix{
X_0 \ar[r]_{h_0} \ar[d]_{a_0} & Y_0 \ar[d]^{b_0} \\
X_{-1} \ar[r]^{h_{-1}} & Y_{-1}
}
$$
soit commutatif. Les constructions décrites ci-dessus
fournissent alors le diagramme commutatif suivant de morphismes de topoi annelés :
$$
\xymatrix{
(\Sh(X_\etale), \mathcal{O}_X) \ar[r]_{h_\etale} \ar[d]_a &
(\Sh(Y_\etale), \mathcal{O}_Y) \ar[d]^b \\
(\Sh(X_{-1}), \mathcal{O}_{X_{-1}}) \ar[r]^{h_{-1}} &
(\Sh(Y_{-1}), \mathcal{O}_{Y_{-1}})
}
$$




\section{Hyperrecouvrements fppf d'espaces algébriques}
\label{section-fppf-hypercovering}

\noindent
Cette section est l'analogue de la section \ref{section-proper-hypercovering}
pour le cas des espaces algébriques et des hyperrecouvrements fppf.
Le lecteur qui le souhaite peut remplacer partout « espace algébrique »
par « schéma » et obtenir des résultats également valables.
Cela présente l'avantage de remplacer les références à
Compléments sur la cohomologie des espaces, section
\ref{spaces-more-cohomology-section-fppf-etale}
par des références à
Cohomologie étale, section \ref{etale-cohomology-section-fppf-etale}.

\medskip\noindent
Nous fixons un schéma de base $S$.
Soit $X$ un espace algébrique sur $S$ et soit $U$ un espace algébrique simplicial
sur $S$. Supposons que nous ayons une augmentation
$$
a : U \to X
$$
Voir la section \ref{section-simplicial-algebraic-spaces}.
On dit que $U$ est un {\it hyperrecouvrement fppf} de $X$ si
\begin{enumerate}
\item $U_0 \to X$ est plat, localement de présentation finie, et surjectif,
\item $U_1 \to U_0 \times_X U_0$ est plat, localement de présentation finie, et
surjectif,
\item $U_{n + 1} \to (\text{cosk}_n\text{sk}_n U)_{n + 1}$
est plat, localement de présentation finie, et surjectif pour $n \geq 1$.
\end{enumerate}
La catégorie des espaces algébriques sur $S$ admet toutes les limites finies; les cosquelettes
utilisés dans la formulation ci-dessus existent donc.
$$
\fbox{Principe : les hyperrecouvrements fppf permettent de calculer la cohomologie étale.}
$$
L'idée essentielle de la démonstration consiste à comparer les topologies
fppf et \'etale sur la catégorie $\textit{Spaces}/S$.
La topologie fppf est en effet plus fine que la topologie \'etale, et l'on a :
(a) tout morphisme plat, localement de présentation finie et surjectif définit
un recouvrement fppf, et
(b) la cohomologie fppf des faisceaux provenant par image inverse du petit site \'etale
coïncide avec la cohomologie \'etale, comme on l'a vu dans
Compléments sur la cohomologie des espaces, section
\ref{spaces-more-cohomology-section-fppf-etale}.

\begin{lemma}
\label{lemma-compare-simplicial-objects-fppf-etale}
Soit $S$ un schéma. Soit $X$ un espace algébrique sur $S$.
Soit $U$ un espace algébrique simplicial sur $S$.
Soit $a : U \to X$ une augmentation. Il existe un diagramme commutatif
$$
\xymatrix{
\Sh((\textit{Spaces}/U)_{fppf, total}) \ar[r]_-h \ar[d]_{a_{fppf}} &
\Sh(U_\etale) \ar[d]^a \\
\Sh((\textit{Spaces}/X)_{fppf}) \ar[r]^-{h_{-1}} &
\Sh(X_\etale)
}
$$
où la flèche verticale gauche est définie dans la section
\ref{section-hypercovering}
et la flèche verticale droite est définie dans la section
\ref{section-simplicial-algebraic-spaces}.
\end{lemma}

\begin{proof}
La notation $(\textit{Spaces}/U)_{fppf, total}$ indique que
nous utilisons la construction de
la section \ref{section-hypercovering}
pour le site $(\textit{Spaces}/S)_{fppf}$ et l'objet simplicial
$U$ de ce site\footnote{On pourrait aussi
utiliser la topologie \'etale; on noterait alors
$(\textit{Spaces}/U)_{\etale, total}$.}.
Nous emploierons les sites $X_{spaces, \etale}$ et $U_{spaces, \etale}$
pour les topoi du membre de droite; cela est permis,
voir la discussion de la section \ref{section-simplicial-algebraic-spaces}.

\medskip\noindent
Observons que $(\textit{Spaces}/U)_{fppf, total}$ et
$U_{spaces, \etale}$
relèvent tous deux du cas (A) de la situation \ref{situation-simplicial-site}.
Cela résulte immédiatement de la construction de
$U_\etale$ dans la section \ref{section-simplicial-algebraic-spaces}
et découle du lemme \ref{lemma-sr-when-fibre-products}
pour $(\textit{Spaces}/U)_{fppf, total}$.
Considérons ensuite les foncteurs
$U_{n, spaces, \etale} \to (\textit{Spaces}/U_n)_{fppf}$, $U \mapsto U/U_n$
et
$X_{spaces, \etale} \to (\textit{Spaces}/X)_{fppf}$, $U \mapsto U/X$.
Nous avons vu que ces foncteurs définissent des morphismes de sites dans
Compléments sur la cohomologie des espaces, section
\ref{spaces-more-cohomology-section-fppf-etale}
où ils étaient notés $a_{U_n} = \epsilon_{U_n} \circ \pi_{u_n}$
et $a_X = \epsilon_X \circ \pi_X$.
On obtient ainsi un morphisme de sites simpliciaux compatible avec les augmentations
comme dans la remarque \ref{remark-morphism-augmentation-simplicial-sites}
et on peut appliquer le lemme
\ref{lemma-morphism-augmentation-simplicial-sites} pour conclure.
\end{proof}

\begin{lemma}
\label{lemma-descent-sheaves-for-fppf-hypercovering}
Soit $S$ un schéma. Soit $X$ un espace algébrique sur $S$.
Soit $U$ un espace algébrique simplicial sur $S$. Soit $a : U \to X$
une augmentation. Si $a : U \to X$ est un hyperrecouvrement fppf de $X$,
alors
$$
a^{-1} : \Sh(X_\etale) \to \Sh(U_\etale)
\quad\text{et}\quad
a^{-1} : \textit{Ab}(X_\etale) \to \textit{Ab}(U_\etale)
$$
sont pleinement fidèles, d'image essentielle les faisceaux cartésiens et
de quasi-inverse $a_*$. Ici, $a : \Sh(U_\etale) \to \Sh(X_\etale)$
est comme dans la section \ref{section-simplicial-algebraic-spaces}.
\end{lemma}

\begin{proof}
Nous allons démontrer l'énoncé pour les faisceaux d'ensembles. Il découlera
presque formellement de résultats déjà établis.
Considérons le diagramme du lemme
\ref{lemma-compare-simplicial-objects-fppf-etale}.
Dans la preuve de ce lemme nous avons vu que
$h_{-1}$ est le morphisme $a_X$ de
Compléments sur la cohomologie des espaces, section
\ref{spaces-more-cohomology-section-fppf-etale}.
Il résulte donc de
Compléments sur la cohomologie des espaces, lemme
\ref{spaces-more-cohomology-lemma-comparison-fppf-etale}
que $(h_{-1})^{-1}$ est pleinement fidèle, de quasi-inverse $h_{-1, *}$.
Il en va de même pour les composantes $h_n$ de $h$.
D'après la description des foncteurs $h^{-1}$ et $h_*$ donnée dans le
lemme \ref{lemma-morphism-simplicial-sites},
on conclut que $h^{-1}$ est pleinement fidèle, de quasi-inverse $h_*$.
Observons que $U$ est un hyperrecouvrement de $X$ dans $(\textit{Spaces}/S)_{fppf}$
au sens de la section \ref{section-hypercovering}.
Par le lemme \ref{lemma-hypercovering-X-simple-descent-sheaves}
on voit que $a_{fppf}^{-1}$ est pleinement fidèle, de quasi-inverse
$a_{fppf, *}$ et d'image essentielle les faisceaux cartésiens
sur $(\textit{Spaces}/U)_{fppf, total}$.
Une chasse formelle dans le diagramme montre alors que
$a^{-1}$ est pleinement fidèle.

\medskip\noindent
Enfin, supposons que $\mathcal{G}$ soit un faisceau cartésien sur $U_\etale$.
Alors $h^{-1}\mathcal{G}$ est un faisceau cartésien sur
$(\textit{Spaces}/U)_{fppf, total}$. Ainsi
$h^{-1}\mathcal{G} = a_{fppf}^{-1}\mathcal{H}$ pour un certain faisceau
$\mathcal{H}$ sur $(\textit{Spaces}/X)_{fppf}$.
En particulier, on obtient
$h_0^{-1}\mathcal{G}_0 = (a_{0, big, fppf})^{-1}\mathcal{H}$.
Comme $h_0 = a_{U_0}$ et que $U_0 \to X$ est
plat, localement de présentation finie et surjectif, le
lemme de Compléments sur la cohomologie des espaces
\ref{spaces-more-cohomology-lemma-descent-sheaf-fppf-etale} montre
qu'il existe un faisceau $\mathcal{F}$ sur $X_\etale$ et un isomorphisme
$\mathcal{H} = (h_{-1})^{-1}\mathcal{F}$.
Puisque $a_{fppf}^{-1}\mathcal{H} = h^{-1}\mathcal{G}$
on en déduit que $h^{-1}\mathcal{G} \cong h^{-1}a^{-1}\mathcal{F}$.
Par la pleine fidélité de $h^{-1}$, nous concluons que
$a^{-1}\mathcal{F} \cong \mathcal{G}$.

\medskip\noindent
Fixons un isomorphisme $\theta : a^{-1}\mathcal{F} \to \mathcal{G}$.
Pour terminer la preuve, il faut montrer $\mathcal{G} = a^{-1}a_*\mathcal{G}$
(afin d'établir que le quasi-inverse est donné par $a_*$; tout le reste
a été prouvé ci-dessus).
Puisque $a^{-1}$ est pleinement fidèle, on a $\text{id} \cong a_*a^{-1}$ d'après
Catégories, lemme \ref{categories-lemma-adjoint-fully-faithful}.
Ainsi $\mathcal{F} \cong a_*a^{-1}\mathcal{F}$ et
$a_*\theta : a_*a^{-1}\mathcal{F} \to a_*\mathcal{G}$
se combinent en un isomorphisme $\mathcal{F} \to a_*\mathcal{G}$.
En prenant l'image inverse par $a$ et en précomposant par $\theta^{-1}$,
on obtient l'isomorphisme souhaité.
\end{proof}

\begin{lemma}
\label{lemma-cohomological-descent-for-fppf-hypercovering}
Soient $S$ un schéma et $X$ un espace algébrique sur $S$.
Soit $U$ un espace algébrique simplicial sur $S$. Soit $a : U \to X$
une augmentation. Si $a : U \to X$ est un hyperrecouvrement fppf de $X$,
alors pour $K \in D^+(X_\etale)$
$$
K \to Ra_*(a^{-1}K)
$$
est un isomorphisme. Ici, $a : \Sh(U_\etale) \to \Sh(X_\etale)$
est comme dans la section \ref{section-simplicial-algebraic-spaces}.
\end{lemma}

\begin{proof}
Considérons le diagramme du lemme \ref{lemma-compare-simplicial-objects-fppf-etale}.
Observons que $Rh_{n, *}h_n^{-1}$ est le foncteur identité
sur $D^+(U_{n, \etale})$ d'après
Compléments sur la cohomologie des espaces, lemme
\ref{spaces-more-cohomology-lemma-cohomological-descent-etale-fppf}.
Par conséquent, $Rh_*h^{-1}$ est le foncteur identité sur
$D^+(U_\etale)$ par
le lemme \ref{lemma-direct-image-morphism-simplicial-sites}.
Nous avons
\begin{align*}
Ra_*(a^{-1}K)
& =
Ra_*Rh_*h^{-1}a^{-1}K \\
& =
Rh_{-1, *}Ra_{fppf, *}a_{fppf}^{-1}(h_{-1})^{-1}K \\
& =
Rh_{-1, *}(h_{-1})^{-1}K \\
& =
K
\end{align*}
La première égalité découle de la discussion ci-dessus; la deuxième,
de la commutativité du diagramme du lemme
\ref{lemma-compare-simplicial-objects}; la troisième, du lemme
\ref{lemma-hypercovering-X-simple-descent-bounded-abelian}
puisque $U$ est un hyperrecouvrement de $X$ dans $(\textit{Spaces}/S)_{fppf}$;
et la dernière, du lemme déjà utilisé de
Compléments sur la cohomologie des espaces
\ref{spaces-more-cohomology-lemma-cohomological-descent-etale-fppf}.
\end{proof}

\begin{lemma}
\label{lemma-compute-via-fppf-hypercovering}
Soient $S$ un schéma et $X$ un espace algébrique sur $S$.
Soit $U$ un espace algébrique simplicial sur $S$. Soit $a : U \to X$
une augmentation. Si $a : U \to X$ est un hyperrecouvrement fppf de $X$, alors
$$
R\Gamma(X_\etale, K) = R\Gamma(U_\etale, a^{-1}K)
$$
pour $K \in D^+(X_\etale)$. Ici, $a : \Sh(U_\etale) \to \Sh(X_\etale)$
est comme dans la section \ref{section-simplicial-algebraic-spaces}.
\end{lemma}

\begin{proof}
Cela découle du lemme
\ref{lemma-cohomological-descent-for-fppf-hypercovering}
car $R\Gamma(U_\etale, -) = R\Gamma(X_\etale, -) \circ Ra_*$ par
Cohomologie sur les sites, remarque \ref{sites-cohomology-remark-before-Leray}.
\end{proof}

\begin{lemma}
\label{lemma-fppf-hypercovering-equivalence-bounded}
Soient $S$ un schéma et $X$ un espace algébrique sur $S$.
Soit $U$ un espace algébrique simplicial sur $S$. Soit $a : U \to X$
une augmentation.
Soit $\mathcal{A} \subset \textit{Ab}(U_\etale)$
la sous-catégorie de Serre faible des faisceaux abéliens cartésiens.
Si $U$ est un hyperrecouvrement fppf de $X$, alors
le foncteur $a^{-1}$ définit une équivalence
$$
D^+(X_\etale) \longrightarrow D_\mathcal{A}^+(U_\etale)
$$
de quasi-inverse $Ra_*$. Ici, $a : \Sh(U_\etale) \to \Sh(X_\etale)$
est comme dans la section \ref{section-simplicial-algebraic-spaces}.
\end{lemma}

\begin{proof}
Observons que $\mathcal{A}$ est une sous-catégorie de Serre faible d'après le
lemme \ref{lemma-Serre-subcat-cartesian-modules}.
L'équivalence est une conséquence formelle
des résultats obtenus jusqu'à présent. On utilise les lemmes
\ref{lemma-descent-sheaves-for-fppf-hypercovering} et
\ref{lemma-cohomological-descent-for-fppf-hypercovering}, ainsi que le lemme de Cohomologie
sur les sites \ref{sites-cohomology-lemma-equivalence-bounded}.
\end{proof}

\begin{lemma}
\label{lemma-spectral-sequence-fppf-hypercovering}
Soient $S$ un schéma et $X$ un espace algébrique sur $S$.
Soit $U$ un espace algébrique simplicial sur $S$. Soit $a : U \to X$
une augmentation. Soit $\mathcal{F}$ un faisceau abélien
sur $X_\etale$. Notons $\mathcal{F}_n$ son image inverse sur $U_{n, \etale}$.
Si $U$ est un hyperrecouvrement fppf de $X$, alors
il existe une suite spectrale canonique
$$
E_1^{p, q} = H^q_\etale(U_p, \mathcal{F}_p)
$$
convergeant vers $H^{p + q}_\etale(X, \mathcal{F})$.
\end{lemma}

\begin{proof}
Conséquence immédiate des lemmes \ref{lemma-compute-via-fppf-hypercovering}
et \ref{lemma-simplicial-sheaf-cohomology-site}.
\end{proof}








\section{Hyperrecouvrements fppf d'espaces algébriques : modules}
\label{section-fppf-hypercovering-modules}

\noindent
Nous poursuivons l'étude de la descente (cohomologique) pour les hyperrecouvrements fppf,
entreprise dans la section \ref{section-fppf-hypercovering},
en examinant ici le cas des faisceaux de modules.
Nous nous intéressons surtout aux modules quasi-cohérents,
pour lesquels une descente cohomologique non bornée est possible.

\begin{lemma}
\label{lemma-compare-simplicial-objects-fppf-etale-modules}
Soit $S$ un schéma. Soit $X$ un espace algébrique sur $S$.
Soit $U$ un espace algébrique simplicial sur $S$.
Soit $a : U \to X$ une augmentation. Il existe un diagramme commutatif
$$
\xymatrix{
(\Sh((\textit{Spaces}/U)_{fppf, total}), \mathcal{O}_{big, total})
\ar[r]_-h \ar[d]_{a_{fppf}} &
(\Sh(U_\etale), \mathcal{O}_U) \ar[d]^a \\
(\Sh((\textit{Spaces}/X)_{fppf}), \mathcal{O}_{big}) \ar[r]^-{h_{-1}} &
(\Sh(X_\etale), \mathcal{O}_X)
}
$$
de topoi annelés, où la flèche verticale gauche est définie dans la section
\ref{section-hypercovering-modules}
et la flèche verticale droite est définie dans la section
\ref{section-simplicial-algebraic-spaces}.
\end{lemma}

\begin{proof}
Pour le diagramme sous-jacent de topoi, voir la discussion de
la preuve du lemme \ref{lemma-compare-simplicial-objects-fppf-etale}.
Le faisceau $\mathcal{O}_U$ est le faisceau structural de l'espace algébrique simplicial
$U$, tel qu'il est défini dans la section
\ref{section-simplicial-algebraic-spaces}.
Le faisceau $\mathcal{O}_X$ est le faisceau structural usuel de l'espace algébrique
$X$. Les faisceaux d'anneaux $\mathcal{O}_{big, total}$ et
$\mathcal{O}_{big}$ proviennent du faisceau structural sur
$(\textit{Spaces}/S)_{fppf}$ de la manière expliquée dans la section
\ref{section-hypercovering-modules}
qui construit aussi $a_{fppf}$ comme morphisme de topoi annelés.
Les morphismes sur les composantes $h_n = a_{U_n}$ et $h_{-1} = a_X$
sont des morphismes de topoi annelés d'après
Compléments sur la cohomologie des espaces, section
\ref{spaces-more-cohomology-section-fppf-etale-modules}.
Enfin, puisque le foncteur continu
$u : U_{spaces, \etale} \to (\textit{Spaces}/U)_{fppf, total}$
utilisé pour définir $h$\footnote{Ce foncteur intervient dans la preuve du
lemme \ref{lemma-compare-simplicial-objects-fppf-etale}
via une application du lemme
\ref{lemma-morphism-augmentation-simplicial-sites}.}
est donné par $V/U_n \mapsto V/U_n$,
on voit que $h_*\mathcal{O}_{big, total} = \mathcal{O}_U$;
c'est ainsi que l'on munit $h$ de la structure d'un morphisme
de sites simpliciaux annelés comme dans
la remarque \ref{remark-morphism-simplicial-sites-modules}.
On obtient alors un morphisme de topoi annelés $h$
par le lemme \ref{lemma-morphism-simplicial-sites-modules}.
Notons que les morphismes $h_n$ coïncident bien
avec les morphismes $a_{U_n}$ décrits ci-dessus.
Nous omettons la vérification
que le diagramme est commutatif en tant que diagramme de topoi annelés
-- nous savons déjà qu'il est commutatif
en tant que diagramme de topoi).
\end{proof}

\begin{lemma}
\label{lemma-descent-qcoh-for-fppf-hypercovering}
Soient $S$ un schéma et $X$ un espace algébrique sur $S$.
Soit $U$ un espace algébrique simplicial sur $S$, et soit $a : U \to X$
une augmentation. Si $a : U \to X$ est un hyperrecouvrement fppf de $X$,
alors
$$
a^* : \QCoh(\mathcal{O}_X) \to \QCoh(\mathcal{O}_U)
$$
est une équivalence, de quasi-inverse $a_*$.
Ici, $a : \Sh(U_\etale) \to \Sh(X_\etale)$
est comme dans la section \ref{section-simplicial-algebraic-spaces}.
\end{lemma}

\begin{proof}
Considérons le diagramme du lemme
\ref{lemma-compare-simplicial-objects-fppf-etale-modules}.
Dans la preuve de ce lemme nous avons vu que
$h_{-1}$ est le morphisme $a_X$ de
Compléments sur la cohomologie des espaces, section
\ref{spaces-more-cohomology-section-fppf-etale-modules}.
Il résulte donc de
Compléments sur la cohomologie des espaces, lemme
\ref{spaces-more-cohomology-lemma-review-quasi-coherent}
que
$$
(h_{-1})^* :
\QCoh(\mathcal{O}_X)
\longrightarrow
\QCoh(\mathcal{O}_{big})
$$
est une équivalence, de quasi-inverse $h_{-1, *}$.
Il en va de même pour les composantes $h_n$ de $h$.
Rappelons que $\QCoh(\mathcal{O}_U)$ et $\QCoh(\mathcal{O}_{big, total})$
sont formées des modules cartésiens dont les composantes sont quasi-cohérentes; voir le
lemme \ref{lemma-quasi-coherent-sheaf}.
Puisque les foncteurs $h^*$ et $h_*$ du
lemme \ref{lemma-morphism-simplicial-sites-modules}
coïncident avec les foncteurs $h_n^*$ et $h_{n, *}$ sur les composantes,
nous concluons que
$$
h^* :
\QCoh(\mathcal{O}_U)
\longrightarrow 
\QCoh(\mathcal{O}_{big, total})
$$
est une équivalence, de quasi-inverse $h_*$.
Observons que $U$ est un hyperrecouvrement de $X$ dans $(\textit{Spaces}/S)_{fppf}$
au sens de la section \ref{section-hypercovering}.
D'après le lemme \ref{lemma-hypercovering-X-simple-descent-modules},
$a_{fppf}^*$ est pleinement fidèle, de quasi-inverse
$a_{fppf, *}$, et son image essentielle est formée des
$\mathcal{O}_{fppf, total}$-modules cartésiens.
Ainsi, par la description de $\QCoh(\mathcal{O}_{big})$ et
$\QCoh(\mathcal{O}_{big, total})$ donnée par le lemme \ref{lemma-quasi-coherent-sheaf},
nous obtenons une équivalence
$$
a_{fppf}^* :
\QCoh(\mathcal{O}_{big})
\longrightarrow
\QCoh(\mathcal{O}_{big, total})
$$
de quasi-inverse $a_{fppf, *}$.
Une chasse formelle dans le diagramme montre alors que
$a^*$ est pleinement fidèle sur $\QCoh(\mathcal{O}_X)$ et que son image est
contenue dans $\QCoh(\mathcal{O}_U)$.

\medskip\noindent
Enfin, supposons que $\mathcal{G}$ appartienne à $\QCoh(\mathcal{O}_U)$.
Alors $h^*\mathcal{G}$ appartient à $\QCoh(\mathcal{O}_{big, total})$.
Ainsi $h^*\mathcal{G} = a_{fppf}^*\mathcal{H}$ avec
$\mathcal{H} = a_{fppf, *}h^*\mathcal{G}$
dans $\QCoh(\mathcal{O}_{big})$ (voir ci-dessus).
On voit ensuite que $\mathcal{H} = (h_{-1})^*\mathcal{F}$
avec $\mathcal{F} = h_{-1, *}\mathcal{H}$ dans $\QCoh(\mathcal{O}_X)$.
En parcourant le diagramme, on en déduit que
$h^*\mathcal{G} \cong h^*a^*\mathcal{F}$.
Par la pleine fidélité de $h^*$, nous concluons que
$a^*\mathcal{F} \cong \mathcal{G}$.
Puisque $\mathcal{F} = h_{-1, *}a_{fppf, *}h^*\mathcal{G} =
a_*h_*h^*\mathcal{G} = a_*\mathcal{G}$, nous obtenons également
l'énoncé selon lequel le quasi-inverse est donné par $a_*$.
\end{proof}

\begin{lemma}
\label{lemma-cohomological-descent-qcoh-for-fppf-hypercovering}
Soient $S$ un schéma et $X$ un espace algébrique sur $S$.
Soit $U$ un espace algébrique simplicial sur $S$. Soit $a : U \to X$
une augmentation. Si $a : U \to X$ est un hyperrecouvrement fppf de $X$,
alors, pour tout $\mathcal{F}$ qui est un $\mathcal{O}_X$-module quasi-cohérent,
le morphisme
$$
\mathcal{F} \to Ra_*(a^*\mathcal{F})
$$
est un isomorphisme. Ici, $a : \Sh(U_\etale) \to \Sh(X_\etale)$
est comme dans la section \ref{section-simplicial-algebraic-spaces}.
\end{lemma}

\begin{proof}
Considérons le diagramme du lemme \ref{lemma-compare-simplicial-objects-fppf-etale}.
Soit $\mathcal{F}_n = a_n^*\mathcal{F}$ la composante de degré $n$ de
$a^*\mathcal{F}$. C'est un $\mathcal{O}_{U_n}$-module quasi-cohérent.
Alors $\mathcal{F}_n = Rh_{n, *}h_n^*\mathcal{F}_n$
par Compléments sur la cohomologie des espaces, lemme
\ref{spaces-more-cohomology-lemma-cohomological-descent-etale-fppf-modules}.
Ainsi $a^*\mathcal{F} = Rh_*h^*a^*\mathcal{F}$ d'après le
lemme \ref{lemma-direct-image-morphism-simplicial-sites-modules}.
Nous avons
\begin{align*}
Ra_*(a^*\mathcal{F})
& =
Ra_*Rh_*h^*a^*\mathcal{F} \\
& =
Rh_{-1, *}Ra_{fppf, *}a_{fppf}^*(h_{-1})^*\mathcal{F} \\
& =
Rh_{-1, *}(h_{-1})^*\mathcal{F} \\
& =
\mathcal{F}
\end{align*}
La première égalité découle de la discussion ci-dessus; la deuxième,
de la commutativité du diagramme du lemme
\ref{lemma-compare-simplicial-objects}; la troisième, du lemme
\ref{lemma-hypercovering-X-simple-descent-bounded-modules}
puisque $U$ est un hyperrecouvrement de $X$ dans $(\textit{Spaces}/S)_{fppf}$
et que $La_{fppf}^* = a_{fppf}^*$, car $a_{fppf}$ est plat
(c'est-à-dire que $a_{fppf}^{-1}\mathcal{O}_{big} = \mathcal{O}_{big, total}$;
voir la remarque \ref{remark-augmentation-ringed}); enfin,
la dernière égalité découle du lemme déjà utilisé de
Compléments sur la cohomologie des espaces
\ref{spaces-more-cohomology-lemma-cohomological-descent-etale-fppf-modules}.
\end{proof}

\begin{lemma}
\label{lemma-coh-descent-qcoh-unbounded-for-fppf-hypercovering}
Soient $S$ un schéma et $X$ un espace algébrique sur $S$.
Soit $U$ un espace algébrique simplicial sur $S$. Soit $a : U \to X$
une augmentation. Supposons que $a : U \to X$ soit un hyperrecouvrement fppf de $X$.
Alors $\QCoh(\mathcal{O}_U)$ est une sous-catégorie de Serre faible de
$\textit{Mod}(\mathcal{O}_U)$ et
$$
a^* : D_\QCoh(\mathcal{O}_X) \longrightarrow D_\QCoh(\mathcal{O}_U)
$$
est une équivalence de catégories, de quasi-inverse
$Ra_*$. Ici, $a : \Sh(U_\etale) \to \Sh(X_\etale)$
est comme dans la section \ref{section-simplicial-algebraic-spaces}.
\end{lemma}

\begin{proof}
Observons d'abord que les morphismes $a_n : U_n \to X$ et $d^n_i : U_n \to U_{n - 1}$
sont plats, localement de présentation finie et surjectifs d'après
Hyperrecouvrements, remarque \ref{hypercovering-remark-P-covering}.

\medskip\noindent
Rappelons qu'un $\mathcal{O}_U$-module $\mathcal{F}$ est quasi-cohérent si et
seulement s'il est cartésien et si $\mathcal{F}_n$ est quasi-cohérent pour tout $n$.
Voir le lemme \ref{lemma-quasi-coherent-sheaf}.
Par le lemme \ref{lemma-Serre-subcat-cartesian-modules}
(et la platitude des morphismes $d^n_i : U_n \to U_{n - 1}$ établie ci-dessus)
les modules cartésiens forment une sous-catégorie de Serre faible de
$\textit{Mod}(\mathcal{O}_U)$. D'autre part,
$\QCoh(\mathcal{O}_{U_n}) \subset \textit{Mod}(\mathcal{O}_{U_n})$
est une sous-catégorie de Serre faible pour tout $n$
(Propriétés des espaces, lemme
\ref{spaces-properties-lemma-properties-quasi-coherent}).
Ces faits montrent ensemble que
$\QCoh(\mathcal{O}_U) \subset \textit{Mod}(\mathcal{O}_U)$
est une sous-catégorie de Serre faible.

\medskip\noindent
Pour achever la preuve, vérifions une à une les conditions (1) -- (5) du lemme
de Cohomologie sur les sites
\ref{sites-cohomology-lemma-equivalence-unbounded-one}.

\medskip\noindent
Pour (1). Comme $a_n$ est plat (voir ci-dessus), $a$ est plat
par le lemme \ref{lemma-flat-augmentation-modules}.

\medskip\noindent
Pour (2). C'est exactement l'énoncé du
lemme \ref{lemma-descent-qcoh-for-fppf-hypercovering}.

\medskip\noindent
Pour (3). C'est exactement l'énoncé du
lemme \ref{lemma-cohomological-descent-qcoh-for-fppf-hypercovering}.

\medskip\noindent
Pour (4). Rappelons que l'on peut utiliser soit le site $U_\etale$, soit le site
$U_{spaces, \etale}$ pour définir le petit topos \'etale
$\Sh(U_\etale)$; voir la section \ref{section-simplicial-algebraic-spaces}.
L'hypothèse de la situation de Cohomologie
sur les sites \ref{sites-cohomology-situation-olsson-laszlo}
est satisfaite pour le triplet
$(U_{spaces, \etale}, \mathcal{O}_U, \QCoh(\mathcal{O}_U))$
et, par le même raisonnement, pour le triplet
$(U_\etale, \mathcal{O}_U, \QCoh(\mathcal{O}_U))$.
Plus précisément, prenons pour
$$
\mathcal{B} \subset \Ob(U_\etale) \subset \Ob(U_{spaces, \etale})
$$
l'ensemble des objets affines. Pour $V/U_n \in \mathcal{B}$,
prenons $d_{V/U_n} = 0$ et pour $\text{Cov}_{V/U_n}$ l'ensemble des
recouvrements étales $\{V_i \to V\}$ tels que les $V_i$ soient affines.
On obtient alors l'annulation voulue, car pour
$\mathcal{F} \in \QCoh(\mathcal{O}_U)$
et tout $V/U_n \in \mathcal{B}$, nous avons
$$
H^p(V/U_n, \mathcal{F}) = H^p(V, \mathcal{F}_n)
$$
d'après le lemme \ref{lemma-sanity-check-modules}. Dans le
membre de droite figure la cohomologie du faisceau quasi-cohérent
$\mathcal{F}_n$ sur $U_n$, évaluée sur l'objet affine $V$
de $U_{n, \etale}$. Ce groupe s'annule pour $p > 0$ d'après
Cohomologie des espaces, section
\ref{spaces-cohomology-section-higher-direct-image} et Cohomologie des schémas,
lemme
\ref{coherent-lemma-quasi-coherent-affine-cohomology-zero}.

\medskip\noindent
Pour (5). Il suffit de prendre pour $\mathcal{B} \subset \Ob(X_{spaces, \etale})$
l'ensemble des objets affines et d'utiliser les références données ci-dessus.
\end{proof}

\begin{lemma}
\label{lemma-compute-via-fppf-hypercovering-modules}
Soit $S$ un schéma. Soit $X$ un espace algébrique sur $S$.
Soit $U$ un espace algébrique simplicial sur $S$. Soit $a : U \to X$
une augmentation. Si $a : U \to X$ est un hyperrecouvrement fppf de $X$, alors
$$
R\Gamma(X_\etale, K) = R\Gamma(U_\etale, a^*K)
$$
pour $K \in D_\QCoh(\mathcal{O}_X)$. Ici, $a : \Sh(U_\etale) \to \Sh(X_\etale)$
est comme dans la section \ref{section-simplicial-algebraic-spaces}.
\end{lemma}

\begin{proof}
Cela découle du lemme
\ref{lemma-coh-descent-qcoh-unbounded-for-fppf-hypercovering}
car $R\Gamma(U_\etale, -) = R\Gamma(X_\etale, -) \circ Ra_*$ par
Cohomologie sur les sites, remarque \ref{sites-cohomology-remark-before-Leray}.
\end{proof}

\begin{lemma}
\label{lemma-spectral-sequence-fppf-hypercovering-modules}
Soient $S$ un schéma et $X$ un espace algébrique sur $S$.
Soit $U$ un espace algébrique simplicial sur $S$. Soit $a : U \to X$
une augmentation. Soit $\mathcal{F}$ un
$\mathcal{O}_X$-module quasi-cohérent. Notons $\mathcal{F}_n$ son image inverse sur
$U_{n, \etale}$. Si $U$ est un hyperrecouvrement fppf de $X$, alors
il existe une suite spectrale canonique
$$
E_1^{p, q} = H^q_\etale(U_p, \mathcal{F}_p)
$$
convergeant vers $H^{p + q}_\etale(X, \mathcal{F})$.
\end{lemma}

\begin{proof}
Conséquence immédiate de
les lemmes \ref{lemma-compute-via-fppf-hypercovering-modules}
et \ref{lemma-simplicial-module-cohomology-site}.
\end{proof}







\section{Descente fppf de complexes}
\label{section-fppf-descent-derived}

\noindent
Dans cette section, on rassemble certains résultats établis plus haut
pour les recouvrements fppf d'espaces algébriques
et les catégories dérivées de modules quasi-cohérents.

\begin{lemma}
\label{lemma-fppf-neg-ext-zero-hom}
Soit $X$ un espace algébrique sur un schéma $S$.
Soient $K, E \in D_\QCoh(\mathcal{O}_X)$.
Soit $a : U \to X$ un hyperrecouvrement fppf.
Supposons que, pour tout $n \geq 0$, on ait
$$
\Ext_{\mathcal{O}_{U_n}}^i(La_n^*K, La_n^*E) = 0
\text{ pour } i < 0
$$
Alors :
\begin{enumerate}
\item $\Ext_{\mathcal{O}_X}^i(K, E) = 0$ pour $i < 0$ et
\item il existe une suite exacte
$$
0
\to
\Hom_{\mathcal{O}_X}(K, E)
\to
\Hom_{\mathcal{O}_{U_0}}(La_0^*K, La_0^*E)
\to
\Hom_{\mathcal{O}_{U_1}}(La_1^*K, La_1^*E)
$$
\end{enumerate}
\end{lemma}

\begin{proof}
Posons $K_n = La_n^*K$ et $E_n = La_n^*E$. Ce sont alors les composantes des
systèmes simpliciaux de la catégorie dérivée de modules
(définition \ref{definition-cartesian-derived-modules})
associés à $La^*K$ et $La^*E$
(lemme \ref{lemma-cartesian-objects-derived-modules}),
où $a : U_\etale \to X_\etale$ est comme dans la section
\ref{section-simplicial-algebraic-spaces}.
Démontrons d'abord (2). D'après le
lemme \ref{lemma-coh-descent-qcoh-unbounded-for-fppf-hypercovering},
nous avons
$$
\Hom_{\mathcal{O}_X}(K, E) =
\Hom_{\mathcal{O}_U}(La^*K, La^*E)
$$
La suite prend donc la forme suivante :
$$
0
\to
\Hom_{\mathcal{O}_U}(La^*K, La^*E)
\to
\Hom_{\mathcal{O}_{U_0}}(K_0, E_0)
\to
\Hom_{\mathcal{O}_{U_1}}(K_1, E_1)
$$
La première flèche est injective par
le lemme \ref{lemma-nullity-cartesian-modules-derived}.
L'image de cette flèche est le noyau de la seconde
par le lemme \ref{lemma-hom-cartesian-modules-derived}.
Cela termine la preuve de (2).
La partie (1) résulte de la partie (2) appliquée à
$K[i]$ et $E$ pour $i > 0$.
\end{proof}

\begin{lemma}
\label{lemma-fppf-glue-neg-ext-zero}
Soit $X$ un espace algébrique sur un schéma $S$.
Soit $a : U \to X$ un hyperrecouvrement fppf.
Supposons que soient donnés $K_0 \in D_\QCoh(U_0)$ et un isomorphisme
$$
\alpha :
L(f_{\delta_1^1})^*K_0
\longrightarrow
L(f_{\delta_0^1})^*K_0
$$
satisfaisant à la condition de cocycle sur $U_1$. Posons
$\tau^n_i : [0] \to [n]$, $0 \mapsto i$ et
posons $K_n = Lf_{\tau^n_n}^*K_0$.
Supposons $\Ext^i_{\mathcal{O}_{U_n}}(K_n, K_n) = 0$ pour $i < 0$.
Il existe alors un objet $K \in D_\QCoh(\mathcal{O}_X)$
et un isomorphisme $La_0^*K \to K$ compatible avec $\alpha$.
\end{lemma}

\begin{proof}
Les objets $K_n$ forment les termes d'un système simplicial de la
catégorie dérivée de modules, d'après le
lemme \ref{lemma-construct-cartesian-objects-derived-modules}.
On obtient alors un objet
$K' \in D_\QCoh(\mathcal{O}_{U_\etale})$
tel que $(K_n, K_\varphi)$ soit le système déduit de $K'$; voir le lemme
\ref{lemma-cartesian-module-derived-from-simplicial}.
Enfin, appliquons le lemme
\ref{lemma-coh-descent-qcoh-unbounded-for-fppf-hypercovering}
pour voir que $K' = La^*K$ pour un certain $K \in D_\QCoh(\mathcal{O}_X)$
ce qui conclut.
\end{proof}











\section{Hyperrecouvrements propres d'espaces algébriques}
\label{section-proper-hypercovering-spaces}

\noindent
Cette section est l'analogue de la section \ref{section-proper-hypercovering}
pour le cas des espaces algébriques.
Le lecteur qui le souhaite peut remplacer partout « espace algébrique »
par « schéma » et obtenir des résultats également valables.
Cela présente l'avantage de remplacer les références à
Compléments sur la cohomologie des espaces, section
\ref{spaces-more-cohomology-section-ph-etale}
par des références à
Cohomologie étale, section \ref{etale-cohomology-section-ph-etale}.

\medskip\noindent
Nous fixons un schéma de base $S$.
Soit $X$ un espace algébrique sur $S$ et soit $U$ un espace algébrique simplicial
sur $S$. Supposons que nous ayons une augmentation
$$
a : U \to X
$$
Voir la section \ref{section-simplicial-algebraic-spaces}.
On dit que $U$ est un {\it hyperrecouvrement propre} de $X$ si
\begin{enumerate}
\item $U_0 \to X$ est propre et surjectif,
\item $U_1 \to U_0 \times_X U_0$ est propre et surjectif,
\item $U_{n + 1} \to (\text{cosk}_n\text{sk}_n U)_{n + 1}$
est propre et surjectif pour $n \geq 1$.
\end{enumerate}
La catégorie des espaces algébriques sur $S$ admet toutes les limites finies; les cosquelettes
utilisés dans la formulation ci-dessus existent donc.
$$
\fbox{Principe : les hyperrecouvrements propres permettent de
calculer la cohomologie étale.}
$$
L'idée essentielle de la démonstration consiste à comparer les topologies
ph et \'etale sur la catégorie $\textit{Spaces}/S$.
La topologie ph est en effet plus fine que la topologie \'etale, et l'on a :
(a) tout morphisme propre surjectif définit un recouvrement ph, et
(b) la cohomologie ph des faisceaux provenant par image inverse du petit site étale
coïncide avec la cohomologie étale, comme on l'a vu dans
Compléments sur la cohomologie des espaces, section
\ref{spaces-more-cohomology-section-ph-etale}.

\medskip\noindent
Tous les résultats de cette section se généralisent au cas
où $U \to X$ n'est qu'un « hyperrecouvrement ph », c'est-à-dire un hyperrecouvrement
de $X$ dans le site $(\textit{Spaces}/S)_{ph}$
au sens de la section \ref{section-hypercovering}. Si nous avons besoin de ce résultat,
nous en donnerons ici un énoncé précis et une démonstration.

\begin{lemma}
\label{lemma-compare-simplicial-objects-ph-etale}
Soit $S$ un schéma. Soit $X$ un espace algébrique sur $S$.
Soit $U$ un espace algébrique simplicial sur $S$.
Soit $a : U \to X$ une augmentation. Il existe un diagramme commutatif
$$
\xymatrix{
\Sh((\textit{Spaces}/U)_{ph, total}) \ar[r]_-h \ar[d]_{a_{ph}} &
\Sh(U_\etale) \ar[d]^a \\
\Sh((\textit{Spaces}/X)_{ph}) \ar[r]^-{h_{-1}} &
\Sh(X_\etale)
}
$$
où la flèche verticale gauche est définie dans la section
\ref{section-hypercovering}
et la flèche verticale droite est définie dans la section
\ref{section-simplicial-algebraic-spaces}.
\end{lemma}

\begin{proof}
La notation $(\textit{Spaces}/U)_{ph, total}$ indique que
nous utilisons la construction de
la section \ref{section-hypercovering}
pour le site $(\textit{Spaces}/S)_{ph}$ et l'objet simplicial
$U$ de ce site\footnote{Pour distinguer
$(\textit{Spaces}/U)_{fppf, total}$, défini à l'aide de la topologie fppf
dans la section \ref{section-fppf-hypercovering}.}.
Nous emploierons les sites $X_{spaces, \etale}$ et $U_{spaces, \etale}$
pour les topoi du membre de droite; cela est permis,
voir la discussion de la section \ref{section-simplicial-algebraic-spaces}.

\medskip\noindent
Observons que $(\textit{Spaces}/U)_{ph, total}$ et
$U_{spaces, \etale}$
relèvent tous deux du cas (A) de la situation \ref{situation-simplicial-site}.
Cela résulte immédiatement de la construction de
$U_\etale$ dans la section \ref{section-simplicial-algebraic-spaces}
et découle du lemme \ref{lemma-sr-when-fibre-products}
pour $(\textit{Spaces}/U)_{ph, total}$.
Considérons ensuite les foncteurs
$U_{n, spaces, \etale} \to (\textit{Spaces}/U_n)_{ph}$, $U \mapsto U/U_n$
et
$X_{spaces, \etale} \to (\textit{Spaces}/X)_{ph}$, $U \mapsto U/X$.
Nous avons vu que ces foncteurs définissent des morphismes de sites dans
Compléments sur la cohomologie des espaces, section
\ref{spaces-more-cohomology-section-ph-etale}
où ils étaient notés $a_{U_n} = \epsilon_{U_n} \circ \pi_{u_n}$
et $a_X = \epsilon_X \circ \pi_X$.
On obtient ainsi un morphisme de sites simpliciaux compatible avec les augmentations
comme dans la remarque \ref{remark-morphism-augmentation-simplicial-sites}
et on peut appliquer le lemme
\ref{lemma-morphism-augmentation-simplicial-sites} pour conclure.
\end{proof}

\begin{lemma}
\label{lemma-descent-sheaves-for-ph-hypercovering}
Soit $S$ un schéma. Soit $X$ un espace algébrique sur $S$.
Soit $U$ un espace algébrique simplicial sur $S$. Soit $a : U \to X$
une augmentation. Si $a : U \to X$ est un hyperrecouvrement propre de $X$,
alors
$$
a^{-1} : \Sh(X_\etale) \to \Sh(U_\etale)
\quad\text{et}\quad
a^{-1} : \textit{Ab}(X_\etale) \to \textit{Ab}(U_\etale)
$$
sont pleinement fidèles, d'image essentielle les faisceaux cartésiens et
de quasi-inverse $a_*$. Ici, $a : \Sh(U_\etale) \to \Sh(X_\etale)$
est comme dans la section \ref{section-simplicial-algebraic-spaces}.
\end{lemma}

\begin{proof}
Nous allons démontrer l'énoncé pour les faisceaux d'ensembles. Il découlera
presque formellement de résultats déjà établis.
Considérons le diagramme du lemme
\ref{lemma-compare-simplicial-objects-ph-etale}.
Dans la preuve de ce lemme nous avons vu que
$h_{-1}$ est le morphisme $a_X$ de
Compléments sur la cohomologie des espaces, section
\ref{spaces-more-cohomology-section-ph-etale}.
Il résulte donc de
Compléments sur la cohomologie des espaces, lemme
\ref{spaces-more-cohomology-lemma-comparison-ph-etale}
que $(h_{-1})^{-1}$ est pleinement fidèle, de quasi-inverse $h_{-1, *}$.
Il en va de même pour les composantes $h_n$ de $h$.
D'après la description des foncteurs $h^{-1}$ et $h_*$ donnée dans le
lemme \ref{lemma-morphism-simplicial-sites},
on conclut que $h^{-1}$ est pleinement fidèle, de quasi-inverse $h_*$.
Observons que $U$ est un hyperrecouvrement de $X$ dans $(\textit{Spaces}/S)_{ph}$
au sens de la section \ref{section-hypercovering}, puisqu'un morphisme propre
surjectif définit un recouvrement ph d'après Topologies sur les espaces, lemme
\ref{spaces-topologies-lemma-surjective-proper-ph}.
Par le lemme \ref{lemma-hypercovering-X-simple-descent-sheaves}
on voit que $a_{ph}^{-1}$ est pleinement fidèle, de quasi-inverse
$a_{ph, *}$ et d'image essentielle les faisceaux cartésiens
sur $(\textit{Spaces}/U)_{ph, total}$.
Une chasse formelle dans le diagramme montre alors que
$a^{-1}$ est pleinement fidèle.

\medskip\noindent
Enfin, supposons que $\mathcal{G}$ soit un faisceau cartésien sur $U_\etale$.
Alors $h^{-1}\mathcal{G}$ est un faisceau cartésien sur
$(\textit{Spaces}/U)_{ph, total}$.
Ainsi $h^{-1}\mathcal{G} = a_{ph}^{-1}\mathcal{H}$ pour un certain faisceau
$\mathcal{H}$ sur $(\textit{Spaces}/X)_{ph}$.
Calculons, en employant des notations quelque peu pédantes :
\begin{align*}
(h_{-1})^{-1}(a_*\mathcal{G})
& =
(h_{-1})^{-1}
\text{Eq}(
\xymatrix{
a_{0, small, *}\mathcal{G}_0
\ar@<1ex>[r] \ar@<-1ex>[r] &
a_{1, small, *}\mathcal{G}_1
}
) \\
& =
\text{Eq}(
\xymatrix{
(h_{-1})^{-1}a_{0, small, *}\mathcal{G}_0
\ar@<1ex>[r] \ar@<-1ex>[r] &
(h_{-1})^{-1}a_{1, small, *}\mathcal{G}_1
}
) \\
& =
\text{Eq}(
\xymatrix{
a_{0, big, ph, *}h_0^{-1}\mathcal{G}_0
\ar@<1ex>[r] \ar@<-1ex>[r] &
a_{1, big, ph, *}h_1^{-1}\mathcal{G}_1
}
) \\
& =
\text{Eq}(
\xymatrix{
a_{0, big, ph, *}(a_{0, big, ph})^{-1}\mathcal{H}
\ar@<1ex>[r] \ar@<-1ex>[r] &
a_{1, big, ph, *}(a_{1, big, ph})^{-1}\mathcal{H}
}
) \\
& =
a_{ph, *}a_{ph}^{-1}\mathcal{H} \\
& =
\mathcal{H}
\end{align*}
La première égalité découle du lemme \ref{lemma-augmentation-site};
la deuxième, de l'exactitude du foncteur $(h_{-1})^{-1}$;
la troisième, du
lemme de Compléments sur la cohomologie des espaces
\ref{spaces-more-cohomology-lemma-proper-push-pull-ph-etale}
(ici nous utilisons que $a_0 : U_0 \to X$ et $a_1: U_1 \to X$ sont propres);
la quatrième, de $a_{ph}^{-1}\mathcal{H} = h^{-1}\mathcal{G}$;
la cinquième, du lemme \ref{lemma-augmentation-site}; enfin la sixième
a été établie ci-dessus. Puisque $a_{ph}^{-1}\mathcal{H} = h^{-1}\mathcal{G}$,
on en déduit que $h^{-1}\mathcal{G} \cong h^{-1}a^{-1}a_*\mathcal{G}$,
ce qui termine la preuve par la pleine fidélité de $h^{-1}$.
\end{proof}

\begin{lemma}
\label{lemma-cohomological-descent-for-ph-hypercovering}
Soient $S$ un schéma et $X$ un espace algébrique sur $S$.
Soit $U$ un espace algébrique simplicial sur $S$. Soit $a : U \to X$
une augmentation. Si $a : U \to X$ est un hyperrecouvrement propre de $X$,
alors pour $K \in D^+(X_\etale)$
$$
K \to Ra_*(a^{-1}K)
$$
est un isomorphisme. Ici, $a : \Sh(U_\etale) \to \Sh(X_\etale)$
est comme dans la section \ref{section-simplicial-algebraic-spaces}.
\end{lemma}

\begin{proof}
Considérons le diagramme du lemme \ref{lemma-compare-simplicial-objects-ph-etale}.
Observons que $Rh_{n, *}h_n^{-1}$ est le foncteur identité
sur $D^+(U_{n, \etale})$ par
Compléments sur la cohomologie des espaces, lemme
\ref{spaces-more-cohomology-lemma-cohomological-descent-etale-ph}.
Par conséquent, $Rh_*h^{-1}$ est le foncteur identité sur
$D^+(U_\etale)$ par
le lemme \ref{lemma-direct-image-morphism-simplicial-sites}.
Nous avons
\begin{align*}
Ra_*(a^{-1}K)
& =
Ra_*Rh_*h^{-1}a^{-1}K \\
& =
Rh_{-1, *}Ra_{ph, *}a_{ph}^{-1}(h_{-1})^{-1}K \\
& =
Rh_{-1, *}(h_{-1})^{-1}K \\
& =
K
\end{align*}
La première égalité découle de la discussion ci-dessus; la deuxième,
de la commutativité du diagramme du lemme
\ref{lemma-compare-simplicial-objects}; la troisième, du lemme
le lemme \ref{lemma-hypercovering-X-simple-descent-bounded-abelian}
puisque $U$ est un hyperrecouvrement de $X$ dans $(\textit{Spaces}/S)_{ph}$
d'après Topologies sur les espaces, lemme
\ref{spaces-topologies-lemma-surjective-proper-ph},
et la dernière, du lemme déjà utilisé de
Compléments sur la cohomologie des espaces, lemme
\ref{spaces-more-cohomology-lemma-cohomological-descent-etale-ph}.
\end{proof}

\begin{lemma}
\label{lemma-compute-via-ph-hypercovering}
Soient $S$ un schéma et $X$ un espace algébrique sur $S$.
Soit $U$ un espace algébrique simplicial sur $S$. Soit $a : U \to X$
une augmentation. Si $a : U \to X$ est un hyperrecouvrement propre de $X$, alors
$$
R\Gamma(X_\etale, K) = R\Gamma(U_\etale, a^{-1}K)
$$
pour $K \in D^+(X_\etale)$. Ici, $a : \Sh(U_\etale) \to \Sh(X_\etale)$
est comme dans la section \ref{section-simplicial-algebraic-spaces}.
\end{lemma}

\begin{proof}
Cela découle du lemme
\ref{lemma-cohomological-descent-for-ph-hypercovering}
car $R\Gamma(U_\etale, -) = R\Gamma(X_\etale, -) \circ Ra_*$ par
Cohomologie sur les sites, remarque \ref{sites-cohomology-remark-before-Leray}.
\end{proof}

\begin{lemma}
\label{lemma-ph-hypercovering-equivalence-bounded}
Soient $S$ un schéma et $X$ un espace algébrique sur $S$.
Soit $U$ un espace algébrique simplicial sur $S$. Soit $a : U \to X$
une augmentation.
Soit $\mathcal{A} \subset \textit{Ab}(U_\etale)$
la sous-catégorie de Serre faible des faisceaux abéliens cartésiens.
Si $U$ est un hyperrecouvrement propre de $X$, alors
le foncteur $a^{-1}$ définit une équivalence
$$
D^+(X_\etale) \longrightarrow D_\mathcal{A}^+(U_\etale)
$$
de quasi-inverse $Ra_*$. Ici, $a : \Sh(U_\etale) \to \Sh(X_\etale)$
est comme dans la section \ref{section-simplicial-algebraic-spaces}.
\end{lemma}

\begin{proof}
Observons que $\mathcal{A}$ est une sous-catégorie de Serre faible d'après le
lemme \ref{lemma-Serre-subcat-cartesian-modules}.
L'équivalence est une conséquence formelle
des résultats obtenus jusqu'à présent. On utilise les lemmes
\ref{lemma-descent-sheaves-for-ph-hypercovering} et
\ref{lemma-cohomological-descent-for-ph-hypercovering}, ainsi que le lemme de Cohomologie
sur les sites \ref{sites-cohomology-lemma-equivalence-bounded}.
\end{proof}

\begin{lemma}
\label{lemma-spectral-sequence-ph-hypercovering}
Soit $S$ un schéma. Soit $X$ un espace algébrique sur $S$.
Soit $U$ un espace algébrique simplicial sur $S$. Soit $a : U \to X$
une augmentation. Soit $\mathcal{F}$ un faisceau abélien
sur $X_\etale$. Notons $\mathcal{F}_n$ son image inverse sur $U_{n, \etale}$.
Si $U$ est un hyperrecouvrement ph de $X$, alors
il existe une suite spectrale canonique
$$
E_1^{p, q} = H^q_\etale(U_p, \mathcal{F}_p)
$$
convergeant vers $H^{p + q}_\etale(X, \mathcal{F})$.
\end{lemma}

\begin{proof}
Conséquence immédiate des lemmes \ref{lemma-compute-via-ph-hypercovering}
et \ref{lemma-simplicial-sheaf-cohomology-site}.
\end{proof}










\begin{multicols}{2}[\section{Autres chapitres}]
\noindent
Préliminaires
\begin{enumerate}
\item \hyperref[introduction-section-phantom]{Introduction}
\item \hyperref[conventions-section-phantom]{Conventions}
\item \hyperref[sets-section-phantom]{Théorie des ensembles}
\item \hyperref[categories-section-phantom]{Catégories}
\item \hyperref[topology-section-phantom]{Topologie}
\item \hyperref[sheaves-section-phantom]{Faisceaux sur les espaces}
\item \hyperref[sites-section-phantom]{Sites et faisceaux}
\item \hyperref[stacks-section-phantom]{Champs}
\item \hyperref[fields-section-phantom]{Corps}
\item \hyperref[algebra-section-phantom]{Algèbre commutative}
\item \hyperref[brauer-section-phantom]{Groupes de Brauer}
\item \hyperref[homology-section-phantom]{Algèbre homologique}
\item \hyperref[derived-section-phantom]{Catégories dérivées}
\item \hyperref[simplicial-section-phantom]{Méthodes simpliciales}
\item \hyperref[more-algebra-section-phantom]{Compléments d'algèbre}
\item \hyperref[smoothing-section-phantom]{Lissage des morphismes d'anneaux}
\item \hyperref[modules-section-phantom]{Faisceaux de modules}
\item \hyperref[sites-modules-section-phantom]{Modules sur les sites}
\item \hyperref[injectives-section-phantom]{Objets injectifs}
\item \hyperref[cohomology-section-phantom]{Cohomologie des faisceaux}
\item \hyperref[sites-cohomology-section-phantom]{Cohomologie sur les sites}
\item \hyperref[dga-section-phantom]{Algèbre différentielle graduée}
\item \hyperref[dpa-section-phantom]{Algèbre à puissances divisées}
\item \hyperref[sdga-section-phantom]{Faisceaux différentiels gradués}
\item \hyperref[hypercovering-section-phantom]{Hyperrecouvrements}
\end{enumerate}
Schémas
\begin{enumerate}
\setcounter{enumi}{25}
\item \hyperref[schemes-section-phantom]{Schémas}
\item \hyperref[constructions-section-phantom]{Constructions de schémas}
\item \hyperref[properties-section-phantom]{Propriétés des schémas}
\item \hyperref[morphisms-section-phantom]{Morphismes de schémas}
\item \hyperref[coherent-section-phantom]{Cohomologie des schémas}
\item \hyperref[divisors-section-phantom]{Diviseurs}
\item \hyperref[limits-section-phantom]{Limites de schémas}
\item \hyperref[varieties-section-phantom]{Variétés}
\item \hyperref[topologies-section-phantom]{Topologies sur les schémas}
\item \hyperref[descent-section-phantom]{Descente}
\item \hyperref[perfect-section-phantom]{Catégories dérivées de schémas}
\item \hyperref[more-morphisms-section-phantom]{Compléments sur les morphismes}
\item \hyperref[flat-section-phantom]{Compléments sur la platitude}
\item \hyperref[groupoids-section-phantom]{Schémas en groupoïdes}
\item \hyperref[more-groupoids-section-phantom]{Compléments sur les schémas en groupoïdes}
\item \hyperref[etale-section-phantom]{Morphismes étales de schémas}
\end{enumerate}
Thèmes de théorie des schémas
\begin{enumerate}
\setcounter{enumi}{41}
\item \hyperref[chow-section-phantom]{Homologie de Chow}
\item \hyperref[intersection-section-phantom]{Théorie de l'intersection}
\item \hyperref[pic-section-phantom]{Schémas de Picard des courbes}
\item \hyperref[weil-section-phantom]{Théories cohomologiques de Weil}
\item \hyperref[adequate-section-phantom]{Modules adéquats}
\item \hyperref[dualizing-section-phantom]{Complexes dualisants}
\item \hyperref[duality-section-phantom]{Dualité pour les schémas}
\item \hyperref[discriminant-section-phantom]{Discriminants et différentes}
\item \hyperref[derham-section-phantom]{Cohomologie de de Rham}
\item \hyperref[local-cohomology-section-phantom]{Cohomologie locale}
\item \hyperref[algebraization-section-phantom]{Géométrie algébrique et formelle}
\item \hyperref[curves-section-phantom]{Courbes algébriques}
\item \hyperref[resolve-section-phantom]{Résolution des surfaces}
\item \hyperref[models-section-phantom]{Réduction semi-stable}
\item \hyperref[functors-section-phantom]{Foncteurs et morphismes}
\item \hyperref[equiv-section-phantom]{Catégories dérivées de variétés}
\item \hyperref[pione-section-phantom]{Groupes fondamentaux de schémas}
\item \hyperref[etale-cohomology-section-phantom]{Cohomologie étale}
\item \hyperref[crystalline-section-phantom]{Cohomologie cristalline}
\item \hyperref[proetale-section-phantom]{Cohomologie pro-étale}
\item \hyperref[relative-cycles-section-phantom]{Cycles relatifs}
\item \hyperref[more-etale-section-phantom]{Compléments de cohomologie étale}
\item \hyperref[trace-section-phantom]{La formule des traces}
\end{enumerate}
Espaces algébriques
\begin{enumerate}
\setcounter{enumi}{64}
\item \hyperref[spaces-section-phantom]{Espaces algébriques}
\item \hyperref[spaces-properties-section-phantom]{Propriétés des espaces algébriques}
\item \hyperref[spaces-morphisms-section-phantom]{Morphismes d'espaces algébriques}
\item \hyperref[decent-spaces-section-phantom]{Espaces algébriques décents}
\item \hyperref[spaces-cohomology-section-phantom]{Cohomologie des espaces algébriques}
\item \hyperref[spaces-limits-section-phantom]{Limites d'espaces algébriques}
\item \hyperref[spaces-divisors-section-phantom]{Diviseurs sur les espaces algébriques}
\item \hyperref[spaces-over-fields-section-phantom]{Espaces algébriques sur des corps}
\item \hyperref[spaces-topologies-section-phantom]{Topologies sur les espaces algébriques}
\item \hyperref[spaces-descent-section-phantom]{Descente et espaces algébriques}
\item \hyperref[spaces-perfect-section-phantom]{Catégories dérivées d'espaces}
\item \hyperref[spaces-more-morphisms-section-phantom]{Compléments sur les morphismes d'espaces}
\item \hyperref[spaces-flat-section-phantom]{Platitude sur les espaces algébriques}
\item \hyperref[spaces-groupoids-section-phantom]{Groupoïdes dans les espaces algébriques}
\item \hyperref[spaces-more-groupoids-section-phantom]{Compléments sur les groupoïdes dans les espaces}
\item \hyperref[bootstrap-section-phantom]{Amorçage}
\item \hyperref[spaces-pushouts-section-phantom]{Sommes amalgamées d'espaces algébriques}
\end{enumerate}
Thèmes de géométrie
\begin{enumerate}
\setcounter{enumi}{81}
\item \hyperref[spaces-chow-section-phantom]{Groupes de Chow des espaces}
\item \hyperref[groupoids-quotients-section-phantom]{Quotients de groupoïdes}
\item \hyperref[spaces-more-cohomology-section-phantom]{Compléments sur la cohomologie des espaces}
\item \hyperref[spaces-simplicial-section-phantom]{Espaces simpliciaux}
\item \hyperref[spaces-duality-section-phantom]{Dualité pour les espaces}
\item \hyperref[formal-spaces-section-phantom]{Espaces algébriques formels}
\item \hyperref[restricted-section-phantom]{Algébrisation des espaces formels}
\item \hyperref[spaces-resolve-section-phantom]{Résolution des surfaces revisitée}
\end{enumerate}
Théorie des déformations
\begin{enumerate}
\setcounter{enumi}{89}
\item \hyperref[formal-defos-section-phantom]{Théorie formelle des déformations}
\item \hyperref[defos-section-phantom]{Théorie des déformations}
\item \hyperref[cotangent-section-phantom]{Le complexe cotangent}
\item \hyperref[examples-defos-section-phantom]{Problèmes de déformation}
\end{enumerate}
Champs algébriques
\begin{enumerate}
\setcounter{enumi}{93}
\item \hyperref[algebraic-section-phantom]{Champs algébriques}
\item \hyperref[examples-stacks-section-phantom]{Exemples de champs}
\item \hyperref[stacks-sheaves-section-phantom]{Faisceaux sur les champs algébriques}
\item \hyperref[criteria-section-phantom]{Critères de représentabilité}
\item \hyperref[artin-section-phantom]{Axiomes d'Artin}
\item \hyperref[quot-section-phantom]{Espaces de Quot et de Hilbert}
\item \hyperref[stacks-properties-section-phantom]{Propriétés des champs algébriques}
\item \hyperref[stacks-morphisms-section-phantom]{Morphismes de champs algébriques}
\item \hyperref[stacks-limits-section-phantom]{Limites de champs algébriques}
\item \hyperref[stacks-cohomology-section-phantom]{Cohomologie des champs algébriques}
\item \hyperref[stacks-perfect-section-phantom]{Catégories dérivées de champs}
\item \hyperref[stacks-introduction-section-phantom]{Introduction aux champs algébriques}
\item \hyperref[stacks-more-morphisms-section-phantom]{Compléments sur les morphismes de champs}
\item \hyperref[stacks-geometry-section-phantom]{La géométrie des champs}
\end{enumerate}
Thèmes de théorie des espaces de modules
\begin{enumerate}
\setcounter{enumi}{107}
\item \hyperref[moduli-section-phantom]{Champs de modules}
\item \hyperref[moduli-curves-section-phantom]{Espaces de modules de courbes}
\end{enumerate}
Divers
\begin{enumerate}
\setcounter{enumi}{109}
\item \hyperref[examples-section-phantom]{Exemples}
\item \hyperref[exercises-section-phantom]{Exercices}
\item \hyperref[guide-section-phantom]{Guide bibliographique}
\item \hyperref[desirables-section-phantom]{Desiderata}
\item \hyperref[coding-section-phantom]{Style de programmation}
\item \hyperref[obsolete-section-phantom]{Obsolète}
\item \hyperref[fdl-section-phantom]{Licence de documentation libre GNU}
\item \hyperref[index-section-phantom]{Index généré automatiquement}
\end{enumerate}
\end{multicols}


\bibliography{my}
\bibliographystyle{amsalpha}

\end{document}
